% Commenti speciali

	% !TEX encoding = UTF-8 Unicode
	% !TEX TS-program = pdflatex
	% !TEX spellcheck = it-IT

% Tipo di documento e impostazioni globali

\documentclass[12 pt, a4paper, openright, titlepage, twoside, leqno]{book}

\usepackage{geometry}
\geometry{a4paper, right= 2.6cm, left = 2.6 cm, bottom=3.4cm, top=3.1cm}

\usepackage{setspace}
\onehalfspacing

% Pacchetti principali

\usepackage[T1]{fontenc}
\usepackage{latexsym,verbatim}
\usepackage[utf8]{inputenc}
\usepackage[italian]{babel}
\usepackage{graphicx}
\usepackage{amsfonts,mathrsfs,amsmath,amssymb,amsthm}
\usepackage{indentfirst}
\usepackage{microtype}
\usepackage{booktabs, caption}
\usepackage{lmodern}
\usepackage{braket}
\usepackage{mathtools}
\usepackage{emptypage}
\usepackage{afterpage}
\usepackage{pdfpages}

\usepackage{stackrel}
\usepackage{pgfplots}
\usepackage{pgf}
\usepackage{multirow}
\usepackage{imakeidx} % pacchetto per indice analitico
\usepackage{bold-extra}
\usepackage{extpfeil}
\usepackage{comment}
\usepackage{fancyhdr}
\usepackage{enumerate}
\usepackage{eso-pic} 
\usepackage{tikz}
\usepackage{graphicx}
\usepackage{wasysym}

% Impostazioni TikZ
\tikzset{every picture/.style={line width=0.75pt}} %set default line width to 0.75pt   

% Pacchetto per integrale barrato
\usepackage{mathtools}
\usepackage{esint}


\usepackage{color}
\definecolor{PrussianBlue}{rgb}{0.0, 0.19, 0.33}
	 
% Funzione utile per valore assoluto e norma

\DeclarePairedDelimiter{\abs}{\lvert}{\rvert}
\DeclarePairedDelimiter{\norm}{\lVert}{\rVert}


%Funzioni utili a scrivere insiemi numerici

\newcommand{\numberset}{\mathbb}
\newcommand{\N}{\numberset{N}}
\newcommand{\R}{\numberset{R}}
\newcommand{\Z}{\numberset{Z}}
\newcommand{\Q}{\numberset{Q}}
\newcommand{\C}{\numberset{C}}
\newcommand{\K}{\numberset{K}}

\newcommand{\B}{\textup{B}}

\renewcommand{\div}{\textup{div}}
	
	
% Definizione ambienti speciali per teoremi e affini

\renewcommand{\thesection}{\arabic{section}}
\numberwithin{equation}{section} 
\renewcommand{\theequation}{\arabic{section}.\arabic{equation}}

\newenvironment{theorframe}[1]{\refstepcounter{equation}\par\medskip\noindent%
	{\bf (\theequation)~#1}\hspace{5pt}\em}%
{\em\par\medskip}
\newenvironment{ossframe}[1]{\refstepcounter{equation}\par\medskip\noindent%
	{\bf (\theequation)~#1}\hspace{5pt}}%
{\par\medskip}

\newenvironment{teorema}{\begin{theorframe}{Teorema}}{\end{theorframe}}
\newenvironment{lemma}{\begin{theorframe}{Lemma}}{\end{theorframe}}
\newenvironment{corollario}{\begin{theorframe}{Corollario}}{\end{theorframe}}
\newenvironment{proposizione}{\begin{theorframe}{Proposizione}}{\end{theorframe}}
\newenvironment{osservazione}{\begin{theorframe}{Osservazione}}{\end{theorframe}}
\newenvironment{notazione}{\begin{ossframe}{Notazione}}{\end{ossframe}}
\newenvironment{definizione}{\begin{theorframe}{Definizione}}{\end{theorframe}}
\newenvironment{esempio}{\begin{theorframe}{Esempio}}{\end{theorframe}}

\newenvironment{problema}{\begin{theorframe}{Problema}}{\end{theorframe}}

\renewcommand{\qed}{~\vrule height5pt depth0pt width5pt}

% Indice analitico (opzioni generali)

\makeindex[name = analitico, intoc, columns = 2]
\makeindex[name = simboli, title = Elenco dei simboli, intoc, columns=2]

% Definizione comando mail

\newcommand{\mail}[1]{\href{mailto:#1}{\texttt{#1}}}

% Proprietà collegamenti ipertestuali

\usepackage[colorlinks]{hyperref}
\hypersetup{hidelinks, colorlinks=true, linkcolor=PrussianBlue, urlcolor=PrussianBlue, linktoc=page}

\pgfplotsset{compat=1.18}

\allowdisplaybreaks

% Documento: il documento è strutturato, per semplicità, in un insieme di sottofile

	\begin{document}
		
		% Commenti speciali

	% !TEX encoding = UTF-8 Unicode
	% !TEX TS-program = pdflatex
	% !TEX spellcheck = it-IT
	% !TEX root = Dispensa/afed.tex

% Frontespizio 

	\begin{titlepage}
		\begin{center}
			{\Large \textbf{UNIVERSITÀ CATTOLICA DEL SACRO CUORE}}
			
			\vspace{20 pt}
			
			{\Large \textbf{Facoltà di Scienze Matematiche, Fisiche e Naturali}}
		\end{center}
		\vspace{180 pt}
		\begin{center}
			{\LARGE \textit{\textbf{ANALISI FUNZIONALE}}}
			
			\vspace{20 pt}
			
			{\LARGE \textbf{ed}}
			
			\vspace{20 pt}
			
			{\LARGE \textit{\textbf{EQUAZIONI DIFFERENZIALI}}}
		\end{center}
		\vspace{190 pt}
		\begin{center}
			{\Large \textbf{Prof. Marco Squassina}}
		\end{center}
		
		\vfill
		
		\begin{center}
			{\Large \textbf{Anno Accademico 2024/2025}}
		\end{center}
	\end{titlepage}
	
	\newpage
	\thispagestyle{empty}
	\hfill
	
	\newpage 	    % Frontespizio 
		
		% Commenti speciali

	% !TEX encoding = UTF-8 Unicode
	% !TEX TS-program = pdflatex
	% !TEX spellcheck = it-IT
	% !TEX root = Dispensa/afed.tex


\thispagestyle{empty}

% Ringraziamenti 


Questa dispensa è stata scritta da Mattia Garatti e raccoglie gli appunti del corsi \emph{Analisi Funzionale} ed \emph{Equazioni Differenziali} tenuti dal Prof. Marco Squassina negli anni accademici 2023/2024 e 2024/2025 (con qualche aggiunta).

Realizzare una dispensa è un'operazione complessa, che richiede numerosi controlli. L'esperienza suggerisce che è praticamente impossibile realizzare un'opera priva di errori. Potete segnalare eventuali errori al seguente indirizzo di posta elettronica:
\[ \mail{mattiagaratti@gmail.com} \]

%\vspace{20 pt}

\noindent \textbf{Ringraziamenti} 
\begin{itemize}
	\item Si ringrazia il Professor Marco Degiovanni per le brillanti osservazioni.
	\item Si ringrazia il Dottor Marco Marzocchi per i consigli sull'utilizzo del linguaggio \LaTeX $\,$ e per la diffusione della passione per l'Analisi Matematica.
	\item Si ringrazia la Dottoressa Chiara Lonati per gli inserti di Fisica Matematica.
	\item Si ringrazia il Dottor Riccardo Moraschi per i disegni del Teorema di estensione, delle immersioni di Sobolev ed del Lemma di Hopf, nonché per le puntuali correzioni e i preziosi consigli.
	\item Si ringrazia il Dottor Gabriele Geroldi per l'aiuto nella revisione di bozze.
	\item Si ringrazia il Dottor Luca Tamanini per le fruttuose discussioni e la revisione di bozze per le parti riguardanti gli spazi riflessivi e la rappresentazione degli elementi del duale di uno spazio di Lebesgue.
\end{itemize}

%\vspace{20 pt}

\noindent \textbf{Nota dell'autore} Non tutte le $C$ sono uguali.

\vfill

\begin{flushright}
	Ultimo aggiornamento: \today
\end{flushright} 	    % Ringraziamenti 
		
		\tableofcontents                                          % Indice
		\thispagestyle{empty}
		
		% Commenti speciali

	% !TEX encoding = UTF-8 Unicode
	% !TEX TS-program = pdflatex
% !TeX spellcheck = it_IT
	% !TEX root = Dispensa/afed.tex

% Capitolo 2

\chapter*{Preliminari}
\addcontentsline{toc}{chapter}{Preliminari}
\markboth{PRELIMINARI}{}
\section{Spazi riflessivi}
\begin{notazione}
	Dato $X$ uno spazio normato su $\K$, indichiamo con $X' = \mathcal{L}(X;\K)$ il suo duale topologico\index[simboli]{$X'$}\index[analitico]{duale topologico} e con $X'' = \mathcal{L}(X';\K)$ il suo biduale topologico\index[simboli]{$X''$}\index[analitico]{biduale topologico}.
\end{notazione}
\begin{definizione}
	Sia $X$ uno spazio normato su $\K$. Chiamiamo \emph{iniezione canonica}\index[simboli]{$J$}\index[analitico]{iniezione canonica} di $X$ in $X''$, l'applicazione $J: X \to X''$ tale che per ogni $x \in X$, per ogni $\varphi \in X'$
	\[ \langle Jx,\varphi\rangle = \langle\varphi, x\rangle. \]
\end{definizione}
\begin{proposizione}
	Sia $X$ uno spazio normato su $\K$ e $J$ l'iniezione canonica di $X$ in $X''$. Valgono i seguenti fatti:
	\begin{enumerate}[$(a)$]
		\item $J$ è ben definita,
		\item $J$ è un'isometria, ovvero per ogni $x \in X$
		\[ \norm{Jx}_{X''} = \norm{x}_{X}, \]
		\item $J$ è lineare e iniettiva\footnote{Il nome è quindi giustificato.}.
	\end{enumerate}
\end{proposizione}
\begin{proof}
	Omettiamo la dimostrazione.
\end{proof}
\begin{definizione}
	Sia $X$ uno spazio normato su $\K$. Diciamo che $X$ è \emph{riflessivo}\index[analitico]{spazio!riflessivo} se $J$ è suriettiva.
\end{definizione}
\begin{proposizione}
	Sia $X$ uno spazio normato su $\K$ di dimensione finita. Allora $J$ è biettiva. In particolare $X$ è riflessivo.
\end{proposizione}
\begin{proof}
	È sufficiente combinare il fatto che $J$ sia iniettiva con il Teorema di Nullità + Rango.
\end{proof}
\begin{teorema}
	Sia $X$ uno spazio normato su $\K$. Allora $X$ è riflessivo se e solo se $X'$ è riflessivo.
\end{teorema}
\begin{proof}
	Omettiamo la dimostrazione.
\end{proof}
\begin{teorema}
	Siano $X$ uno spazio di Banach su $\K$ riflessivo e $Y \le X$ chiuso. Allora $Y$ è riflessivo.
\end{teorema}
\begin{proof}
	Omettiamo la dimostrazione.
\end{proof}
\begin{proposizione}\label{normrifban}
	Sia $X$ uno spazio normato su $\K$ riflessivo. Allora $X$ è di Banach.
\end{proposizione}
\begin{proof}
	Sia $(x_{h})$ di Cauchy in $X$. Allora per ogni $\varepsilon > 0$ esiste $\bar{h} \in \N$ tale che per ogni $h,k \ge \bar{h}$, $\norm{x_{h}- x_{k}}_{X}<\varepsilon$. Detta $J$ l'iniezione canonica associata a $X$,
	\[ \norm{Jx_{h} - Jx_{k}}_{X''} = \norm{J(x_{h} - x_{k})}_{X''} = \norm{x_{h} - x_{k}}_{X} < \varepsilon \]
	ossia $(Jx_{h})$ è di Cauchy in $X''$, che è completo essendo $\K$ completo. Di conseguenza esiste $T \in X''$ tale che $Jx_{h} \to T$. 
	
	Poiché $X$ è riflessivo, esiste $x \in X $ tale che $T = Jx$. Essendo $J$ un'isometria,
	\[ \norm{Jx_{h} - Jx}_{X''} = \norm{x_{h} - x}_{X}\]
	e passando al limite otteniamo che $x_{h} \to x$.
\end{proof}
\begin{esempio}
	Sia 
	\[ C_{00}= \left\lbrace (x_{h}) \in l^{\infty}(\N) \;:\; \textup{esiste } m \in \N \textup{ tale che } x_{h} = 0 \textup{ per ogni } h \ge m\right\rbrace \]
	munito della norma del $\sup$. Allora $(C_{00},\norm{\cdot}_\infty)$ non è riflessivo.\index[simboli]{$C_{00}$} 
\end{esempio}
\begin{proof}
	Consideriamo la successione in $C_{00}$ definita da
	\[ x_{j} = (1,\frac{1}{2},\frac{1}{4}, \dots, \frac{1}{2^{j}}, 0, \dots, 0, \dots). \]
	Siccome $x_{j} \to (1,\frac{1}{2},\frac{1}{4}, \dots, \frac{1}{2^{j}}, \frac{1}{2^{j+1}}, \dots) \notin C_{00}$, possiamo affermare che $C_{00}$ non è chiuso. Non può quindi essere completo. La tesi discende quindi dalla Proposizione \eqref{normrifban}.
\end{proof}
\begin{definizione}
	Sia $X$ uno spazio normato su $\K$. Diciamo che $X$ è \emph{uniformemente convesso}\index[analitico]{spazio!uniformemente convesso} se per ogni $\varepsilon > 0$ esiste $\delta >0$ tale che
	\[ \forall x,y \in X, \norm{x},\norm{y} \le 1 : \norm{x-y} \ge \varepsilon \Longrightarrow \norm*{\frac{x+y}{2}}  < 1-\delta.\]
\end{definizione}
Possiamo immaginare uno spazio normato uniformemente convesso come uno spazio normato la cui palla unitaria sia in un certo senso ``tonda''. Più precisamente, si deve avere che il punto medio di segmenti con estremi sul bordo della palla unitaria sia contenuto nell'interno della palla unitaria.

Il risultato seguente collega una nozione geometrica, l'uniforme convessità, con una nozione topologica, la riflessività. Premettiamo subito che il risultato non è invertibile, i controesempi sono molteplici.
\begin{teorema}\label{uniconvrif}
	Sia $X$ uno spazio di Banach su $\K$ uniformemente convesso. Allora $X$ è riflessivo.
\end{teorema}
\begin{proof}
	Omettiamo la dimostrazione.
\end{proof}

\section{Alcuni complementi sugli spazi di Legesgue} 
Iniziamo con alcune considerazioni di carattere generale.
\begin{proposizione}
	Consideriamo $(X,\norm{\;})$ uno spazio normato su $\K$, $(x_h)$ in $X$ e $x \in X$. I seguenti fatti sono equivalenti:
	\begin{enumerate}[$(a)$]
		\item $x_h \to x$ in $X$,
		\item per ogni sottosuccessione $(x_{h_k})$, esiste una sotto-sottosuccessione $(x_{h_{k_j}})$ tale che $x_{h_{k_j}} \to x$ in $X$.
	\end{enumerate}
\end{proposizione}
\begin{proof}\hfill
	\par\noindent
	$(a) \Longrightarrow (b)$ Evidente.
	\par\noindent
	$(b) \Longrightarrow (a)$ Supponiamo, per assurdo, che esista $\varepsilon > 0$ tale che per ogni $k > 0$ esista $h_k \ge k$ tale che $\norm{x_{h_k}-x} > \varepsilon$. Allora, per ogni sottosuccessione $(x_{h_{k_j}})$ di $(x_{h_k})$, per ogni $j \in \N$, $\norm{x_{h_{k_j}}-x} > \varepsilon$, ossia $(x_{h_{k_j}})$ non converge a $x$, assurdo.
\end{proof}

\begin{proposizione}
Siano $p \in \left[ 1,+\infty\right[ $, $k \in \N$ tale che $k \ge 2$ ed $x_1,\dots, x_k \ge 0$. Si ha
\[ \left( \sum_{i=1}^{k}x_i\right)^p \le k^{p-1}\sum_{i=1}^{k}x_i^p. \]
\end{proposizione}
\begin{proof}
Si tratta di usare la convessità di $\left\lbrace t \mapsto t^p\right\rbrace$ a più punti. Infatti, $\displaystyle\sum_{i=1}^{k}\frac{1}{k}x_i $ è una combinazione convessa di $x_1,\dots, x_k$ essendo $\displaystyle\sum_{i=1}^{k}\frac{1}{k} = 1$, 
quindi
\begin{align*}
	\left( \sum_{i=1}^{k}x_i\right)^p = \left( k\sum_{i=1}^{k}\frac{1}{k}x_i\right)^p = k^p  \left(\sum_{i=1}^{k}\frac{1}{k}x_i\right)^p \le k^p \sum_{i=1}^{k}\frac{1}{k}x_i^p = k^{p-1}\sum_{i=1}^{k}x_i^p.\qedhere
\end{align*}
\end{proof}

Nel seguito $U$ indicherà un generico sottoinsieme aperto di $\R^n$. 

\begin{teorema}\emph{\textbf{(disuguaglianza di interpolazione)}}\index[analitico]{disuguaglianza! di interpolazione}
Siano $1 \le q \le p \le r \le \infty$ e $\vartheta \in \left] 0,1\right[ $ tali che 
\[ \frac{1}{p} = \frac{\vartheta}{q} + \frac{1-\vartheta}{r}.  \]
Se $u \in L^q(U)\cap L^r(U)$, allora $u \in L^p(U)$ e 
\[ \norm{u}_{L^p(U)} \le \norm{u}_{L^q(U)}^\vartheta\norm{u}_{L^r(U)}^{1-\vartheta}.\]

\begin{proof}
Vediamo solo il caso finito. Innanzitutto,
\[ \int_{U} \abs{u}^pd\mathcal{L}^n = \int_{U} \abs{u}^{\vartheta p}\abs{u}^{(1-\vartheta)p}d\mathcal{L}^n  \]
e per la disuguaglianza di Holder\footnote{Essendo \[ 1 = \frac{\vartheta p}{q} + \frac{(1-\vartheta)p}{r}.  \]}
\[ \int_{U} \abs{u}^pd\mathcal{L}^n \le \left( \int_{U}\abs{u}^{q}d\mathcal{L}^n\right) ^{\frac{\vartheta p}{q}} \left( \int_{U}\abs{u}^rd\mathcal{L}^n\right)^{\frac{(1 -\vartheta) p}{r}} = \norm{u}_{L^q(U)}^{\vartheta p}  \norm{u}_{L^r(U)}^{(1- \vartheta) p} \]
da cui la tesi.
\end{proof}
\end{teorema}
\begin{teorema}\textbf{\emph{(disuguaglianza di Holder generalizzata)}}\footnote{Ma è strabiliante! Si potrebbe quasi dire che è \emph{over 9000}. A parte le citazioni di DragonBall (e parodie), si tratta della versione più generale della disuguaglianza di Holder valida per gli spazi normati $\left( L^p(U),\norm{\ }_p\right) $.}\index[analitico]{disuguaglianza! di Holder generalizzata}\label{holdergeneralizzato}
Consideriamo $1\le p_1,\dots, p_N \le \infty$. Sia $1\le r\le \infty$  tale che 
\[ \frac{1}{r} = \frac{1}{p_1} + \dots + \frac{1}{p_N}.\footnote{Di fatto $r$ è completamente determinato dai $p_1,\dots, p_N$.}\]
Se $f_1 \in L^{p_1}(U), \dots, f_N \in L^{p_N}(U)$, allora $f = f_1 \dots f_N \in L^{r}(U)$ e 
\[ \norm{f}_r \le \prod_{i=1}^{N} \norm{f_i}_{p_i}. \] 
\begin{proof}
	Supponiamo che per ogni $k \in \N$ tale che $k < N$ valga la tesi. Mostriamo che vale anche per $N$. Innanzitutto, a meno di riordinare gli indici ed il prodotto delle $f_i$, possiamo supporre $p_1 \le \dots \le p_N$. Se per ogni $i = 1,\dots, N$, $p_i = \infty$, allora anche $r = \infty$ e la tesi è evidente. Se invece esiste $ j \in \N$ tale che $0 <j < N$ e $p_i = \infty$ per $i = j+1, \dots, N$ e $p_i < \infty$ per $i= 1,\dots, j$, allora
	\[ \frac{1}{r} = \frac{1}{p_1} +\dots + \frac{1}{p_{j}} \] 
	ed essendo $j < N$,
	\[ \int_{U} \abs{f}^{r} d\mathcal{L}^n \le \prod_{i=j+1}^{N} \norm{f_i}_\infty^r \int_{U} \prod_{i=1}^{j}\abs{f_i}^r d\mathcal{L}^n \le \prod_{i=j+1}^{N} \norm{f_i}_\infty^r \prod_{i=1}^{j} \norm{f_i}^r_{p_i} = \left( \prod_{i=1}^N \norm{f_i}_{p_i}\right)^r \] 
	da cui 
	\[ \norm{f}_r \le \prod_{i=1}^{N}\norm{f_i}_{p_i}. \]
	Se poi $p_N < \infty$, allora per ogni $i = 1,\dots, N-1$, $p_i < \infty$ e quindi $r < \infty$. Ora, evidentemente $p_N > r$ e quindi, per la disuguaglianza di Holder\footnote{Quella classica, che applichiamo qui con esponente $p = \frac{p_N}{r}$.},
	\[ \int_{U} \abs{f}^{r} d\mathcal{L}^n = \int_{U} \abs{f_N}^r \prod_{i=1}^{N-1} \abs{f_i}^r d\mathcal{L}^n \le \left( \int_{U} \abs{f_N}^{p_N} d\mathcal{L}^n\right) ^{\frac{r}{p_N}} \left(\int_{U} \prod_{i=1}^{N-1} \abs{f_i}^{\frac{p_Nr}{p_N - r}}d\mathcal{L}^n\right) ^{\frac{p_N-r}{p_N}}. \]
	Ora, dal fatto che 
	\begin{align*}
		N-1 &< N, & \frac{p_Nr}{p_N - r} &> 1, & \frac{p_N-r}{p_N r} &= \frac{1}{p_1} + \dots +\frac{1}{p_{N-1}},
	\end{align*}
	possiamo applicare l'ipotesi induttiva forte e ottenere
	\[\int_{U} \prod_{i=1}^{N-1} \abs{f_i}^{\frac{p_Nr}{p_N - r}}d\mathcal{L}^n \le \prod_{i=1}^{N-1} \left( \int_{U} \abs{f_i}^{p_i}d\mathcal{L}^n\right) ^{\dfrac{\frac{p_Nr}{p_N -r}}{p_i}}.\]
	Pertanto,
	\[ \int_{U} \abs{f}^{r} d\mathcal{L}^n \le \left( \int_{U} \abs{f_N}^{p_N} d\mathcal{L}^n\right) ^{\frac{r}{p_N}}  \left(\prod_{i=1}^{N-1} \norm{f_i}_{p_i} \right) ^{r} = \left(\prod_{i=1}^{N} \norm{f_i}_{p_i} \right) ^{r} \]
	da cui la tesi.
\end{proof}
\end{teorema}

\begin{osservazione}
È possibile estendere il Teorema \eqref{holdergeneralizzato} anche per $r, p_1,\dots, p_N \in \left] 0,1\right[ $. Tuttavia si ricorda che in questo caso non abbiamo a che fare con delle norme in quanto non vale la disuguaglianza triangolare, pertanto il Teorema \eqref{holdergeneralizzato} è di fatto la versione più generale se si vuole lavorare con delle norme. In ogni caso la generalizzazione è diretta, lo stesso enunciato può essere scritto per $r, p_1,\dots, p_N \in \left] 0,1\right[ $ intendendo appropriatamente il simbolo $\norm{\ }$.
\end{osservazione}
Presentiamo ora un criterio di compattezza per gli spazi di Lebesgue\footnote{Così come il Teorema di Ascoli--Arzelà per le funzioni continue su un compatto. In effetti la dimostrazione richiede proprio l'uso del Teorema sopra citato.}.
\begin{notazione}\index[simboli]{$\tau_{y}(f)$}
	Per ogni $x \in \R^n$, per ogni $y \in \R^n$, indichiamo con 
	\[ \tau_{y}(f)(x) = f(x+y) \]
	la traslata di vettore $y$ di una qualsiasi applicazione $f$.
\end{notazione}

\begin{teorema}\emph{\textbf{(di Kolmogorov--M. Riesz--Fréchet)}}\index[analitico]{teorema!di Kolmogorov--M. Riesz--Fréchet}
	Siano $1 \le p < \infty$ ed $\mathcal{F}$ un sottoinsieme di $L^p(\R^{n})$. Se $\mathcal{F}$ è limitato e per ogni $\varepsilon > 0$, esiste $\delta > 0$ tale che 
	\[ \forall\; f \in \mathcal{F},\forall\; y \in \R^n: \abs{y} < \delta \Longrightarrow \norm{\tau_{y}(f)-f}_{p} < \varepsilon \]
	allora per ogni sottoinsieme misurabile $\Omega$ di $\R^n$ di misura finita, $\overline{\mathcal{F}_{| \Omega}}$ è compatto in $L^p(\Omega)$.
\end{teorema}
\begin{proof}
	Omettiamo la dimostrazione.
\end{proof}

\begin{corollario}
	Siano $1 \le p < \infty$ ed $\mathcal{F}$ un sottoinsieme di $L^p(U)$. Se $\mathcal{L}^n(U) < +\infty$, $\mathcal{F}$ è limitato e per ogni $\varepsilon > 0$, esiste $\delta > 0$ tale che 
	\[ \forall\; f \in \mathcal{F},\forall\; y \in \R^n: \abs{y} < \delta \Longrightarrow \norm{\tau_{y}(f)-f}_{p} < \varepsilon \]
	allora $\overline{\mathcal{F}}$ è compatto in $L^p(U)$.
\end{corollario}
\begin{proof}
	Si tratta di una riscrittura del Teorema di Kolmogorov.
\end{proof}
\begin{corollario}\label{corollarioKolmogorov}
	Siano $1 \le p < \infty$ ed $\mathcal{F}$ un sottoinsieme di $L^p(\R^n)$. Se $\mathcal{F}$ è limitato, per ogni $\varepsilon > 0$, esiste $\delta > 0$ tale che 
	\[ \forall\; f \in \mathcal{F},\forall\; y \in \R^n: \abs{y} < \delta \Longrightarrow \norm{\tau_{y}(f)-f}_{p} < \varepsilon \]
	ed esiste un sottoinsieme limitato $U$ di $\R^n$ misurabile tale che 
	\[ \forall\; f \in \mathcal{F}: \norm{f}_{L^p(\R^n\setminus U)} < \varepsilon \] 
	allora $\overline{\mathcal{F}}$ è compatto in $L^p(\R^n)$.
\end{corollario}
\begin{proof}
	Omettiamo la dimostrazione.
\end{proof}
\begin{osservazione}
	Il viceversa del Corollario \eqref{corollarioKolmogorov} è anch'esso vero. Pertanto abbiamo una caratterizzazione piena dei sottoinsiemi compatti di $L^p(\R^n)$.
\end{osservazione}

Passiamo quindi ad alcune considerazioni sulla convergenza.
\begin{proposizione}\label{convergenzaLq}
	Supponiamo che $\mathcal{L}^n(U)< +\infty$ e $1 < p < \infty$.  Sia $(f_{h})$ una successione limitata in $L^p(U)$. Se $f_{h} \to f$ \emph{q.o.} in $U$, allora per ogni $q < p$, $f_{h} \to f$ in $L^q(U)$.
\end{proposizione}
\begin{proof}
	Dalla limitatezza di $(f_{h})$ innanzitutto discende l'esistenza di $K > 0$ tale che per ogni $h \in \N$
	\[ \norm{f_{h}}_{p}\le K. \]
	Ora, per il Lemma di Fatou
	\[ \int_{U}\abs{f(x)}^p d\mathcal{L}^n(x) \le \underset{h}{\liminf} \int_{U} \abs{f_h(x)}^pd\mathcal{L}^n(x) = \underset{h}{\liminf} \norm{f_{h}}_{p}^p \le K^p \]
	da cui $f \in L^p(U)$. Mostriamo che per ogni $q \in \left[  1,p \right[$
	\[ \lim_{h} \int_{U}\abs{f_h(x)-f(x)}^qd\mathcal{L}^n(x) = 0. \]
	
	Fissiamo $M > 0 $ e prendiamo 
	\[ U_{h,M} = \Set{ x \in U: \abs{f_h(x) - f(x)} < M}. \]
	Ora,
	\[ \int_{U}\abs{f_h(x)-f(x)}^qd\mathcal{L}^n(x) =  \int_{U_{h,M}}\abs{f_h(x)-f(x)}^qd\mathcal{L}^n(x) + \int_{U\setminus U_{h,M}}\abs{f_h(x)-f(x)}^qd\mathcal{L}^n(x),\]
	e per il Teorema della convergenza dominata\footnote{Indichiamo con $g_h(x) = \abs{f_h(x)-f(x)}^q\chi_{U_{h,M}}$,
		\begin{align*}
			g_{h} &\to 0 \textup{ q.o. in } U & \abs{g_h} \le M^q \in L^1(U)
		\end{align*} 
		in quanto $\mathcal{L}^n(U) < +\infty$.}, a $M$ fissato,
	\[ \lim_{h} \int_{U_{h,M}}\abs{f_h(x)-f(x)}^qd\mathcal{L}^n(x) = 0.\]
	Invece, 
	\[ \begin{split}
		\int_{U\setminus U_{h,M}}\abs{f_h(x)-f(x)}^qd\mathcal{L}^n(x) &= \int_{U\setminus U_{h,M}}\abs{f_h(x)-f(x)}^p \abs{f_h(x)-f(x)}^{q-p}d\mathcal{L}^n(x) \le \\
		&\le \frac{1}{M^{p-q}}\int_{U\setminus U_{h,M}}\abs{f_h(x)-f(x)}^pd\mathcal{L}^n(x) \le \frac{C}{M^{p-q}}.
	\end{split} \]
	Allora, 
	\[ \int_{U}\abs{f_h(x)-f(x)}^qd\mathcal{L}^n(x) \le \int_{U_{h,M}}\abs{f_h(x)-f(x)}^qd\mathcal{L}^n(x) + \frac{C}{M^{p-q}} \]
	e se $h \to +\infty$
	\[ \int_{U}\abs{f_h(x)-f(x)}^qd\mathcal{L}^n(x) \le \frac{C}{M^{p-q}}. \]
	Se ora $M \to +\infty$ otteniamo la tesi.
\end{proof}

\begin{definizione}
	Siano $\mu$ una misura esterna su $\R^n$ ed $E \subseteq \R^n$ misurabile. Siano $(f_{h})$ una successione di applicazioni misurabili da $E$ in $\R$ ed $f: E \to \R$ misurabile. Diciamo che $f_h$ \emph{converge quasi uniformemente} a $f$ in $E$\index[analitico]{convergenza!quasi uniforme} se per ogni $\varepsilon>0$ esiste $E_{\varepsilon}\subseteq E$ tale che 
\begin{align*}
	\mu(E\setminus E_{\varepsilon})&< \varepsilon,  & \lim_{h}\norm{f_{h}-f}_{L^{\infty}(E_{\varepsilon},\mu)} &=0.
\end{align*} 
\end{definizione}
Chiaramente, la convergenza quasi uniforme implica la convergenza q.o. in $E$. Nel caso di misura finita, il seguente risultato realizza l'implicazione contraria.

\begin{teorema}\emph{\textbf{(di Severini--Egorov)}}\index[analitico]{teorema! di Severini--Egorov}
	Siano $\mu$ una misura esterna su $\R^n$ ed $E \subseteq \R^n$ misurabile tale che $\mu(E)< +\infty$. Sia $(f_{h})$ una successione di applicazioni misurabili da $E$ in $\R$. Se $f_{h} \to f$ \emph{q.o.} in $E$, allora $f_h \to f$ quasi uniformemente in $E$.
\end{teorema}
\begin{proof}
Sia $N \subseteq E$ tale che $\mu(N) = 0$ e $f_n \to f$ puntualmente in $E \setminus N$. Per ogni $k \in \N$ consideriamo la successione crescente $(A^{(k)}_m)$ di sottoinsiemi di $E$ tali che
\[ A^{(k)}_m = \bigcap_{n \ge m}\left\lbrace x \in E \setminus N: \abs{f_n(x)-f(x)}\le \frac{1}{k+1}\right\rbrace.  \]
Chiaramente, fissato $k \in \N$
\[ \bigcup_{m \in \N} A^{(k)}_m = E \setminus N.\]
Siccome siamo in misura finita, a $k \in \N$ fissato, $\mu(A^{(k)}_m) \to \mu(E)$, per ogni $\varepsilon >0$ esiste $m_k$ tale che $\mu(E \setminus  A^{(k)}_{m_k}) \le \frac{\varepsilon}{2^{k+2}}$ e 
\[ \forall \; x \in  A^{(k)}_{m_k}, \forall \; n \ge m_k: \abs{f_n(x)-f(x)} \le \frac{1}{k+1}.\]
Posto
\[ E_\varepsilon = \bigcap_{k \in \N} A^{(k)}_{m_k}, \]
allora 
\[ \lim_{h}\norm{f_{h}-f}_{L^{\infty}(E_{\varepsilon},\mu)} =0 \]
e 
\[ \mu(E \setminus E_\varepsilon) \le \sum_{k=0}^{\infty}\frac{\varepsilon}{2^{k+2}} = \frac{\varepsilon}{2} <\varepsilon,\]
da cui la tesi.
\end{proof}
In generale, se togliamo l'ipotesi di misura finita, tuttavia, il risultato non è più vero.
\begin{esempio}
	Sia $E = \R^+$. Allora $\mathcal{L}^1(E) = +\infty$. Consideriamo 
	\[ f_h(x) = \chi_{\left] h,h+1\right[ }(x). \]
	Evidentemente $f_h(x) \to 0$. Supponiamo per assurdo valga il Teorema di Severini--Egorov, allora per ogni $\varepsilon > 0$, esiste $E_{\varepsilon} \subseteq \R^+$ tale che $\mathcal{L}^1(\R^+\setminus E_{\varepsilon}) < \varepsilon$ e 
	\[ \lim_{h} \norm{f_h}_{L^{\infty}(E_{\varepsilon})} = 0. \]
	Allora, definitivamente in $h$, 
	\[ E_{\varepsilon} \cap \left] h,h+1\right[ = \emptyset, \]
	quindi $E_{\varepsilon}$ limitato. Allora $\mathcal{L}^1(\R^+\setminus E_{\varepsilon}) = +\infty$, assurdo.
\end{esempio}


\begin{osservazione}
Mediante il Teorema di Severini--Egorov, possiamo costruire una diversa dimostrazione della Proposizione \eqref{convergenzaLq}. 

Per il Teorema di Severini--Egorov, per ogni $\varepsilon > 0$, esiste $U_{\varepsilon}$ tale che $\mathcal{L}^n(U\setminus U_{\varepsilon})<\varepsilon$ e 
\[ \lim_{h}\norm{f_{h}-f}_{L^{\infty}(U_{\varepsilon})} =0. \]
Ora, per la disuguaglianza di Holder,
\[\begin{split}
	 \int_{U}\abs{f_h(x)-f(x)}^qd\mathcal{L}^n(x) &=   \int_{U_{\varepsilon}}\abs{f_h(x)-f(x)}^qd\mathcal{L}^n(x) + \int_{U \setminus U_{\varepsilon}}\abs{f_h(x)-f(x)}^qd\mathcal{L}^n(x) \le \\
	 &\le \norm{f_h-f}_{L^\infty(U_{\varepsilon})}^q\mathcal{L}^n(U_{\varepsilon}) + \\
	 & + \mathcal{L}^n(U\setminus U_{\varepsilon}))^{\frac{p-q}{p}}\left( \int_{U \setminus U_{\varepsilon}} \abs{f_h(x)-f(x)}^pd\mathcal{L}^n(x)\right) ^{\frac{q}{p}} < \\
	 & < \norm{f_h-f}_{L^\infty(U_{\varepsilon})}^q\mathcal{L}^n(U) + \widehat{C}\varepsilon^{\frac{p-q}{p}}
\end{split} \]
Ora, passando prima al limite per $h \to +\infty$ e poi per $\varepsilon \to 0$ si ottiene la tesi.
\end{osservazione}

\begin{definizione}\index[analitico]{convergenza! in misura}
	Siano $\mu$ una misura esterna su $\R^n$ ed $E \subseteq \R^n$ misurabile. Siano $(f_{h})$ una successione di applicazioni misurabili da $E$ in $\R$ ed $f: E \to \R$ misurabile. Diciamo che $f_h$ \emph{converge in misura} a $f$, e lo indichiamo con $f_h \overset{\mu}{\to} f$\index[simboli]{$\overset{\mu}{\to}$} se, per ogni $\varepsilon > 0 $, detto 
\[ E_{\varepsilon} = \left\lbrace  x \in U: \abs{f_h(x) - f(x)} > \varepsilon\right\rbrace, \]
risulta
\[ \lim_{h} \mu(E_{\varepsilon}) = 0. \]
\end{definizione}
\begin{proposizione}
	Siano $\mu$ una misura esterna su $\R^n$ ed $E \subseteq \R^n$ misurabile. Siano $(f_{h})$ una successione di applicazioni misurabili da $E$ in $\R$ ed $f: E \to \R$ misurabile. Se $f_h \overset{\mu}{\to} f$, allora esiste $(f_{h_k})$ tale che $f_{h_k}\to f$ \emph{q.o.} in $E$.
\end{proposizione}
\begin{proof}
Per ogni $k \in \N$, esiste $n_k$ tale che 
\[ \mu\left( \left\lbrace x \in E: \abs{f_n(x)-f(x)}>\frac{1}{k}\right\rbrace \right) \le \frac{1}{k^2}. \]
Siano
\[ A_k = \left\lbrace x \in E: \abs{f_n(x)-f(x)}>\frac{1}{k}\right\rbrace   \]
e 
\[ A = \bigcap_{m \in \N} \bigcup_{k \ge m} A_k. \]
La tesi segue dal fatto che $\mu(A) \le \displaystyle \sum_{k=m}^{\infty} \frac{1}{k^2} \to 0$ e $f_{n_k}(x) \to f(x)$ in  $E\setminus A$ in quanto 
\[ E \setminus A = \bigcup_{m \in N} \bigcap_{k \ge m} \left\lbrace x \in E: \abs{f_{n_k}(x)-f(x)} \le \frac{1}{k} \right\rbrace. \qedhere\]
\end{proof}

\begin{proposizione}
	Siano $\mu$ una misura esterna su $\R^n$ ed $E \subseteq \R^n$ misurabile tale che $\mathcal{L}^n(E)<+\infty$. Siano $(f_{h})$ una successione di applicazioni misurabili da $E$ in $\R$ ed $f: E \to \R$ misurabile. Se $f_h \to f$ \emph{q.o.} in $E$, allora $f_h \overset{\mu}{\to} f$.
\end{proposizione}
\begin{proof}
Fissiamo $\varepsilon >0$. Per il Teorema di Severini--Egorov, per ogni $\delta >0$ esiste $E_{\delta}\subseteq E$ tale che 
\begin{align*}
	\mu(E\setminus E_{\delta})&<\delta,  & \lim_{h}\norm{f_{h}-f}_{L^{\infty}(E_{\delta},\mu)} &=0.
\end{align*} 
In particolare, esiste $n_\varepsilon$ tale che per ogni $h \ge n_\varepsilon$ e per ogni $x \in E_\delta$
\[ \abs{f_h(x)-f(x)}< \varepsilon. \]
Allora, per $h \ge n_\varepsilon$
\begin{align*}
 \mu\left( \left\lbrace x \in E: \abs{f_h(x)-f(x)}>\varepsilon\right\rbrace \right) &\le \mu\left( \left\lbrace x \in E_\delta: \abs{f_h(x)-f(x)}>\varepsilon \right\rbrace \right) + \\
 &+\mu\left( \left\lbrace x \in E\setminus E_\delta: \abs{f_h(x)-f(x)}>\varepsilon \right\rbrace \right) \le \\
 &\le \mu(\emptyset) + \mu(E \setminus E_\delta) < \delta.\qedhere
\end{align*}
\end{proof}

\begin{proposizione}
Siano $1\le p < \infty$ e $(f_h)$ una successione in $L^p(U,\mu)$. Se 
\[ \lim_{h}\norm{f_h-f}_{p} = 0 \]
allora $f_h \overset{\mu}{\to} f$.
\end{proposizione}
\begin{proof}
Siccome
	\[ \int_{U} \abs{f_h - f}^p d\mu \ge \int_{E_{\varepsilon}}\abs{f_h - f}^p d\mu > \varepsilon^p\mu(E_{\varepsilon}),\]
fissato $\varepsilon$, se $h \to +\infty$ risulta che $\mu(E_{\varepsilon}) \to 0$.
\end{proof}
\begin{esempio}
Sia $1<p<\infty$. Consideriamo la successione definita da $f_h: \R^+ \to \R$ tale che
\[ f_h(x) = h\chi_{\left] 0,\frac{1}{h}\right[ }(x). \]
Evidentemente $f_h \to 0$ \emph{q.o.} ma
\[ \int_{\R^+}\abs{f_h}^pd\mathcal{L}^1 = h^{p-1} \]
per cui $f_h$ non converge in $L^p(\R^+)$. Tuttavia
\[ \mathcal{L}^1 \left( \left\lbrace x \in \R^+: f_h(x) > \varepsilon\right\rbrace \right)  = \mathcal{L}^1\left( \left\lbrace x \in \R^+: h\chi_{\left] 0,\frac{1}{h}\right[ }(x) > \varepsilon\right\rbrace \right)  \le \mathcal{L}^1\left( \left] 0,\frac{1}{h}\right[\right) = \frac{1}{h}\]
per cui  $f_h \overset{\mathcal{L}^1}{\to} 0$.
\end{esempio}

Vediamo ora di completare la panoramica sulle proprietà degli spazi di Lebesgue.
\begin{teorema}\emph{\textbf{(Disuguaglianze di Clarkson)}}\index[analitico]{disuguaglianza!di Clarkson} Sia $p \in \left[ 1, \infty\right[ $. Siano $E \subseteq \R^{n} $ misurabile ed $f,g \in L^{p}(E)$. Valgono i seguenti fatti:
	\begin{enumerate}[$(a)$]
		\item se $p \ge 2$, 
		\[ \norm*{\frac{f+g}{2}}^{p}_{p} + \norm*{\frac{f-g}{2}}^{p}_{p} \le \frac{1}{2}\left( \norm{f}^{p}_{p} +  \norm{g}^{p}_{p}\right) \qquad \textup{(Disuguaglianza di Clarkson I)}\]
		\item se $1 \le p < 2$, 
		\[ \norm*{\frac{f+g}{2}}^{p'}_{p} + \norm*{\frac{f-g}{2}}^{p'}_{p} \le \left( \frac{1}{2}\norm{f}^{p}_{p} +  \frac{1}{2}\norm{g}^{p}_{p}\right)^{\frac{1}{p-1}} \quad \textup{(Disuguaglianza di Clarkson II)}\]
	\end{enumerate}
\end{teorema}
\begin{proof} \hfill
	\par\noindent
	$(a)$ Mostriamo innanzitutto che per ogni $ x,y \ge 0$
	\begin{equation}\label{eq1}
		x^{p} + y^{p} \le (x^{2} + y^2)^{\frac{p}{2}}.
	\end{equation} 
		Se $y = 0$, la disuguaglianza è evidentemente vera. Supponiamo quindi $y \ne 0$ e studiamo la veridicità di 
	\[ \left(\frac{x}{y}\right)^p + 1 \le \left( \left( \frac{x}{y}\right) ^2 + 1\right) ^{\frac{p}{2}}  \]
	o, equivalentemente (introducendo $z = \frac{x}{y}$), di 
	\[ z^p +1 \le (z^2 + 1)^{\frac{p}{2}}. \]
	Posto $\varphi(z) = (z^2 + 1)^{\frac{p}{2}} - z^p - 1 $, è sufficiente mostrare che $\varphi(z) \ge 0$ per $z \ge 0$. Grazie al fatto che $p \ge 2$, 
	\[ \varphi'(z) = p\left( z(z^2 + 1)^{\frac{p-2}{2}} - z^{p-1}\right) \ge 0. \]
	La funzione $\varphi$ è dunque crescente e, dato che $\varphi(0) = 0$, ne deduciamo che $\varphi(z) \ge 0$ per ogni $z \ge 0$.
	
	Siano ora $x = \abs*{\frac{a+b}{2}}, y = \abs*{\frac{a-b}{2}}$, con $a,b \in \R$. Allora applicando \eqref{eq1}, otteniamo 
	\[ \abs*{\frac{a+b}{2}}^{p} + \abs*{\frac{a-b}{2}}^{p} \le \left( \abs*{\frac{a+b}{2}}^{2} + \abs*{\frac{a-b}{2}}^2\right) ^{\frac{p}{2}} = \left( \frac{1}{2}a^{2} + \frac{1}{2}b^{2}\right)^{\frac{p}{2}}\le \frac{1}{2} \left( a^{p} + b^{p}\right), \]
	dove l'ultima maggiorazione è giustificata dalla convessità dell'applicazione $t \mapsto t^{\frac{p}{2}}$ su $t\ge 0$, essendo $p\ge2$.
	
	Non resta che porre $a = f(x)$ e $b= g(x)$ e integrare ambo i membri per ottenere la disuguaglianza cercata.
	\par\noindent
	$(b)$ Omettiamo la dimostrazione.\qedhere
\end{proof}
\begin{teorema}
	Siano $p \in \left[ 2, \infty \right[ $ ed $E \subseteq \R^{n} $ misurabile. Allora $L^{p}(E)$ è riflessivo.
\end{teorema}
\begin{proof}
	Mostriamo che $L^{p}(E)$ è uniformemente convesso; fatto questo, la tesi discende da \eqref{uniconvrif}. Siano $\varepsilon > 0$ e $f,g \in L^{p}(E)$ tali che $\norm{f}_p,\norm{g}_p \leq 1$ e $\norm{f-g}_{p} > \varepsilon$. Per la prima disuguaglianza di Clarkson,
\[ \norm*{\frac{f+g}{2}}^{p}_{p} \le \frac{1}{2}\left( \norm{f}^{p}_{p} +  \norm{g}^{p}_{p}\right) -  \frac{1}{2^{p}}\norm*{f-g}^{p}_{p} < 1 - \frac{\varepsilon^{p}}{2^{p}}. \]
Allora
\[ \norm*{\frac{f+g}{2}}_{p} < \left(1 - \frac{\varepsilon^{p}}{2^{p}}\right) ^{\frac{1}{p}} = 1 - 1 + \left(1 - \frac{\varepsilon^{p}}{2^{p}}\right) ^{\frac{1}{p}} = 1 - \left( 1 - \left(1 - \frac{\varepsilon^{p}}{2^{p}}\right) ^{\frac{1}{p}}\right).\] 
Ponendo $\delta = 1 - \left(1 - \frac{\varepsilon^{p}}{2^{p}}\right) ^{\frac{1}{p}} > 0$ e osservando che non dipende da $f,g$, si conclude la dimostrazione.
\end{proof}

	Per completare il quadro e dimostrare che anche gli spazi $L^p$ con $1 < p < 2$ sono riflessivi, operiamo ora per dualità, come mostrato nel seguente Teorema.

\begin{teorema}\label{1,2rif}
	Siano $p \in \left]  1, 2 \right[ $ ed $E \subseteq \R^{n} $ misurabile. Allora $L^{p}(E)$ è riflessivo. 
\end{teorema}

\begin{proof}
	Innanzitutto, se $p \in \left]  1, 2 \right[ $ allora $p' \in \left]  2, \infty \right[$, sicché $L^{p'}(E)$ è riflessivo. Di conseguenza, anche $\left( L^{p'}(E)\right)'$ è riflessivo per \eqref{teo:reflexive_iff}. Consideriamo poi l'applicazione lineare $T : L^{p}(E) \to \left( L^{p'}(E)\right)' $ tale che per ogni $u \in L^{p}(E) $, per ogni $f \in L^{p'}(E)$,
	\[ \langle Tu,f\rangle = \int_{E}^{} f(x)u(x)d\mathcal{L}^{n}(x).\]
	Risulta, dalla disuguaglianza di H\"older, che
	\[ \abs{\langle Tu,f\rangle} \le \norm{u}_{p}\norm{f}_{p'} \]
	da cui $\norm{Tu}_{\left( L^{p'}(E)\right)'} \le \norm{u}_{p}$. D'altro canto, se $\norm{u}_{p}=0$, allora anche $\norm{Tu}_{\left( L^{p'}(E)\right)'} = 0$; se invece $\norm{u}_{p}>0$, allora prendendo 
	\[ f_{0}(x) = \frac{\abs{u(x)}^{p-2}u(x)}{\norm{u}^{p-1}_{p}}, \]
	risulta $\abs{\langle Tu,f_{0}\rangle} = \norm{u}_{p}$, da cui segue che
	\[ \norm{Tu}_{\left( L^{p'}(E)\right)'} = \sup\left\lbrace \abs{\langle Tu,f\rangle} \; :\; f \in L^{p'}(E) \;,\; \norm{f}_{p'} \le 1\right\rbrace \ge \abs{\langle Tu,f_{0}\rangle} = \norm{u}_{p}. \]
	Se ne deduce che $\norm{Tu}_{\left( L^{p'}(E)\right)'} = \norm{u}_{p} $, ovvero $T$ è un'isometria. %In particolare, è iniettiva.
	
	Dal fatto che $T$ è un'isometria, si verifica facilmente che $T(L^{p}(E)) \le \left( L^{p'}(E)\right)'$ chiuso. Allora, per \eqref{teo:refl_subspace}, $T(L^{p}(E))$ è riflessivo. 
\end{proof}
\begin{osservazione}
	Tramite la seconda disuguaglianza di Clarkson si può ottenere una dimostrazione alternativa del Teorema precedente.
\end{osservazione}

Vediamo ora di andare a studiare la rappresentazione di forme lineari e continue di $L^p$.
\begin{lemma}\emph{\textbf{(di caratterizzazione topologica della densità)}}\index[analitico]{lemma!di caratterizzazione topologica della densità}
	Siano $X$ uno spazio di Banach su $\K$ e $Y \le X$. I seguenti fatti sono equivalenti:
	\begin{enumerate}[$(a)$]
		\item $\overline{Y} = X$, cioè $Y$ è denso in $X$, 
		\item per ogni $\varphi \in X'$
		\[  \left( \langle \varphi, y \rangle = 0 \textup{ per ogni } y \in Y\right) \Longrightarrow \varphi = 0.\]
	\end{enumerate}
\end{lemma}

\begin{proof}\hfill
	\par\noindent
	$(a) \Longrightarrow (b)$ Sia $x \in X$. Per densità esiste $(y_{j})$ in $Y$ tale che $y_{j} \to x$. Per ogni $\varphi \in X'$, per continuità,
	\[ \langle \varphi, x \rangle = \langle \varphi, \lim\limits_{j} y_{j} \rangle = \lim\limits_{j} \langle \varphi,  y_{j} \rangle =  \lim\limits_{j} 0 = 0.\]
	\par\noindent
	$(b) \Longrightarrow (a)$ Omettiamo la dimostrazione.
\end{proof}

\begin{teorema}
	Siano $p \in \left]  1, \infty \right[ $ ed $E \subseteq \R^{n} $ misurabile. Allora $\left( L^{p}(E)\right)' \cong L^{p'}(E)$.
\end{teorema}

\begin{proof}
	Consideriamo l'applicazione lineare $T : L^{p'}(E) \to \left( L^{p}(E)\right)'$ tale che per ogni $u \in L^{p'}(E)$, per ogni $f \in L^{p}(E)$
	\[ \langle Tu,f\rangle = \int_{E}^{} f(x)u(x)d\mathcal{L}^{n}(x).\]
	Con passaggi analoghi a quelli eseguiti nel Teorema \eqref{1,2rif}, possiamo verificare che $T$ è un'isometria e che $T(L^{p'}(E)) \le \left( L^{p}(E)\right)' $ chiuso. Mostriamo che $T$ è anche suriettiva.
	
	Sia $\varphi \in \left( L^{p}(E)\right)'' $ e supponiamo che per ogni $f \in L^{p'}(E)$ valga
	\[ \langle \varphi, Tf \rangle = 0. \]
	Per riflessività di $\left( L^{p}(E)\right)''$ esiste $g \in L^{p}(E)$ tale che $\varphi = Jg$ e quindi, per definizione di iniezione canonica,
	\[ \langle \varphi, Tf \rangle = \langle Jg, Tf \rangle = \langle Tf, g \rangle = \int_{E}^{}f(x)g(x)d\mathcal{L}^{n}(x) = 0 \]
	per ogni $f \in L^{p'}(E)$. In particolare 
	\[ \int_{E}^{}f(x)g(x)d\mathcal{L}^{n}(x) = 0 \]
	per ogni $f \in C^{\infty}_{c}(E)$, da cui $g = 0$ q.o. Ciò significa che anche $\varphi = 0$ e, per il Lemma di caratterizzazione topologica della densità, $T(L^{p'}(E))$ è denso in $X$. Dal fatto che $T(L^{p'}(E))$ è chiuso, si conclude la suriettività, ovvero che
	\[ T(L^{p'}(E)) = \left( L^{p}(E)\right)', \]
	da cui la tesi.
\end{proof}

\begin{teorema}\label{dualediL1}
	Sia $E \subseteq \R^{n} $ misurabile. Allora $\left( L^{1}(E)\right)' \cong L^{\infty}(E)$.
\end{teorema}

\begin{proof}
	Sia $\varphi : L^{1}(E) \to \R$ lineare e continua. Mostriamo che esiste un'unica $u \in L^{\infty}(E)$ tale che 
	\[ \langle \varphi, f \rangle = \int_{E}f(x)u(x)d\mathcal{L}^{n}(x). \]
	Sia $(E_{h})$ una successione di sottoinsiemi misurabili di $E$ tali che 
	\[ E = \bigcup_{h=1}^{\infty} E_{h}, \]
	$\mathcal{L}^{n}(E_{h}) < +\infty$ e $E_{0} = \emptyset$ e costruiamo una funzione $\vartheta \in L^{2}(E)$ tale che per ogni $h \in \N\setminus\left\lbrace 0 \right\rbrace $ esiste $\varepsilon_{h} > 0$ per cui $\vartheta|_{E_{h}} \ge \varepsilon_{h}$. Per fare ciò è sufficiente porre
	\[ \vartheta|_{E_{1}} = \alpha_{1}, \]
	\[ \vartheta|_{E_{i}\setminus E_{i-1}} = \alpha_{i} \textup{ per ogni } i >1,\]
	dove, per ogni $j \in \N\setminus\left\lbrace 0 \right\rbrace$,
	\[ \alpha_{j} \le \dfrac{\left( \mathcal{L}^{n}\left( E_{j}\setminus E_{j-1}\right)\right) ^{-\frac{1}{2}}}{j^{\beta}} \]
	con $\beta > \frac{1}{2}$, di modo che 
	\begin{equation*}
		\begin{split}
			\int_{E} \vartheta^{2}(x)d\mathcal{L}^{n}(x)&=  \int_{\bigcup_{j=1}^{\infty} E_{j}} \vartheta^{2}(x)d\mathcal{L}^{n}(x) \\ 
			&=\sum_{j = 1}^{\infty}\int_{E_{j}\setminus E_{j-1}} \vartheta^{2}(x)d\mathcal{L}^{n}(x) \\
			&=  \sum_{j = 1}^{\infty} \alpha^{2}_{j} \mathcal{L}^{n}\left( E_{j}\setminus E_{j-1}\right) \\
			&= \sum_{j=1}^\infty \frac{1}{j^{2\beta}} < +\infty.
		\end{split}
	\end{equation*}
	Consideriamo ora l'applicazione 
	\[ \begin{cases}
		L^{2}(E) \to \R \\
		f \longmapsto \langle \varphi, \vartheta f \rangle
	\end{cases}, \]
	ben definita in quanto $\vartheta, f \in L^{2}(E)$. Per il Teorema di F. Riesz esiste $v \in L^{2}(E)$ tale che per ogni $f \in L^{2}(E)$
	\begin{equation}\label{Riesz}
		\langle \varphi,\vartheta f \rangle = \int_{E}f(x)v(x)d\mathcal{L}^{n}(x).
	\end{equation}
	Sia $u = \frac{v}{\vartheta}$. Effettuando la scelta $f = \chi_{h}\frac{g}{\vartheta}$ in \eqref{Riesz} con $g \in L^{\infty}(E)$ e $\chi_{h} = \chi_{E_{h}}$, e osservando che la scelta è effettivamente ammissibile, risulta che per ogni $h \in \N$
	\begin{equation}\label{condizioneimportante}
		\forall \; g \in L^{\infty}(E): 	\langle \varphi,\chi_{h} g \rangle = \int_{E}\chi_{h}(x)g(x)u(x)d\mathcal{L}^{n}(x).
	\end{equation}
	Mostriamo che $u \in L^{\infty}(E) $. Consideriamo $C > \norm{\varphi}_{(L^{1}(E))'}$, 
	\[ A = \left\lbrace x \in E \,:\; \abs{u(x)} \ge C \right\rbrace  \]
	e, scegliendo $g = \chi_{A}\frac{u}{\abs{u}} \in L^{\infty}(E)$ in \eqref{condizioneimportante}, osserviamo che
	\[
	\begin{split}
		\langle \varphi, \chi_{h}\chi_{A}\frac{u}{\abs{u}}\rangle & = \int_{E}u(x)\chi_{h}(x)\chi_{A}(x)\frac{u(x)}{\abs{u(x)}}d\mathcal{L}^{n}(x) = \int_{E}\abs{u(x)}\chi_{A \cap E_{h}}(x)d\mathcal{L}^{n}(x) \\
		& \geq C\mathcal{L}^n(A \cap E_h).
	\end{split}
	\]
	D'altra parte,
	\[
	\begin{split}
		\langle \varphi, \chi_{h}\chi_{A}\frac{u}{\abs{u}}\rangle = \langle \varphi, \chi_{A \cap E_{h}}\frac{u}{\abs{u}}\rangle \le \norm{\varphi}_{(L^{1}(E))'}\norm*{\chi_{A \cap E_{h}}\frac{u}{\abs{u}}}_{1} = \norm{\varphi}_{(L^{1}(E))'}\mathcal{L}^{n}\left( A \cap E_{h}\right),
	\end{split}
	\]
	così che
	\[ C\mathcal{L}^{n}(A\cap E_{n}) \le \int_{E}\abs{u(x)}\chi_{A \cap E_{h}}d\mathcal{L}^{n}(x) \le \norm{\varphi}_{(L^{1}(E))'}\mathcal{L}^{n}\left( A \cap E_{h}\right).  \]
	Dato che $C > \norm{\varphi}_{(L^{1}(E))'}$, deduciamo che $\mathcal{L}^{n}\left( A \cap E_{h}\right) = 0$. Dall'arbitrarietà di $h$, risulta che $\mathcal{L}^{n}\left( A \right) = 0$. Quindi $\abs{u(x)} < C$ q.o. in $E$, ovvero $u \in L^{\infty}(E) $.\\
	
	Mostriamo infine che per ogni $f \in L^{1}(E)$
	\begin{equation}\label{eq:representation_L1Linfty}
		\langle \varphi, f \rangle = \int_{E}f(x)u(x)d\mathcal{L}^{n}(x).
	\end{equation}
	Per fare ciò consideriamo l'operatore $T : L^1(E) \to L^\infty(E)$ definito da
	\[ T_{h}f(x) = \begin{cases}
		-h & \textup{se } f(x) \le -h, \\
		f(x) & \textup{se } \abs{f(x)} \le h, \\
		h & \textup{se } f(x) \ge h,
	\end{cases} \]
	e comunemente detto \emph{troncatura\index[analitico]{troncatura} $h$-esima di $f$}. Sostituendo nella \eqref{condizioneimportante} otteniamo
	\begin{equation}\label{dalimitizzare}
		\langle \varphi,\chi_{h} T_{h}f \rangle = \int_{E}\chi_{h}(x)T_{h}f(x)u(x)d\mathcal{L}^{n}(x).
	\end{equation}
	A questo punto, per $h \to +\infty$, $\chi_{h}(x)T_{h}f(x) \to f(x)$ q.o. in $E$ e
	\[ \abs{\chi_{h}(x)T_{h}f(x)}\le \abs{T_{h}f(x)} \le \abs{f(x)} \in L^{1}(E), \]
	così che per il Teorema della convergenza dominata $\chi_{h} T_{h}f \to f$ in $L^{1}(E)$. Possiamo quindi passare al limite nella \eqref{dalimitizzare} e ottenere \eqref{eq:representation_L1Linfty}.
	
	Supponiamo ora esista un'altra $\tilde{u} \in L^\infty(E)$ tale che per ogni $f \in L^{1}(E)$
	\[ \langle \varphi, f \rangle = \int_{E}f(x)\tilde{u}(x)d\mathcal{L}^{n}(x). \]
	Allora per ogni $f \in L^{1}(E)$ vale
	\[ \int_{E}f(x)\tilde{u}(x)d\mathcal{L}^{n}(x) = \int_{E}f(x)u(x)d\mathcal{L}^{n}(x),\]
	ovvero
	\[ \int_{E}(u(x)- \tilde{u}(x))f(x)d\mathcal{L}^{n}(x) =0.\]
	In particolare, la condizione sopra scritta vale per ogni $f \in C^{\infty}_{c}(E)$, così che $u - \tilde{u} = 0 $ q.o in $E$, da cui l'unicità.
\end{proof}

\begin{corollario}
	Siano $E \subseteq \R^{n}$ misurabile, $\varphi : L^{1}(E) \to \R$ lineare e continua e sia $u \in L^{\infty}(E)$ la rappresentazione di $\varphi$ nel senso del Teorema \eqref{dualediL1}. Allora 
	\[ 	\norm{\varphi}_{(L^{1}(E))'} = \norm{u}_{\infty}. \]
\end{corollario}

\begin{proof}
	Innanzitutto, per ogni $f \in L^{1}(E)$
	\[ \abs{\langle \varphi, f\rangle} = \int_{E}f(x)u(x)d\mathcal{L}^{n}(x) \le \norm{u}_{\infty}\norm{f}_{1} \]
	da cui $\norm{\varphi}_{(L^{1}(E))'}\le \norm{u}_{\infty} $. 
	
	Viceversa, dato $C > \norm{\varphi}_{(L^{1}(E))'} $, con passaggi analoghi a quelli del Teorema \eqref{dualediL1}, possiamo affermare che $\norm{u}_{\infty} \le C$. Scelto quindi $C = \norm{\varphi}_{(L^{1}(E))'} + \frac{1}{h}$, con $h \in \N\setminus \left\lbrace 0\right\rbrace $, passando al limite per $h \to +\infty$ otteniamo $\norm{u}_{\infty} \le \norm{\varphi}_{(L^{1}(E))'}$, da cui la tesi.
\end{proof}

\begin{proposizione}
	Sia $E \subseteq \R^{n} $ misurabile. Allora $L^{1}(E) \hookrightarrow (L^{\infty}(E))'$.
\end{proposizione}

\begin{proof}
	Sia $f \in L^{1}(E)$. Consideriamo l'applicazione lineare  $\psi : L^{\infty}(E) \to \R$ definita come
	\[ \psi(g) = \int_{E}^{}f(x)g(x)d\mathcal{L}^{n}(x).\]
	Si vede facilmente che
	\[ \abs{\psi(g)} \le \norm{f}_{1}\norm{g}_{\infty},\]
	così che $\psi$ è continua e di conseguenza $\psi \in (L^{\infty}(E))'$. L'identificazione è quindi fornita dall'applicazione $f \mapsto \psi$. Tuttavia, essa non è suriettiva.
	
	Infatti, si consideri $\xi \in E$ e si definisca l'applicazione $T : (C_{b}(E), \norm{\cdot}_{\infty}) \to \R$ come
	\[ \langle T, f \rangle = f(\xi). \]
	Essa è ben definita, lineare e tale che
	\[ \abs{\langle T, f \rangle} = \abs{f(\xi)} \le \max_E \abs{f} = \norm{f}_{\infty}, \]
	quindi continua. Siccome $C_{b}(E) \le L^{\infty}(E)$, per il Teorema di Hahn -- Banach (forma analitica) esiste $\mathfrak{T}: L^{\infty}(E) \to \R$ lineare e continua tale che 
	\[ \mathfrak{T}|_{C_{b}(E)} = T. \]
	Evidentemente $\mathfrak{T} \in (L^{\infty}(E))'$. Se, per assurdo, esistesse $\tilde{g} \in L^{1}(E)$ tale che $\mathfrak{T} = \psi(\tilde{g})$, allora si avrebbe
	\[ \Braket{\mathfrak{T}, f} = \int_{E}^{}f(x)\tilde{g}(x)d\mathcal{L}^{n}(x) \]
	per ogni $f \in L^{\infty}(E)$. Se $f \in C^{\infty}_{c}(E\setminus \left\lbrace \xi\right\rbrace )$, allora $f \in C_{b}(E)$ e $f(\xi) = 0$; ne segue che
	\[ \int_{E} f(x)\tilde{g}(x)d\mathcal{L}^{n}(x) = \langle T,f \rangle = f(\xi) = 0, \]
	da cui $\tilde{g}(x) = 0$ per q.o. $x \in E\setminus\left\lbrace \xi\right\rbrace $, ovvero $\tilde{g}(x) = 0$ per q.o. $x \in E$. Allora $\mathfrak{T} = 0$, ma $\mathfrak{T}|_{C_{b}(E)} = T \ne 0$, da cui l'assurdo. Pertanto l'inclusione è stretta, ovvero 
	\[ \psi(L^{1}(E)) \subsetneq (L^{\infty}(E))'. \qedhere\]
\end{proof}

\section{Convergenza debole in spazi di Banach}
\begin{definizione}\index[analitico]{convergenza! debole}
Siano $X$ uno spazio di Banach su $\K$ e $(x_h)$ una successione in $X$. Diciamo che \emph{$x_h$ converge debolmente a $x \in X$}, e lo indichiamo con $x_{h} \rightharpoonup x$, se per ogni $\varphi \in X'$
\[ \lim_{h} \braket{\varphi, x_h} = \braket{\varphi, x} \]
\end{definizione}

\begin{proposizione}
Siano $X$ uno spazio di Banach su $\K$ e $(x_h)$ una successione in $X$ ed $x \in X$ tale che $x_h \rightharpoonup x$ in $X$. Allora
\[ \norm{x} \le \underset{h}{\textup{lim inf}} \norm{x_h}. \]
\end{proposizione}
\begin{proof}
Innanzitutto, 
\[ \norm{x} = \sup \Set{\abs{\braket{\varphi,x}} :  \varphi \in X', \norm{\varphi}\le 1}. \]
Inoltre,
\[ \abs{\braket{\varphi,x_h}} \le \norm{\varphi}\norm{x_h} \]
e passando al lim inf, sfruttando la continuità
\[ \abs{\braket{\varphi, x}} = \underset{h}{\textup{lim inf}}\abs{\braket{\varphi,x_h}} \le \underset{h}{\textup{lim inf}}\norm{\varphi}\norm{x_h} \]
e passando al sup per $\norm{\varphi} \le 1$ otteniamo la tesi.
\end{proof}

\begin{proposizione}
Siano $X$ uno spazio di Banach su $\K$ e $(x_h)$ una successione in $X$. Se esiste $x \in X$ tale che $x_h \rightharpoonup x$ in X, allora $x$ è unico.
\end{proposizione}
\begin{proof}
	Omettiamo la dimostrazione.
\end{proof}
\begin{proposizione}
Siano $X$ uno spazio di Banach su $\K$ e $(x_h)$ una successione in $X$. Se esiste $x \in X$ tale che $x_h \to x$ in $X$, allora $x_h \rightharpoonup x$ in $X$.
\end{proposizione}
\begin{proof}
Omettiamo la dimostrazione.
\end{proof}
\begin{proposizione}
Siano $X$ uno spazio di Banach su $\K$ e $(x_h)$ una successione in $X$ limitata. Se
$x_h$ converge debolmente a qualche $x \in X$, allora $(x_{h})$ è limitata.
\end{proposizione}
\begin{proof}
	Omettiamo la dimostrazione.
\end{proof}
Il seguente è un analogo del Teorema di Bolzano--Weierstrass.
\begin{teorema}\emph{\textbf{(di Eberlein--Smulian)}}\index[analitico]{teorema!di Eberlein--Smulian}
	Sia $X$ uno spazio di Banach su $\K$. $X$ è riflessivo se e solo se ogni successione $(x_h)$ in X limitata ammette una sottosuccessione debolmente convergente.
\end{teorema}
\begin{proof}
	Omettiamo la dimostrazione.
\end{proof}
\begin{definizione}\index[analitico]{convergenza!debole*}
	Siano $X$ uno spazio di Banach su $\K$ e $(\varphi_h)$ una successione in $X'$. Diciamo che $\varphi_h$ \emph{converge debolmente*}, e lo indichiamo con $\varphi_h \overset{\ast}{\rightharpoonup} \varphi$\index[simboli]{$\overset{\ast}{\rightharpoonup}$} se per ogni $x \in X$
	\[ \braket{\varphi_h,x} \to \braket{\varphi, x}.\footnote{Praticamente è una convergenza puntuale.} \]
\end{definizione}
\begin{proposizione}\label{separabilitàe*}
	Siano $X$ uno spazio normato separabile su $\K$ e $(\varphi_h)$ una successione in $X'$ limitata. Allora esistono $\varphi \in X'$ ed una sottosuccessione $(\varphi_{h_{k}})$ tali che $\varphi_{h_{k}}\overset{\ast}{\rightharpoonup} \varphi $.
\end{proposizione}
\begin{proof}
	Omettiamo la dimostrazione.
\end{proof}

Veniamo al notevole caso degli spazi di Lebesgue.
\begin{osservazione}
Se $X = L^p(U)$, con $1 \le p < \infty$. Dal fatto che $(L^p(U))'\cong L^{p'}(U)$, per ogni $\varphi \in (L^p(U))'$, esiste $g \in L^{p'}(U)$ tale che 
\[ \braket{\varphi, f} = \int_{U}fgd\mathcal{L}^n. \]
Allora $f_h \rightharpoonup f$ se e solo se per ogni $g \in L^{p'}(U)$
\[ \lim_{h} \int_{U} f_hgd\mathcal{L}^n = \int_{U} fgd\mathcal{L}^n. \]
\end{osservazione}


\begin{proposizione}
Sia $p \in \left]   1,\infty\right[ $. Sia $(f_h)$ una successione in $L^p(U)$ tale che $f_h \to f$ \emph{q.o.} in $U$ e $f_h \rightharpoonup g$ in $L^p(U)$. Allora $g=f$ \emph{q.o.} in $U$.
\end{proposizione}
\begin{proof}
Innanzitutto, siccome $f_h \rightharpoonup g$, risulta che $(f_h)$ è limitata e per ogni $u \in C^\infty_{c}(U)$	
\[ \int_{U}f_h u d\mathcal{L}^n \to \int_U g u d\mathcal{L}^n. \]
Ora, siccome di fatto integriamo solo sul supporto di $u$, per la Proposizione \eqref{convergenzaLq}, per ogni $q < p$, $f_h u \to fu$ in $L^q(U)$. Siccome la convergenza forte implica quella debole, $f_hu \rightharpoonup fu $ in $L^q(U)$, e quindi, a meno di rinominazioni, per ogni $u \in C^\infty_{c}(U)$	
\[ \int_{U}f_h u d\mathcal{L}^n \to \int_U f u d\mathcal{L}^n. \]
A questo punto, per l'unicità del limite, deve essere per ogni $u \in C^\infty_{c}(U)$	
\[ \int_U g u d\mathcal{L}^n = \int_U f u d\mathcal{L}^n \]
ossia
\[ \int_U (g-f) u d\mathcal{L}^n = 0 \]
e quindi, per il Lemma di Du Bois-Reymond, $g = f$ q.o. in $U$.
\end{proof}

\begin{osservazione}\label{convergenzadeboleLp}
Consideriamo una successione $(u_h)$ in $L^p(U)$ limitata.
\begin{itemize}
	\item se $1 < p < \infty$, $L^p(U)$ è riflessivo quindi per il Teorema di Eberlein--Smulian esistono $u \in L^p(U)$ e una sottosuccessione $u_{h_{k}}$ tali che per ogni $\varphi \in L^{p'}(U)$
	\[ \int_{U}u_{h_{k}}\varphi d\mathcal{L}^n \to \int_{U}u\varphi d\mathcal{L}^n \]
	\item se $p = \infty$, siccome $L^\infty(U) \cong (L^1(U))'$, consideriamo $T: L^\infty(U) \to (L^1(U))'$ tale che 
	\[ \braket{Tu_h,\varphi} = \int_{U}u_h\varphi d\mathcal{L}^n. \]
	Essendo $L^1(U)$ separabile, dalla Proposizione \eqref{separabilitàe*} esistono  $G \in (L^1(U))'$ ed $Tu_{h_{k}}$ tali che per ogni $\varphi \in L^{1}(U)$
	\[ \braket{Tu_{h_{k}},\varphi} \to \braket{G,\varphi}\]
	ma deve essere $G = Tu$ per qualche $u \in L^\infty(U)$ e quindi  per ogni $\varphi \in L^{1}(U)$
	\[ \int_{U}u_{h_{k}}\varphi d\mathcal{L}^n \to \int_{U}u\varphi d\mathcal{L}^n \]
	\item se $p=1$, esiste $T: L^1(U) \to \left( C(\overline{U})\right)'$ tale che per ogni $g \in C(\overline{U})$,
	\[ \braket{Tu_h,g} = \int_U u_h g d\mathcal{L}^n. \]
	Essendo $C(\overline{U})$ separabile, per la Proposizione \eqref{separabilitàe*}, esiste $G \in \left( C(\overline{U})\right)'$ ed una sottosuccessione $(u_{h_k})$ tali che per ogni $g \in C(\overline{U})$ 
	\[ \int_U u_{h_k} g d\mathcal{L}^n \to \braket{G,g}. \]
	A questo punto, esiste una ed una sola misura di Radon $\nu$\footnote{Omettiamo la dimostriazione di questo fatto.} tale che
	\[ \braket{G,g} = \int_U gd\nu \]
	e quindi per ogni $g \in C(\overline{U})$
	\[ \int_U u_{h_k} g d\mathcal{L}^n \to \int_U gd\nu.\]
\end{itemize}
\end{osservazione}
\begin{proposizione}
	Sia $E \subseteq \R^{n} $ misurabile. Allora $L^{1}(E)$ non è riflessivo.
\end{proposizione}
\begin{proof}
	Consideriamo innanzitutto $(U_j)$ una successione decrescente di aperti di $E$ tale che $\mathcal{L}^n(U_j) \to 0$ e per ogni $j \in \N$
	\[ u_j = \frac{\chi_{U_j}}{\norm{\chi_{U_j}}_{L^1(E)}}. \]
	Chiaramente, $\norm{u_j}_{L^1(E)} = 1$, quindi $(u_j)$ è limitata in $L^1(E)$. Se, per assurdo, $L^1(E)$ fosse riflessivo, allora, per il Teorema di Eberlein--Smulian, esisterebbero $(u_{j_k})$ in $L^1(E)$ ed $u \in L^1(E)$ tali che $u_{j_k} \rightharpoonup u$ in $L^1(E)$. In altre parole, per ogni $\psi \in L^\infty(E)$
	\[ \int_E u_{j_k}\psi d\mathcal{L}^n \to \int_E u \psi d\mathcal{L}^n. \]
	Per ogni $h \in \N$ si ha $\chi_{U_h} \in L^\infty(E)$, quindi 
	\[ \int_E u_{j_k}\chi_{U_h} d\mathcal{L}^n \to \int_E u \chi_{U_h} d\mathcal{L}^n, \]
	ma $\chi_{U_{j_k}}\chi_{U_h} = \chi_{U_{j_k}\cap U_h}$, perciò, essendo anche $(U_{j_k})$ decrescente, per ogni $h \in \N$ esiste $\overline{k} \in \N$ tale che per ogni $k \ge \overline{k}$ si abbia $\chi_{U_{j_k}}\chi_{U_h} = \chi_{U_{j_k}}$. Allora
	\[  \int_E u_{j_k}\chi_{U_h} d\mathcal{L}^n \to 1 \]
	e, per unicità del limite,
	\[ \int_E u \chi_{U_h} d\mathcal{L}^n = 1. \]
	Ora, però, 
	\[ \lim\limits_{h} \norm{\chi_{U_h}}_{L^1(E)} = \lim\limits_{h} \mathcal{L}^n(U_h) = 0,\]
	quindi $\chi_{U_h} \to 0$ in $L^1(E)$. A meno di sottosuccessioni, per le proprietà degli spazi di Lebesgue, $\chi_{U_h} \to 0$ q.o. in $L^1(E)$. In particolare, $u\chi_{U_h} \to 0$ q.o. in $L^1(E)$ e $u\chi_{U_h} \le u \in L^1(E)$ per cui, applicando il Teorema della convergenza dominata,
	\[  \lim\limits_{h} \int_E u\chi_{U_h} d\mathcal{L}^n = 0, \]
	da cui l'assurdo.
\end{proof}

\section*{Riepilogo proprietà degli Spazi di Lebesgue}
Sia $E$ un sottoinsieme misurabile di $\R^{n}$. Possiamo così riepilogare le proprietà di $(L^{p}(E), \norm{\;}_{p})$.

\hfill

\begin{center}
	\begin{tabular}{lccc}
		
		& $p=1$ & $1<p<\infty$ & $p=\infty$ \\
		
		Spazio di Banach & \checked & \checked & \checked \\
		
		Spazio di Banach separabile & \checked & \checked &  \\
		
		Spazio di Banach riflessivo &  & \checked &  \\
		
	\end{tabular}
\end{center}

\section{Approssimazione polinomiale di funzioni continue}
Il seguente risultato è stato dimostrato da Karl Weierstrass nel 1855.
\begin{teorema}\emph{\textbf{(di approssimazione di Weierstrass)}}\index[analitico]{teorema!di approssimazione di Weierstrass}
Siano $a,b \in \R$, con $a<b$, ed $f : [a,b] \to \R$ un'applicazione continua. Allora esiste una successione di polinomi $(p_h)$ tale che $p_h \to f$ uniformemente su $[a,b]$.
\end{teorema}
\begin{proof}
Innanzitutto, a meno di una traslazione ed un riscalamento possiamo supporre $[a,b] = [0,1]$. Inoltre, senza perdita di generalità possiamo supporre che $f(0) = f(1) = 0$.\footnote{Consideriamo l'applicazione 
	\[ \tilde{f}(x) = f(x) - f(0) -x(f(1) - f(0)).  \]
	Chiaramente, $\tilde{f}$ realizza l'ipotesi richiesta. Se dimostriamo che la tesi vale per $\tilde{f}$, essendo quest'ultima $f$ privata di un polinomio, allora sarà vera anche per $f$. Non è restrittivo considerare quindi fin da subito $f$ con l'ipotesi menzionata.} Possiamo inoltre, a meno di un prolungamento continuo, pensare che $f = 0$ fuori da $[0,1]$. In particolare, $f$ è uniformemente continua.
	
	Consideriamo il polinomio $Q_n$ tale che
	\[ Q_n(x) = c_n(1-x^2)^n, \]
	dove $c_n \in \R$ tale che
	\[ \int_{-1}^1 Q_n d\mathcal{L}^1 = 1. \]
	In particolare, siccome $(1-x^2)^n \ge 1-nx^2$ su $[0,1]$,\footnote{Sia $u:[0,1]\to \R$ tale che
	\[ u(x) = (1-x^2)^n - 1+nx^2. \]
	Sicuramente $u$ è continua su $[0,1]$ e derivabile su $\left] 0,1\right[ $. Inoltre $u(0) = 0$ e 
	\[ u'(x) = -2nx (1-x^2)^{n-1} +2nx = 2nx(1-(1-x^2)^{n-1}) \ge 0, \]
	quindi $u(x) \ge 0$.}
	\begin{align*}
	\int_{-1}^1 (1-x^2)^n dx &= 2 \int_0^1 (1-x^2)^ndx \ge 2\int_{0}^{\frac{1}{\sqrt{n}}}(1-x^2)^n dx \ge 2\int_{0}^{\frac{1}{\sqrt{n}}}(1-nx^2) dx =  \\ 
	&=\frac{4}{3\sqrt{n}} > \frac{1}{\sqrt{n}},
	\end{align*}
	quindi $c_n < \sqrt{n}$. Inoltre, dato $\delta \in \left] 0,1\right[ $, su $[-1,-\delta]\cup [\delta,1]$
	\[ Q_n(x) < \sqrt{n} (1-\delta^2)^n, \]
	quindi $Q_n \to 0$ uniformemente su $[-1,-\delta]\cup [\delta,1]$.
	
	Consideriamo l'applicazione $p_n: [0,1]\to \R$ tale che
	\[ p_n(x) = \int_{-1}^1 f(x+t)Q_n(t)dt. \]
	Innanzitutto, siccome $f$ è nulla fuori da $[0,1]$,
	\[ p_n(x) = \int_{-x}^{1-x}f(x+t)Q_n(t) dt = \int_0^1 f(t) Q_n(t-x) dt, \]
	ossia $p_n$ è un polinomio in $x$. Sia $M = \max f$. Per ogni $\varepsilon>0$, sia $\delta >0$ tale che per ogni $x,y \in [0,1]$ con $\abs{x-y} <\delta$, si abbia $\abs{f(x)-f(y)}<\frac{\varepsilon}{2}$. Senza perdita di generalità possiamo supporre $\delta <1$. A questo punto, per ogni $x \in [0,1]$, se $t \in \left] -\delta,\delta\right[ $ si ha $\abs{x+t-x} = \abs{t}<\delta$, allora
	\begin{align*}
		\abs{p_n(x)-f(x)}&\le \int_{-1}^1 \abs{f(x+t)-f(x)}Q_n(t) dt = \int_{-1}^{-\delta}\abs{f(x+t)-f(x)}Q_n(t) dt + \\
		&+ \int_{-\delta}^{\delta}\abs{f(x+t)-f(x)}Q_n(t) dt+ \int_{\delta}^{1}\abs{f(x+t)-f(x)}Q_n(t) dt <\\
		&<2M \int_{-1}^{-\delta} Q_n d\mathcal{L}^1 + \frac{\varepsilon}{2}\int_{-\delta}^{\delta}Q_nd\mathcal{L}^1 +2M \int_{\delta}^{1} Q_n d\mathcal{L}^1 \le \\
		&\le 2M\left( \int_{-1}^{-\delta} Q_n d\mathcal{L}^1 + \int_{\delta}^{1} Q_n d\mathcal{L}^1\right)  + \frac{\varepsilon}{2}\int_{-1}^{1}Q_nd\mathcal{L}^1 = \\
		&=2M\left( \int_{-1}^{-\delta} Q_n d\mathcal{L}^1 + \int_{\delta}^{1} Q_n d\mathcal{L}^1\right)  + \frac{\varepsilon}{2}
	\end{align*}
	e, per il fatto che $Q_n \to 0$ uniformemente su $[-1,-\delta]\cup [\delta,1]$, esiste $N >0$ tale che per ogni $n \ge N$ si abbia $Q_n\le \frac{\varepsilon}{8M}$ su $[-1,-\delta]\cup [\delta,1]$, da cui 
	\[ 	\abs{p_n(x)-f(x)}< \varepsilon, \]
	quindi $p_n \to f$ uniformemente su $[0,1]$, da cui la tesi.
\end{proof}

\begin{esempio}
Consideriamo $f: [0,1]\to \R$ continua tale che per ogni $n \in \N$
\[ \int_{0}^{1}f(x)x^n dx = 0. \]
Mostriamo che $f = 0$.

Innanzitutto, per linearità, per ogni polinomio $p$
\[ \int_{0}^{1}f p d\mathcal{L}^1 = 0. \]
Ora, per il Teorema di approssimazione di Weierstrass, esiste $(p_h)$ tale che $p_h \to f$ uniformemente su $[0,1]$, quindi 
\[ \int_{0}^{1} f^2 d\mathcal{L}^1 = \lim\limits_{h} \int_{0}^{1} fp_h d\mathcal{L}^1 = \lim\limits_{h} 0 = 0, \]
da cui, essendo $f$ continua, $f=0$.
\end{esempio}

Miriamo ora a presentare la generalizzazione del risultato di Weierstrass proposta da Mashall Stone nel 1937: a malincuore dobbiamo parlare di algebre.
\begin{definizione}
Consideriamo $(X,\norm{\;})$ uno spazio normato su $\R$ ed $E \subseteq X$. Diciamo che $\mathcal{A} \subseteq C_b(E; \R^n)$ è un'\emph{algebra}\index[analitico]{algebra} (su $\R$) se valgono i seguenti fatti:
\begin{enumerate}[$(a)$]
\item se $f,g \in \mathcal{A}$, allora $f+g \in \mathcal{A}$,
\item se $f,g \in \mathcal{A}$, allora $fg \in \mathcal{A}$,
\item se $k \in \R$ e $f \in \mathcal{A}$, allora $kf \in \mathcal{A}$.
\end{enumerate}
\end{definizione}

\begin{lemma}
Consideriamo $(X,\norm{\;})$ uno spazio normato su $\R$ ed $E \subseteq X$. Se $\mathcal{A} \subseteq C_b(E; \R^n)$ è un'algebra, allora $\overline{\mathcal{A}}$ è un'algebra.
\end{lemma}
\begin{proof}
Siano $f,g \in \overline{\mathcal{A}}$ e $k \in \R$. Esistono $(f_h),(g_h)$ in $\mathcal{A} \subseteq C_b(E; \R^n)$ tali che $f_h \to f$, $g_h \to g$ uniformemente su $E$. Essendo $C_b(E;\R^n)$ chiuso in $\mathcal{B}(E;\R^n)$, ne deduciamo che $f,g \in C_b(E; \R^n)$. A questo punto, chiaramente $f_h + g_h \to f+g$ e $kf_h \to kf$ uniformemente su $E$, quindi $f+g, kf \in \overline{\mathcal{A}}$. Inoltre, essendo tutte le applicazioni limitate, anche $f_hg_h \to fg$ uniformemente su $E$, da cui $fg \in \overline{\mathcal{A}}$.
\end{proof}

\begin{definizione}
Consideriamo $(X,\norm{\;})$ uno spazio normato su $\R$ ed $E \subseteq X$. Diciamo che un'algebra $\mathcal{A} \subseteq C_b(E; \R^n)$ \emph{separa i punti}\index[analitico]{algebra!che separa i punti} se per ogni $x,y \in E$ esiste $f \in \mathcal{A}$ tale che $f(x) \ne f(y)$.
\end{definizione}

\begin{definizione}
Consideriamo $(X,\norm{\;})$ uno spazio normato su $\R$ ed $E \subseteq X$. Diciamo che un'algebra $\mathcal{A} \subseteq C_b(E; \R^n)$ \emph{non si annulla}\index[analitico]{algebra!che non si annulla} su $E$ se per ogni $x \in E$ esiste $f \in \mathcal{A}$ tale che $f(x) \ne 0$.
\end{definizione}

\begin{lemma}\label{tecnicoperSW}
Consideriamo $(X,\norm{\;})$ uno spazio normato su $\R$ ed $E \subseteq X$. Se $\mathcal{A} \subseteq C_b(E; \R)$ è un'algebra che separa i punti e non si annulla su $E$, allora per ogni $x,y \in E$ e per ogni $c_1,c_2 \in \R$ esiste $f \in \mathcal{A}$ tale che $f(x) = c_1$ e $f(y) = c_2$.
\end{lemma}
\begin{proof}
Dal fatto che $\mathcal{A}$ separa i punti, esiste $g \in \mathcal{A}$ tale che $g(x) \ne g(y)$. Inoltre, dal fatto che $\mathcal{A}$ non si annulla su $E$, esistono $h,k \in \mathcal{A}$ tale che $h(x),k(y) \ne 0$. Posto
\begin{align*}
 u &= gk -g(x)k, & v= gh-g(y)h,
\end{align*}
basta considerare 
\[ f = \frac{c_2u}{u(y)} + \frac{c_1v}{v(x)}.\qedhere \]
\end{proof}

\begin{teorema}\emph{\textbf{(di Stone--Weierstrass)}}\index[analitico]{teorema!di Stone--Weierstrass}
Consideriamo $K \subseteq \R^n$ compatto. Se $\mathcal{A} \subseteq C(K; \R^m)$ è un'algebra che separa i punti e non si annulla su $K$, allora $\overline{\mathcal{A}} = C(K; \R^m)$.
\end{teorema}
\begin{proof}
Innanzitutto, a meno di ragionare per componenti, è sufficiente considerare il caso $\mathcal{A} \subseteq C(K; \R)$. Sicuramente, $\overline{\mathcal{A}} \subseteq C(K; \R^m)$. 

Mostriamo, prima di tutto, che se $f \in \overline{\mathcal{A}}$, allora $\abs{f} \in \overline{\mathcal{A}}$. In primo luogo, per il Teorema di Weierstrass, $f$ è limitata, ossia esiste $a \in \R$ tale che per ogni $x \in K$ si abbia $\abs{f(x)}\le a$. Per il Teorema di approssimazione di Weierstrass, esiste $(p_n)$ convergente uniformemente a $\abs{x}$ su $[-a,a]$: in altre parole, per ogni $\varepsilon > 0$ esistono $c_1,\dots,c_{n_\varepsilon} \in \R$ tali che per ogni $x \in [-a,a]$
\[ \abs*{\sum_{i=1}^{n_\varepsilon}c_ix^i- \abs{x}}<\varepsilon.\footnote{Siccome l'applicazione valore assoluto si annulla in zero, la successione di polinomi approssimanti può essere presa senza termini noti.} \]
Siccome $g_\varepsilon = \sum_{i=1}^{n_\varepsilon}c_if^i \in \overline{\mathcal{A}}$ e $\abs{g_\varepsilon-\abs{f}}<\varepsilon$ su $K$, ne deduciamo che $\abs{f} \in \overline{\mathcal{A}}$. In particolare, date $f,g \in \overline{\mathcal{A}}$, 
\begin{align*}
	\max\left\lbrace f,g\right\rbrace  &= \frac{1}{2}(f+g + \abs{f-g})\in \overline{\mathcal{A}}, & \min\left\lbrace f,g\right\rbrace  &= \frac{1}{2}(f+g - \abs{f-g})\in \overline{\mathcal{A}}
\end{align*}
e, per induzione, lo stesso vale per il massimo ed il minimo tra un numero finito di elementi di $\overline{\mathcal{A}}$.

Fissiamo $f \in C(K;\R)$ ed $\varepsilon>0$. Mostriamo che per ogni $x \in K$ esiste $g_x \in \overline{\mathcal{A}}$ tale che $g_x(x) = f(x)$ e per ogni $t \in K$ si abbia $g_x(t)\ge f(t)-\frac{\varepsilon}{2}$. Per il Lemma \eqref{tecnicoperSW}, per ogni $y \in K$ esiste $h_y \in \mathcal{A}\subseteq \overline{\mathcal{A}}$ tale che $h_y(x) = f(x)$ e $h_y(y) = f(y)$. Siccome $h_y - f$ è continua, per il Teorema di di permanenza del segno, esiste $U_y$ intorno aperto di $y$ tale che $h_y \ge f -\frac{\varepsilon}{2}$ su $U_y$. Siccome 
\[ K \subseteq \bigcup_{y \in K} U_y, \]
per compattezza, esistono $y_1,\dots,y_m$ tali che 
\[ K \subseteq \bigcup_{j=1}^m U_{y_j}. \]
Consideriamo $g_x = \max \left\lbrace h_{y_1},\dots, h_{y_m}\right\rbrace \in \overline{\mathcal{A}}$. Allora, 
\[ g_x(x) = \max \left\lbrace h_{y_1}(x),\dots, h_{y_m}(x)\right\rbrace = f(x) \]
e per ogni $t \in K$ esiste $i =1,\dots, m$ tale che $t \in U_{y_i}$, quindi 
\[ g_x(t) \ge h_{y_i}(t) \ge f(t)-\frac{\varepsilon}{2}, \]
quindi $g_x$ soddisfa i requisti richiesti.

A questo punto, sempre per il Teorema di permanenza del segno, per ogni $x \in K$ esiste un intorno aperto $V_x$ di $x$ tale che $g_x \le f+\frac{\varepsilon}{2}$ su $V_x$. Siccome 
\[ K \subseteq \bigcup_{x \in K} V_x, \]
per compattezza, esistono $x_1,\dots,x_N$ tali che 
\[ K \subseteq \bigcup_{j=1}^N V_{x_j}. \]
Preso $h = \min \left\lbrace g_{x_1},\dots g_{x_N}\right\rbrace \in \overline{\mathcal{A}}$, si ottiene che $\abs{h - f} \le \frac{\varepsilon}{2} <\varepsilon$, da cui la tesi.
\end{proof}

\begin{osservazione}
A ben vedere, non serve conoscere il Teorema di approssimazione di Weierstrass per dimostrare il risultato generale. Nella dimostrazione utilizziamo solo il fatto che l'approssimazione polinomiale valga per l'applicazione valore assoluto: in particolare, è sufficiente sapere che, per ogni $a >0$, la successione $(p_n)$ definita ricorsivamente da 
\begin{align*}
		p_0 &= 0, & p_{n+1}(x) &= p_n(x) + \frac{1}{2}(x^2 -p_n(x)^2),
\end{align*}
converge uniformemente a $\abs{x}$ su $[-a,a]$.
\end{osservazione}

\begin{osservazione}
Consideriamo $E \subseteq \R^n$ compatto. L'algebra delle applicazioni a componenti polinomiali separa i punti\footnote{Basta considerare l'applicazione identica, che ha componenti polinomiali.} e non si annulla su $E$\footnote{Basta considerare un'applicazione costante con almeno una componente non nulla.}, quindi per ogni $f \in C(E; \R^n)$ esiste una successione $(p_h)$ di applicazioni a componenti polinomiali tale che $p_h \to f$ uniformemente su $E$.
\end{osservazione}

\section{Spazi di Lebesgue che coinvolgono il tempo}
Prima di tutto necessitiamo di generalizzare la teoria dell'integrazione di Lebesgue.
\begin{definizione}
Consideriamo $(X,\norm{\;})$ uno spazio di Banach su $\R$ e $T>0$. Diciamo che un'applicazione $s: [0,T] \to X$ è \emph{semplice}\index[analitico]{applicazione!semplice} se
\[ s(t) = \sum_{i=1}^{m}\chi_{E_i}(t)u_i,  \]
dove $E_1,\dots,E_m \subseteq [0,T]$ misurabili e $u_1,\dots,u_m \in X$. 
\end{definizione}

\begin{definizione}
Consideriamo $(X,\norm{\;})$ uno spazio di Banach su $\R$ e $T>0$. Diciamo che un'applicazione $f: [0,T] \to X$ è \emph{fortemente misurabile}\index[analitico]{applicazione!misurabile!fortemente} se esiste una successione $(s_h)$ di funzioni semplici tali che $s_h \to f$ \emph{q.o.} in $[0,T]$.
\end{definizione}

\begin{definizione}
Consideriamo $(X,\norm{\;})$ uno spazio di Banach su $\R$ e $T>0$. Se $s: [0,T] \to X$ è un'applicazione semplice, definiamo 
\[ \int_0^T sd\mathcal{L}^1 = \sum_{i=1}^{m} \mathcal{L}^1(E_i)u_i. \]
\end{definizione}

\begin{definizione}
Consideriamo $(X,\norm{\;})$ uno spazio di Banach su $\R$ e $T>0$. Diciamo che un'applicazione $f: [0,T] \to X$ fortemente misurabile è \emph{sommabile}\index[analitico]{applicazione!sommabile} se esiste una successione $(s_h)$ di funzioni semplici tali che 
\[ \lim\limits_{h} \int_{0}^{T} \norm{s_h(t)-f(t)}d\mathcal{L}^1(t) = 0.  \]
\end{definizione}

\begin{definizione}
	Consideriamo $(X,\norm{\;})$ uno spazio normato su $\R$ e $T>0$. Se $f: [0,T] \to X$ è un'applicazione sommabile, definiamo 
	\[ \int_0^T fd\mathcal{L}^1 = \lim\limits_{h} \int_0^T s_hd\mathcal{L}^1 \]
\end{definizione}

\begin{definizione}
Consideriamo $(X,\norm{\;})$ uno spazio di Banach su $\R$, $p \in [1,\infty]$ e $T>0$. Chiamiamo $L^p(0,T;X)$\index[simboli]{$L^p(0,T;X)$} l'insieme delle (classi di equivalenza) di applicazioni fortemente misurabili $u: [0,T]\to X$ tali che, se $p \in \left[ 1,\infty\right[$, 
\[ \int_0^T \norm{u(t)}^p d\mathcal{L}^1 <\infty \]
oppure, nel caso $p=\infty$,
\[ \underset{t \in [0,T]}{\textup{ess sup}}\; \norm{u(t)} <+\infty. \]
\end{definizione}

\begin{definizione}
Consideriamo $(X,\norm{\;})$ uno spazio di Banach su $\R$, $p \in [1,\infty]$ e $T>0$. Data $u \in L^p(0,T;X)$, poniamo 
	\[ \norm{u}_{L^p(0,T;X)} = \begin{cases}
		\left( \displaystyle \int_0^T \norm{u(t)}^p d\mathcal{L}^1\right) ^{\frac{1}{p}} & 1 \le p < \infty, \\
		\underset{t \in [0,T]}{\textup{ess sup}}\; \norm{u(t)} & p=\infty.
	\end{cases} \]
	
\end{definizione} 
		
		% Commenti speciali

% !TEX encoding = UTF-8 Unicode
% !TEX TS-program = pdflatex
% !TeX spellcheck = it_IT
% !TEX root = Dispensa/afed.tex

% Capitolo 1
\chapter{Spazi di Sobolev}
\section{Prime proprietà degli spazi di Sobolev}
\begin{definizione}
	Chiamiamo $W^{1,1}_{loc}(U)$\index[simboli]{$W^{1,1}_{loc}(U)$} l’insieme delle $u \in L^{1}_{loc}(U)$ per cui esistono $v_{1},\dots,v_{n} \in L^{1}_{loc}(U)$ tali che per ogni $j = 1,\dots, n$, per ogni $\varphi \in C^{\infty}_{c}(U)$ risulta
	\[ \int_{U} uD_{j}\varphi d\mathcal{L}^n = - \int_{U}v_{j}\varphi d\mathcal{L}^n.\]
\end{definizione}
\begin{osservazione}\label{derivatadebolecasoregolare}
	Se $u \in C^{1}(U)$, allora per ogni $j = 1,\dots, n$ $D_{j}u \in C(U)$ e per la Formula di Gauss--Green
	\[ \int_{U} uD_{j}\varphi d\mathcal{L}^n = - \int_{U}D_{j}u\varphi d\mathcal{L}^n.\]
\end{osservazione}
\begin{proposizione}
	Sia $u \in L^{1}_{loc}(U)$ tale che esistono $v_{1},\dots,v_{n} \in L^{1}_{loc}(U)$ tali che per ogni $j = 1,\dots, n$, per ogni $\varphi \in C^{\infty}_{c}(U)$ risulta
	\[ \int_{U} uD_{j}\varphi d\mathcal{L}^n = - \int_{U}v_{j}\varphi d\mathcal{L}^n.\]
	Allora per ogni $j = 1,\dots, n$, $v_{j}$ è unica in $L^{1}_{loc}(U)$.
\end{proposizione}
\begin{proof}
	Fissato $j$, supponiamo esistano $v_{j}, \hat{v}_{j}$ soddisfacenti le ipotesi. Allora per ogni $\varphi \in C^{\infty}_{c}(U)$
	\[ \int_{U}v_{j}\varphi d\mathcal{L}^n = \int_{U}\hat{v}_{j}\varphi d\mathcal{L}^n \]
	da cui 
	\[ \int_{U}(v_{j} - \hat{v}_{j})\varphi d\mathcal{L}^n = 0 \]
	e quindi $v_{j} = \hat{v}_{j}$ q.o. da cui l'unicità.
\end{proof}

\begin{definizione}\index[analitico]{derivata!debole}
	Chiamiamo le $v_{1},\dots,v_{n}$ sopra introdotte \emph{derivate deboli} o derivate distribuzionali, o derivate variazionali.
\end{definizione}
In particolare, per l'Osservazione \eqref{derivatadebolecasoregolare} denotiamo le derivate deboli con lo stesso simbolo delle derivate ordinarie.
\begin{esempio}
	Sia $U = \left] 0,2\right[ $ e 
	\[ u(x) = \begin{cases}
		x & \textup{se } 0 < x \leq 1, \\
		1 & \textup{se } 1 < x < 2.
	\end{cases} \]
	Mostriamo che $u \in W^{1,1}_{loc}(0,2)$. Consideriamo l'applicazione
	\[ v(x) = \begin{cases}
		1 & \textup{se } 0 < x \le 1, \\
		0 & \textup{se } 1 < x < 2.
	\end{cases} \]
	Per ogni $\varphi \in C^{\infty}_{c}(0,2)$, osservando che $\varphi(2) = 0 $,
	\[ \int_{0}^{2} u \varphi' d\mathcal{L}^1 = \int_{0}^{1}x\varphi'(x)dx + \int_{1}^{2}\varphi'(x)dx = - \int_{0}^{1}\varphi(x)dx = - \int_{0}^{2}v(x)\varphi(x)dx,  \]
	allora $v$ è la derivata debole di $u$.
\end{esempio}
\begin{esempio}
	Sia $U = \left] 0,2\right[ $ e
	\[ u(x) = \begin{cases}
		x & \textup{se } 0 < x \leq 1, \\
		2 & \textup{se } 1 < x < 2.
	\end{cases} \]
	Supponiamo, per assurdo, che esista $v \in L^1_{loc}(0,2)$ tale che per ogni $\varphi \in C^{\infty}_{c}(0,2)$
	\[ \int_{0}^{2}u\varphi'd\mathcal{L}^1 = - \int_{0}^{2}v\varphi d\mathcal{L}^1. \]
	Innanzitutto,
	\[ \int_{0}^{2}u\varphi'd\mathcal{L}^1 = -\int_{0}^{1}\varphi d\mathcal{L}^1 - \varphi(1). \]
	Si può costruire\footnote{Data $\eta \in C^{\infty}_{c}(0,2)$, basta considerare
		\[ \varphi_h(x) = \eta(x)e^{-\abs{x-1}h}. \]} $(\varphi_h) \in C^{\infty}_{c}(0,2)$ tale che $0 \le \varphi_h \le 1$, $\varphi_h(1) = 1$ e $\varphi_h(x) \to 0$ se $x \ne 1$. Allora,
	\[ \int_{0}^{2}v\varphi_h d\mathcal{L}^1 = \int_{0}^{1}\varphi_hd\mathcal{L}^1 + \varphi_h(1)\]
	e dal Teorema della convergenza dominata otteniamo $0=1$, assurdo.
\end{esempio}
\begin{notazione}
	Sia $p \in \left[ 1,\infty\right] $. Indichiamo con  $W^{1,p}(U)$\index[simboli]{$W^{1,p}(U)$} l'insieme delle $u \in W^{1,1}_{loc}(U)$ tali che $u \in L^p(U)$ e per ogni $j = 1,\dots,n$ si ha $D_{j}u \in L^p(U)$.
\end{notazione}

\begin{esempio}\label{Esempiopalla}
	Consideriamo $U= \textup{B}(0,1)$ in $\R^n$, $\alpha > 0$, $p \in \left]
	1,n\right[$ e 
	\[ u(x)=\frac{1}{\abs{x}^\alpha}. \]
	Mostriamo che $u \in W^{1,p}(U)$ se e solo se $\alpha < \frac{n-p}{p}$.
	
	Innanzitutto, per $x \ne 0$, $u$ è regolare e 
	\[ D_ju(x) = -\alpha\frac{x_i}{\abs{x}^{\alpha+2}}, \]
	per cui candidiamo come derivata debole proprio $D_ju$ opportunamente estesa.
	
	Sia ora $\varepsilon > 0$. Senza perdita di generalità, possiamo supporre $\varepsilon < 1$. Consideriamo l'insieme $U_{\varepsilon} = U\setminus \textup{B}(0,\varepsilon) $. Se $\varphi \in C^\infty_{c}(U)$,
	\[ \int_{U_{\varepsilon}} uD_j\varphi d\mathcal{L}^n = - \int_{U_{\varepsilon}} D_ju \varphi d\mathcal{L}^n + \int_{\partial\textup{B}(0,\varepsilon)} u\varphi \nu^j d\mathcal{H}^{n-1}. \]
	Osserviamo che 
	\[ \int_{U} \abs{u}^pd\mathcal{L}^n = K \int_{0}^{1}\frac{1}{\rho^{\alpha p}}\rho^{n-1}d\rho < +\infty \]
	e 
	\[ \int_{U} \abs{Du}^pd\mathcal{L}^n = K \int_{0}^{1}\frac{1}{\rho^{(\alpha+1) p}}\rho^{n-1}d\rho < +\infty \]
	se e solo se $\alpha < \frac{n-p}{p}$. 
	Inoltre, per $\alpha < n-1$, 
	\begin{align*}
		\abs*{\int_{\partial\textup{B}(0,\varepsilon)} u\varphi \nu^j d\mathcal{H}^{n-1}} &\le \varepsilon^{-\alpha}\int_{\partial\textup{B}(0,\varepsilon)} \abs{\varphi} d\mathcal{H}^{n-1} \le \\
		&\le \norm{\varphi}_{\infty}\varepsilon^{-\alpha}\mathcal{H}^{n-1}(\partial\textup{B}(0,\varepsilon)) \le C \varepsilon^{n-1-\alpha} \to 0.
	\end{align*} 
	A questo punto, per il Teorema della convergenza dominata combinato con il fatto che $\chi_{U_{\varepsilon}} \to \chi_{U}$ \emph{q.o.} in $U$,
	\[ \int_{U_{\varepsilon}} uD_j\varphi d\mathcal{L}^n \to  \int_{U} uD_j\varphi d\mathcal{L}^n \]
	e 
	\[ \int_{U} D_ju \varphi d\mathcal{L}^n \]
	che permettono di concludere che $u \in W^{1,p}(U)$ se e solo se $\alpha < \frac{n-p}{p}$.
\end{esempio}

\begin{definizione}
	Siano $k \in \N\setminus\left\lbrace 0\right\rbrace $ e $1\le p \le \infty$. Chiamiamo \emph{spazio di Sobolev}\index[analitico]{spazio!di Sobolev} l'insieme $W^{k,p}(U)$\index[simboli]{$W^{k,p}(U)$} delle $u \in L^1_{loc}(U)$ tali che per ogni multi indice $\alpha \in \N^n$, esiste $v_{\alpha} \in L^p(U)$, detta \emph{$\alpha$-esima derivata debole}, tale che per ogni $\varphi \in C^\infty_c(U)$
	\[ \int_{U}uD^\alpha\varphi d\mathcal{L}^n = (-1)^{\abs{\alpha}}\int_{U}v_{\alpha}\varphi d\mathcal{L}^n. \]
\end{definizione}

\begin{notazione}
	Indichiamo con $H^{1}(U) = W^{1,2}(U)$\index[simboli]{$H^{1}(U)$}.
\end{notazione}

\begin{proposizione}\label{propderivatedeboli}
	Siano $k \in \N\setminus\left\lbrace 0\right\rbrace $, $1 \le p \le \infty$ ed $u,v \in W^{k,p}(U)$.  Valgono i seguenti fatti:
	\begin{enumerate}[$(a)$]
		\item se $\alpha, \beta \in \N^n$ tali che $\abs{\alpha} + \abs{\beta} \le k$, allora $D^\alpha u \in W^{k - \abs{\alpha},p}(U)$ e 
		\[ D^\beta (D^\alpha u) = D^\alpha (D^\beta u) = D^{\alpha + \beta} u;  \]
		\item per ogni $\lambda,\mu \in \R$, $\lambda u + \mu v \in W^{k,p}(U)$ e se $\alpha \in \N^n$ tale che $\abs{\alpha}\le k$,
		\[ D^\alpha (\lambda u + \mu v )= \lambda D^\alpha u + \mu D^\alpha v. \]
		In particolare, $W^{k,p}(U)$ è uno spazio vettoriale su $\R$;
		\item se $V$ è un sottoinsieme aperto di $U$, allora $u \in W^{k,p}(V)$;
		\item  se $\eta \in C^\infty_c(U)$, allora $\eta u \in W^{k,p}(U)$ e se $\alpha \in \N^n$ tale che $\abs{\alpha}\le k$,
		\[ D^\alpha(\eta u) = \sum_{\beta \le \alpha}^{} \binom{\alpha}{\beta} D^\beta \eta D^{\alpha - \beta} u. \qquad \textup{(Formula di Leibniz)} \]
		In particolare, se $\abs{\alpha} = 1$, 
		\[ D^\alpha(\eta u) = u D^\alpha \eta + \eta D^\alpha u. \]
	\end{enumerate}
\end{proposizione}
\begin{proof}\hfill
	\par\noindent
	$(a)$ Innanzitutto, per definizione, per ogni $\eta \in C^\infty_c(U)$
	\[ \int_{U} u D^\alpha \eta d\mathcal{L}^n = (-1)^{\abs{\alpha}} \int_{U} D^\alpha u \eta d\mathcal{L}^n. \]
	Se $\varphi \in C^\infty_c(U)$, allora $D^\beta \varphi \in C^\infty_{c}(U)$ e quindi
	\begin{align*}
		\int_{U} D^\alpha u D^\beta \varphi d\mathcal{L}^n &= (-1)^{\abs{\alpha}} \int_{U}  u D^{\alpha + \beta} \varphi d\mathcal{L}^n = \\ 
		& =(-1)^{\abs{\alpha}} (-1)^{\abs{\alpha} + \abs{\beta}} \int_{U} D^{\alpha + \beta} u \varphi  d\mathcal{L}^n = (-1)^{\abs{\beta}} \int_{U} D^{\alpha + \beta} u \varphi  d\mathcal{L}^n
	\end{align*}
	da cui $D^\beta(D^\alpha u) = D^{\alpha + \beta} u$. Scambiando $\alpha$ e $\beta$ si giunge alla conclusione.
	\par\noindent
	$(b)$ Discende dalla linearità dell'integrale.
	\par\noindent
	$(c)$ Evidente.
	\par\noindent
	$(d)$ Vediamo il caso $\abs{\alpha} = 1$, il caso generale può essere ricavato per induzione. Sia $\varphi \in C^\infty_c(U)$, allora
	
	\[ \int_{U} \eta u D^\alpha \varphi d\mathcal{L}^n  =\int_{U} \left( u D^\alpha(\eta \varphi) - u (D^\alpha \eta)\varphi\right) d\mathcal{L}^n =- \int_{U}\left( \eta D^\alpha u + uD^\alpha \eta\right) \varphi d\mathcal{L}^n.  \qedhere \]
	
\end{proof}

\begin{definizione}
	Siano $k \in \N\setminus\left\lbrace 0\right\rbrace $, $1 \le p \le \infty$ ed $u \in W^{k,p}(U)$. Poniamo 
	\[ \norm{u}_{W^{k,p}(U)} = \begin{cases}
		\left( \displaystyle \sum_{\abs{\alpha} \le k} \norm{D^\alpha u}_{L^p(U)}^p\right) ^{\frac{1}{p}} & 1 \le p < \infty, \\
		\displaystyle \sum_{\abs{\alpha} \le k} \norm{D^\alpha u}_{L^\infty(U)} & p=\infty.
	\end{cases} \]
	In particolare, se $k = 1$,\footnote{Alle volte useremo invece la norma equivalente
		\[ \norm{u}_{W^1,p(U)} = \begin{cases}
			\left( \norm{u}_{L^p(U)}^p + \norm{D u}_{L^p(U)}^p\right) ^{\frac{1}{p}} & 1 \le p < \infty, \\
			\norm{u}_{L^\infty(U)} + \norm{D u}_{L^\infty(U)} & p=\infty,
		\end{cases} \]
		dove con $\norm{D u}_p$ si intende, per ogni $p \in [1,\infty]$, la norma della funzione scalare $\abs{Du}$.} 
	\[ \norm{u}_{W^{1,p}(U)} = \begin{cases}
		\left( \norm{u}_{L^p(U)}^p + \displaystyle \sum_{j = 1}^{n}\norm{D_j u}_{L^p(U)}^p\right) ^{\frac{1}{p}} & 1 \le p < \infty, \\
		\norm{u}_{L^\infty(U)} + \displaystyle \sum_{j = 1}^{n}\norm{D_j u}_{L^\infty(U)} & p=\infty. 
	\end{cases} \]
	
\end{definizione}

\begin{teorema}
	Siano $k \in \N\setminus\left\lbrace 0\right\rbrace $ e $1\le p \le \infty$.	Valgono i seguenti fatti:
	\begin{enumerate}[$(a)$]
		\item $(W^{k,p}(U), \norm{\;}_{W^{k,p}(U)})$ è uno spazio normato su $\R$;
		\item $(W^{k,p}(U), \norm{\;}_{W^{k,p}(U)})$ è uno spazio di Banach su $\R$.
	\end{enumerate}
\end{teorema}
\begin{proof} Limitiamo la dimostrazione al caso $k = 1$. Il caso generale si ottiene aggiustando gli indici.
	\par\noindent
	$(a)$ Vediamo innanzitutto il caso $1 \le p < \infty$. Evidentemente $\norm{u}_{1,p} \ge 0$ per ogni $u \in W^{1,p}(U)$ e $\norm{u}_{W^{1,p}(U)} = 0 $ se e solo se $u = 0$ e per ogni $\lambda \in \R$, per ogni $u \in W^{1,p}(U)$, 
	\[ \norm{\lambda u}_{W^{1,p}(U)} = \abs{\lambda} \norm{u}_{W^{1,p}(U)}. \]
	Mostriamo che vale anche la disuguaglianza triangolare. Siano $u,v \in W^{1,p}(U)$, allora per la disuguaglianza di Minkowski nel caso discreto,
	\begin{align*}
		\norm{u+v}_{W^{1,p}(U)} &= \left( \sum_{j = 0}^{n} \norm{D_j u + D_j v}_{L^p(U)}^p\right) ^{\frac{1}{p}} \le \\
		&\le \left( \sum_{j=0}^{n} \left( \norm{D_ju}_{L^p(U)}+ \norm{D_ju}_{L^p(U)}\right) ^p\right)^{\frac{1}{p}} \le \left( \sum_{j = 0}^{n} \norm{D_j u}_{L^p(U)}^p\right) ^{\frac{1}{p}} +  \\
		&+\left( \sum_{j = 0}^{n} \norm{D_j v}_{L^p(U)}^p\right) ^{\frac{1}{p}} = \norm{u}_{W^{1,p}(U)} + \norm{v}_{W^{1,p}(U)}.
	\end{align*}
	Il caso $p = \infty$ può essere trattato in modo simile.
	\par\noindent
	$(b)$ Vediamo innanzitutto il caso $1 \le p < \infty$. Sia $(u_h)$ una successione di Cauchy in $W^{1,p}(U)$, allora per ogni $\varepsilon > 0$, esiste $\bar{h} \in \N$ tale che per ogni $h,k \ge \bar{h}$
	\[ \norm{u_h - u_k}_{W^{1,p}(U)} < \varepsilon. \]
	In particolare, per definizione della norma di Sobolev, fissato $j$,
	\[ \norm{u_h - u_k}_{L^p(U)} < \varepsilon \]
	e 
	\[ \norm{D_ju_h - D_ju_k}_{L^p(U)} < \varepsilon, \]
	da cui $(u_h), (D_ju_h)$ sono di Cauchy in $L^p(U)$. Essendo $L^p(U)$ completo, esistono $u, w_j \in L^p$ tali che $u_h \to u$ in $L^p(U)$ e $D_ju_h \to w_j$ in $L^p(U)$. 
	
	Mostriamo che le $w_j$ sono le derivate deboli di $u$. Siccome $u_h \in W^{1,p}(U)$, per ogni $\varphi \in C^\infty_c(U)$,
	\[ \int_U u_h D_j\varphi d\mathcal{L}^n = - \int_{U}D_j u_h \varphi d\mathcal{L}^n. \]
	Ora, esistono $(u_{h_{k}})$ in $L^p(U)$ e $g \in L^p(U)$ tali che $u_{h_{k}}(x) \to u(x)$ q.o. in $U$ e $\abs{u_{h_{k}}(x)} \le g(x)$ q.o. in $U$ e quindi per il Teorema della convergenza dominata
	\[ \int_U uD_j\varphi d\mathcal{L}^n = - \int_{U}w_j u \varphi d\mathcal{L}^n. \]
	da cui $D_j u = w_j$ per l'unicità della derivata debole.
	
	Il caso $p = \infty$ si tratta in modo simile.
\end{proof}
Dal Teorema precedente risulta in particolare che $(H^{1}(U), \norm{\;}_{1,2})$ è uno spazio di Hilbert. Per gli spazi di Sobolev $W^{k,p}(U)$ valgono le stesse proprietà di riflessività e separabilità delle loro controparti $L^p(U)$.
\begin{definizione}
	Siano $k \in \N\setminus\left\lbrace 0\right\rbrace $ e $1\le p \le \infty$. Chiamiamo $W^{k,p}_{0}(U)$\index[simboli]{$W^{k,p}_{0}(U)$} la chiusura rispetto alla norma di Sobolev di $C^\infty_c(U)$. Più esplicitamente, $u \in W^{k,p}_{0}(U)$ se e solo se esiste $(u_h)$ in $C^\infty_c(U)$ tale che $u_h \to u$ in $ W^{k,p}_{0}(U)$. In particolare $H^{1}_0(U) = W^{1,2}_{0}(U)$\index[simboli]{$H^{1}_0(U)$}.
\end{definizione}
Possiamo intuitivamente interpretare $ W^{k,p}_{0}(U)$ come lo spazio delle applicazioni tali che $D^\alpha u = 0$ su $\partial U$ per ogni $\abs{\alpha} \le k-1$ anche se allo stato attuale non ha senso quanto scritto siccome il bordo ha misura nulla. Questa idea intuitiva verrà resa rigorosa nel seguito.

Concludiamo la sezione con un esempio che contiene una morale profonda: anche se una funzione appartenente ad uno spazio di Sobolev possiede certe proprietà di regolarità, potrebbe comunque comportarsi piuttosto male in altri modi. A voi poi il compito di portare l'analogia fuori dall'ambito dell'Analisi Matematica.
\begin{esempio}
	Sia $U = \textup{B}(0,1) \subseteq \R^n$. Consideriamo $(q_h)$ un sottoinsieme numerabile denso di $U$, ossia $(q_h)$ in $\Q^n \cap U$. Prendiamo $p \in \left] 1,n\right[ $, $\alpha \in \left] 0,\frac{n-p}{p}\right[ $ e 
	\[ u(x) = \sum_{h= 0 }^{+\infty} \frac{2^{-h}}{\abs{x-q_h}^\alpha}. \] 
	Consideriamo $(u_j)$ tale che
	\[ u_j(x) = \sum_{h= 0 }^{j} \frac{2^{-h}}{\abs{x-q_h}^\alpha}. \]
	Chiaramente
	\[ \abs{u_j(x)}^p \to \abs{u(x)}^p \textup{ q.o.}, \]
	\[ 0 \le \abs{u_{j-1}}^p\le \abs{u_{j}}^p \textup{ per ogni  } j \in \N,\]
	quindi dal Teorema di Beppo--Levi
	\[ \norm{u}_{L^p(U)}^p = \int_U \abs{u}^pd\mathcal{L}^n = \lim_j \int_U \abs{u_j}^pd\mathcal{L}^n. \] 
	Ora, per la disuguaglianza di Minkowski,
	\[ \int_U \abs{u_j}^pd\mathcal{L}^n = \norm{u_j}_{L^p(U)}^p \le \left( \sum_{h= 0 }^j \norm*{\frac{2^{-h}}{\abs{x-q_h}^\alpha}}_p \right) ^p = \left( \sum_{h= 0 }^j 2^{-h}\left( \int_{U} \frac{1}{\abs{x-q_h}^{\alpha p}}d\mathcal{L}^n(x)\right) ^{\frac{1}{p}}\right) ^{p}. \]
	A questo punto però, con un cambio di variabile $y = x - q_h$, otteniamo, ricordando che la misura di Lebesgue è invariante per traslazione,
	\begin{align*}
		&\int_{U} \frac{1}{\abs{x-q_h}^{\alpha p}}d\mathcal{L}^n(x) = \int_{\textup{B}(q_h,1-\abs{q_h})} \frac{1}{\abs{x-q_h}^{\alpha p}}d\mathcal{L}^n(x) + \int_{U \setminus \textup{B}(q_h,1-\abs{q_h})} \frac{1}{\abs{x-q_h}^{\alpha p}}d\mathcal{L}^n(x) = \\ 
		&= \int_{\textup{B}(0,1-\abs{q_h})} \frac{1}{\abs{y}^{\alpha p}}d\mathcal{L}^n(y) + \int_{\tilde{U}\setminus \textup{B}(0,1-\abs{q_h})} \frac{1}{\abs{y}^{\alpha p}}d\mathcal{L}^n(y)\le 2 \int_{U} \frac{1}{\abs{y}^{\alpha p}}d\mathcal{L}^n(y)
	\end{align*}
	perciò, essendo $y$ una variabile muta, sfruttando la continuità della funzione $t \mapsto t^p$ e per quanto visto nell'Esempio \eqref{Esempiopalla},
	\[  \norm{u}_{L^p(U)}^p  \le 2 \left( \int_{U} \frac{1}{\abs{x}^{\alpha p}}d\mathcal{L}^n(x)\right) \left( \lim_{j} \sum_{h=0}^{j} 2^{-h}\right) ^p < +\infty  \]
	da cui $u \in L^p(U)$. In modo simile si mostra che $u \in W^{1,p}(U)$.
	
	Tuttavia, preso un qualsiasi compatto $K$ in $U$, esiste per densità un certo $q_{\bar{h}} \in K$. Presa una successione $(x_i)$ in $U\setminus \Q^n$, per non incappare in altri $q_h$, tale che $x_i \to q_{\bar{h}}$, risulta che 
	\[ \lim_{i} u(x_i) = +\infty, \]
	che ci permette di affermare che non esiste alcun compatto $K$ in $U$ tale che $u \in L^\infty(K)$, ossia $u$ è illimitata su ogni compatto di $U$.
\end{esempio}

\section{Rappresentazione di forme lineari e continue}
\begin{lemma}\label{generappresentazioneLp}
	Siano $m \in \N\setminus \left\lbrace 0\right\rbrace $, $p \in \left[ 1,\infty\right[  $ e $T: \left( L^p(U)\right) ^m \to \R$ un'applicazione lineare e continua. Allora esiste $g \in \left(L^{p'}(U) \right)^m $ tale che per ogni $f \in \left( L^p(U)\right) ^m$ si abbia
	\[ \braket{T,f} = \int_U f \cdot g d\mathcal{L}^n = \sum_{i=1}^{m}\int_U f_i g_i d\mathcal{L}^n. \]
\end{lemma}
\begin{proof}
	Si tratta di una variante del caso $m=1$ già noto.
\end{proof}

Il seguente risultato fornisce una rappresentazione degli elementi di $(W^{1,p}(U))'$.

\begin{teorema}
	Siano $p \in \left[ 1,\infty\right[  $ e $T \in (W^{1,p}(U))'$. Allora esistono $f_0,\dots,f_n \in L^{p'}(U)$ tali che per ogni $u \in W^{1,p}(U)$
	\[ \braket{T,u} = \int_U uf_0d\mathcal{L}^n + \sum_{j=1}^{n} \int_U D_j u f_j d\mathcal{L}^n.\]
	In particolare, 
	\[ \norm{T} = \max \left\lbrace \norm{f_0}_{p'},\dots,\norm{f_n}_{p'}\right\rbrace . \]
\end{teorema}
\begin{proof}
	Innanzitutto indichiamo con $X = \left( L^p(U)\right)^{n+1}$ e lo muniamo della norma
	\[ \norm{v}_X = \left( \sum_{j=1}^{n+1}\norm{v_j}_p^p\right) ^{\frac{1}{p}}. \]
	Consideriamo l'operatore lineare $\Pi: W^{1,p}(U) \to X$ tale che
	\[ \Pi u = (u,D_1u,\dots,D_n u).\]
	Chiaramente $\Pi$ è un'isometria. Sia $Y = \Pi(W^{1,p}(U))\le X$. Definiamo l'applicazione lineare $\lambda : Y \to \R$ tale che
	\[ \braket{\lambda, y} = \braket{T,y_1}. \]
	Siccome
	\[ \abs{\braket{\lambda, y}} =\abs{\braket{T,y_1}} \le \norm{T}\norm{y_1}_{1,p} =\footnote{Qua si usa il fatto che $\Pi$ è un'isometria.}  \norm{T}\norm{y}_{Y}, \]
	allora $\lambda$ è continua. Per il Teorema di Hahn--Banach esiste quindi un'applicazione lineare e continua $\Lambda : X \to \R$ tale che $\Lambda|_Y = \lambda$. Per il Lemma \eqref{generappresentazioneLp}, esistono $g_0,\dots,g_n \in L^{p'}(U)$ tali che
	\[ \braket{\Lambda, x} = \sum_{j=0}^{n} \int_U g_j x_j d\mathcal{L}^n. \]
	Se quindi $u \in W^{1,p}(U)$, 
	\[ \braket{T, u} = \braket{\lambda, (u,D_1u,\dots,D_n u)} = \braket{\Lambda, (u,D_1u,\dots,D_n u)} = \int_U ug_0d\mathcal{L}^n + \sum_{j=1}^{n} \int_U D_j u g_j d\mathcal{L}^n.  \]
	
	Omettiamo il calcolo sulla norma di $T$.
\end{proof}

\begin{osservazione}
	Mentre, per il Lemma di Du Bois-Reymond, la rappresentazione degli elementi del duale di $L^p(U)$ è unica, ciò non accade per il duale di $W^{1,p}(U)$.
\end{osservazione}

\begin{osservazione}\label{convergenzadeboleW1,p}
	Se $1 < p < \infty$, consideriamo una successione $(u_h)$ in $W^{1,p}(U)$ limitata. Essendo $W^{1,p}(U)$ riflessivo, per il Teorema di Eberlein--Smulian, esistono $u \in W^{1,p}(U)$ ed una sottosuccessione $u_{h_{k}}$ tali che per ogni $\varphi \in L^{p'}(U)$ e per ogni $j=1,\dots n$ 
	\begin{align*}
		\int_{U}u_{h_{k}}\varphi d\mathcal{L}^n &\to \int_{U}u\varphi d\mathcal{L}^n \\
		\int_{U}D_j u_{h_{k}}\varphi d\mathcal{L}^n &\to \int_{U}D_ju\varphi d\mathcal{L}^n.
	\end{align*}
\end{osservazione}

\begin{notazione}
	Sia $p \in \left[ 1,\infty\right[  $. Indichiamo con $W^{-1,p'}(U)$ il duale topologico di $W^{1,p}_0(U)$. In altre parole, $W^{-1,p'}(U) = (W^{1,p}_0(U))'$. In particolare, $H^{-1}(U) = (H^1_0(U))'$.
\end{notazione}

\section{Approssimazione con funzioni lisce}
\begin{teorema}\emph{\textbf{(sui punti di Lebesgue)}}\index[analitico]{teorema!sui punti di Lebesgue}
	Sia $f \in L^1_{loc}(\R^n)$. Valgono i seguenti fatti:
	\begin{enumerate}[$(a)$]
		\item per \emph{q.o.} $x \in \R^n$, 
		\[ \lim\limits_{r \to 0} \fint_{\textup{B}(x,r)} f d\mathcal{L}^n = f(x);\]
		\item per \emph{q.o.} $x \in \R^n$, 
		\[ \lim\limits_{r \to 0} \fint_{\textup{B}(x,r)} \abs{f(y)- f(x)}d\mathcal{L}^n(y)= 0.\]
	\end{enumerate}
\end{teorema}
\begin{proof}
	Omettiamo la dimostrazione.
\end{proof}
\begin{definizione}
	I punti $x$ tali che soddisfano il Teorema sui punti di Lebesgue sono detti \emph{punti di Lebesgue di $f$}\index[analitico]{punto!di Lebesgue}.
\end{definizione}
\begin{notazione}
	Dato $\varepsilon>0$, indichiamo con $U_{\varepsilon}$, l'insieme
	\[U_{\varepsilon} = \Set{ x \in U: \textup{dist}(x,\partial U) > \varepsilon}.\]
\end{notazione}
\begin{definizione}\index[analitico]{mollificatori}
	Chiamiamo \emph{mollificatore standard}, l'applicazione $\varrho \in C^\infty(\R^n)$ tale che
	\[ \varrho(x) = \begin{cases}
		C e^{\frac{1}{\abs{x}^2-1}} & \textup{se } \abs{x}^2 < 1, \\
		0 & \textup{altrimenti},
	\end{cases} \]
	dove $C > 0$ tale che 
	\[ \int \varrho d\mathcal{L}^n = 1.\]
	Per ogni $\varepsilon > 0$, chiamiamo poi $\varepsilon$\emph{-mollificatore} l'applicazione $\varrho_\varepsilon \in C^\infty_c(\textup{B}(0,\varepsilon))$ tale che
	\[ \varrho_\varepsilon(x) = \frac{1}{\varepsilon^n} \varrho\left( \frac{x}{\varepsilon}\right) .\footnote{Chiaramente, per ogni $\varepsilon> 0 $
		\[ \int \varrho_\varepsilon d\mathcal{L}^n = 1. \]} \]
\end{definizione}
\begin{definizione}
	Sia $u \in L^1_{loc}(U)$ ed $\varepsilon > 0$. Chiamiamo \emph{$\varepsilon$-mollificata di $u$}, l'applicazione definita su $U_\varepsilon$ da
	\[ u^\varepsilon = \varrho \ast u,\]
	ossia tale che per ogni $x \in U_{\varepsilon}$
	\[ u^\varepsilon(x) = \int_{U} \varrho_\varepsilon(x-y)u(y)d\mathcal{L}^n(y) = \int_{\textup{B}(0,\varepsilon)} \varrho_\varepsilon(y) u(x-y)d\mathcal{L}^n(y). \]
\end{definizione}
\begin{proposizione}\label{propmollificatori}
	Sia $u \in L^1_{loc}(U)$ ed $\varepsilon >0$. Valgono i seguenti fatti:
	\begin{enumerate}[$(a)$]
		\item $u^\varepsilon \in C^\infty(U_\varepsilon)$ e per ogni $j =1,\dots,n$ si ha
		\[ D_j u^\varepsilon (x) = \int_{U} u(y)(D_{x_j}\rho_\varepsilon)(x-y)d\mathcal{L}^n(y);\]
		\item $u^\varepsilon(x) \to u(x)$ \emph{q.o.} in $U_\varepsilon$;
		\item se $u \in C(U)$, allora $u^\varepsilon \to u$ per $\varepsilon \to 0$,  uniformemente sui compatti di $U$;
		\item se $1 \le p < \infty$ e $u \in L^p_{loc}(U)$, allora $u^\varepsilon \to u$ in $L^p_{loc}(U)$ per $\varepsilon \to 0$.
	\end{enumerate}
\end{proposizione}
\begin{proof}\hfill
	\par\noindent 
	$(a)$ La dimostrazione può essere svolta per esercizio.
	\par\noindent
	$(b)$ Per il Teorema di differenziabilità di Lebesgue, per q.o. $x \in U_{\varepsilon}$,
	\begin{align*}
		\abs{u_\varepsilon(x) - u(x)} &\le \frac{1}{\varepsilon^n} \int_{\textup{B}(x,\varepsilon)} \varrho\left( \frac{x-y}{\varepsilon}\right) \abs{u(y)-u(x)}d\mathcal{L}^n(y) \le \\
		&\le C \fint_{\textup{B}(x,\varepsilon)} \abs{u(y)-u(x)}d\mathcal{L}^n(y) \to 0
	\end{align*}
	per $\varepsilon \to 0$. 
	\par\noindent
	$(c)$ La dimostrazione può essere svolta per esercizio.
	\par\noindent
	$(d)$ Innanzitutto, dal fatto che $u \in L^p_{loc}(U)$, per ogni $K$ compatto in $U$ si ha $u \in L^p(K)$. Ora, esiste $W$ compatto in $U$ tale che $K \subseteq \textup{int}(W)$. Mostriamo che se $\varepsilon$ è tale che, per ogni $x \in K$, $\textup{B}(x,\varepsilon)\subseteq W$, risulta $\norm{u^\varepsilon}_{L^p(K)} \le \norm{u}_{L^p(W)}$. Per ogni $x \in K$, per la disuguaglianza di Holder, 
	\begin{align*}
		\abs{u^\varepsilon(x)} &\le \int_{\textup{B}(x,\varepsilon)} \varrho_\varepsilon(x-y)\abs{u(y)}d\mathcal{L}^n(y) = \int_{\textup{B}(x,\varepsilon)}\varrho^{\frac{1}{p'}}_\varepsilon(x-y)\varrho^{\frac{1}{p}}_\varepsilon(x-y)\abs{u(y)}d\mathcal{L}^n(y) \le \\
		&\le \left( \int_{\textup{B}(x,\varepsilon)}\varrho_\varepsilon(x-y)\abs{u(y)}^pd\mathcal{L}^n(y)\right) ^{\frac{1}{p}}.
	\end{align*} 
	Ora, per il Teorema di Fubini--Tonelli,
	\begin{align*}
		\int_K \abs{u^\varepsilon(x)}^pd\mathcal{L}^n(x) &\le \int_K \int_{\textup{B}(x,\varepsilon)}\varrho_\varepsilon(x-y)\abs{u(y)}^pd\mathcal{L}^n(y) d\mathcal{L}^n(x) \le \\
		& \le \int_K \int_{W}\varrho_\varepsilon(x-y)\abs{u(y)}^pd\mathcal{L}^n(y)  d\mathcal{L}^n(x) \le \\
		& \le \int_W \int_{K}\varrho_\varepsilon(x-y)\abs{u(y)}^pd\mathcal{L}^n(x)  d\mathcal{L}^n(y) = \norm{u}_{L^p(W)}^p.
	\end{align*}
	
	Dato ora $\delta > 0$, esiste, per densità, $g \in C_c(W)$ tale che 
	\[ \norm{g-u}_{L^p(W)} < \frac{\delta}{3} \]
	ed esiste $\bar{\varepsilon}>0$ tale che per ogni $\varepsilon \in \left] 0,\bar{\varepsilon}\right] $
	\[ \norm{g^\varepsilon - g }_{L^p(K)} <   \frac{\delta}{3}.\]
	
	Siccome $\norm{g^\varepsilon - u^\varepsilon}_{L^p(K)} \le \norm{g - u}_{L^p(W)}$, risulta che 
	\[ \norm{u^\varepsilon - u}_{L^p(K)} < \delta.\qedhere \]
\end{proof}
\begin{teorema}\emph{\textbf{(di locale approssimazione con funzioni lisce)}}\label{locappr}
	Siano $k \in \N\setminus\left\lbrace 0\right\rbrace $, $1\le p < \infty$ e $u \in W^{k,p}(U)$. Valgono i seguenti fatti:
	\begin{enumerate}[$(a)$]
		\item per ogni $\varepsilon > 0$, $u^\varepsilon \in C^\infty(U_\varepsilon)$;
		\item per $\varepsilon \to 0$, $u^\varepsilon \to u$ in $W^{k,p}_{loc}(U)$.
	\end{enumerate}
\end{teorema}
\begin{proof}\hfill
	\par\noindent
	$(a)$ Discende direttamente dalla $(a)$ della Proposizione \eqref{propmollificatori}.
	\par\noindent
	$(b)$ Limitiamo la dimostrazione al caso $k = 1$. Il caso generale si ottiene aggiustando gli indici. 
	
	Per la $(a)$ della Proposizione \eqref{propmollificatori} combinata con la definizione di derivata debole,
	\begin{align*}
		D_j u^\varepsilon(x) &= \int_{U} D_{x_j}\varrho_\varepsilon(x-y) u(y) d\mathcal{L}^n(y) =\\
		&= - \int_{U} D_{y_j}\varrho_\varepsilon(x-y) u(y) d\mathcal{L}^n(y) = \int_{U} \varrho_\varepsilon(x-y) D_{y_j}u(y) d\mathcal{L}^n(y)
	\end{align*}
	da cui,
	\[ D_j u^\varepsilon = \varrho_\varepsilon \ast D_j u.\footnote{Per quanto riguarda $u$, la derivata è da intendere in senso debole.} \]
	
	Ora, dal fatto che $u \in W^{1,p}(U)$, per la $(d)$ della Proposizione \eqref{propmollificatori}, $u^\varepsilon \to u$ in $L^p_{loc}(U)$. Analogamente, per quanto visto sopra, anche $D_ju^\varepsilon \to D_ju$ in $L^p_{loc}(U)$. Allora per ogni compatto $K \subseteq U$
	\[ \norm{u^\varepsilon -u}_{W^{1,p}(K)}^p = \norm{u^\varepsilon -u}_{L^{p}(K)} + \sum_{j=1}^{n}\norm{D_ju^\varepsilon -D_ju}_{L^{p}(K)}<\varepsilon.\qedhere \]
\end{proof}

\begin{definizione}
	Sia $(X,\tau)$ uno spazio topologico. Consideriamo una successione $(\xi_i)$ di applicazioni continue da $X$ in $\R$. Diciamo che $(\xi_i)$ è una \emph{partizione dell'unità}\index[analitico]{partizione dell'unità}, se valgono i seguenti fatti:
	\begin{enumerate}[$(a)$]
		\item per ogni $i \in \N$, $0 \le \xi_i \le 1$,
		\item per ogni $x \in X$, solo un numero finito di $\xi_i(x)$ è non nullo,
		\item per ogni $x \in X$,
		\[ \sum_{i=0}^{\infty} \xi_i(x) = 1. \]
	\end{enumerate}
\end{definizione}

\begin{osservazione}
	La somma introdotta nella $(c)$ è finita in ogni punto per la $(b)$.
\end{osservazione}

\begin{proposizione}
	Sia $(X,\tau)$ uno spazio topologico e $(V_i)$ un suo ricoprimento aperto. Allora esiste una partizione dell'unità $(\xi_i)$ tale che $\textup{supt}(\xi_i) \subseteq V_i$.
\end{proposizione}
\begin{proof}
	Omettiamo la dimostrazione.
\end{proof}

In particolare, se $X = U$, $\xi_i \in C^\infty_c(V_i)$.

\begin{teorema}\emph{\textbf{(di globale approssimazione con funzioni lisce)}}\label{globappr}
	Siano $U$ limitato, $k \in \N\setminus\left\lbrace 0\right\rbrace $, $1\le p < \infty$ e $u \in W^{k,p}(U)$. Allora esiste $(u_h)$ in $C^\infty(U) \cap W^{k,p}(U)$ tale che $u_h \to u$ in $W^{k,p}(U)$.
\end{teorema}
\begin{proof}
	Limitiamo la dimostrazione al caso $k=1$. Per ogni $i \in \N\setminus \left\lbrace 0\right\rbrace  $, sia
	\[ U_i = \Set{ x \in U: \textup{dist}(x,\partial U) >\frac{1}{i}}. \]
	Evidentemente,
	\[  U = \bigcup_{i=1}^{+\infty} U_i.\]
	
	Consideriamo per ogni $i \in \N\setminus \left\lbrace 0\right\rbrace  $, gli insiemi $V_i = U_{i+3} \setminus \overline{U}_{i+1}$, evidentemente aperti, ed un aperto $V_0$ in $U$ tale che $\overline{V}_0 \subseteq U$ e
	\[ U =  \bigcup_{i=0}^{+\infty} V_i.\]
	
	Sia $(\xi_i)$ una partizione dell'unità relativa a $(V_i)$. Per la Proposizione \eqref{propderivatedeboli}, $\xi_iu \in W^{1,p}(U)$. Inoltre $\textup{supt}(\xi_iu) \subseteq V_i$.
	
	Fissiamo $\delta > 0$. Per il Teorema \eqref{locappr}, per ogni $V$ tale che $\overline{V} \subseteq U$ esiste $\varepsilon_i > 0$ tale che, detto
	\[ u^i = \varrho_{\varepsilon_{i}} \ast (\xi_i u), \]
	si abbia
	\begin{align*}
		&\norm{u^i - \xi_iu}_{W^{1,p}(V)} < \frac{\delta}{2^{i+1}}  & \textup{ per } i &\in \N, \\
		&\textup{supt}(u^i) \subseteq W_i = U_{i+4} \setminus \overline{U}_{i} & \textup{ per } i &\in \N\setminus\left\lbrace 0\right\rbrace.
	\end{align*}
	Chiaramente $V_i \subseteq W_i $ per ogni $i \in \N\setminus\left\lbrace 0\right\rbrace$.
	
	Posto 
	\[ v = \sum_{i=0}^{\infty} u^i, \]
	si ha che $v \in C^\infty_c(U)$, ma dal fatto che 
	\[ \sum_{i=0}^\infty \xi_i(x) = 1, \]
	possiamo scrivere 
	\[ u = \sum_{i=0}^\infty \xi_i u \]
	e quindi per ogni $V$ tale che $\overline{V}\subseteq U$, si ha
	\[ \norm{u-v}_{W^{1,p}(V)} \le \sum_{i=0}^\infty \norm{\xi_iu-u^i}_{W^{1,p}(V)} \le \sum_{i=0}^\infty \frac{\delta}{2^{i+1}} = \delta. \]
	Posto ora $\delta = \frac{1}{m+1}$, per ogni $m \in \N$ è sufficiente porre $u_m = v.$ La tesi segue passando al sup su $V$. 
\end{proof}

\begin{osservazione}
	Nel Teorema \eqref{globappr}, non si afferma che $u_h \in C^\infty(\overline{U})$.
\end{osservazione}

\begin{teorema}\emph{\textbf{(di globale approssimazione fino al bordo con funzioni lisce)}}\label{globapprbordo}
	Siano $U$ limitato con $\partial U$ di classe $C^1$, $k \in \N\setminus\left\lbrace 0\right\rbrace $, $1\le p < \infty$ e $u \in W^{k,p}(U)$. Allora esiste $(u_h)$ in $C^\infty(\overline{U})$ tale che $u_h \to u$ in $W^{k,p}(U)$.
\end{teorema}
\begin{proof}
	Omettiamo la dimostrazione.
\end{proof}

\begin{definizione}\index[analitico]{rapporto incrementale}\index[simboli]{$D_{i}^{h}$}\index[simboli]{$D^{h}$}
	Siano $u:U\to\R$ un'applicazione localmente sommabile, $V$ tale che $\overline{V} \subseteq U$, $i = 1,\dots,n$ ed $h \in \R$ tale che $0 < \abs{h}< \textup{dist}(V,\partial U)$. Chiamiamo \emph{$i$-esimo rapporto incrementale di ampiezza $h$}\index[analitico]{rapporto incrementale}\index[simboli]{$D_{i}^{h}$} l'applicazione $D_{i}^{h}u : V \to \R$ tale che
	\[D_{i}^{h}u(x) = \dfrac{u(x+he_i)- u(x)}{h}.\]
	In particolare, poniamo $ D^{h}u = \left( D^h_1u,\dots,D^h_nu\right)$.  
\end{definizione}

\begin{teorema}\textbf{\emph{(sui rapporti incrementali)}}
	Valgono i seguenti fatti:
	\begin{enumerate}[$(a)$]
		\item se $1 \le p < \infty$ e $u \in W^{1,p}(U)$, allora per ogni $V$ tale che $\overline{V} \subseteq U$, esiste $C > 0$ tale che per ogni $h \in \R$ tale che $0 < \abs{h}< \frac{1}{2}\textup{dist}(V,\partial U)$ risulti 
		\[ \norm{D^{h}u}_{L^p(V)} \le C \norm{Du}_{L^p(U)}, \]
		\item se $1 < p <\infty$, $V$ tale che $\overline{V} \subseteq U$, $u \in L^p(V)$ ed esiste $C>0$ tale che per ogni $h \in \R$ tale che $0 < \abs{h}< \frac{1}{2}\textup{dist}(V,\partial U)$
		\[ \norm{D^hu}_{L^p(V)} \le C, \]
		allora $u \in W^{1,p}(V)$ e $\norm{Du}_{L^p(V)} \le C$. 
	\end{enumerate}
	\begin{proof}\hfill
		\par\noindent
		$(a)$ Supponiamo inizialmente $u$ liscia. Allora per il Teorema fondamentale del calcolo integrale
		\[ u(x+he_i) - u(x) = \int_{0}^{1} D_{x_i}u(x+the_i)dt \cdot he_i,\]
		da cui 
		\[ D^h_iu(x) = C \int_{0}^{1} D_{x_i}u(x+the_i) \cdot e_i dt \]
		e quindi 
		\[ \abs{ D^h_iu(x) } \le C \int_{0}^{1} \abs{Du(x+the_i)} dt.\]
		Applicando la disuguaglianza di Holder, otteniamo 
		\[ \abs{ D^h_iu(x) }^p \le C\int_{0}^{1} \abs{Du(x+the_i)}^p dt\]
		che integrata su $V$ restituisce
		\[ \int_V \abs{ D^h_iu }^p d\mathcal{L}^n\le C\int_V \int_{0}^{1} \abs{Du(x+the_i)}^p dt d\mathcal{L}^n(x).\]
		Tramite il cambio di variabile $y = x+the_i$, possiamo scrivere
		\[ \int_V \abs{ D^h_iu }^p d\mathcal{L}^n\le C\int_{\tilde{V}} \int_{0}^{1} \abs{Du}^p dt d\mathcal{L}^n \le C\int_U \int_{0}^{1} \abs{Du}^p dt d\mathcal{L}^n = C\int_U \abs{Du}^p d\mathcal{L}^n.\]
		Se invece $u \in W^{1,p}(U)$, costruiamo una successione di funzioni lisce che la approssima e poi con il Teorema della convergenza dominata otteniamo ancora la stessa stima. In ogni caso,
		\[ \norm{D^{h}u}_{L^p(V)} \le C \norm{Du}_{L^p(U)}. \]
		\par\noindent
		$(b)$ Sia $\varphi \in C^\infty_c(V)$. Osserviamo che 
		\[ \int_V u(x) \dfrac{\varphi(x+he_i) - \varphi(x)}{h} d\mathcal{L}^n(x) = - \int_V \dfrac{u(x)-u(x-he_i)}{h} \varphi(x) d\mathcal{L}^n(x)\]
		ossia 
		\begin{equation}\label{integrazioneperpartirapportoincrementale}
			\int_V u D^h_i\varphi d\mathcal{L}^n =-\int_V D^{-h}_iu \varphi d\mathcal{L}^n.
		\end{equation}
		Ora, siccome 
		\[ \underset{0 < \abs{h}< \frac{1}{2}\textup{dist}(V,\partial U)}{\textup{sup}} \norm{D^hu}_{L^p(V)} \le C, \]
		considerando una successione $(h_k)$ tale che $h_k \to 0$ possiamo scrivere che definitivamente
		\[ \norm{D^{-h_k}_iu}_{L^p(V)} \le C. \]
		Essendo $L^p(V)$ riflessivo, per il Teorema di Eberlein--Smulian esiste $w_i \in L^p(V)$ tale che, a meno di sottosuccessioni, $D^{-h_k}_iu \rightharpoonup w_i$ in $L^p(V)$. Passando al limite in \eqref{integrazioneperpartirapportoincrementale}, tramite il Teorema della convergenza dominata, otteniamo 
		\[ \int_V uD_i \varphi d\mathcal{L}^n = -\int_V w_i \varphi d\mathcal{L}^n \]
		da cui $w_i = D_i u \in L^p(V)$. Allora $u \in W^{1,p}(V)$ e 
		\[ \int_V \abs{Du}^pd\mathcal{L}^n \le \underset{k}{\textup{lim inf}} \int_V \abs{D^{-h_k}u}^pd\mathcal{L}^n \le C.\qedhere\]
	\end{proof}
\end{teorema}
\section{Estensioni}
Nel seguito, salvo diversa specificazione, ci limiteremo a studiare il caso $k=1$.

Data $u \in L^p(U)$, è sempre possibile estenderla a tutto $\R^n$ considerando il suo prolungamento triviale. Questo procedimento non può però essere in generale applicato anche ad applicazioni in $W^{1,p}(U)$ in quanto dobbiamo anche preservare l'esistenza delle derivate deboli su $\partial U$. Un caso però in cui il prolungamento triviale restituisce un'estensione in $W^{1,p}(\R^n)$ è il seguente.
\begin{proposizione}
	Sia $u \in W^{1,p}_0(U)$. Allora, posto
	\[ \tilde{u} = \begin{cases}
		u & \textup{in } U, \\
		0 & \textup{in } \R^n\setminus U,
	\end{cases} \]
	si ha $\tilde{u} \in W^{1,p}(\R^n)$.
\end{proposizione}
\begin{proof}
	Consideriamo 
	\[ v = \begin{cases}
		D_j u & \textup{se } x \in U,\\
		0 & \textup{se } x \in \R^n\setminus U.
	\end{cases} \]
	Evidentemente $v \in L^p(U)$, inoltre per ogni $\varphi \in C^\infty_c(U)$,
	
	\[ \int \tilde{u} D_j \varphi d\mathcal{L}^n = \int_{U} u D_j \varphi d\mathcal{L}^n. \]
	Ora, siccome $u \in W^{1,p}_0(U)$, esiste $(\psi_h) $ in $C^\infty_c(U)$ tale che $\psi_h \to u$ in $W^{1,p}_0(U)$, ed in particolare in $L^p(U)$. Allora, per il Teorema della convergenza dominata, a meno di sottosuccessioni,
	\[ \int \tilde{u} D_j \varphi d\mathcal{L}^n = \int_{U} \lim_h \psi_h D_j \varphi d\mathcal{L}^n = \lim_h  \int_{U} \psi_h D_j \varphi d\mathcal{L}^n\]
	ma sull'ultimo integrale è possibile applicare la formula di Gauss--Green, in quanto $\psi_h \in C^\infty_c(U)$, per cui
	\[ \int \tilde{u} D_j \varphi d\mathcal{L}^n = - \lim_h \int_U D_j \psi_h \varphi d\mathcal{L}^n.\]
	Sempre dal fatto che $\psi_h \to u$ in $W^{1,p}_0(U)$, possiamo anche dedurre, tramite il Teorema della convergenza dominata, che $D_j \psi_h \to D_j u$ in $L^p(U)$ e quindi
	\[  \int \tilde{u} D_j \varphi d\mathcal{L}^n = - \int_U D_j u \varphi d\mathcal{L}^n = - \int v \varphi d\mathcal{L}^n. \]
	Allora $D_j \tilde{u} = v$.
\end{proof}

\begin{teorema}\emph{\textbf{(di estensione)}}\index[analitico]{teorema!di estensione}
	Siano $1 \le p \le \infty$, $U$ limitato con $\partial U$ di classe $C^1$. Allora per ogni $V$ aperto limitato tale che $\overline{U} \subseteq V$ esiste un operatore $E: W^{1,p}(U) \to W^{1,p}(\R^n)$ lineare e continuo tale che per ogni $u \in W^{1,p}(U)$ valgano i seguenti fatti:
	\begin{enumerate}[$(a)$]
		\item $Eu (x) = u(x)$ \emph{q.o.} in $U$,
		\item $Eu = 0$ \emph{q.o.} in $\R^n\setminus V$, ossia $Eu$ ha supporto in $V$,
		\item esiste $C$, dipendente solo da $U,V,n$ e $p$, tale che
		\[ \norm{Eu}_{W^{1,p}(\R^n)} \le C \norm{u}_{W^{1,p}(U)}. \]
	\end{enumerate}
\end{teorema}
\begin{proof}
	Limitiamo la dimostrazione al caso $1\le p < \infty$. Vediamo dapprima il caso $u \in C^1(\overline{U})$. Fissiamo $x_0 \in \partial U$ e supponiamo che, in un intorno di $x_0$, $\partial U$ giaccia sul piano $x_n = 0$, ossia che esista $r > 0 $ tale che, detta $\textup{B} = \textup{B}(x_0,r)$,
	\begin{align*}
		\textup{B}^+ &:= \textup{B}(x_0,r) \cap \left\lbrace x_n \ge 0\right\rbrace \subseteq \overline{U}, \\
		\textup{B}^- &:= \textup{B}(x_0,r) \cap \left\lbrace x_n \le 0\right\rbrace \subseteq \R^n \setminus U.
	\end{align*}
	Consideriamo l'applicazione\footnote{Detta riflessione di ordine superiore da $\textup{B}^+$ a $\textup{B}^-$.}
	\[ \bar{u}(x) = \begin{cases}
		u\left( x\right)  & \textup{se } x \in \textup{B}^+,\\
		-3u\left( x_1,\dots, x_{n-1}, -x_n\right)  + 4u\left( x_1,\dots, x_{n-1},-\frac{x_n}{2}\right)  & \textup{se } x \in \textup{B}^-.
	\end{cases} \]
	Mostriamo che $\bar{u} \in C^1(\textup{B})$. La continuità di $\bar{u}$ è evidente. Se $j = 1,\dots,n-1$ si ha
	\[ \frac{\partial \bar{u}}{\partial x_j}(x) = \begin{cases}
		\frac{\partial u}{\partial x_j}(x) & \textup{se } x \in \textup{B}^+,\\
		-3\frac{\partial u}{\partial x_j}(x_1, \dots, x_{n-1},-x_n) + 4\frac{\partial u}{\partial x_j}\left( x_1,\dots, x_{n-1},-\frac{x_n}{2}\right)   & \textup{se } x \in \textup{B}^-
	\end{cases} \]
	e si vede subito che $\frac{\partial \bar{u}}{\partial x_j}$ esiste ed è continua. Invece
	\[ \frac{\partial \bar{u}}{\partial x_n}(x) = \begin{cases}
		\frac{\partial u}{\partial x_n}(x) & \textup{se } x \in \textup{B}^+,\\
		3\frac{\partial u}{\partial x_n}(x_1, \dots, x_{n-1},-x_n) - 2\frac{\partial u}{\partial x_n}(x_1, \dots, x_{n-1},-\frac{x_n}{2})  & \textup{se } x \in \textup{B}^-,
	\end{cases} \]
	da cui $\frac{\partial \bar{u}}{\partial x_n}$ esiste ed è continua. Per il Teorema del differenziale totale possiamo affermare che $\bar{u} \in C^1(\textup{B})$. 
	
	Mostriamo ora che esiste $C > 0$ tale che
	\[ \norm{\bar{u}}_{W^{1,p}(\textup{B})} \le C \norm{u}_{W^{1,p}(\textup{B}^+)}. \]
	Innanzitutto, 
	\begin{align*}
		\norm{\bar{u}}_{W^{1,p}(\textup{B})}^p &= \norm{\bar{u}}_{L^{p}(\textup{B})}^p + \sum_{j=1}^{n}\norm{D_j\bar{u}}_{L^{p}(\textup{B})}^p = \\
		& = \int_{\textup{B}^+} \abs{u}^p  d\mathcal{L}^n +  \sum_{j=1}^{n}\int_{\textup{B}^+}\abs{D_ju}^p d\mathcal{L}^n + \int_{\textup{B}^-} \abs{\bar{u}}^p  d\mathcal{L}^n +  \sum_{j=1}^{n}\int_{\textup{B}^-}\abs{D_j\bar{u}}^p d\mathcal{L}^n.
	\end{align*}  
	Ora, 
	\begin{align*}
		\int_{\textup{B}^-}\abs{\bar{u}}^pd\mathcal{L}^n &\le 2^{p-1} \left[  3^p\int_{\textup{B}^-} \abs{u(x_1,\dots,x_{n-1},-x_n)}^pd\mathcal{L}^n(x_1,\dots,x_n) + \right. \\
		&\left. + 4^p\int_{\textup{B}^-} \abs*{u\left( x_1,\dots, x_{n-1},-\frac{x_n}{2}\right) }^pd\mathcal{L}^n(x_1,\dots,x_n) \right] 
	\end{align*}
	dove, considerando il cambio di variabile, con jacobiano $1$,
	\[ (x_1,\dots,x_n) \mapsto (x_1,\dots,-x_n) \]
	si ottiene facilmente che 
	\[ \int_{\textup{B}^-} \abs{u(x_1,\dots,x_{n-1},-x_n)}^pd\mathcal{L}^n(x_1,\dots,x_n) = \norm{u}_{L^p(\textup{B}^+)}^p \]
	e considerando il cambio di variabile, con jacobiano $2$,
	\[ 	(x_1,\dots,x_n) \mapsto \left( x_1,\dots,-\frac{x_n}{2}\right) \]
	si ottiene analogamente che
	\[ \int_{\textup{B}^-} \abs*{u\left( x_1,\dots, x_{n-1},-\frac{x_n}{2}\right) }^pd\mathcal{L}^n(x_1,\dots,x_n) \le 2\norm{u}_{L^p(\textup{B}^+)}^p, \]
	perciò abbiamo che 
	\[ 	\int_{\textup{B}^-}\abs{\bar{u}}^pd\mathcal{L}^n \le 2^{p-1}\left( 3^p+2^{2p+1}\right)\norm{u}_{L^p(\textup{B}^+)}^p.  \]
	Invece, con passaggi analoghi, si mostra che 
	\[ \sum_{j=1}^{n}\int_{\textup{B}^-}\abs{D_j\bar{u}}^p d\mathcal{L}^n \le 2^{p-1}\left( 3^p + 2^{2p+1}+2^{p+2}\right) \sum_{j=1}^{n} \norm{D_j u}_{L^p(\textup{B}^+)}^{p}. \]
	In conclusione
	\[ \norm{\bar{u}}_{W^{1,p}(\textup{B})}^p \le \left( 1 + 2^{p-1}(3^{p}+2^{2p+2})\right)  \norm{u}_{W^{1,p}(\textup{B}^+)}^p \]
	e prendendo la radice $p$-esima si ha la disuguaglianza cercata. 
	
	Consideriamo ora il caso in cui $\partial U$ non sia necessariamente piatto vicino ad $x_0$. Consideriamo il diffeomorfismo\footnote{La Geometria Differenziale ce ne garantisce l'esistenza e non ce ne preoccupiamo.} $\Phi$, con inversa $\Psi$, che appiattisce $\partial U$ in un intorno di $x_0$. Introduciamo le notazioni $y = \Phi(x)$, $x = \Psi(y)$ e $u'(y) = u(\Psi(y))$. Operando come nel passo precedente otteniamo un'estensione di $u'$ da $\textup{B}^+$ a $\textup{B}$ tale che esiste $C > 0 $ che fornisce la stima
	\[ \norm{\bar{u}'}_{W^{1,p}(\textup{B})} \le C \norm{u'}_{W^{1,p}(\textup{B}^+)}. \]
	
	\begin{center}
		\begin{tikzpicture}[x=0.65pt,y=0.65pt,yscale=-1,xscale=1]
			%uncomment if require: \path (0,300); %set diagram left start at 0, and has height of 300
			
			%Shape: Ellipse [id:dp0023267519166336736] 
			\draw   (457.13,131.8) .. controls (457.13,91.37) and (489.98,58.6) .. (530.5,58.6) .. controls (571.02,58.6) and (603.87,91.37) .. (603.87,131.8) .. controls (603.87,172.23) and (571.02,205) .. (530.5,205) .. controls (489.98,205) and (457.13,172.23) .. (457.13,131.8) -- cycle ;
			%Shape: Polygon Curved [id:ds9277975115378949] 
			\draw   (127,90.6) .. controls (147,80.6) and (162,72.6) .. (211,72.6) .. controls (260,72.6) and (225,112.6) .. (246,128.6) .. controls (267,144.6) and (212.41,187.27) .. (185,184.6) .. controls (157.59,181.93) and (110,187.6) .. (90,157.6) .. controls (70,127.6) and (107,100.6) .. (127,90.6) -- cycle ;
			%Curve Lines [id:da9936244162893775] 
			\draw    (28,169.6) .. controls (18,159.6) and (93.25,142.66) .. (111,150.6) .. controls (128.75,158.54) and (181.67,136.96) .. (193.74,119.69) .. controls (205.81,102.42) and (233,48.6) .. (197,25.6) ;
			
			%Straight Lines [id:da09958792944365125] 
			\draw    (429,130.6) -- (632,130.58) ;
			
			%Curve Lines [id:da5907621652570023] 
			\draw    (259,214.8) .. controls (343.15,263.71) and (409.66,245.18) .. (449.79,215.7) ;
			\draw [shift={(451,214.8)}, rotate = 143.13] [color={rgb, 255:red, 0; green, 0; blue, 0 }  ][line width=0.75]    (10.93,-3.29) .. controls (6.95,-1.4) and (3.31,-0.3) .. (0,0) .. controls (3.31,0.3) and (6.95,1.4) .. (10.93,3.29)   ;
			%Curve Lines [id:da5283825799891098] 
			\draw    (449,54.4) .. controls (412.37,12.82) and (301.25,11.82) .. (262.16,44.79) ;
			\draw [shift={(261,45.8)}, rotate = 318.18] [color={rgb, 255:red, 0; green, 0; blue, 0 }  ][line width=0.75]    (10.93,-3.29) .. controls (6.95,-1.4) and (3.31,-0.3) .. (0,0) .. controls (3.31,0.3) and (6.95,1.4) .. (10.93,3.29)   ;
			%Shape: Path Data [id:dp48216478983696454] 
			\draw  [fill={rgb, 255:red, 155; green, 155; blue, 155 }  ,fill opacity=0.5 ] (211,72.6) .. controls (211.97,72.6) and (212.9,72.62) .. (213.81,72.65) .. controls (210.35,91.89) and (200.01,110.72) .. (193.74,119.69) .. controls (181.67,136.96) and (128.75,158.54) .. (111,150.6) .. controls (106.04,148.38) and (96.59,148.1) .. (85.76,149.08) .. controls (76.93,122.59) and (108.92,99.64) .. (127,90.6) .. controls (147,80.6) and (162,72.6) .. (211,72.6) -- cycle ;
			%Shape: Path Data [id:dp8186232193666636] 
			\draw  [fill={rgb, 255:red, 155; green, 155; blue, 155 }  ,fill opacity=0.5 ] (530.5,58.6) .. controls (570.62,58.6) and (603.21,90.72) .. (603.86,130.58) -- (457.14,130.6) .. controls (457.78,90.72) and (490.38,58.6) .. (530.5,58.6) -- cycle ;
			
			%Shape: Circle [id:dp6915413941378992] 
			\draw  [fill={rgb, 255:red, 0; green, 0; blue, 0 }  ,fill opacity=1 ] (175,132.6) .. controls (175,130.39) and (176.79,128.6) .. (179,128.6) .. controls (181.21,128.6) and (183,130.39) .. (183,132.6) .. controls (183,134.81) and (181.21,136.6) .. (179,136.6) .. controls (176.79,136.6) and (175,134.81) .. (175,132.6) -- cycle ;
			
			%Shape: Circle [id:dp5855646505789538] 
			\draw  [fill={rgb, 255:red, 0; green, 0; blue, 0 }  ,fill opacity=1 ] (526.5,130.59) .. controls (526.5,128.38) and (528.29,126.59) .. (530.5,126.59) .. controls (532.71,126.59) and (534.5,128.38) .. (534.5,130.59) .. controls (534.5,132.8) and (532.71,134.59) .. (530.5,134.59) .. controls (528.29,134.59) and (526.5,132.8) .. (526.5,130.59) -- cycle ;
			
			% Text Node
			\draw (534,103.4) node [anchor=north west][inner sep=0.75pt]    {$y_{0}$};
			% Text Node
			\draw (345,38.4) node [anchor=north west][inner sep=0.75pt]  [font=\large]  {$\Psi $};
			% Text Node
			\draw (342,210.4) node [anchor=north west][inner sep=0.75pt]  [font=\large]  {$\Phi $};
			% Text Node
			\draw (188,124) node [anchor=north west][inner sep=0.75pt]    {$x_{0}$};
			% Text Node
			\draw (94,78.4) node [anchor=north west][inner sep=0.75pt]    {$W$};
			% Text Node
			\draw (600,80) node [anchor=north west][inner sep=0.75pt]    {B};
			% Text Node
			\draw (500,90) node [anchor=north west][inner sep=0.75pt]    {$\textup{B}^+$};
		\end{tikzpicture}
	\end{center}
	
	Posto ora $W = \Psi(\textup{B})$, che non sarà una palla, ma comunque sarà un intorno aperto di $x_0$, convertendo quanto detto nelle variabili $x$, essendo lo jacobiano limitato, otteniamo un'estensione di $u$ a tutto $W$ con stima
	\[ \norm{\bar{u}}_{W^{1,p}(W)} \le C \norm{u}_{W^{1,p}(U)}. \]
	
	A questo punto, essendo $\partial U$ compatto, esisterà un numero finito di $x_j \in \partial U$ e di $W_j$ aperti tali che $u$ si estenda a $\bar{u}_j$ su $W_j$ e tali che 
	\[ \partial U \subseteq \bigcup_{j=1}^m W_j. \]
	Sia poi $W_0$ un aperto in $U$ tale che $\overline{W}_0 \subseteq U$ e 
	\[ U \subseteq \bigcup_{j=0}^m W_j \]
	e consideriamo una partizione dell'unità $(\xi_j)$ associata a questo ricoprimento. Posto $\bar{u}_0 = u$, definiamo
	\[ \bar{u} = \sum_{j=0}^{m}\xi_j \bar{u}_j. \]
	Potendo scegliere $\textup{B}$ di raggio arbitrariamente piccolo nei passaggi precedenti, si ha che $\bar{u}$ può sempre essere pensata con supporto in $V$. Inoltre, da un non proprio velocissimo calcolo, combinando le stime già ottenute, si vede che esiste $C > 0$ tale che
	\[ \norm{\bar{u}}_{W^{1,p}(\R^n)} \le C \norm{u}_{W^{1,p}(U)}. \]
	Da ciò possiamo concludere che se $u \in C^1(\overline{U})$, $Eu= \bar{u}$, con $E$ lineare per costruzione.
	
	Consideriamo ora il caso generale, ossia $u \in W^{1,p}(U)$. Per il Teorema \eqref{globapprbordo}, esiste $(u_m)$ in $C^\infty(\overline{U})$ tale che $u_m \to u$ in $W^{1,p}(U)$. Essendo convergente, allora $(u_m)$ è di Cauchy. Da fatto che 
	\[ \norm{Eu_m- Eu_l}_{W^{1,p}(\R^n)} = \norm{E(u_m- u_l)}_{W^{1,p}(\R^n)} \le C \norm{u_m - u_l}_{W^{1,p}(U)} \]
	otteniamo che  anche $(Eu_m)$ è una successione di Cauchy. Ma dal fatto che $W^{1,p}(\R^n)$ è di Banach risulta che $(Eu_m)$ è convergente. Risulta quindi ben definita l'estensione
	\[ Eu = \lim\limits_{m} Eu_m, \]
	infatti se si avesse anche $(v_m)$ in $C^\infty(\overline{U})$ tale che $v_m \to u$ in $W^{1,p}(U)$, risulterebbe
	\[ \norm{Eu_m- Ev_l}_{W^{1,p}(\R^n)} \le C \norm{u_m - v_l}_{W^{1,p}(U)} \]
	e dal fatto che 
	\[ \lim\limits_{m} u_m  = \lim\limits_{l} v_l = u\]
	si ha pertanto che 
	\[ \lim\limits_{l} Ev_l = Eu. \]
	Questo fornisce l'operatore richiesto. 
\end{proof}

\begin{definizione}
	Siano $1 \le p \le \infty$, $U$ limitato con $\partial U$ di classe $C^1$ e $V$ un aperto limitato tale che $\overline{U} \subseteq V$. Chiamiamo $Eu$\index[simboli]{$Eu$} un'\emph{estensione di $u$ a $\R^n$}\index[analitico]{estensione}.
\end{definizione}
\begin{osservazione}
	Se supponiamo inoltre che $\partial U$ sia di classe $C^2$, allora la stessa tecnica utilizzata nella dimostrazione del Teorema di estensione permette di ottenere un operatore $E: W^{2,p}(U) \to W^{2,p}(\R^n)$ lineare e continuo con le stesse proprietà viste nel Teorema di estensione. In particolare, esiste $C > 0$, dipendente solo da $U,V,n$ e $p$, tale che 
	\[ \norm{Eu}_{W^{2,p}(\R^n)}\le C\norm{u}_{W^{2,p}(U)}.  \] 
\end{osservazione}
\begin{osservazione}
	Se $k > 2$, è sempre possibile ottenere un Teorema di estensione. La dimostrazione è tuttavia più articolata e sfrutta delle tecniche di riflessione ad ordine superiore più complesse.
\end{osservazione}

\section{Tracce}
Se $u \in C(\overline{U})$, chiaramente ha senso parlare di valori su $\partial U$ per $u$. Invece, se $u \in W^{1,p}(U)$, in generale non è continua ed in più è definita solo quasi ovunque su $U$. Essendo $\partial U$ trascurabile, non c'è un modo diretto per parlare di valori al bordo. Introduciamo in questa sezione un operatore che gestisce queste situazioni.

\begin{teorema} \emph{\textbf{(di traccia)}}\index[analitico]{teorema!di traccia}
	Siano $1 \le p < \infty$ e $U$ limitato con $\partial U$ di classe $C^1$. Allora esiste un operatore lineare e continuo $T: W^{1,p}(U) \to L^p(\partial U)$ tale che
	\begin{enumerate}[$(a)$]
		\item $Tu = u |_{\partial U}$ ogniqualvolta $u \in W^{1,p}(U) \cap C(\overline{U})$;
		\item esiste $C >0$, dipendente solo da $p$ ed $U$, tale che
		\[ \norm{Tu}_{L^p(\partial U)} \le C \norm{u}_{W^{1,p}(U)}. \]
	\end{enumerate}
\end{teorema}
\begin{proof}
	Supponiamo inizialmente $u \in C^1(\overline{U})$. Fissiamo $x_0 \in \partial U$ e supponiamo anche che, in un intorno di $x_0$, $\partial U$ giaccia sul piano $x_n = 0$, ossia che esista $r > 0 $ tale che, detta $\textup{B} = \textup{B}(x_0,r)$,
	\begin{align*}
		\textup{B}^+ &:= \textup{B}(x_0,r) \cap \left\lbrace x_n \ge 0\right\rbrace \subseteq \overline{U}, \\
		\textup{B}^- &:= \textup{B}(x_0,r) \cap \left\lbrace x_n \le 0\right\rbrace \subseteq \R^n \setminus U.
	\end{align*}
	Sia anche $\hat{\textup{B}} = \textup{B}(x_0,\frac{r}{2}) $. Denotiamo con $\Gamma = \hat{\textup{B}} \cap \partial U$ e con $x' = \left( x_1,\dots, x_{n-1}\right) $, ossia $x = (x',x_n)$. Prendiamo $\zeta \in C^{\infty}_c(\textup{B})$ tale che $\zeta \ge 0$ in $\textup{B}$ e $\zeta = 1$ in $\hat{\textup{B}}$. Si ha, dalla Formula di Gauss--Green e dalla disuguaglianza di Young,
	\begin{align*}
		\int_\Gamma \abs{u}^pdx' &\le \int_{\left\lbrace x_n = 0\right\rbrace } \zeta \abs{u}^pdx' = \\
		& = - \int_{\textup{B}^+} D_{x_n}(\zeta \abs{u}^p)dx \le C \int_{\textup{B}^+} \left( \abs{u}^p + \abs{u}^{p-1}D_{x_n}u\right)  dx \le \\
		&\le C \int_{\textup{B}^+} \left( \abs{u}^p + \sum_{j=1}^{n} \abs{D_ju}^p\right)  dx \le C \norm{u}^p_{W^{1,p}(U)}.
	\end{align*}
	Consideriamo ora il caso in cui $\partial U$ non sia necessariamente piatto vicino ad $x_0$. Sia $\Phi$ il diffeomorfismo, con inversa $\Psi$, che appiattisce $\partial U$ in un intorno di $x_0$. Introduciamo le notazioni $y = \Phi(x)$, $x = \Psi(y)$ e $u'(y) = u(\Psi(y))$. Applicando la stima fatta nel caso precedente e cambiando variabili si ottiene la maggiorazione
	\[ 	\int_\Gamma \abs{u}^pdS \le C \norm{u}^p_{W^{1,p}(U)} \]
	dove ora $\Gamma$ è un qualche aperto di $\partial U$ contenente $x_0$.
	
	A questo punto, dalla compattezza di $\partial U$, segue che esisterà un numero finito di $x_j \in \partial U$ e di $\Gamma_j$ aperti in $\partial U$ tali che $\partial U = \bigcup_{j=1}^{m}\Gamma_j$ e per ogni $j=1,\dots,m$
	\[\norm{u}_{L^p(\Gamma_j)} \le C\norm{u}_{W^{1,p}(U)}. \]
	Detto quindi $Tu = u|_{\partial U}$, lineare per costruzione, si ha 
	\[ \norm{Tu}_{L^p(\partial U)} \le \sum_{j=1}^{m} \norm{Tu}_{L^p(\Gamma_j)} \le \tilde{C} \norm{u}_{W^{1,p}(U)} \]
	che fornisce la tesi in questo caso.
	
	Consideriamo ora il caso generale, ossia $u \in W^{1,p}(U)$. Per il Teorema \eqref{globapprbordo}, esiste $(u_m)$ in $C^\infty(\overline{U})$ tale che $u_m \to u$ in $W^{1,p}(U)$. Essendo convergente, allora $(u_m)$ è di Cauchy. Da fatto che 
	\[ \norm{Tu_m- Tu_l}_{L^{p}(\partial U)} = \norm{T(u_m- u_l)}_{L^{p}(\partial U)} \le C \norm{u_m - u_l}_{W^{1,p}(U)} \]
	otteniamo che  anche $(Tu_m)$ è una successione di Cauchy. Ma dal fatto che $L^{p}(\partial U)$ è di Banach risulta che $(Tu_m)$ è convergente. Risulta quindi ben definita l'estensione
	\[ Tu = \lim\limits_{m} Tu_m, \]
	infatti se si avesse anche $(v_m)$ in $C^\infty(\overline{U})$ tale che $v_m \to u$ in $W^{1,p}(U)$, risulterebbe
	\[ \norm{Tu_m- Tv_l}_{L^{p}(\partial U)} \le C \norm{u_m - v_l}_{W^{1,p}(U)} \]
	e dal fatto che 
	\[ \lim\limits_{m} u_m  = \lim\limits_{l} v_l = u\]
	si ha pertanto che 
	\[ \lim\limits_{l} Tv_l = Tu. \]
	Infine, se $u \in W^{1,p}(U)\cap C(\overline{U})$, allora $u_m \to u$ uniformemente e quindi dal fatto che il limite uniforme di applicazioni continue è continuo si ha
	\[ Tu = \lim_{m} Tu_m = \lim_{m} u_m|_{\partial U} = u|_{\partial U}\]
	e la dimostrazione è conclusa.
\end{proof}
\begin{definizione}
	Siano $1 \le p <\infty$ ed $U$ limitato con $\partial U$ di classe $C^1$. Chiamiamo $Tu$\index[simboli]{$Tu$} la \emph{traccia di $u$ su $\partial U$}\index[analitico]{traccia}.
\end{definizione}
Il seguente risultato permette di caratterizzare pienamente gli elementi di $W^{1,p}_0(U)$ come gli zeri dell'operatore di traccia. In questo senso risulta ora ragionevole l'espressione \emph{fare zero al bordo} per queste applicazioni.
\begin{teorema}
	Siano $1 \le p <\infty$, $U$ limitato con $\partial U$ di classe $C^1$ ed $u \in W^{1,p}(U)$. I seguenti fatti sono equivalenti:
	\begin{enumerate}[$(a)$]
		\item $u \in  W^{1,p}_0(U)$,
		\item $Tu = 0$ su $\partial U$.
	\end{enumerate}
\end{teorema}
\begin{proof}\hfill
	\par\noindent
	$(a) \Longrightarrow (b)$ Supponiamo $u \in W^{1,p}_0(U)$, allora esiste una successione $(u_m)$ in $C^\infty_c(U)$ tale che $u_m \to u$ in $W^{1,p}(U)$. Per continuità dell'operatore di traccia si ha che su $\partial U$ 
	\[ Tu = \lim\limits_{m} Tu_m = \lim\limits_{m} 0 = 0.  \]
	\par\noindent
	$(b) \Longrightarrow (a)$ Sia $u \in W^{1,p}(U)$ tale che $Tu = 0$ su $\partial U$. 
	
	Senza perdita di generalità possiamo limitarci al caso $u \in W^{1,p}(\R^n_{+})$ con $Tu = 0$ su $\partial\R^n_{+} = \R^{n-1} $.\footnote{Si tratta di appiattire il bordo di $U$, considerare, come nel teorema di estensione, una partizione dell'unità $(\xi_i)$ e scrivere $u = \sum_{i=0}^{\infty}u\xi_i$. A questo punto possiamo ragionare sui singoli termini della serie, che sono a supporto compatto e quindi definibili su tutto $\R^n_+$.}
	
	Dal fatto che $Tu = 0$ su $\R^{n-1} $, segue che esiste una successione $(u_m)$ in $C^1(\R^n_{+})$, tale che $u_m \to u$ in $W^{1,p}(\R^n_{+})$ e $Tu_m = u_m|_{\R^{n-1}} \to 0$ in $L^p(\R^{n-1})$.
	
	Indicando $x = (x,x_n)$, dove possiamo orientare gli assi in modo da avere $x_n \ge 0$, dal Teorema fondamentale del calcolo integrale segue
	\[ \abs{u_m(x',x_n)} \le \abs{u_m(x',0)}+ \int_{0}^{x_n} \abs*{D_{x_n}u_m(x',t)}dt \]
	da cui
	\[ \abs{u_m(x',x_n)}^p \le C\left(  \abs{u_m(x',0)}^p+ \left( \int_{0}^{x_n} \abs*{D_{x_n}u_m(x',t)}dt\right) ^p\right).  \]
	Integrando in $x'$, otteniamo
	\[ \int_{\R^{n-1}}\abs{u_m(x',x_n)}^p dx' \le C \left( \int_{\R^{n-1}} \abs{u_m(x',0)}^pdx'+ \int_{\R^{n-1}}\left( \int_{0}^{x_n} \abs*{Du_m(x',t)}dt\right) ^pdx'\right)  \]
	ma per la disuguaglianza di Holder,
	\[ \int_{0}^{x_n} \abs*{Du_m(x',t)}dt \le x_n^{\frac{p-1}{p}}\left( \int_{0}^{x_n}\abs*{Du_m(x',t)}^pdt\right)^{\frac{1}{p}},  \]
	perciò
	\[ 	\int_{\R^{n-1}}\abs{u_m(x',x_n)}^p dx' \le C \left( \int_{\R^{n-1}} \abs{u_m(x',0)}^pdx' + x_n^{p-1}\int_{\R^{n-1}} \int_{0}^{x_n}\abs*{Du_m(x',t)}^pdt dx'\right).  \]
	A questo punto, dal fatto che $u_m \to u$ in $W^{1,p}(\R^n_{+})$ e $Tu_m = u_m|_{\R^{n-1}} \to 0$ in $L^p(\R^{n-1})$, combinato con il Lemma di Fatou ed il Teorema di Fubini--Tonelli, risulta che per q.o. $x_n \ge 0$
	\begin{equation}\label{maggiorazionetraccia0}
		\int_{\R^{n-1}}\abs{u(x',x_n)}^p dx' \le C x_n^{p-1} \int_{0}^{x_n}\int_{\R^{n-1}}\abs*{Du(x',t)}^p dx'dt.
	\end{equation}
	Sia ora $\zeta \in C^\infty(\R^+)$ tale che $\zeta = 1$ in $[0,1]$, $\zeta = 0$ in $\R\setminus [0,2]$ e $0\le \zeta \le 1$. Consideriamo le successioni $(\zeta_m)$ e $(w_m)$ in $R^n_{+}$ tali che
	\[ \zeta_m(x) = \zeta(mx_n), \]
	\[ w_m(x) = u(x)(1-\zeta_m(x)). \]
	Chiaramente
	\[ D_{x_n}w_m(x) =  \left( 1-\zeta_m(x)\right)D_{x_n}u(x)  -mu \zeta'_m(x), \]
	\[ D_{x'}w_m(x) = \left( 1-\zeta_m(x)\right)D_{x'}u(x). \]
	Di conseguenza,
	\[ \int_{\R^n_+}\abs{Dw_n -Du}^p d\mathcal{L}^n \le C \int_{\R^n_+}\abs{\zeta_m}^p\abs{Du}^pd\mathcal{L}^n + Cm^p\int_{0}^{\frac{2}{m}}\int_{\R^{n-1}}\abs{u}^p dx'dt.  \]
	A questo punto, per \eqref{maggiorazionetraccia0},
	\begin{align*}
		Cm^p\int_{0}^{2m}\int_{\R^{n-1}}\abs{u}^p dx'dt &\le C m^p \int_{0}^{\frac{2}{m}}t^{p-1}dt \int_{0}^{\frac{2}{m}}\int_{\R^{n-1}}\abs{Du}^pdx'dt \le \\
		&\le C \int_{0}^{\frac{2}{m}}\int_{\R^{n-1}}\abs{Du}^pdx'dt
	\end{align*}
	e quindi dal Teorema della convergenza dominata si ha che $Dw_m \to Du$ in $L^{p}(\R^n_{+})$. Dal fatto che, chiaramente, $w_m \to u$ in $L^{p}(\R^n_{+})$ possiamo quindi affermare che $w_m \to u$ in $W^{1,p}(\R^n_{+})$. Dal fatto che $w_m = 0 $ se $0 < x_n < \frac{1}{m}$, segue che possiamo regolarizzare la successione $(w_m)$ producendo una successione $(u_m)$ in $C^\infty_c(\R^n_{+})$ tale che $u_m \to u$ in $W^{1,p}(\R^n_{+})$, ossia $u \in W^{1,p}_0(\R^n_{+})$.
\end{proof}
\section{Teoremi di immersione}
Se $u \in W^{1,p}(U)$, per definizione $u \in L^p(U)$. In questa sezione vediamo come il fatto che anche $D_ju \in L^p(U)$, fornisca un guadagno sulla sommabilità di $u$.

Iniziamo facendo un piccolo esperimento.
\begin{osservazione}
	Sia $1 \le p < n$. Supponiamo che esistano $C > 0$ e $q \in \left[ 1,\infty \right[ $ tali che per ogni $u \in C^\infty_c(\R^n)$ si abbia
	\[ \norm{u}_{L^q(\R^n)} \le C \norm{Du}_{L^p(\R^n)}. \]
	
	Sia $\lambda > 0$ e definiamo $u_\lambda(x) = u(\lambda x)$. Applicando la disuguaglianza a $u_\lambda$ otteniamo 
	\[ \frac{1}{\lambda^{\frac{n}{q}}}\norm{u}_{L^q(\R^n)} \le C \frac{\lambda}{\lambda^{\frac{n}{p}}}\norm{Du}_{L^p(\R^n)} \]
	da cui
	\[\norm{u}_{L^q(\R^n)} \le C \lambda^{1-\frac{n}{p}+\frac{n}{q}}\norm{Du}_{L^p(\R^n)}.  \]
	Se, per assurdo, $1-\frac{n}{p}+\frac{n}{q} \ne 0$ siamo di fronte a due casi:
	\begin{itemize}
		\item se $1-\frac{n}{p}+\frac{n}{q} > 0$, mandando $\lambda \to 0$, si ottiene una contraddizione,
		\item se $1-\frac{n}{p}+\frac{n}{q} < 0$, mandando $\lambda \to \infty$, si ottiene ancora una contraddizione.
	\end{itemize} 
	Perciò l'esponente deve essere necessariamente nullo, ossia $q = \frac{np}{n-p}$.
\end{osservazione}
La precedente osservazione motiva la seguente definizione.
\begin{definizione}
	Per ogni $p \in \left[ 1,\infty\right] $, poniamo 
	\[ p^* := \begin{cases}
		\infty & \textup{se }p=n, \\
		\frac{np}{n-p} & \textup{se } p < n.
	\end{cases} \]
	Diciamo che $p^*$ è l'esponente \emph{coniugato di Sobolev}\index[analitico]{coniugato di Sobolev}\index[simboli]{$p^*$} di $p$.
\end{definizione}
È evidente che $p^* > p$ e 
\[ \frac{1}{p^*} = \frac{1}{p} - \frac{1}{n}. \]
Vedremo più avanti il motivo per cui non abbiamo speso fatica cercando di estendere la definizione di esponente coniugato di Sobolev anche per $p>n$. In ogni caso, iniziamo studiando il caso $p < n$.
\begin{teorema}\emph{\textbf{(disuguaglianza di Gagliardo--Nirenberg--Sobolev)}}\label{GNS}\index[analitico]{disuguaglianza! di Gagliardo--Nirenberg--Sobolev}
	Sia $p \in \left[ 1, n\right[ $. Allora esiste $C > 0$, dipendente solo da $p$ e $n$, tale che per ogni $u \in C^1_c(\R^n)$ si abbia
	\[  \norm{u}_{L^{p^*}(\R^n)}\le C \norm{Du}_{L^p(\R^n)}. \]
\end{teorema}
\begin{proof}
	Consideriamo innanzitutto il caso $p = 1$. Allora $p^\ast = \frac{n}{n-1}$ e dal Teorema fondamentale del calcolo integrale, siccome $u(y) = 0$ per ogni $y \notin \textup{supt}(u)$,
	\[ u(x) = \int_{-\infty}^{x_i} D_{x_i}u(x_1,\dots,x_{i-1},y_i,x_{i+1},\dots,x_n)dy_i, \]
	da cui
	\[ \abs{u(x)} \le \int_{-\infty}^{+\infty} \abs{Du(x_1,\dots,x_{i-1},y_i,x_{i+1},\dots,x_n)}dy_i. \]
	Allora,
	\[ \abs{u(x)}^{\frac{n}{n-1}} \le \prod_{i=1}^{n}\left( \int_{-\infty}^{+\infty} \abs{Du(x_1,\dots,x_{i-1},y_i,x_{i+1},\dots,x_n)}dy_i\right) ^{\frac{1}{n-1}} \]
	per ogni $x \in \R^n$. Effettuiamo ora una prima integrazione rispetto ad $x_1$, 
	\begin{align*}
		\int_{-\infty}^{+\infty} \abs{u}^{\frac{n}{n-1}}dx_1 &\le \int_{-\infty}^{+\infty} \prod_{i=1}^{n}\left( \int_{-\infty}^{+\infty} \abs{Du(x_1,\dots,x_{i-1},y_i,x_{i+1},\dots,x_n)}dy_i\right) ^{\frac{1}{n-1}} dx_1 = \\
		&= \int_{-\infty}^{+\infty} \left( \int_{-\infty}^{+\infty} \abs{Du}dy_1\right) ^{\frac{1}{n-1}} \prod_{i=2}^{n}\left( \int_{-\infty}^{+\infty} \abs{Du}dy_i\right) ^{\frac{1}{n-1}} dx_1 = \\
		& = \left( \int_{-\infty}^{+\infty} \abs{Du}dy_1\right) ^{\frac{1}{n-1}} \int_{-\infty}^{+\infty} \prod_{i=2}^{n}\left( \int_{-\infty}^{+\infty} \abs{Du}dy_i\right) ^{\frac{1}{n-1}} dx_1
	\end{align*}
	e per la disuguaglianza di Holder generalizzata otteniamo 
	\[ \int_{-\infty}^{+\infty} \abs{u}^{\frac{n}{n-1}}dx_1 \le \left( \int_{-\infty}^{+\infty} \abs{Du}dy_1\right) ^{\frac{1}{n-1}} \left(\prod_{i=2}^{n}\int_{-\infty}^{+\infty}  \int_{-\infty}^{+\infty} \abs{Du}dy_idx_1\right) ^{\frac{1}{n-1}}. \]
	Procediamo ora con una integrazione rispetto ad $x_2$,
	\begin{align*}
		\int_{-\infty}^{+\infty}&\int_{-\infty}^{+\infty} \abs{u}^{\frac{n}{n-1}}dx_1 dx_2 \le \\
		&\le \int_{-\infty}^{+\infty}\left( \int_{-\infty}^{+\infty} \abs{Du}dy_1\right) ^{\frac{1}{n-1}} \left(\prod_{i=2}^{n}\int_{-\infty}^{+\infty}  \int_{-\infty}^{+\infty} \abs{Du}dy_idx_1\right) ^{\frac{1}{n-1}}dx_2,
	\end{align*}
	quindi
	\begin{align*}
		\int_{-\infty}^{+\infty}&\int_{-\infty}^{+\infty} \abs{u}^{\frac{n}{n-1}}dx_1 dx_2 \le \\
		&\le \left( \int_{-\infty}^{+\infty} \int_{-\infty}^{+\infty} \abs{Du}dx_1dy_2\right) ^{\frac{1}{n-1}} \\
		&\int_{-\infty}^{+\infty} \left( \int_{-\infty}^{+\infty}\abs{Du}dy_1\right)^{\frac{1}{n-1}}  \prod_{i=3}^{n} \left( \int_{-\infty}^{+\infty}\int_{-\infty}^{+\infty}\abs{Du}dx_1 dy_i\right) ^{\frac{1}{n-1}} dx_2
	\end{align*}
	ed applicando nuovamente la disuguaglianza di Holder generalizzata si ha
	\begin{align*}
		\int_{-\infty}^{+\infty}&\int_{-\infty}^{+\infty} \abs{u}^{\frac{n}{n-1}}dx_1 dx_2 \le \\
		&\le \left( \int_{-\infty}^{+\infty}\int_{-\infty}^{+\infty}\abs{Du}dx_1dy_2\right) ^{\frac{1}{n-1}}\left( \int_{-\infty}^{+\infty}\int_{-\infty}^{+\infty}\abs{Du}dy_1dx_2\right) ^{\frac{1}{n-1}} \\
		& \prod_{i=3}^{n} \left(\int_{-\infty}^{+\infty} \int_{-\infty}^{+\infty}\int_{-\infty}^{+\infty}\abs{Du}dx_1dx_2 dy_i\right) ^{\frac{1}{n-1}}
	\end{align*}
	dove i primi due integrali sono uguali per il Teorema di Fubini--Tonelli. A questo punto, procedendo ricorsivamente ed applicando il Teorema di Fubini--Tonelli si ottiene
	\[ 	\int_{\R^n}\abs{u}^{\frac{n}{n-1}}d\mathcal{L}^n \le \prod_{i=1}^n \left( \int_{\R^n} \abs{Du}d\mathcal{L}^n\right) ^{\frac{1}{n-1}} = \left( \int_{\R^n} \abs{Du}d\mathcal{L}^n\right) ^{\frac{n}{n-1}} \]
	ossia $\norm{u}_{\frac{n}{n-1}} \le \norm{Du}_1$. Consideriamo ora il caso $1 < p < n$. Dato $\gamma > 1$, applichiamo la stima precedente a $v = \abs{u}^\gamma$
	\[ \left(	\int_{\R^n}\abs{u}^{\frac{\gamma n}{n-1}}d\mathcal{L}^n\right) ^{\frac{n-1}{n}} \le \int_{\R^n} \gamma \abs{u}^{\gamma -1}\abs{Du}d\mathcal{L}^n \]
	ed utilizzando la disuguaglianza di Holder
	\[ \left(\int_{\R^n}\abs{u}^{\frac{\gamma n}{n-1}}d\mathcal{L}^n\right) ^{\frac{n-1}{n}} \le C \left( \int_{\R^n} \abs{u}^{\frac{(\gamma -1)p}{p-1}}d\mathcal{L}^n\right) ^{\frac{p-1}{p}} \left(\int_{\R^n} \abs{Du}^p d\mathcal{L}^n\right) ^{\frac{1}{p}}.  \]
	La posizione\footnote{It's a kind of magic!} $ \gamma = \frac{p(n-1)}{n-p} $ che corrisponde alla condizione
	\[ \frac{\gamma n}{n-1} = \left( \gamma -1\right) \frac{p}{p-1} \]
	fornisce 
	\[ \frac{\gamma n}{n-1} = p^{\ast} \]
	e 
	\[ \left(\int_{\R^n}\abs{u}^{p^{\ast}}d\mathcal{L}^n\right) ^{\frac{1}{p^{\ast}}} \le C \left(\int_{\R^n} \abs{Du}^p d\mathcal{L}^n\right) ^{\frac{1}{p}} \]
	che è la tesi.
\end{proof}
\begin{osservazione}
	L'ipotesi di supporto compatto nel Teorema \eqref{GNS} è necessaria: basti pensare alle funzioni costanti. Se per assurdo si potesse eliminare l'ipotesi, otterremmo $+\infty \le 0$. Comunque $C$ non dipende dall'ampiezza del supporto in quanto non dipende da $u$.
\end{osservazione}
\begin{teorema}\textbf{\emph{(di Sobolev)}}\label{thdiSobolev}\index[analitico]{teorema! di Sobolev}
	Siano $1 \le p < n$, $U$ limitato con $\partial U$ di classe $C^1$ e $u \in W^{1,p}(U)$. Allora $u \in L^{p^\ast}(U)$ ed esiste $C$, dipendente solo da $p,n$ ed $U$, tale che 
	\[ \norm{u}_{L^{p^\ast}(U)} \le C\norm{u}_{W^{1,p}(U)}. \]
\end{teorema}
\begin{proof}
	Innanzitutto per il Teorema di estensione, sia $\overline{u} = Eu \in W^{1,p}(\R^{n})$ tale che $\overline{u} = u$ q.o. in $U$, $\overline{u}$ è a supporto compatto in $\R^n$ e 
	\[ \norm{\overline{u}}_{W^{1,p}(\R^n)} \le C \norm{u}_{W^{1,p}(U)}. \]
	Siccome $\overline{u}$ è a supporto compatto, per il Teorema di locale approssimazione con funzione lisce, esiste $(u_h)$ in $C^\infty_c(\R^n)$ tale che $u_h \to \overline{u}$ in $W^{1,p}(\R^n)$. Ora, dalla Disuguaglianza di Sobolev, 
	\[ \norm{u_h - u_k}_{L^{p^\ast}(\R^n)} \le C\norm{Du_h - Du_k}_{L^p(\R^n)} \]
	allora $(u_h)$ è di Cauchy in $L^{p^\ast}(\R^n)$, per cui è convergente ad un certo $v$ in $L^{p^\ast}(\R^n)$. Per unicità del limite\footnote{Basta passare a delle sottosuccessioni per verificarlo esplicitamente.} si ha che $\overline{u}=v$ q.o. in $\R^n$. Ora, per il Teorema \eqref{GNS},
	\[ \norm{u_h}_{L^{p^\ast}(\R^n)} \le C \norm{Du_h}_{L^p(\R^n)} \]
	e se $h \to +\infty$
	\[ \norm{\overline{u}}_{L^{p^\ast}(\R^n)} \le C \norm{D\overline{u}}_{L^p(\R^n)} \le C\norm{\overline{u}}_{W^{1,p}(\R^n)}, \]
	ma 
	\[ \norm{\overline{u}}_{W^{1,p}(\R^n)} \le C \norm{u}_{W^{1,p}(U)}\]
	e sicuramente
	\[ \norm{u}_{L^{p^\ast}(U)} \le C \norm{\overline{u}}_{L^{p^\ast}(\R^n)} \]
	da cui 
	\[ \norm{u}_{L^{p^\ast}(U)} \le C\norm{u}_{W^{1,p}(U)}.\qedhere \]
\end{proof}

\begin{proposizione}
	Siano $1 \le p < n$, $U$ limitato con $\partial U$ di classe $C^1$ e $u \in W^{1,p}(U)$. Allora per ogni $q \in \left[ 1,p^\ast\right] $, $u \in L^{q}(U)$ ed esiste $C$, dipendente solo da $p,q,n$ ed $U$, tale che 
	\[ \norm{u}_{L^q(U)} \le C\norm{u}_{W^{1,p}(U)}. \]
\end{proposizione}
\begin{proof}
	Discende da Teorema \eqref{thdiSobolev} e dall'incapsulamento degli spazi di Lebesgue per $U$ limitato.
\end{proof}
Vediamo ora una disuguaglianza che, limitatamente alle $u \in W^{1,p}_0(U)$, estende al caso di un dominio limitato la disuguaglianza di Gagliardo--Nirenberg--Sobolev.
\begin{teorema}\textbf{\emph{(disuguaglianza di Poincaré)}}\index[analitico]{disuguaglianza! di Poincaré}\label{disPoincaré}
	Siano $1 \le p < +\infty$, $U$ limitato e $u \in W^{1,p}_0(U)$. Allora per ogni $q \in \left[ 1,p^\ast\right] $, $u \in L^{q}(U)$ ed esiste $C$, dipendente solo da $p,q,n$ ed $U$, tale che 
	\[ \norm{u}_{L^q(U)} \le C\norm{Du}_{L^p(U)}. \]
\end{teorema}
\begin{proof}
	Vediamo solo il caso $1 \le p < n$. Per definizione esiste $(u_m)$ in $C^\infty_c(U)$ tale che $u_m \to u$ in $W^{1,p}(U)$. Per il Teorema \eqref{GNS}, 
	\[ \norm{u_m}_{L^{p^\ast}(\R^n)} \le C \norm{Du_m}_{L^p(\R^n)} = C\norm{Du_m}_{L^p(U)} \]
	e analogamente $\norm{u_m}_{L^{p^\ast}(\R^n)} = \norm{u_m}_{L^{p^\ast}(U)}$, così che
	\begin{equation}\label{eqPoincaré}
		\norm{u_m}_{L^{p^\ast}(U)}\le C\norm{Du_m}_{L^p(U)}.
	\end{equation}
	Il secondo membro converge per definizione a $C\norm{Du}_{L^p(U)}$. Per quanto riguarda il primo membro è sufficiente osservare che esiste una sottosuccessione $(u_{m_k})$ convergente a $u$ q.o. in $U$, per cui, per il Lemma di Fatou,
	\[ \norm{u}_{L^{p^\ast}(U)} \le \underset{k}{\textup{lim inf }} \norm{u_{m_k}}_{L^{p^\ast}(U)} \]
	e, quindi, scrivendo la \eqref{eqPoincaré} per la sottosuccessione $(u_{m_k})$ e passando al lim inf si ha 
	\[ \norm{u}_{L^{p^\ast}(U)} \le C\norm{Du_m}_{L^p(U)}.\]
	Dall'incapsulamento degli spazi di Lebesgue per $U$ limitato segue la tesi.
\end{proof}
\begin{osservazione}
Vale la pena osservare che, limitatamente al caso $p<n$ e $q=p^\ast$, è possibile rimuovere le ipotesi di limitatezza su $U$. In particolare, in questo caso la costante $C$ dipende solo da $p$ ed $n$.
\end{osservazione}

\begin{osservazione}
	In $W^{1,p}_0(U)$ la norma $ \norm{Du}_{L^p(U)} $ è equivalente alla norma $\norm{u}_{W^{1,p}(U)}$. Infatti già sappiamo che
	\[ \norm{Du}_{L^p(U)} \le c_1 \norm{u}_{W^{1,p}(U)} . \]
	Per la Disuguaglianza di Poincaré otteniamo proprio l'altra stima. In particolare $H^{1}_0(U)$ munito della norma del gradiente è uno spazio di Hilbert.
\end{osservazione}

\begin{proposizione}\textbf{\emph{(disuguaglianza di Poincaré duale)}}\index[analitico]{disuguaglianza! di Poincaré!duale}\label{disPoincaréduale}
	Siano $\frac{n}{n-1} < q \le +\infty$ e $u \in L^{\frac{nq}{n+q}}(U)$. Allora , $u \in W^{-1,q}(U)$ ed esiste $C$, dipendente solo da $n$ e $q$, tale che 
	\[ \norm{u}_{W^{-1,q}(U)} \le C\norm{u}_{L^{\frac{nq}{n+q}}(U)}, \]
	dove si pone $\frac{nq}{n+q} = n$ ogniqualvolta $q=+\infty$.
\end{proposizione}
\begin{proof}
	Data $u \in  L^{\frac{nq}{n+q}}(U)$, consideriamo $T_u: W^{1,q'}_0(U) \to \R$ tale che 
	\[ \braket{T_u, v} = \int_U uvd\mathcal{L}^n. \]
	A meno di identificare $u$ e $T_u$, per la disuguaglianza di Holder\footnote{Siccome $\left( \frac{nq}{n+q}\right)' =(q')^\ast$.},
	\[ \abs{\braket{u,v}} \le \int_U \abs{u} \abs{v} d\mathcal{L}^n \le \norm{u}_{L^{\frac{nq}{n+q}}(U)} \norm{v}_{L^{(q')^\ast}(U)}\]
	e, per la disuguaglianza di Poincaré,
	\[ \abs{\braket{u,v}} \le  \norm{u}_{L^{\frac{nq}{n+q}}(U)} C \norm{v}_{W^{1,q'}_0(U)}, \]
	da cui la tesi.
\end{proof}
\begin{osservazione}
	La disuguaglianza di Poincaré duale può essere dimostrata anche semplicemente osservando il seguente fatto generale: consideriamo $X,Y$ due spazi di Banach su $\R$ tali che $X\subseteq Y$. Se l'immersione $X \hookrightarrow Y$ è continua, allora l'immersione $Y' \hookrightarrow X'$ è continua.\footnote{Questo deriva dal fatto che, essendo $X \subseteq Y$, a meno di identificare $\varphi \in Y'$ con $\varphi|_X$, l'immersione è ben definita e lineare. La continuità segue osservando che per ogni $\varphi \in Y'$
		\[ \norm{\varphi}_{X'} = \underset{\underset{z \in X}{\abs{z}\le 1}}{\sup}\; \abs{\braket{\varphi,x}} \le  \underset{\underset{z \in Y}{\abs{z}\le 1}}{\sup}\; \abs{\braket{\varphi,x}} = \norm{\varphi}_{Y'}.\]}
	
	Dalla disuguaglianza di Poincaré, sappiamo che l'immersione $W^{1,q'}_0(U)  \hookrightarrow L^{(q')^\ast}(U)$ è continua, quindi lo è anche l'immersione $(L^{(q')^\ast}(U))' \hookrightarrow W^{-1,q}(U)$. In altre parole, l'immersione $(L^{((q')^\ast)'}(U)) \hookrightarrow W^{-1,q}(U)$ è continua. A questo punto è sufficiente osservare che $((q')^\ast)' =\frac{nq}{n+q}$.
	
\end{osservazione}



\begin{osservazione}\label{probleman=p}
	Dal fatto che 
	\[ \lim\limits_{n \to p} \frac{np}{n-p} = +\infty, \]
	parrebbe ragionevole aspettarsi che se $u \in W^{1,n}(U)$, allora $u \in L^\infty(U)$. Tuttavia per $n > 1$ ciò è falso.
\end{osservazione}

Vediamo un controesempio che giustifica l'Osservazione \eqref{probleman=p}.
\begin{esempio}\footnote{Con leggere varianti si può fare anche in $\R^n$.}
	Sia $U = \textup{B}(0,1)$ in $\R^2$ e 
	\[ u(x) = \ln \left( \ln \left( 1 + \frac{1}{\abs{x}}\right)\right)  . \]
	Chiaramente $u \notin L^\infty(U)$. Mostriamo che $u \in W^{1,2}(U)$. Passando alle coordinate polari si vede subito che
	\[ \int_{U} \ln^2 \left( \ln \left( 1 + \frac{1}{\abs{x}}\right)\right)d\mathcal{L}^2(x) \sim \int_{0}^{1} \ln^2 \left( \ln \left( 1 + \frac{1}{\rho}\right)\right)\rho d\rho.  \]
	Ora, posto $y =\ln \left( 1 + \frac{1}{\rho}\right)$,
	\[ \int_{0}^{1} \ln^2 \left( \ln \left( 1 + \frac{1}{\rho}\right)\right)\rho d\rho \sim \int_{\ln 2}^{+\infty} \dfrac{e^y \ln^2 y}{(e^y-1)^3} dy \sim \int_{\ln 2}^{+\infty} \frac{\ln^2 y}{e^{2y}} dy. \]
	Ricordando che $\ln y \le y-1$ per ogni $y > 0$, 
	\[ \int_{\ln 2}^{+\infty} \frac{\ln^2 y}{e^{2y}} dy \le  \int_{\ln 2}^{+\infty} \frac{(y-1)^2}{e^{2y}} dy  \sim \int_{\ln 2}^{+\infty} \frac{y^2}{e^{2y}} dy. \]
	Se effettuiamo il cambio di variabile $y = \ln k$, otteniamo 
	\[ \int_{\ln 2}^{+\infty} \frac{y^2}{e^{2y}} dy = \int_{2}^{+\infty} \frac{1}{k^3\abs{\ln k}^{-2}}dk < +\infty \]
	da cui $u \in L^2(U)$. Invece, sempre utilizzando le coordinate polari osserviamo che
	\[ \int_{U} \abs{Du(x)}^2d\mathcal{L}^2(x) \sim \int_{0}^{1} \dfrac{1}{\rho(\rho + 1)^2\ln^2(1+\frac{1}{\rho})} d\rho \sim \int_{0}^{1} \frac{1}{\rho \ln^2(1+\frac{1}{\rho})} d\rho. \]
	Ora, posto $y =\ln \left( 1 + \frac{1}{\rho}\right)$,
	\[ \int_{0}^{1} \frac{1}{\rho \ln^2(1+\frac{1}{\rho})} d\rho \sim \int_{\ln 2}^{+\infty} \frac{1}{y^2} dy < +\infty \]
	da cui anche $Du \in L^2(U)$.  
\end{esempio}

Il caso $p=n$ è gestito dalla \emph{Disuguaglianza di Trudinger--Moser}, che non trattiamo nei dettagli. Passiamo quindi al caso $p > n$.
\begin{notazione}\index[simboli]{$\norm{u}_{C(\overline{U})}$}\index[simboli]{$\left[ u\right] _{C^{0,\gamma}(\overline{U})}$}\index[simboli]{$\norm{u}_{C(\overline{U})}$}\index[simboli]{$ \norm{u}_{C^{0,\gamma}(\overline{U})}$}
	Se $u \in C_b(U)$ e $\gamma \in \left] 0,1\right] $, indichiamo con
	\[ \left[ u\right] _{C^{0,\gamma}(\overline{U})} = \underset{\underset{x\ne y}{x, y \in U}}{\textup{sup}}\frac{\abs{u(x)-u(y)}}{\abs{x-y}^\gamma} \]
	e
	\[ \norm{u}_{C^{0,\gamma}(\overline{U})} = \norm{u}_{\infty} + \left[ u\right] _{C^{0,\gamma}(\overline{U})}.\]
\end{notazione}
\begin{definizione}
	Sia $\gamma \in \left] 0,1\right] $ e $k\ge 0$. Chiamiamo \emph{spazio di Holder}\index[analitico]{spazio! di Holder}, denotato con $C^{k,\gamma}(\overline{U})$\index[simboli]{$C^{k,\gamma}(\overline{U})$}, l'insieme delle $u \in C^k(\overline{U})$ tali che 
	\[ \norm{u}_{C^{k,\gamma}(\overline{U})} = \sum_{\abs{\alpha}\le k}\norm{D^\alpha u}_{\infty} + \sum_{\abs{\alpha}\le k} \left[ D^\alpha u\right] _{C^{0,\gamma}(\overline{U})} < +\infty.\]
\end{definizione}
\begin{proposizione}
	Sia $\gamma \in \left] 0,1\right] $ e $k\ge 0$. Allora $\norm{\ }_{C^{k,\gamma}(\overline{U})}$ è una norma su  $C^{k,\gamma}(\overline{U})$. Inoltre $\left(C^{k,\gamma}(\overline{U}),\norm{\ }_{C^{k,\gamma}(\overline{U})}\right)$ è uno spazio di Banach su $\R$.
\end{proposizione}
\begin{proof}
	Omettiamo la dimostrazione.
\end{proof}

Non è invece una norma $ \left[\ \right] _{C^{0,\gamma}(\overline{U})}$ in quanto nulla su tutte le funzioni costanti.

\begin{teorema}\textbf{\emph{(disuguaglianza di Morrey)}}\index[analitico]{disuguaglianza! di Morrey}
	Siano $n < p \le \infty$ e $\gamma = 1 - \frac{n}{p}$. Allora esiste $C$, dipendente solo da $p$ e $n$, tale che per ogni $u \in C^1(\R^n)$ si abbia
	\[ \norm{u}_{C^{0,\gamma}(\R^n)} \le C \norm{u}_{W^{1,p}(\R^n)}. \]
	\begin{proof}
		Limitiamo la dimostrazione al caso $p \ne \infty$. Sia $x \in \R^n$ ed $r >0$. Mostriamo che esiste $C$, dipendente solo da $n$ tale che
		\[ \fint_{\textup{B}(x,r)} \abs{u(y) - u(x)}d\mathcal{L}^n(y) \le C \int_{\textup{B}(x,r)} \frac{\abs{Du(y)}}{\abs{y-x}^{n-1}}d\mathcal{L}^n(y). \]
		Siano $w \in \partial \textup{B}(0,1)$ e $s \in \left] 0,r\right[ $. Chiaramente $\abs{w}=1$ e $x + sw \in \partial \textup{B}(x,s)$. Inoltre, per il Teorema fondamentale del calcolo integrale, 
		\[ \abs{u(x+sw) - u(x)} = \abs*{\int_{0}^{s}\frac{d}{dt}u(x+tw)dt} = \abs*{\int_{0}^{s}Du(x+tw) \cdot w dt} \le \int_{0}^{s} \abs*{Du(x+tw)}dt.  \]
		Integriamo ora rispetto a $w$, ottenendo
		\[ \int_{\partial\textup{B}(0,1)} \abs{u(x+sw) - u(x)} d\mathcal{H}^{n-1}(w) \le \int_{0}^{s} \int_{\partial\textup{B}(0,1)} \abs*{Du(x+tw)} d\mathcal{H}^{n-1}(w) dt. \]
		Consideriamo l'integrale
		\[ \int_{\textup{B}(x,s)} \frac{\abs{Du(y)}}{\abs{x-y}^{n-1}}d\mathcal{L}^n(y) \]
		ed applichiamogli il cambio di variabile $\varphi: \left] 0,s\right[ \times \partial \textup{B}(0,1) \to \textup{B}(x,s)$ tale che $y = \varphi(t,w) = x+tw$ ottenendo 
		\[ \int_{\textup{B}(x,s)} \frac{\abs{Du(y)}}{\abs{x-y}^{n-1}}d\mathcal{L}^n(y) =  \int_{0}^{s} \int_{\partial\textup{B}(0,1)} \abs*{Du(x+tw)} \frac{t^{n-1}}{t^{n-1}} d\mathcal{H}^{n-1}(w) dt.\]
		Pertanto,
		\begin{align*}
		\int_{\partial\textup{B}(0,1)} \abs{u(x+sw) - u(x)} d\mathcal{H}^{n-1}(w) &\le \int_{\textup{B}(x,s)} \frac{\abs{Du(y)}}{\abs{x-y}^{n-1}}d\mathcal{L}^n(y) \le \\
		&\le \int_{\textup{B}(x,r)} \frac{\abs{Du(y)}}{\abs{x-y}^{n-1}}d\mathcal{L}^n(y).
		\end{align*}
		Se ora moltiplichiamo ambo i membri per $s^{n-1}$ ed integriamo rispetto ad $s$ otteniamo 
		\[ 	\int_{0}^{r} s^{n-1} \int_{\partial\textup{B}(0,1)} \abs{u(x+sw) - u(x)} d\mathcal{H}^{n-1}(w) ds  \le \int_{0}^{r} s^{n-1} \int_{\textup{B}(x,r)} \frac{\abs{Du(y)}}{\abs{x-y}^{n-1}}d\mathcal{L}^n(y) ds \]
		ossia
		\[ \int_{\textup{B}(x,r)} \abs{u(y) - u(x)} d\mathcal{L}^{n}(y) \le \frac{r^n}{n} \int_{\textup{B}(x,r)} \frac{\abs{Du(y)}}{\abs{x-y}^{n-1}}d\mathcal{L}^n(y).  \]
		Dividendo ambo i membri per la misura di $\textup{B}(x,r)$ otteniamo proprio
		\[ \fint_{\textup{B}(x,r)} \abs{u(y) - u(x)}d\mathcal{L}^n(y) \le C \int_{\textup{B}(x,r)} \frac{\abs{Du(y)}}{\abs{y-x}^{n-1}}d\mathcal{L}^n(y).\footnote{A voler essere precisi, fino ad ora non abbiamo usato nessuna delle ipotesi del Teorema.} \]
		
		Mostriamo ora che 
		\[ \int_{\textup{B}(x,r)} \dfrac{1}{\abs{x-y}^{(n-1)\frac{p}{p-1}}}d\mathcal{L}^n(y) = C r^{\frac{p-n}{p-1}} < +\infty. \]
		Per fare ciò effettuiamo prima di tutto il cambio di variabile $z = x-y$,
		\[ \int_{\textup{B}(x,r)} \dfrac{1}{\abs{x-y}^{(n-1)\frac{p}{p-1}}}d\mathcal{L}^n(y) = \int_{\textup{B}(0,r)} \dfrac{1}{\abs{z}^{(n-1)\frac{p}{p-1}}}d\mathcal{L}^n(z)\]
		ed utilizzando poi le coordinate ipersferiche, ricordando che $n<p$,
		\[ \int_{\textup{B}(0,r)} \dfrac{1}{\abs{z}^{(n-1)\frac{p}{p-1}}}d\mathcal{L}^n(z) = C\int_{0}^{r} \rho^{\frac{1-n}{p-1}}d\rho = Cr^{\frac{p-n}{p-1}} < +\infty.  \]
		
		A questo punto, dato $x \in \R^n$ e $y \in \textup{B}(x,1)$,
		\[ \abs{u(x)} \le \abs{u(x) - u(y)} + \abs{u(y)} \]
		e quindi, integrando in media rispetto a $y$,
		\[ \abs{u(x)} \le \fint_{\textup{B}(x,1)}\abs{u(x) - u(y)} d\mathcal{L}^n(y) +  \fint_{\textup{B}(x,1)}\abs{u(y)} d\mathcal{L}^n(y)\]
		e, se nel primo integrale applichiamo la stima ricavata sopra e nel secondo la disuguaglianza di Holder otteniamo 
		\[ \abs{u(x)} \le C\int_{\textup{B}(x,1)} \frac{\abs{Du(y)}}{\abs{y-x}^{n-1}}d\mathcal{L}^n(y) + C\norm{u}_{L^{p}(\textup{B}(x,1))}. \]
		Se applichiamo la disuguaglianza di Holder al primo integrale possiamo scrivere
		\[ \int_{\textup{B}(x,1)} \frac{\abs{Du(y)}}{\abs{y-x}^{n-1}}d\mathcal{L}^n(y) \le \left( \int_{\textup{B}(x,1)} \abs{Du}^pd\mathcal{L}^n\right) ^{\frac{1}{p}}\left( \int_{\textup{B}(x,1)} \dfrac{1}{\abs{x-y}^{(n-1)\frac{p}{p-1}}} d\mathcal{L}^n(y) \right)^{\frac{p-1}{p}}.\]
		Allora possiamo scrivere
		\[ \abs{u(x)} \le C\norm{u}_{W^{1,p}(\R^n)} \]
		e dall'arbitrarietà di $x$, possiamo passare al $\textup{sup}$ ottenendo
		\begin{equation}\label{stimanormainfty}
			\norm{u}_{\infty} \le C\norm{u}_{W^{1,p}(\R^n)}. 
		\end{equation}
		\indent Siano ora $x,y \in \R^n$ ed indichiamo con $r = \abs{x-y}$ e con $W = \textup{B}(x,r) \cap \textup{B}(y,r)$. Se $w \in W$,
		\[ \abs{u(x)-u(y)} \le \abs{u(x)-u(w)} + \abs{u(y)-u(w)}\]
		e quindi procedendo con un'integrazione mediata rispetto a $w$
		\begin{align*}
			\abs{u(x)-u(y)} &\le \fint_{W}\abs{u(x)-u(w)}d\mathcal{L}^n(w) + \fint_{W}\abs{u(y)-u(w)}d\mathcal{L}^n(w) \le \\
			& \le \fint_{\textup{B}(x,r)}\abs{u(x)-u(w)}d\mathcal{L}^n(w) + \fint_{\textup{B}(y,r)}\abs{u(y)-u(w)}d\mathcal{L}^n(w).
		\end{align*}
		Siccome
		\[\fint_{\textup{B}(x,r)}\abs{u(x)-u(w)}d\mathcal{L}^n(w) \le C\int_{\textup{B}(x,r)} \frac{\abs{Du(y)}}{\abs{y-x}^{n-1}}d\mathcal{L}^n(y)   \]
		e applicando la disuguaglianza di Holder
		\begin{align*}
			\int_{\textup{B}(x,r)} \frac{\abs{Du(y)}}{\abs{y-x}^{n-1}}d\mathcal{L}^n(y) &\le C \left( \int_{\textup{B}(x,r)} \abs{Du}^pd\mathcal{L}^n\right) ^{\frac{1}{p}}\left( \int_{\textup{B}(x,r)} \dfrac{1}{\abs{x-y}^{(n-1)\frac{p}{p-1}}} d\mathcal{L}^n(y) \right)^{\frac{p-1}{p}} \le \\
			&\le C r^{1-\frac{n}{p}}\norm{Du}_{L^p(\R^n)}.
		\end{align*}
		Analogamente vale una stima simile per il secondo addendo, da cui
		\[ \abs{u(x)-u(y)} \le  C r^{1-\frac{n}{p}}\norm{Du}_{L^p(\R^n)} = C \abs{x-y}^{1-\frac{n}{p}}\norm{Du}_{L^p(\R^n)}\]
		e, per l'arbitrarietà di $x,y$,
		\begin{equation}\label{stimaseminormaholder}
			\left[ u\right] _{C^{0,\gamma}(\R^n)} \le C\norm{Du}_{L^p(\R^n)}.
		\end{equation}
		La tesi segue combinando \eqref{stimanormainfty} e \eqref{stimaseminormaholder}.
	\end{proof}
\end{teorema}
\begin{teorema}\textbf{\emph{(di Morrey)}}\index[analitico]{teorema! di Morrey}
	Siano $n < p \le \infty$, $\gamma = 1 - \frac{n}{p}$, $U$ limitato con $\partial U$ di classe $C^1$ e $u \in W^{1,p}(U)$. Allora esiste $C$, dipendente solo da $p,n$ e $U$, e un rappresentante della classe di equivalenza $u$, indicato con $u^\ast$, tale che $u^\ast \in C^{0,\gamma}(\overline{U})$
	\[ \norm{u^\ast}_{ C^{0,\gamma}(\overline{U})} \le C \norm{u}_{W^{1,p}(U)}.\]
\end{teorema}
\begin{proof}
	Innanzitutto per il Teorema di estensione, sia $\overline{u} = Eu \in W^{1,p}(\R^{n})$ tale che $\overline{u} = u$ q.o. in $U$, $\overline{u}$ è a supporto compatto in $\R^n$ e 
	\[ \norm{\overline{u}}_{W^{1,p}(\R^n)} \le C \norm{u}_{W^{1,p}(U)}. \]
	Dal fatto che $\overline{u}$ è a supporto compatto, esiste $(u_h)$ in $C^\infty_c(\R^n)$ tale che $u_h \to \overline{u}$ in $W^{1,p}(\R^n)$. Ora, dalla Disuguaglianza di Morrey, 
	\[ \norm{u_h - u_k}_{C^{0,\gamma}(\R^n)} \le C\norm{u_h - u_k}_{W^{1,p}(\R^n)} \]
	allora $(u_h)$ è di Cauchy in $C^{0,\gamma}(\R^n)$, per cui è convergente ad un certo $u^\ast$ in $C^{0,\gamma}(\R^n)$. Per unicità del limite puntuale si ha che $u=u^\ast$  q.o. in $U$.  Sempre per la Disuguaglianza di Morrey,
	\[ \norm{u_h}_{C^{0,\gamma}(\overline{U})} \le \norm{u_h}_{C^{0,\gamma}(\R^n)} \le C \norm{u_h}_{W^{1,p}(\R^n)}\]
	e se $h \to +\infty$
	\[ \norm{u^\ast}_{C^{0,\gamma}(\overline{U})} \le C \norm{\overline{u}}_{W^{1,p}(\R^n)}, \]
	da cui 
	\[ \norm{u^\ast}_{ C^{0,\gamma}(\overline{U})} \le C \norm{u}_{W^{1,p}(U)}.\qedhere\]
\end{proof}
Ricapitolando, se $U$ limitato con $\partial U$ di classe $C^1$, le seguenti immersioni sono continue:
\begin{description}
	\item[\textbf{Sobolev} ($\mathbf{1\le p< n}$)] $W^{1,p}(U) \hookrightarrow L^{q}(U)$ per ogni $1 \le q \le p^\ast$,
	\item[\textbf{Trudinger--Moser} ($\mathbf{p=n}$)] $W^{1,n}(U) \hookrightarrow L^{q}(U)$ per ogni $1 \le q < \infty$,
	\item[\textbf{Morrey} ($\mathbf{p > n}$)] $W^{1,p}(U) \hookrightarrow C^{0,1 - \frac{n}{p}}(\overline{U})$, quindi, in particolare, $W^{1,p}(U) \hookrightarrow L^{q}(U)$ per ogni $1 \le q \le \infty$.
\end{description}
In altre parole, all'aumentare di $p$ rispetto ad $n$ si ha un guadagno sempre maggiore sulla regolarità.

Vediamo alcune interessanti conseguenze nel caso $p>n$.
\begin{osservazione}
	In $\R$, se $U$ limitato con $\partial U$ di classe $C^1$, ogni $u \in H^1(U)$ ammette un rappresentante continuo.
\end{osservazione}

\begin{proposizione}
	Siano $U$ limitato con $\partial U$ di classe $C^1$ ed $u : U \to \R$. I seguenti fatti sono equivalenti:
	\begin{enumerate}[$(a)$]
		\item $u$ è lipschitziana,
		\item $u \in W^{1,\infty}(U)$.
	\end{enumerate}
	
	\begin{proof} Supponiamo $U = \R^n$ ed $u$ a supporto compatto in U. Il caso generale segue dal Teorema di estensione.
		\par\noindent
		$(b) \Longrightarrow (a)$ Se $u \in W^{1,\infty}(\R^n)$, allora $u^\varepsilon = \varrho_\varepsilon \ast u$ è tale che $u^\varepsilon \to u$ uniformemente per $\varepsilon \to 0$, in quanto esiste un rappresentante holderiano, quindi continuo, e 
		\[ \norm{Du^\varepsilon}_{L^\infty(\R^n)} \le \norm{Du}_{L^\infty(\R^n)}. \]
		Dati $x,y \in \R^n$ tali che $x \ne y$, per il Teorema fondamentale del calcolo integrale,
		\[\abs*{u^\varepsilon(x) - u^\varepsilon(y)} \le \int_{0}^{1} \abs*{Du^\varepsilon(tx + (1-t)y)}dt \abs{x-y} \le \norm{Du^\varepsilon}_{L^\infty(\R^n)} \abs{x-y} \]
		da cui 
		\[ \abs*{u^\varepsilon(x) - u^\varepsilon(y)} \le \norm{Du}_{L^\infty(\R^n)}\abs{x-y} \]
		e, per $\varepsilon \to 0$, dal fatto che il limite $u^\varepsilon \to u$ è uniforme, segue la lipschitzianità di $u$.
		\par\noindent
		$(a) \Longrightarrow (b)$ Supponiamo $u$ lipschitziana. Allora esiste $C>0$ tale che per ogni $x,y \in \R^n$
		\[ \abs{u(x) -u(y)}\le C\abs{x-y}. \]
		Consideriamo $(h_k)$ tale che $h_k \to 0$. Allora
		\[ \abs{u(x-h_ke_j) -u(x)}\le Ch_k \]
		da cui 
		\[ D^{-h_k}_ju(x) \le C \]
		e quindi $\norm{ D^{-h_k}_ju}_\infty \le C$ ossia $(D^{-h_k}_ju)$ limitata in $L^\infty(\R^n)$. Per l'Osservazione \eqref{convergenzadeboleLp}, esiste $w \in L^\infty$ tale che, a meno di sottosuccessioni, per ogni $\varphi \in L^1(\R^n)$ si abbia
		\[ \int D^{-h_k}_j u \varphi d\mathcal{L}^n \to \int w \varphi d\mathcal{L}^n \]
		e quindi in particolare per ogni $\varphi \in C^\infty_c(\R^n)$. Ma per \eqref{integrazioneperpartirapportoincrementale}, 
		\[ \int D_j^{-h_k}u\varphi d\mathcal{L}^n = -\int u D^{h_k}_j\varphi d\mathcal{L}^n. \]
		Allora per il Teorema della convergenza dominata, per ogni $\varphi \in C^\infty_c(\R^n)$
		\[ \int uD_j\varphi d\mathcal{L}^n = -\int w \varphi d\mathcal{L}^n \]
		da cui $D_j u = w \in L^\infty(\R^n)$, ossia $u \in W^{1,\infty}(\R^n)$.
	\end{proof}
\end{proposizione}
\begin{osservazione}
	In modo analogo $u \in W^{1,\infty}_{loc}(U)$ se e solo se $u$ è localmente lipschitziana in $U$. Inoltre, si può dimostrare che dati $U$ un aperto qualunque e $p \in \left] n,\infty\right[ $, se $u \in W^{1,p}(U)$ allora $u \in C^{0,1-\frac{n}{p}}(U)$.
\end{osservazione}
\begin{teorema}
	Sia $n < p \le \infty.$ Se $u \in W^{1,p}_{loc}(U)$ allora $u$ è differenziabile \emph{q.o.} in $U$. 
	
	In particolare, il gradiente di $u$ è uguale al gradiente debole di $u$ \emph{q.o.} in $U$.
	\begin{proof} Consideriamo dapprima il caso $p\ne \infty.$ Innanzitutto, fissati $x,y \in U$ e posto $r=\abs{x-y}$, per ogni $v \in C^1(\textup{B}(x,2r))$, in modo analogo a quanto fatto per dimostrare la disuguaglianza di Morrey, si ha che
		\[ \abs{v(y)-v(x)} \le C r^{1-\frac{n}{p}}\left( \int_{\textup{B}(x,2r)}\abs{Dv}^pd\mathcal{L}^n\right) ^{\frac{1}{p}}.\]
		La precedente relazione, mediante approssimazione, vale poi per ogni $v \in W^{1,p}(\textup{B}(x,2r))$, identificata con il suo rappresentante holderiano, quindi continuo. Sia ora $u \in W^{1,p}_{loc}(U)$, identificata con il suo rappresentante continuo. Per una variante del Teorema sui punti di Lebesgue, per q.o. $x \in U$
		\[ \lim\limits_{r \to 0} \fint_{\textup{B}(x,2r)} \abs{Du(x)-Du(y)}^pd\mathcal{L}^n(y) = 0. \]
		Scelta $v(y) = u(y)-u(x)-Du(x)\cdot(y-x)$, osservando che $v(x) = 0$,
		\[ \abs{u(y)-u(x)-Du(x)\cdot(y-x)} = \abs{v(y)-v(x)} \le Cr^{1-\frac{n}{p}}\left( \int_{\textup{B}(x,2r)}\abs{Dv}^pd\mathcal{L}^n\right) ^{\frac{1}{p}}  \]
		ossia 
		\[ \abs{u(y)-u(x)-Du(x)\cdot(y-x)} \le C r\left( \fint_{\textup{B}(x,2r)} \abs{Du(x)-Du(y)}^pd\mathcal{L}^n(y)\right)^{\frac{1}{p}} \]
		e quindi
		\[ \dfrac{\abs{u(y)-u(x)-Du(x)\cdot(y-x)}}{\abs{y-x}} \le C \left( \fint_{\textup{B}(x,2r)} \abs{Du(x)-Du(y)}^pd\mathcal{L}^n(y)\right)^{\frac{1}{p}}. \]
		Dal Teorema del confronto si ha la tesi. 
		
		Il caso $p=\infty$ è simile in quanto $W^{1,\infty}_{loc}(U) \subseteq W^{1,p}_{loc}(U)$ per ogni $1\le p<\infty$.
	\end{proof}
\end{teorema}
\begin{teorema}\emph{\textbf{(di Rademacher)}}\index[analitico]{teorema!di Rademacher}
	Se $u:U\to \R$ è localmente lipschitziana allora è differenziabile \emph{q.o.} in $U$. 
	\begin{proof}
		Si tratta di combinare i risultati precedenti.
	\end{proof}
\end{teorema}
Di frequente utilità in situazioni non lineari è il seguente risultato.
\begin{teorema}\emph{\textbf{(disuguaglianza di Hardy)}}\index[analitico]{disuguaglianza!di Hardy}
	Siano $n \ge 3$, $r > 0$ ed $u \in H^1(\textup{B}(0,r))$. Allora $\frac{u}{\abs{x}} \in L^2(\textup{B}(0,r))$. In particolare,
	\[ \int_{\textup{B}(0,r)} \frac{u(x)}{\abs{x}} d\mathcal{L}^n(x) \le C  \int_{\textup{B}(0,r)} \left( \abs{Du}^2 + \frac{u^2}{r^2}\right)  d\mathcal{L}^n. \]
\end{teorema}
\begin{proof}
	Supponiamo $u \in C^\infty(\textup{B}(0,r))$, il caso generale segue dal Teorema di globale approssimazione con funzioni lisce.
	
	Innanzitutto,
	\[ D\left( \frac{1}{\abs{x}}\right) = -\frac{x}{\abs{x}^3}, \]
	allora, utilizzando la Formula di Gauss--Green,
	\begin{align*}
		\int_{\textup{B}(0,r)}\frac{u(x)^2}{\abs{x}^2} d\mathcal{L}^n(x) &= \int_{\textup{B}(0,r)}\frac{u(x)^2}{\abs{x}^2} \left( -\frac{x}{\abs{x}}\right)\cdot \left( -\frac{x}{\abs{x}}\right)   d\mathcal{L}^n(x) = \\
		&= - \int_{\textup{B}(0,r)} u(x)^2 D\left( \frac{1}{\abs{x}}\right)  \cdot \frac{x}{\abs{x}} d\mathcal{L}^n(x),
	\end{align*}
	quindi
	\begin{align*}
		\int_{\textup{B}(0,r)}\frac{u(x)^2}{\abs{x}^2} d\mathcal{L}^n(x) &= -\sum_{j=1}^{n} \int_{\textup{B}(0,r)} u(x)^2 D_j\left( \frac{1}{\abs{x}}\right)\frac{x_j}{\abs{x}} d\mathcal{L}^n(x) = \\
		&= -\sum_{j=1}^{n}\left[  \int_{\partial \textup{B}(0,r)}\!\!\!\!  \frac{u(x)^2}{\abs{x}^2}x_j\nu_j d\mathcal{H}^{n-1}(x) -\! \int_{ \textup{B}(0,r)}  \!\frac{1}{\abs{x}} D_j\left( x_j\frac{u(x)^2}{\abs{x}}\right)\! d\mathcal{L}^n(x)\right]\!.
	\end{align*}
	A questo punto, essendo 
	\[ D_j\left( x_j\frac{u^2}{\abs{x}}\right) = \frac{u(x)^2}{\abs{x}}D_j x_j + x_j D_j \left( \frac{u(x)^2}{\abs{x}}\right)   = \frac{u(x)^2}{\abs{x}} + 2u(x)D_j u(x) \frac{x_j}{\abs{x}} + x_j u(x)^2D_j\left( \frac{1}{\abs{x}}\right), \]
	possiamo scrivere
	\begin{align*}
		\int_{\textup{B}(0,r)}\!\!\!\frac{u(x)^2}{\abs{x}^2} d\mathcal{L}^n(x) &= \int_{\textup{B}(0,r)} \sum_{j=1}^{n}\left( \frac{u(x)^2}{\abs{x}^2} +2u(x) D_j u(x) \frac{x_j}{\abs{x}^2} + \frac{x_j}{\abs{x}}u(x)^2D_j\left( \frac{1}{\abs{x}}\right)\right) d\mathcal{L}^n(x) + \\
		&- \int_{\partial \textup{B}(0,r)}x \cdot \nu \frac{u(x)^2}{\abs{x}^2}d\mathcal{H}^{n-1}(x)
	\end{align*}
	\begin{align*}
		\int_{\textup{B}(0,r)}\!\!\!\frac{u(x)^2}{\abs{x}^2} d\mathcal{L}^n(x) &= \!\!\int_{\textup{B}(0,r)} \left( n\frac{u(x)^2}{\abs{x}^2} + 2u(x)Du(x) \cdot \frac{x}{\abs{x}^2} + u(x)^2 \frac{x}{\abs{x}}\cdot D\left( \frac{1}{\abs{x}}\right)\right) d\mathcal{L}^n(x) + \\
		&- \int_{\partial \textup{B}(0,r)}x \cdot \nu \frac{u(x)^2}{\abs{x}^2} d\mathcal{H}^{n-1}(x) = \\
		&= \int_{\textup{B}(0,r)} \left( (n-1)\frac{u(x)^2}{\abs{x}^2} + 2u(x)Du(x) \cdot \frac{x}{\abs{x}^2}\right)  d\mathcal{L}^n(x)+ \\
		&- \int_{\partial \textup{B}(0,r)}x \cdot \nu \frac{u(x)^2}{\abs{x}^2} d\mathcal{H}^{n-1}(x).
	\end{align*}
	Allora, 
	\[ (2-n)\int_{\textup{B}(0,r)}\frac{u(x)^2}{\abs{x}^2} d\mathcal{L}^n(x) = 2\int_{\textup{B}(0,r)}u(x)Du(x) \cdot \frac{x}{\abs{x}^2}  d\mathcal{L}^n(x)  - \int_{\partial \textup{B}(0,r)}x \cdot \nu \frac{u(x)^2}{\abs{x}^2} d\mathcal{H}^{n-1}(x) \]
	ed essendo $\nu = \frac{x}{\abs{x}} = \frac{x}{r}$ su $\partial \textup{B}(0,r)$, si ha
	\[ (2-n)\int_{\textup{B}(0,r)}\frac{u(x)^2}{\abs{x}^2} d\mathcal{L}^n(x) = 2\int_{\textup{B}(0,r)}u(x)Du(x) \cdot \frac{x}{\abs{x}^2}  d\mathcal{L}^n(x)  - \frac{1}{r}\int_{\partial \textup{B}(0,r)}u^2 d\mathcal{H}^{n-1}. \]
	A questo punto, utilizzando la disuguaglianza di Young,
	\begin{align*}
		\int_{\textup{B}(0,r)}\frac{u(x)^2}{\abs{x}^2} d\mathcal{L}^n(x) &= \frac{2}{2-n}\int_{\textup{B}(0,r)}u(x)Du(x) \cdot \frac{x}{\abs{x}^2}  d\mathcal{L}^n(x)  - \frac{1}{2-n}\frac{1}{r}\int_{\partial \textup{B}(0,r)}u^2 d\mathcal{H}^{n-1} \le \\
		&\le \frac{2}{2-n}\int_{\textup{B}(0,r)}\frac{\abs{u(x)}}{\abs{x}}\abs{Du(x)} d\mathcal{L}^n(x)  - \frac{1}{2-n}\frac{1}{r}\int_{\partial \textup{B}(0,r)}u^2 d\mathcal{H}^{n-1} \le \\
		&\le\frac{1}{2-n}\int_{\textup{B}(0,r)}\frac{\abs{u(x)}^2}{\abs{x}^2} d\mathcal{L}^n(x) + \frac{1}{2-n}\int_{\textup{B}(0,r)}\abs{Du}^2 d\mathcal{L}^n +\\
		&-\frac{1}{2-n}\frac{1}{r}\int_{\partial \textup{B}(0,r)}u^2 d\mathcal{H}^{n-1},
	\end{align*}
	e quindi
	\[ \frac{2n}{n-2}\int_{\textup{B}(0,r)}\frac{u(x)^2}{\abs{x}^2} d\mathcal{L}^n(x) \le \frac{1}{2-n}\int_{\textup{B}(0,r)}\abs{Du}^2 d\mathcal{L}^n -\frac{1}{2-n}\frac{1}{r}\int_{\partial \textup{B}(0,r)}u^2 d\mathcal{H}^{n-1}. \]
	Allora, 
	\[ \int_{\textup{B}(0,r)}\frac{u(x)^2}{\abs{x}^2} d\mathcal{L}^n(x) \le -\frac{1}{2n}\int_{\textup{B}(0,r)}\abs{Du}^2 d\mathcal{L}^n + \frac{1}{2n}\frac{1}{r}\int_{\partial \textup{B}(0,r)}u^2 d\mathcal{H}^{n-1} \]
	ed esisterà $C> 0$ tale che
	\[ \int_{\textup{B}(0,r)}\frac{u(x)^2}{\abs{x}^2} d\mathcal{L}^n(x) \le C\int_{\textup{B}(0,r)}\abs{Du}^2 d\mathcal{L}^n + \frac{C}{r}\int_{\partial \textup{B}(0,r)}u^2 d\mathcal{H}^{n-1}. \]
	Osserviamo ora che, essendo $r = x \cdot \nu$,
	\begin{align*}
		r^2\int_{\partial \textup{B}(0,r)}u^2 d\mathcal{H}^{n-1} &= \int_{ \textup{B}(0,r)} \textup{div} \left( xu^2\right)  d\mathcal{L}^n(x) = \int_{ \textup{B}(0,r)}\left( nu^2 + 2uDu \cdot x \right) d\mathcal{L}^n(x)
	\end{align*}
	ed applicando la disuguaglianza di Young,
	\[ 	r^2\int_{\partial \textup{B}(0,r)}u^2 d\mathcal{H}^{n-1} \le C\int_{\textup{B}(0,r)} \left( u^2 + r^2\abs{Du}^2 \right) d\mathcal{L}^n. \]
	Dividendo per $r^2$ otteniamo la disuguaglianza
	\[ \frac{1}{r}\int_{\partial \textup{B}(0,r)}u^2 d\mathcal{H}^{n-1} \le C\int_{\textup{B}(0,r)} \left( \frac{u^2}{r^2} + \abs{Du}^2 \right) d\mathcal{L}^n. \]
	In conclusione,
	\[ \int_{\textup{B}(0,r)} \frac{u(x)}{\abs{x}} d\mathcal{L}^n(x) \le C  \int_{\textup{B}(0,r)} \left( \abs{Du}^2 + \frac{u^2}{r^2}\right)  d\mathcal{L}^n \]
	e la dimostrazione è conclusa.
\end{proof}

Volendo generalizzare tutti i ragionamenti fatti nel caso $k=1$ alla situazione generale, si arriva al seguente risultato.
\begin{teorema}\textbf{\emph{(immersioni di Sobolev)}}\index[analitico]{immersioni! di Sobolev}
	Sia $U$ limitato con $\partial U$ di classe $C^1$ e $u \in W^{k,p}(U)$. Valgono i seguenti fatti:
	\begin{enumerate}[$(a)$]
		\item se $kp < n$, allora se $q$ è tale che
		\[ \frac{1}{q} = \frac{1}{p}-\frac{k}{n} \]
		si ha che $u \in L^q(U)$ ed esiste $C$, dipendente solo da $k,p,n$ ed $U$, tale che
		\[ \norm{u}_{L^q(U)} \le C\norm{u}_{W^{k,p}(U)}, \]
		\item se $kp > n$, allora se $\gamma$ è tale che
		\[ \gamma = \begin{cases}
			\lfloor\frac{n}{p}\rfloor +1 - \frac{n}{p} & \textup{se } \frac{n}{p} \notin \Z, \\
			\textup{qualsiasi } k \in \left] 0,1\right[ & \textup{altrimenti},
		\end{cases} \]
		si ha che $u \in C^{k - \lfloor\frac{n}{p}\rfloor  -1,\gamma}(\overline{U})$ ed esiste $C$, dipendente solo da $k,p,n,\gamma$ ed $U$ tale che
		\[ \norm{u}_{C^{k - \lfloor\frac{n}{p}\rfloor  -1,\gamma}(\overline{U})} \le C\norm{u}_{W^{k,p}(U)}.\]
	\end{enumerate}
\end{teorema}
\begin{proof}
	Omettiamo la dimostrazione.
\end{proof}
Possiamo schematizzare il precedente risultato nel seguente modo.
\begin{center}
	\begin{tikzpicture}[x=0.8pt,y=0.8pt,yscale=-1,xscale=1]
		%uncomment if require: \path (0,300); %set diagram left start at 0, and has height of 300
		
		%Shape: Axis 2D [id:dp27904890581218256] 
		\draw  (50,248.5) -- (554,248.5)(100.4,13.6) -- (100.4,274.6) (547,243.5) -- (554,248.5) -- (547,253.5) (95.4,20.6) -- (100.4,13.6) -- (105.4,20.6)  ;
		%Straight Lines [id:da651292971601942] 
		\draw  [dash pattern={on 0.84pt off 2.51pt}]  (100,50.6) -- (161,50.6) -- (531,50.6) ;
		%Shape: Circle [id:dp21548707964318892] 
		\draw  [fill={rgb, 255:red, 0; green, 0; blue, 0 }  ,fill opacity=1 ] (321.5,130.1) .. controls (321.5,128.44) and (320.16,127.1) .. (318.5,127.1) .. controls (316.84,127.1) and (315.5,128.44) .. (315.5,130.1) .. controls (315.5,131.76) and (316.84,133.1) .. (318.5,133.1) .. controls (320.16,133.1) and (321.5,131.76) .. (321.5,130.1) -- cycle ;
		%Straight Lines [id:da32979052674260245] 
		\draw  [dash pattern={on 0.84pt off 2.51pt}]  (103,129.6) -- (534,130.6) ;
		%Shape: Circle [id:dp751169393115473] 
		\draw  [fill={rgb, 255:red, 0; green, 0; blue, 0 }  ,fill opacity=1 ] (213,130) .. controls (213,128.34) and (211.66,127) .. (210,127) .. controls (208.34,127) and (207,128.34) .. (207,130) .. controls (207,131.66) and (208.34,133) .. (210,133) .. controls (211.66,133) and (213,131.66) .. (213,130) -- cycle ;
		%Shape: Circle [id:dp8645473255346661] 
		\draw  [fill={rgb, 255:red, 0; green, 0; blue, 0 }  ,fill opacity=1 ] (321.5,72.1) .. controls (321.5,70.44) and (320.16,69.1) .. (318.5,69.1) .. controls (316.84,69.1) and (315.5,70.44) .. (315.5,72.1) .. controls (315.5,73.76) and (316.84,75.1) .. (318.5,75.1) .. controls (320.16,75.1) and (321.5,73.76) .. (321.5,72.1) -- cycle ;
		%Shape: Circle [id:dp29280488987845565] 
		\draw  [fill={rgb, 255:red, 0; green, 0; blue, 0 }  ,fill opacity=1 ] (212.5,201.1) .. controls (212.5,199.44) and (211.16,198.1) .. (209.5,198.1) .. controls (207.84,198.1) and (206.5,199.44) .. (206.5,201.1) .. controls (206.5,202.76) and (207.84,204.1) .. (209.5,204.1) .. controls (211.16,204.1) and (212.5,202.76) .. (212.5,201.1) -- cycle ;
		%Shape: Circle [id:dp40222412538059893] 
		\draw  [fill={rgb, 255:red, 0; green, 0; blue, 0 }  ,fill opacity=1 ] (503.5,92.1) .. controls (503.5,90.44) and (502.16,89.1) .. (500.5,89.1) .. controls (498.84,89.1) and (497.5,90.44) .. (497.5,92.1) .. controls (497.5,93.76) and (498.84,95.1) .. (500.5,95.1) .. controls (502.16,95.1) and (503.5,93.76) .. (503.5,92.1) -- cycle ;
		%Shape: Circle [id:dp31999878516510627] 
		\draw  [fill={rgb, 255:red, 0; green, 0; blue, 0 }  ,fill opacity=1 ] (423.5,172.1) .. controls (423.5,170.44) and (422.16,169.1) .. (420.5,169.1) .. controls (418.84,169.1) and (417.5,170.44) .. (417.5,172.1) .. controls (417.5,173.76) and (418.84,175.1) .. (420.5,175.1) .. controls (422.16,175.1) and (423.5,173.76) .. (423.5,172.1) -- cycle ;
		%Shape: Boxed Bezier Curve [id:dp9244206607185501] 
		\draw    (331,109.48) .. controls (331,120.43) and (325,126.41) .. (325,112.47) .. controls (325,98.8) and (325,109.62) .. (325,82.32) ;
		\draw [shift={(325,80.6)}, rotate = 90] [color={rgb, 255:red, 0; green, 0; blue, 0 }  ][line width=0.75]    (10.93,-3.29) .. controls (6.95,-1.4) and (3.31,-0.3) .. (0,0) .. controls (3.31,0.3) and (6.95,1.4) .. (10.93,3.29)   ;
		%Shape: Boxed Bezier Curve [id:dp5390642291075387] 
		\draw    (288.96,136.53) .. controls (303.8,136.45) and (311.92,141.44) .. (293.04,141.55) .. controls (274.44,141.66) and (289.34,141.57) .. (251.64,141.79) ;
		\draw [shift={(249.88,141.8)}, rotate = 359.66] [color={rgb, 255:red, 0; green, 0; blue, 0 }  ][line width=0.75]    (10.93,-3.29) .. controls (6.95,-1.4) and (3.31,-0.3) .. (0,0) .. controls (3.31,0.3) and (6.95,1.4) .. (10.93,3.29)   ;
		%Shape: Boxed Bezier Curve [id:dp6379681515373254] 
		\draw    (476.6,120.17) .. controls (491.3,101.56) and (493.44,86.78) .. (474.74,110.46) .. controls (456.23,133.91) and (471.23,114.91) .. (433.17,163.12) ;
		\draw [shift={(432,164.6)}, rotate = 308.29] [color={rgb, 255:red, 0; green, 0; blue, 0 }  ][line width=0.75]    (10.93,-3.29) .. controls (6.95,-1.4) and (3.31,-0.3) .. (0,0) .. controls (3.31,0.3) and (6.95,1.4) .. (10.93,3.29)   ;
		%Straight Lines [id:da8211205559412464] 
		\draw  [dash pattern={on 0.84pt off 2.51pt}]  (210,20) -- (210,73.6) -- (210,246) ;
		%Straight Lines [id:da7008636267010666] 
		\draw  [dash pattern={on 0.84pt off 2.51pt}]  (318.5,20.1) -- (318.5,246.1) ;
		
		% Text Node
		\draw (77,40) node [anchor=north west][inner sep=0.75pt]    {$L^{1}$};
		% Text Node
		\draw (78,122) node [anchor=north west][inner sep=0.75pt]    {$L^{2}$};
		% Text Node
		\draw (76,227) node [anchor=north west][inner sep=0.75pt]    {$L^{\infty }$};
		% Text Node
		\draw (331,62.4) node [anchor=north west][inner sep=0.75pt]    {$W^{2,p}$};
		% Text Node
		\draw (215,133.4) node [anchor=north west][inner sep=0.75pt]    {$H^{1}$};
		% Text Node
		\draw (323.5,133.5) node [anchor=north west][inner sep=0.75pt]    {$H^{2}$};
		% Text Node
		\draw (220,191.4) node [anchor=north west][inner sep=0.75pt]    {$W^{1,p}$};
		% Text Node
		\draw (511,82.4) node [anchor=north west][inner sep=0.75pt]    {$W^{k,p}$};
		% Text Node
		\draw (363,177.4) node [anchor=north west][inner sep=0.75pt]    {$W^{k-1,p^{*}}$};
		% Text Node
		\draw (59,9.4) node [anchor=north west][inner sep=0.75pt]    {$1/p$};
		% Text Node
		\draw (563,243.4) node [anchor=north west][inner sep=0.75pt]    {$k$};
	\end{tikzpicture}
\end{center}
\section{Teoremi di compattezza}
Presentiamo un risultato di compattezza per gli spazi di Sobolev nel caso $1 \le p < n$.
\begin{teorema}\textbf{\emph{(di Rellich--Kondrachov, $\mathbf{1\le p<n}$)}}\index[analitico]{teorema! di Rellich--Kondrachov}
	Siano $1 \le p < n$, $U$ limitato e con $\partial U$ di classe $C^1$. Allora, per ogni $1 \le q < p^\ast$, l'immersione $W^{1,p}(U) \hookrightarrow L^q(U)$ è compatta.
\end{teorema}
\begin{proof}
	Sia $(u_h)$ una successione limitata in $W^{1,p}(U)$. Allora esiste $M>0$ tale che per ogni $h \in \N$,
	\[ \norm{u_h}_{W^{1,p}(U)} \le M. \]
	In particolare, 
	\[ \underset{h}{\textup{sup}}\norm{u_h}_{W^{1,p}(U)} \le M < +\infty. \]
	In conformità al Teorema di estensione siano $V, W$ aperti limitati tali che $\overline{U} \subseteq V$ e $\overline{V} \subseteq W$ e, per ogni $h \in \N$, $\overline{u}_h = Eu_h$ tali che $\overline{u}_h = 0$ in $\R^n\setminus V$. In particolare, $(\overline{u}_h)$ limitata in $W^{1,p}(U)$ e
	\[ \underset{h}{\textup{sup}}\norm{\overline{u}_h}_{W^{1,p}(V)} < +\infty. \]
	Osserviamo subito che, essendo $V$ limitato, 
	\[ \underset{h}{\textup{sup}}\norm{\overline{u}_h}_{L^1(V)} < +\infty.\]
	Consideriamo le $\varepsilon$-mollificate delle $\overline{u}_h$, $\overline{u}_h^\varepsilon$. Allora, per $\varepsilon$ sufficientemente piccolo, le $\overline{u}_h^\varepsilon$ avranno supporto compatto in $W$. In particolare $\overline{u}_h^\varepsilon = 0$ in $\R^n\setminus W$.
	
	Mostriamo che 
	\[ \lim\limits_{\varepsilon \to 0} \left(\underset{h}{\textup{sup}}\norm*{\overline{u}_h^\varepsilon -\overline{u}_h}_{L^q(W)} \right) = 0  \]
	ossia che $\overline{u}_h^\varepsilon \to \overline{u}_h$ uniformemente in $L^q(W)$. Supponiamo inizialmente $\overline{u}_h$ liscia. Si ha che
	\[ \abs{\overline{u}_h^\varepsilon(x) - \overline{u}_h(x)} = \abs*{\frac{1}{\varepsilon^n} \int_{\textup{B}(0,\varepsilon)} \varrho\left( \frac{y}{\varepsilon}\right) \overline{u}_h(x-y) d\mathcal{L}^n(y) - \overline{u}_h(x)\int_{\textup{B}(0,1)}\varrho(y)d\mathcal{L}^n(y)}. \]
	Operando un cambio di variabile nel primo integrale del tipo $y' = \frac{y}{\varepsilon}$ possiamo scrivere, denotando comunque la variabile di integrazione con $y$, 
	\[ 	\abs{\overline{u}_h^\varepsilon(x) - \overline{u}_h(x)} \le \int_{\textup{B}(0,1)}\varrho(y)\abs*{\overline{u}_h(x-\varepsilon y) - \overline{u}_h(x)}d\mathcal{L}^n(y)\]
	e per il Teorema fondamentale del calcolo integrale
	\begin{align*}
		\abs*{\overline{u}_h(x-\varepsilon y) - \overline{u}_h(x)} &= \abs*{ \int_{0}^{1} \frac{d}{dt}\overline{u}_h(x-\varepsilon t y)dt} \le \\
		& \le \int_{0}^{1} \abs{D\overline{u}_h(x-\varepsilon t y) \cdot (-\varepsilon y)}dt \le \varepsilon \int_{0}^{1} \abs{D\overline{u}_h(x-\varepsilon t y)}dt,
	\end{align*}
	quindi
	\[ \abs{\overline{u}_h^\varepsilon(x) - \overline{u}_h(x)} \le \varepsilon \int_{\textup{B}(0,1)} \varrho(y) \int_{0}^{1} \abs{D\overline{u}_h(x-\varepsilon t y)}dt d\mathcal{L}^n(y), \]
	che integrata in $W$ fornisce
	\begin{align*}
		\int_W \abs{\overline{u}_h^\varepsilon(x) - \overline{u}_h(x)} d\mathcal{L}^n(x) &\le \varepsilon \int_W \int_{\textup{B}(0,1)} \varrho(y) \int_{0}^{1} \abs{D\overline{u}_h(x-\varepsilon t y)}dt d\mathcal{L}^n(y)d\mathcal{L}^n(x) = \\
		&= \varepsilon \int_{\textup{B}(0,1)} \varrho(y) \int_{0}^{1} \int_W \abs{D\overline{u}_h(x-\varepsilon t y)} d\mathcal{L}^n(x) dt d\mathcal{L}^n(y).
	\end{align*}
	Operando il cambio di variabile $z = x - \varepsilon t y$, otteniamo 
	\[ \int_W \abs{D\overline{u}_h(x-\varepsilon t y)} d\mathcal{L}^n(x) = \int_{W_{\varepsilon,t}} \abs{D\overline{u}_h(z)}d\mathcal{L}^n(z) \]
	ma fuori da $W$ si ha $\overline{u}_h = 0$, quindi 
	\[ \int_W \abs{D\overline{u}_h(x-\varepsilon t y)} d\mathcal{L}^n(x) = \int_{W} \abs{D\overline{u}_h(z)}d\mathcal{L}^n(z). \]
	Riassumendo,
	\[ \int_W \abs{\overline{u}_h^\varepsilon(x) - \overline{u}_h(x)} d\mathcal{L}^n(x)\le \varepsilon \int_{W} \abs{D\overline{u}_h(z)}d\mathcal{L}^n(z). \]
	Se invece $\overline{u}_h \in W^{1,p}(V)$, costruiamo una successione di funzioni lisce che la approssima e poi con il Teorema della convergenza dominata otteniamo ancora la stessa stima, perciò per la disuguaglianza di Holder si ha
	\begin{align*}
		\norm{\overline{u}_h^\varepsilon - \overline{u}_h}_{L^1(W)} &\le \varepsilon \int_{W} \abs{D\overline{u}_h(z)}d\mathcal{L}^n(z) \le \varepsilon C \norm{D\overline{u}_h}_{L^p(W)}\le \\
		&\le \varepsilon C\norm{\overline{u}_h}_{W^{1,p}(V)} \le \varepsilon C\underset{h}{\textup{sup}} \norm{\overline{u}_h}_{W^{1,p}(V)}. 
	\end{align*}
	A questo punto, per la disuguaglianza di interpolazione, detto $\vartheta \in \left] 0,1\right[ $ tale che
	\[ \frac{1}{q} = \vartheta + \frac{(1-\vartheta)}{p^\ast}, \]
	\[  \norm{\overline{u}_h^\varepsilon - \overline{u}_h}_{L^q(W)} \le \norm{\overline{u}_h^\varepsilon - \overline{u}_h}_{L^1(W)}^{\vartheta}\norm{\overline{u}_h^\varepsilon - \overline{u}_h}_{L^{p^\ast}(W)}^{(1-\vartheta)} \]
	e per il Teorema di Sobolev,
	\[ \norm{\overline{u}_h^\varepsilon - \overline{u}_h}_{L^q(W)} \le C \norm{\overline{u}_h^\varepsilon - \overline{u}_h}_{L^1(W)}^{\vartheta} \le \varepsilon^\vartheta C. \]
	Passando al sup possiamo scrivere
	\[ \underset{h}{\textup{sup}}\norm*{\overline{u}_h^\varepsilon -\overline{u}_h}_{L^q(W)} \le \varepsilon^\vartheta C \]
	ed il Teorema del confronto restituisce il limite desiderato. Mostriamo ora che, ad $\varepsilon >0$ fissato, la successione $(\overline{u}_h^\varepsilon)$ è limitata ed equi-uniformemente continua. Innanzitutto, 
	\[ \abs{\varrho_\varepsilon(x)} = \abs*{\frac{1}{\varepsilon^n}\varrho\left( \frac{x}{\varepsilon}\right) } \le \frac{1}{\varepsilon^n}, \]
	quindi $\norm{\varrho_\varepsilon}_{L^\infty(\R^n)} \le \frac{1}{\varepsilon^n}$ e per ogni $h \in \N$,
	\[ \abs{\overline{u}_h^\varepsilon(x)} \le \int_{\textup{B}(x,\varepsilon)} \varrho_\varepsilon(x-y)\abs*{\overline{u}_h(y)}d\mathcal{L}^n(y) \le \norm{\varrho_\varepsilon}_{L^\infty(\R^n)} \norm{\overline{u}_h}_{L^1(V)} \le \frac{C}{\varepsilon^n} < +\infty \]
	e quindi 
	\[ \norm{\overline{u}_h^\varepsilon}_\infty \le \frac{C}{\varepsilon^n} \]
	da cui la limitatezza di $(\overline{u}_h^\varepsilon)$. Inoltre, in modo analogo si vede che
	\[ \abs{D\overline{u}_h^\varepsilon(x)} \le \frac{C}{\varepsilon^{n+1}}< +\infty.  \]
	Ora, dati $x,y \in \R^n$, per la disuguaglianza di Lagrange, esiste $\xi$ nel segmento congiungente $x,y$ tale che
	\[ \abs{\overline{u}_h^\varepsilon(x)- \overline{u}_h^\varepsilon(y)} \le \abs{D\overline{u}_h^\varepsilon(\xi)} \abs{x-y}. \]
	Dato $\eta >0$, sia $\zeta > 0$ tale che
	\[  \zeta  < \frac{\eta \varepsilon^{n+1}}{C}. \]
	Allora per ogni $h \in \N$, per ogni $x,y \in \R^n$ tali che $\abs{x-y} < \zeta $, 
	\[ \abs{\overline{u}_h^\varepsilon(x)- \overline{u}_h^\varepsilon(y)} < \eta \]
	per cui è equi-uniformemente continua su $\R^n$. In particolare lo è su $W$.
	
	Fissato $\delta > 0$, mostriamo che esiste una sottosuccessione $(\overline{u}_{h_j})$ tale che 
	\[ \lim\limits_{j,k} \norm{\overline{u}_{h_j}-\overline{u}_{h_k}}_{L^q(W)} \le \delta. \]
	Per fare ciò, sia $\varepsilon >0$ sufficientemente piccolo, tale che 
	\[ \norm{\overline{u}_h^\varepsilon - \overline{u}_h}_{L^q(W)} \le \frac{\delta}{2}. \]
	Essendo poi le $\overline{u}_h^\varepsilon$ a supporto compatto in $W$\footnote{Si usa quindi qui il fatto di aver esteso.}, possiamo applicare il Teorema di Ascoli--Arzelà ottenendo una sottosuccessione $(\overline{u}_{h_j}^\varepsilon)$ convergente uniformemente ad un qualche $w$ in $W$. In particolare, per ogni $\eta > 0$, esiste $i \in \N$ tale che per ogni $j \ge i$, 
	\[ \norm{\overline{u}_{h_j}^\varepsilon - w}_\infty < \eta \]
	e quindi 
	\[  \abs{\overline{u}_{h_j}^\varepsilon(x) - w(x)}^q < \eta^q. \]
	Allora, per ogni $j,k \ge i$,
	\[ \int_W \abs{\overline{u}_{h_j}^\varepsilon(x) - \overline{u}_{h_k}^\varepsilon(x)}^qd\mathcal{L}^n(x) < C\eta^q, \]
	ossia
	\[ \underset{j,k}{\textup{lim}} \norm{\overline{u}_{h_j}^\varepsilon-\overline{u}_{h_k}^\varepsilon}_{L^q(W)} = 0.\]
	Allora, 
	\begin{align*}
		\norm{\overline{u}_{h_j} - \overline{u}_{h_k}}_{L^q(W)} &\le \norm{\overline{u}_{h_j} - \overline{u}_{h_j}^\varepsilon}_{L^q(W)} +\norm{\overline{u}_{h_k} - \overline{u}_{h_k}^\varepsilon}_{L^q(W)} + \norm{\overline{u}_{h_j}^\varepsilon - \overline{u}_{h_k}^\varepsilon}_{L^q(W)} \le \\
		& \le \delta + \norm{\overline{u}_{h_j}^\varepsilon - \overline{u}_{h_k}^\varepsilon}_{L^q(W)}
	\end{align*}
	e passando al limite per $j,k \to +\infty,$
	\[ \lim\limits_{j,k} \norm{\overline{u}_{h_j} - \overline{u}_{h_k}}_{L^q(W)} \le \delta.\]
	Dato $\delta_1 = 1$, esiste quindi una sottosuccessione $(\overline{u}_{h_j})$ tale che 
	\[ \lim\limits_{j,k} \norm{\overline{u}_{h_j} - \overline{u}_{h_k}}_{L^q(W)} \le 1.\]
	Ora, dato $\delta_2 = \frac{1}{2}$, applicando lo stesso ragionamento a $(\overline{u}_{h_j})$, possiamo estrarre una sottosuccessione $(\overline{u}_{h_{j_l}})$ tale che 
	\[ \lim\limits_{l,k} \norm{\overline{u}_{h_{j_l}} - \overline{u}_{h_{j_k}}}_{L^q(W)} \le \frac{1}{2}.\]
	Procedendo ricorsivamente con $\delta_n = \frac{1}{n}$, si ottiene applicando il metodo della diagonale di Cantor, a meno di sottosuccessioni, una sottosuccessione $(\overline{u}_{h_l})$ tale che per ogni $\varepsilon > 0$ esiste $\overline{n} \in \N$ tale che per ogni $h_l,h_k \ge \overline{n}$ si abbia
	\[ \norm{\overline{u}_{h_l}-\overline{u}_{h_k}}_{L^q(W)}<\varepsilon, \]
	ossia $(\overline{u}_{h_l})$ è di Cauchy. A maggior ragione,
	\[ \norm{\overline{u}_{h_l}-\overline{u}_{h_k}}_{L^q(U)}<\varepsilon, \]
	quindi $(u_{h_l})$ è di Cauchy in $L^q(U)$ e, per completezza di $L^q(U)$, è convergente.
\end{proof}

Vediamo un controesempio alla validità del Teorema di Rellich per $q = p^\ast$.
\begin{esempio}
	Sia $\varphi \in C^\infty_c(\textup{B}(0,1))$. Consideriamo la successione $(\varphi_h)$ tale che
	\[ \varphi_h (x) = h^{\frac{n-p}{p}}\varphi(hx).  \]
	Evidentemente ha limite puntuale $0$. Effettuando poi il cambio di variabile $y = hx$,
	\[ \norm{\varphi_h}_{L^{p^\ast}(\textup{B}(0,1))}^{p^\ast} =\int_{\textup{B}(0,1)}h^n\abs*{\varphi(hx)}^{p^\ast} d\mathcal{L}^n(x) = \int_{\textup{B}(0,1)} \abs{\varphi(y)}^{p^\ast}d\mathcal{L}^n(y) =: C > 0. \]
	Invece, con lo stesso cambio di variabile,
	\begin{align*}
		\norm{\varphi_h}_{W^{1,p}(\textup{B}(0,1))}^{p} &= h^{n-p}\int_{\textup{B}(0,1)}\abs*{\varphi(hx)}^{p} d\mathcal{L}^n(x) + h^n\int_{\textup{B}(0,1)}\abs*{D\varphi(hx)}^{p} d\mathcal{L}^n(x) = \\
		&= \frac{1}{h^p} \int_{\textup{B}(0,1)}\abs*{\varphi\left(y \right)}^{p}d\mathcal{L}^n(y) + \int_{\textup{B}(0,1)}\abs*{D\varphi\left(y \right)}^{p}d\mathcal{L}^n(y). 
	\end{align*}
	Chiaramente, 
	\[ \lim\limits_{h} \frac{1}{h^p} \int_{\textup{B}(0,1)}\abs*{\varphi\left(y \right)}^{p}d\mathcal{L}^n(y) = 0,\]
	allora esiste $M > 0 $ tale che per ogni $h \in \N$,  
	\[ \frac{1}{h^p} \int_{\textup{B}(0,1)}\abs*{\varphi\left(y \right)}^{p}d\mathcal{L}^n(y) \le M. \]
	Allora 
	\[ \norm{\varphi_h}_{W^{1,p}(\textup{B}(0,1))}^{p} \le M + \int_{\textup{B}(0,1)}\abs*{D\varphi\left(y \right)}^{p}d\mathcal{L}^n(y) < +\infty. \]
	In particolare $(\varphi_h)$ è limitata in $W^{1,p}(\textup{B}(0,1))$. 
	
	Se, per assurdo, l'immersione fosse compatta, esisterebbe una sottosuccessione $(\varphi_{h_k})$ convergente ad un certo $w$ in $L^{p^\ast}(\textup{B}(0,1))$. 
	In particolare, da
	\[ \norm{\varphi_{h_k}}_{L^{p^\ast}(\textup{B}(0,1))}^{p^\ast} = C > 0 \]
	si ha che 
	\[ \norm{w}_{L^{p^\ast}(\textup{B}(0,1))}^{p^\ast} = C > 0. \]
	Tuttavia esisterà un'altra sottosuccessione $(\varphi_{h_{k_j}})$ tale che $\varphi_{h_{k_j}}(x) \to w(x)$ \emph{q.o.} in $\textup{B}(0,1)$. Ma allora $w(x) = 0$ \emph{q.o.} in $\textup{B}(0,1)$ perché il limite puntuale di $(\varphi_h)$ è $0$. Allora
	\[ \norm{w}_{L^{p^\ast}(\textup{B}(0,1))}^{p^\ast} = 0, \]
	assurdo.
\end{esempio}
Vediamo un controesempio alla validità del Teorema di Rellich senza ipotesi di limitatezza su $U$.
\begin{esempio}
	Sia $\varphi \in C^\infty_c(\R^n)$. Consideriamo la successione $(\varphi_h)$ tale che
	\[ \varphi_h (x) = \varphi(x+h).  \]
	Evidentemente ha limite puntuale $0$. Inoltre, per ogni $q \in \left[ 1,p^\ast\right[$, 
	\[ \norm{\varphi_h}_{L^{q}(\R^n)}^{q} = \int \abs{\varphi}^{q^\ast}d\mathcal{L}^n =: C > 0. \]
	Invece,
	\begin{align*}
		\norm{\varphi_h}_{W^{1,p}(\R^n)}^{p} \le (\norm{\varphi}_\infty^p + \norm{D\varphi}_\infty^p) \mathcal{L}^n(\textup{supt}(\varphi)) <+\infty
	\end{align*}
	allora $(\varphi_h)$ è limitata in $W^{1,p}(\R^n)$. 
	
	Se, per assurdo, l'immersione fosse compatta, esisterebbe una sottosuccessione $(\varphi_{h_k})$, dipendente dalla scelta di $q$, convergente ad un certo $w$ in $L^{q}(\R^n)$. 
	In particolare, da
	\[ \norm{\varphi_{h_k}}_{L^{q}(\R^n)}^{q} = C > 0 \]
	si ha che 
	\[ \norm{w}_{L^{q}(\R^n)}^{q} = C > 0. \]
	Tuttavia esisterà un'altra sottosuccessione $(\varphi_{h_{k_j}})$ tale che $\varphi_{h_{k_j}}(x) \to w(x)$ \emph{q.o.} in $\R^n$. Ma allora $w(x) = 0$ \emph{q.o.} in $\R^n$ perché il limite puntuale di $(\varphi_h)$ è $0$. Allora
	\[ \norm{w}_{L^{q}(\R^n)}^{q} = 0, \]
	assurdo.
\end{esempio}

Grazie al Teorema di Rellich possiamo dimostrare una variante della disuguaglianza di Poincaré.
\begin{teorema}\emph{\textbf{(di Poincaré a media nulla)}}\label{poincaréamedianulla}
	Siano $1 \le p < n$, $U$ limitato, connesso e con $\partial U$ di classe $C^1$. Allora esiste $C> 0$, dipendente solo da $n,p$ ed $U$, tale che per ogni $u \in W^{1,p}(U)$
	\[ \norm{u - (u)_{U}}_{L^p(U)} \le C \norm{Du}_{L^p(U)}\]
	dove $(u)_{U} = \fint_U u d\mathcal{L}^n$.
\end{teorema}
\begin{proof}
	Per assurdo, per ogni $h \in \N$, esiste $u_h \in W^{1,p}(U)$ tale che
	\[\norm{u_h - (u_h)_{U}}_{L^p(U)} > h \norm{Du_h}_{L^p(U)}.\]
	Chiamiamo poi 
	\[ v_h = \dfrac{u_h - (u_h)_{U}}{\norm{u_h - (u_h)_{U}}_{L^p(U)}}. \]
	Evidentemente, per ogni $h \in \N$
	\begin{align*}
		(v_h)_{U} &= 0, & \norm{v_h}_{L^p(U)} &= 1, & \norm{Dv_h}_{L^p(U)} &< \frac{1}{h}.
	\end{align*}
	Allora $(v_h)$ è limitata in $W^{1,p}(U)$. Per il Teorema di Rellich, esisterà $(v_{h_k})$ convergente ad un qualche $w$ in $L^p(U)$. In particolare, 
	\begin{align*}
		(w)_{U} &= 0, & \norm{w}_{L^p(U)} &= 1,
	\end{align*}
	quindi $D_j v_{h_k} \to 0$ in $L^p(U)$. 
	Ora, per ogni $\varphi \in C^\infty_c(U)$, per ogni $j = 1,\dots,n$,
	\[ \int_U v_{h_k} D_j \varphi d\mathcal{L}^n= - \int_U  D_jv_{h_k} \varphi d\mathcal{L}^n \]
	e per il Teorema della convergenza dominata,
	\[ \int_U w D_j \varphi d\mathcal{L}^n= \int_U0 D_j\varphi d\mathcal{L}^n. \]
	Allora $w \in W^{1,p}(U)$ con $Dw = 0$. Dal fatto che $U$ è connesso, segue $w$ è costante q.o. in $U$. Tuttavia, essendo $(w)_{U} = 0$ , dovrebbe essere allora $\norm{w}_{L^p(U)} = 0$, in contraddizione con $\norm{w}_{L^p(U)} = 1$.
\end{proof}

Non approfondiamo nel dettaglio il caso $p=n$. Per il caso $p > n$, il risultato di compattezza è il seguente.
\begin{teorema}\textbf{\emph{(di Rellich--Kondrachov, $\mathbf{p>n}$)}}\index[analitico]{teorema! di Rellich--Kondrachov}
	Siano $p > n$, $U$ limitato e con $\partial U$ di classe $C^1$. Allora, per ogni $1 \le q < p^\ast$, l'immersione $W^{1,p}(U) \hookrightarrow C^{0,1-\frac{n}{p}}$ è compatta.
\end{teorema}
\begin{proof}
	Omettiamo la dimostrazione.
\end{proof}
Comunque, ricapitolando, se $U$ limitato con $\partial U$ di classe $C^1$, le seguenti immersioni sono compatte:
\begin{description}
	\item[\textbf{Sobolev} ($\mathbf{1\le p< n}$)] $W^{1,p}(U) \hookrightarrow L^{q}(U)$ per ogni $1 \le q < p^\ast$,
	\item[\textbf{Trudinger--Moser} ($\mathbf{p=n}$)] $W^{1,n}(U) \hookrightarrow L^{q}(U)$ per ogni $1 \le q < \infty$,
	\item[\textbf{Morrey} ($\mathbf{p > n}$)] $W^{1,p}(U) \hookrightarrow C^{0,1 - \frac{n}{p}}(\overline{U})$, quindi, in particolare, $W^{1,p}(U) \hookrightarrow L^{q}(U)$ per ogni $1 \le q \le \infty$.
\end{description}

\begin{osservazione}
	Quanto visto in questa sezione può essere generalizzato al caso $W^{k,p}(U)$. Per i nostri scopi è importante tenere presente che se $U$ limitato con $\partial U$ di classe $C^1$, allora l'immersione $H^2(U) \hookrightarrow H^1(U)$ è compatta.
\end{osservazione}

Per completezza specifichiamo che, in un caso particolare, è possibile dare un risultato con ipotesi di struttura su $U$ meno stringenti.
\begin{proposizione}\label{Rellichsottosteroidi}
	Siano $1 \le p < +\infty$ ed $U$ limitato. L'immersione $W^{1,p}_0(U) \hookrightarrow L^p(U)$ è compatta.
\end{proposizione}
\begin{proof}
Omettiamo la dimostrazione.
\end{proof}
\begin{proposizione}\label{Rellichsottosteroididuale}
	Siano $1 < q \le +\infty$ ed $U$ limitato. L'immersione $L^q(U) \hookrightarrow W^{-1,q}(U)$ è compatta.
\end{proposizione}
\begin{proof}
Omettiamo la dimostrazione.
\end{proof}

\section{Spazi di Sobolev che coinvolgono il tempo}
\begin{definizione}
	Consideriamo $(X,\norm{\;})$ uno spazio di Banach su $\R$ e $T>0$. Data $u \in L^1(0,T;X)$, diciamo che $u' \in L^1(0,T;X)$ è la \emph{derivata debole} di $u$ \index[analitico]{derivata!debole} se per ogni $\varphi \in C^\infty_c(0,T)$
	\[ \int_0^T u \varphi' d\mathcal{L}^1 = - \int_0^T u' \varphi d\mathcal{L}^1. \]
\end{definizione}

\begin{notazione}
	Consideriamo $(X,\norm{\;})$ uno spazio di Banach su $\R$, $p \in [1,\infty]$ e $T>0$. Indichiamo con  $W^{1,p}(0,T;X)$\index[simboli]{$W^{1,p}(0,T;X)$} l'insieme delle $u \in L^p(0,T;X)$ tali che $u' \in L^p(0,T;X)$. In particolare, poniamo $H^{1}(0,T;X) = W^{1,2}(0,T;X)$\index[simboli]{$H^{1}(0,T;X)$}.
\end{notazione}

\begin{definizione}
	Consideriamo $(X,\norm{\;})$ uno spazio di Banach su $\R$, $p \in [1,\infty]$ e $T>0$. Data $u \in W^{1,p}(0,T;X)$, poniamo 
	\[ \norm{u}_{W^{1,p}(0,T;X)} = \begin{cases}
		\left( \displaystyle \int_0^T \left( \norm{u(t)}^p +  \norm{u'(t)}^p\right)  d\mathcal{L}^1\right) ^{\frac{1}{p}} & 1 \le p < \infty, \\
		\underset{t \in [0,T]}{\textup{ess sup}}\;\left(  \norm{u(t)} + \norm{u'(t)}\right)  & p=\infty.
	\end{cases} \]
\end{definizione}

\begin{teorema}
	Consideriamo $(X,\norm{\;})$ uno spazio di Banach su $\R$, $p \in [1,\infty]$ e $T>0$. Data $u \in W^{1,p}(0,T;X)$, valgono i seguenti fatti:
	\begin{enumerate}[$(a)$]
		\item esiste un rappresentante $u^\ast$ di $u$ tale che $u^\ast \in C([0,T];X)$,
		\item per ogni $s\le t \in [0,T]$ risulta
		\[ u(t) = u(s) + \int_s^t u'd\mathcal{L}^1, \]
		\item esiste $C$, dipendente solo da $T$, tale che 
		\[ \underset{t \in [0,T]}{\max}\; \norm{u(t)} \le C \norm{u}_{W^{1,p}(0,T;X)}. \]
	\end{enumerate}
\end{teorema}
\begin{proof}
	Omettiamo la dimostrazione.
\end{proof}

\begin{teorema}
	Data $u \in L^2(0,T;H^1_0(U))$ con $u' \in L^2(0,T;H^{-1}(U))$, valgono i seguenti fatti:
	\begin{enumerate}[$(a)$]
		\item esiste un rappresentante $u^\ast$ di $u$ tale che $u^\ast \in C([0,T];L^2(U))$,
		\item l'applicazione $\left\lbrace t \mapsto \norm{u(t)}_{L^2(U)}^2\right\rbrace $ è assolutamente continua e per \emph{q.o.} $t \in [0,T]$
		\[ \frac{d}{dt}\norm{u(t)}_{L^2(U)}^2 = 2\braket{u'(t),u(t)}, \]
		\item esiste $C$, dipendente solo da $T$, tale che 
		\[ \underset{t \in [0,T]}{\max}\; \norm{u(t)}_{L^2(U)} \le C\left(  \norm{u}_{L^2(0,T;H^1_0(U))} + \norm{u'}_{L^2(0,T;H^{-1}(U))}\right) . \]
	\end{enumerate}
\end{teorema}
\begin{proof}
	Omettiamo la dimostrazione.
\end{proof}

\section{Operatore di Nemytskij}
\begin{lemma}\label{compattezzaelimitiuniformi}
	Consideriamo $X,Y$ due spazi di Banach su $\R$ ed $E \subseteq X$. Se $(f_h)$ in $C(E;Y)\cap \mathcal{K}(E;Y)$ e $f: E \to Y$ tale che per ogni $B \subseteq E$ limitato 
	\[ \lim\limits_{h} \norm{f_h-f}_{L^\infty(B;Y)} = 0, \]
	allora $f \in C(E;Y)\cap \mathcal{K}(E;Y)$.
\end{lemma}
\begin{proof}
	Innanzitutto, essendo limite uniforme di applicazioni continue, per ogni $B \subseteq E$ limitato $f$ è continua su $B$. In particolare, per ogni $x \in E$ possiamo affermare che $f$ è continua su $\B(x,1)$, quindi $f$ è continua in $x$. Da ciò si deduce $f \in C(E;Y)$. 
	
	Poiché $Y$ è completo, basta dimostrare che per ogni $B \subseteq E$ limitato si ha che $f(B)$ è totalmente limitato in $Y$. Fissiamo $B \subseteq E$ limitato. Dato $\varepsilon >0$, sia $h \in \N$ tale che 
	\[ \norm{f_h-f}_{L^\infty(B;Y)} < \frac{\varepsilon}{3}. \]
	Per compattezza di $f_h$, si ha che $f_h(B)$ è compatto, quindi può essere ricoperto da un numeri finito di palle di raggio minore di $\frac{\varepsilon}{6}$, ossia esistono $y_1,\dots,y_k$ tali che
	\[ f_h(B) \subseteq \bigcup_{j=1}^k \B(y_j,\frac{\varepsilon}{6}).\]
	Per ogni $j=1,\dots, k$ consideriamo 
	\[ C_j = \left\lbrace y \in Y: \textup{dist}(y, \B(y_j,\frac{\varepsilon}{6})) \le \frac{\varepsilon}{3} \right\rbrace. \]
	Applicando la disuguaglianza triangolare a $z_1,z_4 \in C_j$ e $z_2,z_3 \in \B(y_j,\frac{\varepsilon}{6})$, si ottiene $\textup{diam}(C_j) \le \varepsilon$.
	\begin{center}
		\begin{tikzpicture}[x=0.5pt,y=0.5pt,yscale=-1,xscale=1]
			%uncomment if require: \path (0,367); %set diagram left start at 0, and has height of 367
			
			%Shape: Circle [id:dp7038724886842289] 
			\draw   (270,171) .. controls (270,143.39) and (292.39,121) .. (320,121) .. controls (347.61,121) and (370,143.39) .. (370,171) .. controls (370,198.61) and (347.61,221) .. (320,221) .. controls (292.39,221) and (270,198.61) .. (270,171) -- cycle ;
			%Shape: Circle [id:dp20408854625505946] 
			\draw   (169.95,171) .. controls (169.95,88.13) and (237.13,20.95) .. (320,20.95) .. controls (402.87,20.95) and (470.05,88.13) .. (470.05,171) .. controls (470.05,253.87) and (402.87,321.05) .. (320,321.05) .. controls (237.13,321.05) and (169.95,253.87) .. (169.95,171) -- cycle ;
			%Straight Lines [id:da8755098905110549] 
			\draw    (169.95,171) -- (470.05,171) ;
			%Shape: Circle [id:dp34245043590360025] 
			\draw   (232,237.5) .. controls (232,237.22) and (232.22,237) .. (232.5,237) .. controls (232.78,237) and (233,237.22) .. (233,237.5) .. controls (233,237.78) and (232.78,238) .. (232.5,238) .. controls (232.22,238) and (232,237.78) .. (232,237.5) -- cycle ;
			%Shape: Circle [id:dp3226357363041006] 
			\draw   (299.5,200) .. controls (299.5,199.72) and (299.72,199.5) .. (300,199.5) .. controls (300.28,199.5) and (300.5,199.72) .. (300.5,200) .. controls (300.5,200.28) and (300.28,200.5) .. (300,200.5) .. controls (299.72,200.5) and (299.5,200.28) .. (299.5,200) -- cycle ;
			%Shape: Circle [id:dp8931755100052371] 
			\draw   (348.5,194) .. controls (348.5,193.72) and (348.72,193.5) .. (349,193.5) .. controls (349.28,193.5) and (349.5,193.72) .. (349.5,194) .. controls (349.5,194.28) and (349.28,194.5) .. (349,194.5) .. controls (348.72,194.5) and (348.5,194.28) .. (348.5,194) -- cycle ;
			%Shape: Circle [id:dp78063200787119] 
			\draw   (417,227.5) .. controls (417,227.22) and (417.22,227) .. (417.5,227) .. controls (417.78,227) and (418,227.22) .. (418,227.5) .. controls (418,227.78) and (417.78,228) .. (417.5,228) .. controls (417.22,228) and (417,227.78) .. (417,227.5) -- cycle ;
			%Straight Lines [id:da9144843916038057] 
			\draw  [dash pattern={on 0.84pt off 2.51pt}]  (299.5,200) -- (233,237.5) ;
			%Straight Lines [id:da2444506606827901] 
			\draw  [dash pattern={on 0.84pt off 2.51pt}]  (348.5,194) -- (300.5,200) ;
			%Straight Lines [id:da16195075964841843] 
			\draw  [dash pattern={on 0.84pt off 2.51pt}]  (417.5,228) -- (349,194.5) ;
			
			% Text Node
			\draw (312,140) node [anchor=north west][inner sep=0.75pt]    {$\frac{\varepsilon }{3}$};
			% Text Node
			\draw (406,140) node [anchor=north west][inner sep=0.75pt]    {$\frac{\varepsilon }{3}$};
			% Text Node
			\draw (216,140) node [anchor=north west][inner sep=0.75pt]    {$\frac{\varepsilon }{3}$};
			% Text Node
			\draw (234,234.9) node [anchor=north west][inner sep=0.75pt]    {$z_{1}$};
			% Text Node
			\draw (301.5,200) node [anchor=north west][inner sep=0.75pt]    {$z_{2}$};
			% Text Node
			\draw (345.5,175.4) node [anchor=north west][inner sep=0.75pt]    {$z_{3}$};
			% Text Node
			\draw (419.5,230.4) node [anchor=north west][inner sep=0.75pt]    {$z_{4}$};
		\end{tikzpicture}
	\end{center}
	Se $y \in f(B)$, allora esiste $x \in B$ tale che $y = f(x)$ e 
	\[ \norm{y- f_h(x)}_Y = \norm{f(x)-f_h(x)}_Y <\frac{\varepsilon}{3}. \]
	Inoltre, esiste $j =1,\dots,k$ tale che $f_h(x) \in \B(y_j,\frac{\varepsilon}{6})$, quindi $y \in C_j$, da cui 
	\[ f(B) \subseteq \bigcup_{j=1}^k C_j. \qedhere \]
\end{proof}
Nel corso di questa sezione, $E$ denoterà un sottoinsieme misurabile di $\R^n$.

\begin{definizione}
	Diciamo che un'applicazione $g : E \times \R^k \to \R$ è \emph{di Carathéodory} \index[analitico]{applicazione!di Carathéodory}, se valgono i seguenti fatti:
	\begin{enumerate}[$(a)$]
		\item per ogni $s \in \R^k$ l'applicazione $\left\lbrace x \mapsto g(x,s) \right\rbrace$ è misurabile su $E$,
		\item per \emph{q.o.} $x \in E$ l'applicazione $\left\lbrace s \mapsto g(x,s) \right\rbrace$ è continua su $\R^k$.
	\end{enumerate}
\end{definizione}
\begin{notazione}
	Data un'applicazione $u: E \to \R^k$, denotiamo con $g(x, u): E \to \R$ l'applicazione $\left\lbrace x \longmapsto g(x,u(x))\right\rbrace$.
\end{notazione}
\begin{teorema}\label{tecnicoperNemytskij}
	Se $g : E \times \R^k \to \R$ un'applicazione di Carathéodory, allora per ogni applicazione misurabile $u : E \to \R^k$, l'applicazione $g(x,u): E \to \R$ è
	misurabile.
	
	In particolare, se $u,v: E \to \R^k$ sono misurabili ed uguali \emph{q.o.} in $E$, anche $g(x, u)$ e $g(x, v)$ sono uguali \emph{q.o.} in $E$.
\end{teorema}
\begin{proof}
	Omettiamo la dimostrazione.
\end{proof}
\begin{teorema}\label{fondamentalepercontinuitàNemytskij}
	Sia $g: E \times \R^k \to \R$ un'applicazione di Carathéodory e siano $p_1,\dots, p_k,q \in \left[ 1,+\infty\right[$. Supponiamo che esistano $a \in L^q(E)$ e $b \in \R$ tali che
	\[ \abs{g(x,s_1,\dots,s_k)} \le a(x) + b\abs{s_1}^{\frac{p_1}{q}} + \dots + b\abs{s_k}^{\frac{p_k}{q}} \]
	per \emph{q.o.} $x \in E$ e per ogni $s \in \R^k$.
	Allora per ogni $u_1 \in L^{p_1}(E), \dots, u_k \in L^{p_k}(E)$ si ha $g(x,u_1,\dots , u_k) \in L^q(E)$ e l’applicazione $\mathcal{G}: L^{p_1}(E) \times \dots L^{p_k}(E) \to L^q(E)$ tale che
	\[ \mathcal{G}(u_1,\dots,u_k) = g(x,u_1,\dots,u_k) \]
	è continua.
\end{teorema}
\begin{proof}
	Siano $u_1 \in L^{p_1}(E),\dots ,u_k \in L^{p_k}(E)$. Per il Teorema \eqref{tecnicoperNemytskij}, $g(x, u_1, \dots, u_k)$ è ben definita, come classe di equivalenza, ed è misurabile.	
	
	Dal momento che
	\begin{align*}
		\abs{g(x,u_1,\dots,u_k)}^q &\le \left( a(x) + b\abs{u_1}^{\frac{p_1}{q}}+ \dots + b\abs{u_k}^{\frac{p_k}{q}}\right) \le \\
		&\le (k+1)^{q-1}\left( a(x)^q + b^q\abs{u_1}^{p_1}+ \dots + b^q\abs{u_k}^{p_k}\right), 
	\end{align*}
	risulta
	\[ \int_E \abs{g(x,u_1,\dots,u_k)}^q d \mathcal{L}^n \le (k+1)^{q-1} \left( \norm{a}^{q}_{L^{q}(E)} + b^q\norm{u_1}^{p_1}_{L^{p_1}(E)} +\dots + b^q\norm{u_k}^{p_k}_{L^{p_k}(E)}\right) .\]
	Ne segue $g(x,u_1,\dots,u_k) \in L^q(E)$.
	
	Sia ora, per $j = 1, . . . , k$, $(u_j^{(h)})$ una successione convergente ad $u_j$ in $L^{p_j}(E)$. Per ogni sottosuccessione $(u_j^{(h_l)})$, esistono una sottosuccessione $(u_j^{(h_{l_r})})$ e $w_j \in L^{p_j}(E)$ tali che
	\begin{align*}
		\lim\limits_{h} u_j^{(h_{l_r})}(x) &= u_j(x) \textup{ per q.o. } x \in E,\\
		\abs{u_j^{(h_{l_r})}} &\le w_j  \textup{ q.o. in }E.
	\end{align*}
	Per continuità di $g$, a meno di ulteriori sottosuccessioni,
	\begin{align*}
		\lim\limits_{h} g(x,u_1^{(h_{l_r})}(x),\dots,u_k^{(h_{l_r})}(x)) &= g(x,u_1(x),\dots,u(x)) \textup{ per q.o. } x \in E,\\
		\abs{g(x,u_1^{(h_{l_r})},\dots,u_k^{(h_{l_r})})}^q &\le (k+1)^{q-1} \left( a^q + b^qw_1^{p_1} + \dots + b^qw_k^{p_k} \right)  \textup{ q.o. in }E.
	\end{align*}
	e, per il Teorema della convergenza dominata,
	\[ \lim\limits_{m} \int_E \abs{g(x,u_1^{(h_{l_r})},\dots,u_k^{(h_{l_r})})- g(x,u_1,\dots,u_k)}^q d \mathcal{L}^n =0, \]
	ossia $g(x,u_1^{(h_{l_r})},\dots,u_k^{(h_{l_r})}) \to g(x,u_1,\dots,u_k)$ in $L^q(E)$. Allora $g(x,u_1^{(h)},\dots,u_k^{(h)}) \to g(x,u_1,\dots,u_k)$ in $L^q(E)$, ossia $\mathcal{G}$ è continua.
\end{proof}

\begin{definizione}
	Sia $g : E \times \R^k \to \R$ un'applicazione di Carathéodory. L'applicazione $\mathcal{G}: L^{p_1}(E) \times \dots L^{p_k}(E) \to L^q(E)$ tale che
	\[ \mathcal{G}(u_1,\dots,u_k) = g(x,u_1,\dots,u_k) \]
	si chiama \emph{operatore di Nemytskij}\index[analitico]{operatore!di Nemytskij}\index[simboli]{$\mathcal{G}$} o \emph{operatore di superposizione} associato a $g$.
\end{definizione}

Iniziamo dallo studio della continuità.
\begin{corollario}
	Siano $g : U \times (\R \times \R^n) \to \R$ un'applicazione di Carathéodory, $p \in \left[ 1,n\right[$ e $q \in \left[ 1,+\infty\right[ $. Supponiamo che esistano $a \in L^q(U)$ e $b \in \R$ tali
	che
	\[ \abs{g(x,s,\xi)} \le a(x) + b\abs{s}^{\frac{np}{(n-p)q}} + b\abs{\xi}^{\frac{p}{q}} \]
	per \emph{q.o.} $x \in U$ e per ogni $(s,\xi) \in \R \times \R^n$. Allora per ogni $u \in W^{1,p}_0(U)$ si ha $g(x, u, Du) \in L^q(U)$ e l’applicazione $\mathcal{T}: W^{1,p}_0(U) \to L^q(U)$ tale che
	\[ \mathcal{T}(u) = g(x,u,Du) \]
	è continua.
\end{corollario}
\begin{proof}
	Per il Teorema \eqref{fondamentalepercontinuitàNemytskij}, l'applicazione $\Psi_1: L^{\frac{np}{n-p}}(U) \times \left( L^p(U)\right) ^n \to L^q(U)$ tale che
	\[ \Psi_1(u,v_1,\dots,v_n) = g(x,u,v_1,\dots,v_n) \]
	è ben definita e continua. Inoltre, per ogni $j =1,\dots,n$ l'applicazione lineare $\Phi_j : W^{1,p}_0(U) \to L^p(U)$ tale che
	\[ \Phi_ju = D_ju \]
	è continua, infatti 
	\[ \norm{D_ju}_{L^p(U)} \le \norm{Du}_{L^p(U)} = \norm{u}_{W^{1,p}_0(U)}. \]
	Dal momento che, per la disuguaglianza di Poincaré, l'immersione $W^{1,p}_0(U) \hookrightarrow L^{\frac{np}{n-p}}(U)$ è continua, l'applicazione lineare $\Psi_2: W^{1,p}_0(U) \to L^{\frac{np}{n-p}}(U) \times \left( L^p(U)\right) ^n$ tale che
	\[ \Psi_2u = (u,D_1u,\dots,D_nu) \]
	è continua, infatti
	\[ \norm{\Psi_2u}_{L^{\frac{np}{n-p}}(U) \times \left( L^p(U)\right) ^n}^2  = \norm{u}_{L^{\frac{np}{n-p}}(U)}^2 + \norm{D_1u}^2_{L^p(U)} + \dots + \norm{D_nu}^2_{L^p(U)} \le C \norm{u}_{W^{1,p}_0(U)}^2. \]
	La tesi segue dal fatto che $\mathcal{T} = \Psi_1 \circ \Psi_2$.
\end{proof}

\begin{corollario}\label{continuitàNemytskij}
	Siano $g : U \times (\R \times \R^n) \to \R$ un'applicazione di Carathéodory, $p \in \left[ 1,n\right[$ e $q \in \left] \frac{n}{n-1},+\infty\right] $. Supponiamo che esistano $a \in L^{\frac{nq}{n+q}}(U)$ e $b \in \R$ tali
	che
	\[ \abs{g(x,s,\xi)} \le a(x) + b\abs{s}^{\frac{p(n+q)}{(n-p)q}} + b\abs{\xi}^{\frac{p(n+q)}{nq}} \]
	per \emph{q.o.} $x \in U$ e per ogni $(s,\xi) \in \R \times \R^n$. Allora per ogni $u \in W^{1,p}_0(U)$ si ha $g(x, u, Du) \in W^{-1,q}(U)$ e l’applicazione $\mathcal{T}: W^{1,p}_0(U) \to W^{-1,q}(U)$ tale che
	\[ \mathcal{T}(u) = g(x,u,Du) \]
	è continua.
\end{corollario}
\begin{proof}
	Per il Corollario precedente, l'applicazione $\Psi: W^{1,p}_0(U) \to L^{\frac{nq}{n+q}}(U)$ tale che
	\[ \Psi u = g(x,u,Du) \]
	è ben definita e continua. La tesi segue dunque dal fatto che, per la disuguaglianza di Poincaré duale, l'immersione $ L^{\frac{nq}{n+q}}(U) \hookrightarrow W^{-1,q}(U)$ è continua.
\end{proof}
\begin{osservazione}
	Il Corollario \eqref{continuitàNemytskij} appena dimostrato vale, in particolare, per $q = p'$.
\end{osservazione}

Affrontiamo ora lo studio della compattezza.

\begin{teorema}\label{compattezzaNemytskij}
Siano $g : U \times (\R \times \R^n) \to \R$ un'applicazione di Carathéodory, $p \in \left[ 1,n\right[$ e $q \in \left] \frac{n}{n-1},+\infty\right] $. Supponiamo che per ogni $\varepsilon>0$ esista $a_\varepsilon \in L^{\frac{nq}{n+q}}(U)$ tale che 
	\[ \abs{g(x,s,\xi)} \le a_\varepsilon(x) + \varepsilon\abs{s}^{\frac{p(n+q)}{(n-p)q}} + \varepsilon \abs{\xi}^{\frac{p(n+q)}{nq}} \]
per \emph{q.o.} $x \in U$ e per ogni $(s,\xi) \in \R \times \R^n$. Allora l’applicazione $\mathcal{T}: W^{1,p}_0(U) \to W^{-1,q}(U)$ tale che
\[ \mathcal{T}(u) = g(x,u,Du) \]
è continua e compatta.
\end{teorema}
\begin{proof}
La continuità segue dal Corollario \eqref{continuitàNemytskij}.

Vediamo, innanzitutto, il caso in cui esistono $R>0$ ed $a \in L^\infty(U)$ tali che 
\[ \abs{g(x,s,\xi)} \le a(x) \textup{ per q.o. } x \in U \textup{ ed ogni } (s,\xi) \in \R \times \R^n, \]
\[ a(x) = 0 \textup{ per q.o. } x \in U \setminus \B(0,R). \]
Consideriamo $(u_h)$ in $W^{1,p}_0(U)$ limitata. Siccome
\[ \norm{g(x,u_h,Du_h)}^q_{L^q(U)} \le \norm{a}_{L^\infty(U)}^q\mathcal{L}^n(\B(0,R)) <+\infty, \]
allora $(g(x,u_h,Du_h))$ è limitata in $L^q(U)$ e $g(x,u_h,Du_h) = 0$ q.o. in $U\setminus \B(0,R)$. Sia $\vartheta \in C^\infty_c(\B(0,R+1))$ tale che $0\le \vartheta \le 1$ e $\vartheta = 1$ in $\overline{\B(0,R)}$.

Per la Proposizione \eqref{Rellichsottosteroididuale}, a meno di sottosuccessioni, $(g(x,u_h,Du_h))$ converge in $W^{-1,q}(U \cap \B(0,R+1))$, quindi è di Cauchy in $W^{-1,q}(U \cap \B(0,R+1))$. Per ogni $v \in W^{1,q'}_0(U)$ si ha che $\vartheta v \in W^{1,q'}_0(U\cap B(0,R+1))$, quindi, per la disuguaglianza di Poincaré,
\begin{align*}
&\abs*{\int_U (g(x,u_h,Du_h)- g(x,u_k,Du_k))vd\mathcal{L}^n} = \abs*{\int_{U \cap \B(0,R)}\hspace{-25pt} (g(x,u_h,Du_h)- g(x,u_k,Du_k))vd\mathcal{L}^n} = \\
&= \abs*{\int_{U \cap \B(0,R)}\hspace{-25pt} (g(x,u_h,Du_h)- g(x,u_k,Du_k))\vartheta vd\mathcal{L}^n} = \\
&= \abs*{\int_{U \cap \B(0,R+1)}\hspace{-35pt} (g(x,u_h,Du_h)- g(x,u_k,Du_k))\vartheta vd\mathcal{L}^n} \le \\
&\le \norm{g(x,u_h,Du_h)- g(x,u_k,Du_k)}_{W^{-1,q}(U\cap \B(0,R+1))} \norm{\vartheta v}_{W^{1,q'}_0(U\cap \B(0,R+1))} \le \\
&\le C\norm{g(x,u_h,Du_h)- g(x,u_k,Du_k)}_{W^{-1,q}(U\cap \B(0,R+1))}\norm{v}_{W^{1,q'}_0(U)},
\end{align*}
da cui
\begin{align*}
&\norm{g(x,u_h,Du_h)- g(x,u_k,Du_k)}_{W^{-1,q}(U)} \le  \\
&\le C\norm{g(x,u_h,Du_h)- g(x,u_k,Du_k)}_{W^{-1,q}(U\cap \B(0,R+1))}
\end{align*}
per cui $(g(x,u_h,Du_h))$ è di Cauchy in $W^{-1,q}(U)$, quindi convergente.

Vediamo ora il caso in cui esista $a \in L^{\frac{nq}{n+q}}(U)$ tale che 
\[ \abs{g(x,s,\xi)} \le a(x) \textup{ per q.o. } x \in U \textup{ ed ogni } (s,\xi) \in \R\times \R^n. \]
Consideriamo 
\[ g_h(x,s,\xi) = \begin{cases}
	g(x,s,\xi) & \textup{se } a(x) \le h \textup{ e } \abs{x}<h,\\
	0 & \textup{ altrimenti},
\end{cases} \]
e 
\[ a_h(x) = \begin{cases}
	a(x) & \textup{se } a(x) \le h \textup{ e } \abs{x}<h,\\
	0 & \textup{ altrimenti},
\end{cases} \]
Evidentemente $g_h$ è di Carathéodory e 
\[ \abs{g_h(x,s,\xi)} \le a_h(x) \textup{ per q.o. } x \in U \textup{ ed ogni } (s,\xi) \in \R\times \R^n. \]
Allora, per il passo precedente, l'applicazione $\mathcal{T}_h: W^{1,p}_0(U) \to W^{-1,q}(U)$ tale che
\[ \mathcal{T}_h(u) = g_h(x,u,Du) \]
è continua e compatta. Siccome per ogni $u \in W^{1,p}_0(U)$ si ha
\[ \int_U \abs{g(x,u,Du) - g_h(x,u,Du)}^{\frac{nq}{n+q}}d\mathcal{L}^n \le \int_U a^{\frac{nq}{n+q}}(1-\chi_{\left\lbrace x \in \R^n: a(x)\le h\right\rbrace}\chi_{\B(0,h)})d\mathcal{L}^n. \]
Dal Teorema della convergenza dominata, si deduce che per ogni $B \subseteq W^{1,p}_0(U)$ limitato,
\[ \lim\limits_{h} \underset{u \in B}{\sup}\; \norm{g_h(x,u,Du) -g(x,u,Du)}_{L^{\frac{nq}{n+q}}(U)} = 0.  \]
Per la Disuguaglianza di Poincaré duale, 
\[ \lim\limits_{h} \underset{u \in B}{\sup}\; \norm{g_h(x,u,Du) -g(x,u,Du)}_{W^{-1,q}(U)} = 0,  \]
che, combinata con il Lemma \eqref{compattezzaelimitiuniformi}, da la compattezza desiderata.

Nel caso generale, per $\varepsilon = \frac{1}{h}$, con $h \in N\setminus\left\lbrace 0\right\rbrace$, esiste $a_h \in L^{\frac{nq}{n+q}}(U)$ tale che 
	\[ \abs{g(x,s,\xi)} \le a_h(x) + \frac{1}{h}\abs{s}^{\frac{p(n+q)}{(n-p)q}} + \frac{1}{h} \abs{\xi}^{\frac{p(n+q)}{nq}} \]
per q.o. $x \in U$ e per ogni $(s,\xi) \in \R \times \R^n$. Poniamo 
\[ g_h(x,s,\xi) = \min \left\lbrace \max\left\lbrace g(x,s,\xi),-a_h(x)\right\rbrace a_h(x)\right\rbrace.  \]
Evidentemente $ g_h$ è di Carathéodory e 
\[ \abs{g_h(x,s,\xi)} \le a_h(x) \textup{ per q.o. } x \in U \textup{ ed ogni } (s,\xi) \in \R\times \R^n. \]
Allora, per il passo precedente, l'applicazione $\mathcal{T}_h: W^{1,p}_0(U) \to W^{-1,q}(U)$ tale che
\[ \mathcal{T}_h(u) = g_h(x,u,Du) \]
è continua e compatta. Osservando che 
\[ \abs{g(x,s,\xi) - g_h(x,s,\xi)} \le \frac{1}{h}\abs{s}^{\frac{p(n+q)}{(n-p)q}} + \frac{1}{h} \abs{\xi}^{\frac{p(n+q)}{nq}} \]
per q.o. $x \in U$ e per ogni $(s,\xi) \in \R \times \R^n$, per ogni $u \in W^{1,p}_0(U)$ si ha
\[ \int_U \abs{g(x,s,\xi) - g_h(x,s\xi)}^{\frac{nq}{n+q}} d\mathcal{L}^n \le 2^{\frac{nq-n-q}{n+q}}\left(\frac{1}{h} \right)^{\frac{nq}{n+q}} \left(\int_U \abs{u}^\frac{np}{n-p} d\mathcal{L}^n + \int_U \abs{Du}^p d\mathcal{L}^n \right)  \]
e, per la disuguaglianza di Poincaré,
\[ \int_U \abs{g(x,s,\xi) - g_h(x,s\xi)}^{\frac{nq}{n+q}} d\mathcal{L}^n \le C \left(\frac{1}{h} \right)^{\frac{nq}{n+q}}\left( \norm{u}_{W^{1,p}_0(U)}^{\frac{np}{n-p}} + \norm{u}_{W^{1,p}_0(U)}^p\right).   \]
Pertanto per ogni $B \subseteq W^{1,p}_0(U)$ limitato si ha
\[ \lim\limits_{h} \underset{u \in B}{\sup}\; \norm{g_h(x,u,Du) -g(x,u,Du)}_{L^{\frac{nq}{n+q}}(U)} = 0.  \]
Per la Disuguaglianza di Poincaré duale, 
\[ \lim\limits_{h} \underset{u \in B}{\sup}\; \norm{g_h(x,u,Du) -g(x,u,Du)}_{W^{-1,q}(U)} = 0,  \]
che, in combinazione con il Lemma \eqref{compattezzaelimitiuniformi}, da la compattezza desiderata.
\end{proof}

\begin{osservazione}
	Il Teorema \eqref{compattezzaNemytskij} appena dimostrato vale, in particolare, per $q = p'$. In questo caso la stima di crescita su $g$ assume la forma
		\[ \abs{g(x,s,\xi)} \le a_\varepsilon(x) + \varepsilon\abs{s}^{\frac{n(p-1)+p}{n-p}} + \varepsilon \abs{\xi}^{\frac{n(p-1)+p}{n}}. \]
\end{osservazione}

\begin{osservazione}
E opportuno confrontare la condizione di crescita del Teorema \eqref{compattezzaNemytskij} con quella del Corollario \eqref{continuitàNemytskij}. Gli esponenti $\frac{p(n+q)}{(n-p)q}$ e $\frac{p(n+q)}{nq}$ rappresentano gli andamenti critici con cui si può garantire la continuità, ma non la completa continuità dell’operatore. Non appena si rimane al di sotto di tali crescite, come nel Teorema \eqref{compattezzaNemytskij}, si ottiene la completa continuità dell’operatore. Non occorre invece imporre una diversa
sommabilità sul termine $a_\varepsilon$. Si osservi anche che nel Teorema \eqref{compattezzaNemytskij} non viene fatta alcuna ipotesi sull’aperto $U$.
In particolare, non si richiede che $U$ sia limitato.
\end{osservazione}

\begin{corollario}
Siano $g : U \times (\R \times \R^n) \to \R$ un'applicazione di Carathéodory, $p \in \left[ 1,n\right[$ e $q \in \left] \frac{n}{n-1},+\infty\right] $. Supponiamo $\mathcal{L}^n(U)<+\infty$ e che esistano $a \in L^{\frac{nq}{n+q}}(U)$, $b \in \R$, $r_0 \in \left[ 1,\frac{p(n+q)}{(n-p)q}\right[$ e $r_1 \in \left[ 1,\frac{p(n+q)}{nq}\right[$  tali che 
\[ \abs{g(x,s,\xi)} \le a(x) + b\abs{s}^{r_0} + b \abs{\xi}^{r_1} \]
per \emph{q.o.} $x \in U$ e per ogni $(s,\xi) \in \R \times \R^n$. Allora l’applicazione $\mathcal{T}: W^{1,p}_0(U) \to W^{-1,q}(U)$ tale che
\[ \mathcal{T}(u) = g(x,u,Du) \]
è continua e compatta.
\end{corollario}
\begin{proof}
Ricordiamo che, per la Disuguaglianza di Young pesata, per ogni $\alpha \in \left] 1,+\infty\right[, s,t \in \R$ ed $\varepsilon >0$ si ha
\[ \abs{st} \le \varepsilon \abs{s}^{\alpha} + \frac{\alpha-1}{\alpha^{\frac{\alpha}{\alpha-1}}\varepsilon^{\frac{1}{\alpha -1}}}\abs{t}^{\alpha'}. \]
Scelti 
\begin{align*}
	\alpha &= \frac{p(n+q)}{(n-p)qr_0}, & \beta &=\frac{p(n+q)}{(nqr_1},
\end{align*}
per ogni $\varepsilon > 0$ risulta
\[ \abs{g(x,s,\xi)} \le a(x) + \frac{\alpha -1}{\alpha^{\frac{\alpha}{\alpha-1}}\varepsilon^{\frac{1}{\alpha -1}}}b^{\alpha'} + \frac{\beta -1}{\beta^{\frac{\beta}{\beta-1}}\varepsilon^{\frac{1}{\beta -1}}}b^{\beta'} + \varepsilon\abs{s}^{\frac{p(n+q)}{(n-p)q}} + \varepsilon \abs{\xi}^{\frac{p(n+q)}{nq}}. \]
D’altronde, avendo $U$ misura finita, ogni funzione costante appartiene a $L^{\frac{nq}{n+q}}(U)$. La tesi discende allora dal Teorema \eqref{compattezzaNemytskij}.
\end{proof}

\begin{osservazione}
	Il Corollario appena dimostrato vale, in particolare, per $q = p'$.
\end{osservazione}    
		
		% Commenti speciali

% !TEX encoding = UTF-8 Unicode
% !TEX TS-program = pdflatex
% !TEX spellcheck = it-IT
% !TEX root = Dispensa/af.tex

% Capitolo 2
\chapter{Equazioni ellittiche lineari del secondo ordine}\label{Equazioni ellittiche lineari del secondo ordine}
\section{Introduzione}
Nel seguito indicheremo con $\Delta = \sum_{j=1}^{n}\frac{\partial^2}{\partial x_j^2}$ l'operatore di Laplace, o Laplaciano. Altro non è che la traccia della matrice hessiana.

Iniziamo presentando l'esempio capostipite di tutta la teoria che intendiamo costruire.
\begin{esempio}\textbf{\emph{(problema di Poisson)}}\label{esempiopoisson}
	Siano $U$ limitato con $\partial U$ di classe $C^1$ ed $f \in L^2(U)$. Siamo interessati alle applicazioni $u$ che soddisfano 
	\begin{equation}\label{Poisson}
		\begin{cases}
			-\Delta u = f & \textup{in } U,\\
			u= 0 & \textup{su } \partial U.
		\end{cases}
	\end{equation}
	Se $f = 0$, chiamiamo tali applicazioni \emph{armoniche}.
	
	Data $f \in C(U)$, diciamo che $u$ è una \emph{soluzione classica del problema di Poisson} \eqref{Poisson}, se $u \in C^2(U)\cap C(\overline{U})$, $-\Delta u(x) = f(x)$ per ogni $x \in U$ e $u(x) = 0$ per ogni $x \in \partial U$. 
	
	Supponiamo dunque di avere una soluzione classica $u$. Allora per ogni $x \in U$,
	\[ -\Delta u(x) = f(x). \]
	Moltiplicando ambo i membri per una certa $\varphi \in C^\infty_c(U)$ ed integrando otteniamo 
	\[ -\int_U\Delta u \varphi d\mathcal{L}^n = \int_U f \varphi d\mathcal{L}^n \] 
	e per le formule di Gauss--Green abbiamo quindi, per ogni $\varphi \in C^\infty_c(U)$
	\begin{equation}\label{poissondebole}
		\int_U Du \cdot D\varphi d\mathcal{L}^n = \int_U f \varphi d\mathcal{L}^n
	\end{equation}
	L'equazione \eqref{poissondebole} ha significato in condizioni molto più deboli rispetto a quelle classiche. In particolare ha senso per $u \in H^1(U)$. Data $f \in L^2(U)$, diciamo quindi che $u$ è una \emph{soluzione debole del problema di Poisson} \eqref{Poisson}, se $u \in H^{1}_0(U)$, in modo che sia $u=0$ al bordo nel senso delle tracce, e per ogni $\varphi \in H^{1}_0(U)$ valga \eqref{poissondebole}. 
	
	Osserviamo subito che, mentre una soluzione classica è anche una soluzione debole, il viceversa non è scontato.
	
	Infine, diciamo che $u$ è una \emph{soluzione forte del problema di Poisson} \eqref{Poisson}, se $u \in H^{1}_0(U)$ e $-\Delta u(x) = f(x)$ per \emph{q.o.} $x \in U$.
	
	Osserviamo che, se $u$ è una soluzione debole e $u \in H^2_{loc}(U)$, allora $u$ è una soluzione forte. Infatti dal fatto che per ogni $\varphi \in C^\infty_c(U)$
	\[ \int_U Du \cdot D\varphi d\mathcal{L}^n = \int_U f \varphi d\mathcal{L}^n, \]
	per definizione di derivata debole abbiamo subito, essendo di fatto il dominio di integrazione il supporto di $\varphi$, che è compatto,
	\[ -\int_U \Delta u \varphi d\mathcal{L}^n = \int_U f \varphi d\mathcal{L}^n\]
	ossia 
	\[ \int_U (-\Delta u - f)\varphi d\mathcal{L}^n = 0\]
	e per il lemma di Du Bois--Reymond $-\Delta u - f= 0$ \emph{q.o.} in $U$.\footnote{Qui il laplaciano è costituito dalle derivate deboli, quindi di fatto è un laplaciano debole.} 
\end{esempio}

Nel corso di questo capitolo studieremo principalmente il problema differenziale con condizione di Dirichlet
\begin{equation}
	\begin{cases}
		Lu = f & \textup{in } U,\\
		u = 0 & \textup{su } \partial U,
	\end{cases}
\end{equation}
dove $U$ è un aperto limitato di $\R^n$ con $\partial U$ di classe $C^1$, $f: U \to \R$ è un dato, $u: \overline{U} \to \R$ è l'incognita e $L$ denota un operatore differenziale del secondo ordine a coefficienti funzioni che può essere \emph{non in forma di divergenza}
\[ Lu = -\sum_{i,j=1}^{n}a_{ij}D^2_{i,j}u + \sum_{i=1}^{n}b_i D_iu + cu \]
oppure \emph{in forma di divergenza}
\[ Lu = -\sum_{i,j=1}^{n}D_{j}(a_{ij}D_iu) + \sum_{i=1}^{n}b_i D_iu + cu. \]

In particolare, $L$ è lineare.

\begin{osservazione}
	Se $a_{ij} \in C^1(U)$, le due forme sono equivalenti. Infatti, l'operatore può essere scritto come
	\[ Lu = -\sum_{i,j=1}^{n}a_{ij}D^2_{i,j}u + \sum_{i=1}^{n}\tilde{b}_i D_iu + cu \]
	con $\tilde{b}_i = b_i - \sum_{j=1}^{n}D_j a_{ij}$, che è chiaramente non in forma di divergenza.
\end{osservazione}

\begin{definizione}
	Diciamo che l'operatore differenziale $L$ è \emph{uniformemente ellittico}\index[analitico]{uniforme ellitticità}, se esiste $\nu > 0$ tale che per \emph{q.o.} $ x \in U$ e per ogni $\xi \in \R^n$ si ha
	\[ \sum_{i,j=1}^n a_{ij}(x)\xi_i\xi_j \ge \nu \abs{\xi}^2. \]
	In altre parole, per \emph{q.o.} $x \in U$, la matrice $A = (a_{i,j}(x))$ è definita positiva ed il più piccolo autovalore è maggiore uguale a $\nu$.
\end{definizione}

Nel caso $a_{ij} = \delta_{ij}, b_i = 0, c=0$, abbiamo $L = - \Delta$. Quindi $L$ può essere visto come una generalizzazione dell'operatore di Laplace.

\begin{osservazione}
	Possiamo dare un'interpretazione fisica dell'operatore $L$:
	\begin{itemize}
		\item il termine di secondo ordine $A : D^2 u = \sum_{i,j=1}^na_{ij}D^2_{i,j}u$, rappresenta la \emph{diffusione} di $u$ in $U$. I coefficienti $a_{ij}$ descrivono l'anisotropa ed eterogenea natura del mezzo;
		\item il vettore $\mathbf{f} = -ADu$ rappresenta la \emph{densità di flusso diffusivo}. La condizione di ellitticità, $\mathbf{f} \cdot Du \le 0$, può essere interpretata come la tendenza del flusso a passare dalle zone a concentrazione maggiore verso quelle a concentrazione inferiore;
		\item il termine di primo ordine $\mathbf{b} \cdot Du = \sum_{i=1}^n b_iD_iu$ rappresenta il trasporto attraverso $U$;
		\item il termine di ordine zero $cu$ descrive le \emph{sorgenti} (risp. i \emph{pozzi}) \emph{locali}.
	\end{itemize}
\end{osservazione}

\section{Formulazione debole}
\begin{definizione}
	Siano $U$ un aperto limitato di $\R^n$ con $\partial U$ di classe $C^1$, $f: U \to \R$ un'applicazione continua ed $L$ un operatore differenziale in forma di divergenza con $a_{ij} \in C^1(U), b_i,c \in C(U)$ per ogni $i,j=1,\dots,n$. Diciamo che $u$ è una \emph{soluzione classica}\index[analitico]{soluzione!classica} del problema di Dirichlet
	\begin{equation}
		\begin{cases}
			Lu = f & \textup{in } U,\\
			u = 0 & \textup{su } \partial U,
		\end{cases}
	\end{equation}
	se $u \in C^2(U)\cap C(\overline{U})$, $Lu(x) = f(x)$ per ogni $x \in U$ e $u(x) = 0$ per ogni $x \in \partial U$.
\end{definizione}

Se $u$ è una soluzione classica di
\[ \begin{cases}
	Lu = f & \textup{in } U,\\
	u = 0 & \textup{su } \partial U,
\end{cases} \]
allora per ogni $x \in U$,
\[ Lu(x) = f(x). \]
Moltiplicando ambo i membri per una certa $\varphi \in C^\infty_c(U)$ ed integrando otteniamo 
\[ \int_U L u \varphi d\mathcal{L}^n = \int_U f \varphi d\mathcal{L}^n, \]
ossia
\[ \int_U \left( -\sum_{i,j=1}^{n}D_{j}(a_{ij}D_iu) + \sum_{i=1}^{n}b_i D_iu + cu\right)  \varphi d\mathcal{L}^n = \int_U f \varphi d\mathcal{L}^n, \]
che corrisponde a 
\[ \int_U \left( -\sum_{i,j=1}^{n}D_{j}(a_{ij}D_iu)\varphi + \sum_{i=1}^{n}b_i D_iu\varphi + cu\varphi\right) d\mathcal{L}^n = \int_U f \varphi d\mathcal{L}^n. \]
Per le formule di Gauss--Green abbiamo quindi, per ogni $\varphi \in C^\infty_c(U)$
\[ \int_U \left( \sum_{i,j=1}^{n}a_{ij}D_iuD_{j}\varphi + \sum_{i=1}^{n}b_i D_iu\varphi + cu\varphi\right) d\mathcal{L}^n = \int_U f \varphi d\mathcal{L}^n. \]
Osserviamo che l'equazione precedente ha significato con $f \in L^2(U)$ e $u,\varphi \in H^1_0(U)$.

Sulla base della precedente osservazione, possiamo quindi dare un nuovo significato all'espressione "risolvere un problema di Dirichlet".

Nel seguito, salvo diversa specificazione, supporremo sempre che $U$ sia un aperto limitato di $\R^n$ con $\partial U$ di classe $C^1$ e che siano assegnati $f \in L^2(U)$ ed un operatore differenziale $L$ uniformemente ellittico in forma di divergenza tale che
\begin{itemize}
	\item per ogni $i,j = 1,\dots, n$, 
	\[ a_{ij}=a_{ji}. \]
	In altre parole, $A$ è simmetrica,
	\item per ogni $i,j = 1,\dots, n$, 
	\[ a_{ij},b_i,c \in L^\infty(U). \]
\end{itemize}

\begin{definizione}
	Sia $B: H^1_0(U)\times H^1_0(U)\to \R$ l'applicazione bilineare tale che
	\[ B(u,v) = \int_U \left( \sum_{i,j=1}^{n}a_{ij}D_iuD_{j}v + \sum_{i=1}^{n}b_i D_iuv + cuv\right) d\mathcal{L}^n. \]
	Chiamiamo $B$ la \emph{forma bilineare associata ad $L$}\index[analitico]{forma bilineare!associata ad un operatore in forma di divergenza}.
\end{definizione}
La definizione precedente è ben posta.
\begin{definizione}
	Diciamo che $u \in H^1_0(U)$ è una \emph{soluzione debole}\index[analitico]{soluzione!debole} del problema di Dirichlet
	\begin{equation}
		\begin{cases}
			Lu = f & \textup{in } U,\\
			u = 0 & \textup{su } \partial U,
		\end{cases}
	\end{equation}
	se per ogni $v \in H^1_0(U)$, 
	\[ 	 B(u,v) = (f|v)_{2} \qquad \textup{(Formulazione variazionale)} \]
	dove $(\;|\;)_2$ denota il prodotto scalare in $L^2(U)$.
\end{definizione}
\begin{definizione}
	Diciamo che $u$ è una \emph{soluzione forte}\index[analitico]{soluzione!forte} del problema di Dirichlet
	\begin{equation}
		\begin{cases}
			Lu = f & \textup{in } U,\\
			u = 0 & \textup{su } \partial U,
		\end{cases}
	\end{equation}
	se $u \in H^{1}_0(U)$ e $L u(x) = f(x)$ per \emph{q.o.} $x \in U$.
\end{definizione}
\begin{proposizione}
	Se $u$ è una soluzione debole del problema di Dirichlet
	\begin{equation}
		\begin{cases}
			Lu = f & \textup{in } U,\\
			u = 0 & \textup{su } \partial U,
		\end{cases}
	\end{equation}
	e $u \in H^2_{loc}(U)$, allora $u$ è una soluzione forte dello stesso problema.	
\end{proposizione}
\begin{proof}
	Si tratta di una variante di quanto visto nell'esempio \eqref{esempiopoisson}.
\end{proof}
\begin{osservazione}
	Il lettore potrebbe essere ora sconfortato, in quanto abbiamo definito solo soluzioni di problemi con condizione al bordo nulla. In realtà questo non è restrittivo. Consideriamo il problema
	\[ \begin{cases}
		Lu = f & \textup{in } U,\\
		u = g & \textup{su } \partial U.
	\end{cases} \]
	Se cerchiamo soluzioni deboli, la seconda condizione significa che $Tu = g$, ossia che $g$ è la traccia di una qualche funzione $w \in H^1(U)$. Ma allora $\tilde{u} = u-w \in H^1_0(U)$ è soluzione del problema
	\[ \begin{cases}
		L\tilde{u} = \tilde{f}& \textup{in } U,\\
		\tilde{u} = 0 & \textup{su } \partial U.
	\end{cases} \]
	dove $\tilde{f} = f - Lw \in H^{-1}(U) = (H^1(U))'$.
\end{osservazione}

Concludiamo la sezione con un'osservazione riguardo una possibile generalizzazione che non tratteremo.
\begin{osservazione}
	La forma bilineare $B$ è ben definita già con $a_{ij}\in L^\infty(U), b_i \in L^n(U)$ e $c \in L^{\frac{n}{2}}(U)$. La motivazione è una semplice applicazione della disuguaglianza di Holder generalizzata, infatti per ogni $i,j = 1,\dots,n$, essendo $u,v \in H^1_0(U)$, $D_iu, D_jv \in L^2(U)$ e quindi
	\[ \int_U \abs*{a_{ij}D_iuD_{j}v}d\mathcal{L}^n \le \norm{a_{ij}}_\infty \norm{D_iu}_{2}\norm{D_jv}_{2}<+\infty, \]
	\[ \int_U \abs*{b_i D_iuv}d\mathcal{L}^n \le \norm{b_i}_{n}\norm{D_iu}_{2}\norm{v}_{2}<+\infty \]
	e
	\[  \int_U \abs*{cuv}d\mathcal{L}^n \le \norm{c}_{\frac{n}{2}}\norm{u}_{2}\norm{v}_{2}<+\infty. \]
\end{osservazione}

\section{Esistenza di soluzioni deboli}
Affrontiamo in questa sezione la teoria dell'esistenza delle nostre simpatiche soluzioni deboli, con la speranza che si rivelino in seguito più regolari di quanto assunto.
\begin{notazione}
	Indichiamo con $\norm{u}_{H^1_0(U)} = \norm{Du}_{L^2(U)}$\index[simboli]{$\norm{u}_{H^1_0(U)}$}.
\end{notazione} 
\begin{proposizione}\label{stimesuB}
	Valgono i seguenti fatti:
	\begin{enumerate}[$(a)$]
		\item $B$ è continua, ossia esiste $\alpha>0$ tale che per ogni $u,v \in H^1_0(U)$,
		\[ \abs*{B(u,v)} \le \alpha\norm{u}_{H^1_0(U)}\norm{v}_{H^1_0(U)}, \]
		\item esistono $\beta > 0$ e $\gamma \ge 0$ tali che per ogni $u \in H^1_0(U)$,
		\[ 	\beta\norm{u}_{H^1_0(U)}^2 \le B(u,u) + \gamma\norm{u}_{L^2(U)}^2. \]
	\end{enumerate}
	\begin{proof}\hfill
		\par\noindent
		$(a)$ Innanzitutto,
		\begin{align*}
			\abs*{B(u,v)} &= \abs*{\int_U \left( \sum_{i,j=1}^{n}a_{ij}D_iuD_{j}v + \sum_{i=1}^{n}b_i D_iuv + cuv\right) d\mathcal{L}^n} \le \\
			&\le \sum_{i,j=1}^{n} \norm{a_{ij}}_\infty \int_U \abs*{D_iu} \abs*{D_jv}d\mathcal{L}^n + \sum_{i=1}^{n}\norm{b_{i}}_\infty\int_U\abs*{D_iu} \abs*{v}d\mathcal{L}^n + \norm{c}_\infty\int_U\abs{u}\abs{v}d\mathcal{L}^n.
		\end{align*}
		Ora, per la disuguaglianza di Holder,
		\[	\abs*{B(u,v)}  \le C_1 \norm{Du}_2\norm{Dv}_2 + C_2 \norm{Du}_2\norm{v}_2 + C_3 \norm{u}_2\norm{v}_2  \]
		e, per la disuguaglianza di Poincaré,
		\[	\abs*{B(u,v)}  \le  C_1 \norm{Du}_2\norm{Dv}_2 + C_2 \norm{Du}_2\norm{Dv}_2 + C_3 \norm{Du}_2\norm{Dv}_2 \le C\norm{u}_{H^1_0(U)}\norm{v}_{H^1_0(U)}.  \]
		Si tratta quindi di considerare $\alpha > C$.
		\par\noindent
		$(b)$ Per la condizione di uniforme ellitticità, esiste $\nu> 0$ tale che per q.o. $ x \in U$ e per ogni $\xi \in \R^n$ si ha
		\[ \sum_{i,j=1}^n a_{ij}(x)\xi_i\xi_j \ge \nu\abs{\xi}^2. \]
		Allora preso $\xi = Du$,
		\begin{align*}
			\nu \int_U \abs*{Du}^2 d\mathcal{L}^n &\le \int_U \sum_{i,j=1}^n a_{ij}D_iuD_ju d\mathcal{L}^n = B(u,u) - \int_U\sum_{i=1}^n b_iD_iu \, ud\mathcal{L}^n -\int_Ucu^2 d\mathcal{L}^n \le \\
			&\le B(u,u) + \abs*{\int_U\sum_{i=1}^n b_iD_iu ud\mathcal{L}^n} + \abs*{\int_Ucu^2 d\mathcal{L}^n} \le \\
			&\le B(u,u) + \sum_{i=1}^n \int_U \abs{b_i}\abs{D_iu}\abs{u}d\mathcal{L}^n + \int_U \abs{c}\abs{u}^2d\mathcal{L}^n \le \\
			&\le B(u,u) +  \sum_{i=1}^n  \norm{b_i}_\infty\int_U \abs{D_iu}\abs{u}d\mathcal{L}^n +  \norm{c}_\infty\int_U \abs{u}^2d\mathcal{L}^n.
		\end{align*}
		Ora, per la disuguaglianza di Young pesata, dato $\varepsilon > 0 $ tale che
		\[ \varepsilon \sum_{i=1}^n  \norm{b_i}_\infty < \frac{\nu}{2},\]
		possiamo scrivere
		\[ \int_U \abs{D_iu}\abs{u}d\mathcal{L}^n \le \varepsilon \int_U\abs{Du}^2d\mathcal{L}^n + \frac{1}{4\varepsilon}\int_U\abs{u}^2d\mathcal{L}^n \]
		e quindi esiste $C > 0$ tale che
		\[ \frac{\nu}{2} \int_U \abs*{Du}^2 d\mathcal{L}^n \le B(u,u) + C\int_U\abs{u}^2d\mathcal{L}^n = B(u,u) + C\norm{u}_2^2,\]
		da cui la tesi.
	\end{proof}
\end{proposizione}
\begin{teorema}\emph{\textbf{(di esistenza di soluzioni deboli, I)}}\index[analitico]{teorema! di esistenza di soluzioni deboli, I}
	Esiste $\gamma \ge 0$ tale che per ogni $\mu \ge \gamma$ esiste un'unica soluzione debole $u \in H^1_0(U)$ del problema di Dirichlet
	\[ \begin{cases}
		Lu + \mu u = f& \textup{in } U,\\
		u = 0 & \textup{su } \partial U.
	\end{cases} \]
\end{teorema}
\begin{proof}
	Sia $\gamma \ge 0$ conforme alla Proposizione \eqref{stimesuB}. Dato $\mu \ge \gamma$, consideriamo la forma bilineare
	\[ B_\mu(u,v) = B(u,v) + \mu(u|v)_2.  \]
	Per la Proposizione \eqref{stimesuB}, $B_\mu$ è continua e coercitiva. Allora, per il Teorema di Lax--Milgram, per ogni $f \in L^2(U)$, esiste ed è unica $u \in H^1_0(U)$ tale che per ogni $v \in H^1_0(U)$
	\[ B_\mu(u,v) = (f|v)_2, \]
	da cui la tesi.
\end{proof}

\begin{definizione}
	Chiamiamo \emph{operatore aggiunto di $L$}\index[analitico]{operatore!aggiunto} l'applicazione $L'$ formalmente definita da
	\[ L'u = -\sum_{i,j=1}^{n}D_i\left( a_{ij}D_{j}u\right) -\sum_{i=1}^{n}b_i D_i u + \left( c-\sum_{i=1}^{n}D_ib_i\right)u.\]
	Chiamiamo poi \emph{forma bilineare aggiunta di $L$}\index[analitico]{forma bilineare!aggiunta} l'applicazione bilineare $B':H^1_0(U)\times H^1_0(U) \to \R$ tale che per ogni $u,v \in H^1_0(U)$
	\[ B'(u,v) = B(v,u). \]
	Diciamo infine che $u \in H^1_0(U)$ è soluzione debole del problema
	\[ \begin{cases}
		L'u=f & \textup{in } U,\\
		u= 0 & \textup{su } \partial U,
	\end{cases} \]
	se per ogni $v \in H^1_0(U)$
	\[ B'(u,v) = (f|v)_2. \]
	
	\begin{osservazione}
		La definizione precedente nasce in maniera molto naturale.
		\[ B'(u,v) = B(v,u) = \int_U \left( \sum_{i,j=1}^{n}a_{ij}D_ivD_{j}u + \sum_{i=1}^{n}b_i D_ivu + cvu\right) d\mathcal{L}^n, \]
		ma per ogni $i = 1,\dots, n$, nel caso regolare,
		\[ \int_U b_i D_i v u d\mathcal{L}^n = -\int_U D_i(b_iu)v d\mathcal{L}^n =  -\int_U D_ib_i uv d\mathcal{L}^n -\int_U b_iD_iuv d\mathcal{L}^n.  \]
		Quindi
		\[ B'(u,v) = \int_U \left( \sum_{i,j=1}^{n}a_{ij}D_ivD_{j}u + \sum_{i=1}^{n}b_i D_iuv + (c- \sum_{i=1}^nD_ib_i)uv\right) d\mathcal{L}^n \]
		ed è quindi naturale porre
		\[ L'u = -\sum_{i,j=1}^{n}D_i\left( a_{ij}D_{j}u\right) -\sum_{i=1}^{n}b_i D_i u + \left( c-\sum_{i=1}^{n}D_ib_i\right)u.\]
	\end{osservazione}
	
\end{definizione}
\begin{teorema}\emph{\textbf{(di esistenza di soluzioni deboli, II)}}\index[analitico]{teorema! di esistenza di soluzioni deboli, II}
	Una delle seguenti affermazioni è vera\footnote{In mutua esclusione.}:
	\begin{enumerate}[$(a)$]
		\item esiste un'unica soluzione debole $u \in H^1_0(U)$ del problema di Dirichlet
		\[ \begin{cases}
			Lu = f& \textup{in } U,\\
			u = 0 & \textup{su } \partial U,
		\end{cases} \]
		\item esiste una soluzione debole $u \in H^1_0(U)$, non nulla, del problema di Dirichlet
		\[ \begin{cases}
			Lu = 0 & \textup{in } U,\\
			u = 0 & \textup{su } \partial U.
		\end{cases} \]
		In particolare, si ha che $\dim{\mathcal{N}(L)} = \dim{\mathcal{N}(L')} < +\infty. $
	\end{enumerate}
	\par\noindent
	In ogni caso, il problema di Dirichlet
	\[ \begin{cases}
		Lu = f& \textup{in } U,\\
		u = 0 & \textup{su } \partial U,
	\end{cases} \]
	ammette una soluzione debole se e solo se per ogni $v \in \mathcal{N}(L')$
	\[ (f|v)_2 = 0. \]
\end{teorema}
\begin{proof}
	Innanzitutto, per il primo Teorema di esistenza di soluzioni deboli, scelto $\mu = \gamma$, per ogni $g \in L^2(U)$ esiste una ed una sola $u \in H^1_0(U)$ tale che per ogni $v \in H^1_0(U)$
	\[ B_\gamma(u,v) = (g|v)_2. \]	
	In particolare, $B_\gamma$ è biettivo in $u$ allora possiamo porre formalmente
	\[ u = L_{\gamma}^{-1}g, \]
	intendendo che $u$ è l'unico elemento di $H^1_0(U)$ che soddisfa la formulazione variazionale del problema avente come operatore $L_\gamma = L + \gamma Id$.
	
	Osserviamo che $u \in H^1_0(U)$ è una soluzione debole del problema
	\[ \begin{cases}
		Lu = f& \textup{in } U,\\
		u = 0 & \textup{su } \partial U,
	\end{cases} \]
	se e solo se per ogni $v \in H^1_0(U)$
	\[ B_\gamma(u,v) = (\gamma u + f|v), \]
	ossia se e solo se 
	\[ u = L_\gamma^{-1}(\gamma u + f).\footnote{Infatti, 
		\[ Lu = f \Longleftrightarrow L_\gamma u = Lu+\gamma u = f + \gamma u. \]} \]
	Poniamo $Ku = \gamma L_\gamma^{-1}u$ e $h = L_\gamma^{-1} f$. Allora l'ultima uguaglianza può essere riscritta come
	\[ u - Ku = h. \]
	
	Mostriamo che $K: L^2(U) \to L^2(U)$ è un operatore compatto. Per la Proposizione \eqref{stimesuB}
	\[ \beta \norm{u}_{H^1_0(U)}^2\le B_\gamma[u,u] =(g|u)_2,  \]
	ma per la disuguaglianza di Holder,
	\[ (g|u)_2 \le \norm{g}_2 \norm{u}_2 \]
	e per la disuguaglianza di Poincaré,
	\[ \norm{g}_2 \norm{u}_2 \le C \norm{g}_2 \norm{u}_{H^1_0(U)}, \]
	pertanto $ \beta \norm{u}_{H^1_0(U)}\le C \norm{g}_2$ e quindi, dal fatto che per ogni $g \in L^2(U)$
	\[ \norm{Kg}_{H^1_0(U)} = \norm{\gamma L_\gamma^{-1} g}_{H^1_0(U)} =  \norm{\gamma u}_{H^1_0(U)} \]
	si ha $ \norm{Kg}_{H^1_0(U)} \le C\norm{g}_2$. Sia quindi $(g_h)$ limitata in $L^2(U)$. Allora, per la disuguaglianza precedente, $(Kg_h)$ è limitata in $H^1_0(U)$. Essendo l'immersione $H^1_0(U)\hookrightarrow L^2(U)$ compatta per il Teorema di Rellich, si ha che $(Kg_h)$ ammette una sottosuccessione convergente in $L^2(U)$. Allora $K$ è compatto. La tesi dunque segue dal Teorema dell'alternativa di Fredholm.
\end{proof}

\begin{teorema}\emph{\textbf{(di esistenza di soluzioni deboli, III)}}\index[analitico]{teorema! di esistenza di soluzioni deboli, III}
	Esiste un insieme $\Sigma \subseteq \R$ al più numerabile con le seguenti proprietà:
	\begin{enumerate}[$(a)$]
		\item esiste ed è unica la soluzione debole del problema di Dirichlet
		\[ \begin{cases}
			Lu = \lambda u + f & \textup{in } U,\\
			u= 0 & \textup{su } \partial U,
		\end{cases} \]
		se e solo se $\lambda \notin \Sigma$,
		\item se $\Sigma$ è infinito, allora $\Sigma = \left\lbrace \lambda_h: h \in \N\right\rbrace $ tale che $\lambda_h \to +\infty$.
	\end{enumerate}
\end{teorema}
\begin{proof}
	Sia, innanzitutto, $\gamma$ definito dal primo Teorema di esistenza di soluzioni deboli.	Supponiamo senza perdita di generalità che $\lambda > -\gamma $ e $\gamma > 0$. 
	
	In base al secondo Teorema di esistenza di soluzioni deboli, il problema dato ha un'unica soluzione se e solo se la soluzione nulla è l'unica soluzione del problema omogeneo associato.
	
	Quest'ultimo fatto è equivalente al fatto che la soluzione nulla sia l'unica soluzione del problema di Dirichlet
	\[ \begin{cases}
		Lu + \gamma u = (\gamma + \lambda) u & \textup{in } U,\\
		u= 0 & \textup{su } \partial U,
	\end{cases} \]
	che può essere riscritto come $u = L_\gamma^{-1}(\gamma + \lambda)u = \frac{\gamma+\lambda}{\gamma}Ku$, dove al solito  $K: L^2(U) \to L^2(U)$ tale che $Ku= \gamma L^{-1}_\gamma u $ è un operatore compatto. 
	
	Quindi, avere come unica soluzione del problema omogeneo associato la soluzione nulla corrisponde a richiedere che $\frac{\gamma}{\gamma+\lambda}$ non sia un autovalore di $K$. Ossia $\lambda$ non deve essere del tipo successione che tende a $+\infty$ per la ben nota forma degli autovalori di un operatore compatto.
\end{proof}

Osserviamo che il problema di Dirichlet
\[ \begin{cases}
	Lu = \lambda u & \textup{in } U,\\
	u= 0 & \textup{su } \partial U,
\end{cases} \]
ha soluzione non nulla $w$ se e solo se $\lambda \in \Sigma$. In questo caso $\lambda$ è detto \emph{autovalore} di $L$ e $w$ $\emph{autofunzione}$. Ciò giustifica la seguente definizione.
\begin{definizione}
	L'insieme $\Sigma$ del Teorema precedente è detto \emph{spettro di} $L$.
\end{definizione}


\begin{corollario}
	Sia $\lambda \notin \Sigma$. Allora esiste una costante $C>0$, dipendente solo da $\lambda, U$ ed $L$, tale che per ogni $f \in L^2(U)$, detta $u \in H^1_0(U)$ la soluzione debole del problema di Dirichlet
	\[ \begin{cases}
		Lu= \lambda u + f & \textup{in } U,\\
		u=0 & \textup{su } \partial U, 
	\end{cases} \]
	risulta
	\[ \norm{u}_2 \le C\norm{f}_2. \]
\end{corollario}
\begin{proof}
	Supponiamo, per assurdo, che esistano $(f_h)$ in $L^2(U)$ e $(u_h)$ in $H^1_0(U)$ tali che per ogni $h \in \N$ $u_h$ sia la soluzione debole del problema dato e 
	\[ \norm{u_h}_2 > h\norm{f_h}_2. \]
	Poniamo $\tilde{u}_h = \frac{u_h}{\norm{u_h}_2}$ e $\tilde{f}_h = \frac{f_h}{\norm{u_h}_2}$. Allora $\norm{\tilde{u}_h}_2=1$ e, per linearità, $\tilde{u}_h$ è la soluzione debole del problema di Dirichlet
	\[ \begin{cases}
		L\tilde{u}_h= \lambda \tilde{u}_h + \tilde{f}_h & \textup{in } U,\\
		\tilde{u}_h=0 & \textup{su } \partial U, 
	\end{cases} \]
	con $1 > h\norm{\tilde{f}_h }_2$. Allora per $h \to +\infty$, $\tilde{f}_h \to 0$ in $L^2(U)$. 
	
	Inoltre, dalla Proposizione \eqref{stimesuB} combinata con la formulazione variazionale,
	\[ \norm{\tilde{u}_h}_{H^1_0(U)}^2 \le \dfrac{B(\tilde{u}_h,\tilde{u}_h)+\gamma\norm{\tilde{u}_h}_{2}^2}{\beta} =  \frac{(\lambda\tilde{u}_h+\tilde{f}_h|\tilde{u}_h)_2 + \gamma}{\beta}\]
	ed applicando la disuguaglianza di Holder, 
	\[ \norm{\tilde{u}_h}_{H^1_0(U)}^2 \le \frac{\lambda + \norm{\tilde{f}_h}_2 + \gamma}{\beta}.\]
	Ricordando che $(\tilde{f}_h)$ è convergente, quindi limitata, in $L^2(U)$, si ha
	\[ \norm{\tilde{u}_h}_{H^1_0(U)} \le C, \]
	ossia $(\tilde{u}_h)$ limitata in $H^1_0(U)$. Essendo $H^1_0(U)$ riflessivo, per il Teorema di Eberlein--Smulian, a meno di sottosuccessioni $\tilde{u}_h \rightharpoonup \tilde{u} $ in $H^1_0(U)$. In particolare,  $\tilde{u}_h \rightharpoonup  \tilde{u}$ in $L^2(U)$ e $D_j\tilde{u}_h \rightharpoonup D_j \tilde{u}$ in $L^2(U)$. Ora, per ogni $v \in H^1_0(U)$ e per ogni $i,j=1,\dots, n$, si ha che $a_{ij}D_jv, b_iv, cv, v \in L^2(U)$ e per la disuguaglianza di Holder $(\tilde{f}_h|v)_2 \to 0$. Allora passando al limite nella formulazione variazionale
	\[ B(\tilde{u}_h,v) = (\lambda\tilde{u}_h +\tilde{f}_h|v)_2, \]
	ossia, 
	\[  \int_U \left( \sum_{i,j=1}^{n}a_{ij}D_i\tilde{u}_hD_{j}v + \sum_{i=1}^{n}b_i D_i\tilde{u}_hv + c\tilde{u}_hv\right) d\mathcal{L}^n = \lambda \int_U \tilde{u}_h v d\mathcal{L}^n + \int_U \tilde{f}_h v d\mathcal{L}^n\]
	per definizione di convergenza debole,
	\[  \int_U \left( \sum_{i,j=1}^{n}a_{ij}D_i\tilde{u}D_{j}v + \sum_{i=1}^{n}b_i D_i\tilde{u}v + c\tilde{u}v\right) d\mathcal{L}^n = \lambda \int_U \tilde{u} v d\mathcal{L}^n\]
	ossia $\tilde{u}$ è soluzione debole del problema di Dirichlet
	\[ \begin{cases}
		L\tilde{u} = \lambda \tilde{u} & \textup{in } U,\\
		u = 0 & \textup{su } \partial U.
	\end{cases} \]
	Essendo $\lambda \notin \Sigma$, deve essere $\tilde{u} = 0$. L'assurdo deriva dal fatto che, per il Teorema di Rellich, le convergenze deboli trovate, a meno di sottosuccessioni, sono in realtà convergenze forti. Essendo la norma continua si ha $\norm{\tilde{u}}_2 = 1$.
\end{proof}
\begin{osservazione}
	La costante $C$ del teorema precedente tende a $+\infty$ all'avvicinarsi di $\lambda$ ad un autovalore.
\end{osservazione}
Concludiamo la sezione con un breve affondo sui problemi di Neumann. Analizziamo solo l'operatore di Laplace con condizioni al bordo nulle
\begin{equation*}
	\begin{cases}
		-\Delta u = f & \textup{in } U,\\
		\frac{\partial u}{\partial \nu} = 0 & \textup{su } \partial U
	\end{cases} 
\end{equation*}
e $U$ connesso. 
\begin{definizione}
	Sia $U$ connesso. Data $f \in C(U)$, diciamo che $u \in C^2(\overline{U})$ è una soluzione classica del problema di Neumann
	\[ \begin{cases}
		-\Delta u = f & \textup{in } U,\\
		\frac{\partial u}{\partial \nu} = 0 & \textup{su } \partial U,
	\end{cases} \]
	se per ogni per ogni $x \in U$ $-\Delta u(x) = f(x)$ e per ogni $x \in \partial U$ $\frac{\partial u}{\partial \nu}(x) = 0$.
\end{definizione}
Supponiamo dunque di avere una soluzione classica $u$. Allora per ogni $x \in U$,
\[ -\Delta u(x) = f(x). \]
Moltiplicando ambo i membri per una certa $\varphi \in C^\infty(U)$ ed integrando otteniamo 
\[ -\int_U\Delta u \varphi d\mathcal{L}^n = \int_U f \varphi d\mathcal{L}^n \] 
e per le formule di Gauss--Green abbiamo quindi, per ogni $\varphi \in C^\infty(U)$
\[ \int_U Du \cdot D\varphi d\mathcal{L}^n - \int_{\partial U} \frac{\partial u}{\partial \nu} \varphi d\mathcal{H}^{n-1} = \int_U f \varphi d\mathcal{L}^n\]
ed imponendo la condizione sulla derivata normale si ha
\[ \int_U Du \cdot D\varphi d\mathcal{L}^n = \int_U f \varphi d\mathcal{L}^n.\]
L'equazione precedente ha significato in condizioni molto più deboli rispetto a quelle classiche. In particolare ha senso per $u \in H^1(U)$. Risulta naturale quindi porre la seguente definizione.
\begin{definizione}
	Sia $U$ connesso. Diciamo che $u \in H^1(U)$ è soluzione debole del problema di Neumann
	\[ \begin{cases}
		-\Delta u = f & \textup{in } U,\\
		\frac{\partial u}{\partial \nu} = 0 & \textup{su } \partial U,
	\end{cases} \]
	se per ogni $v \in H^1(U)$
	\[ \int_U Du \cdot Dv d\mathcal{L}^n = \int_U fv d\mathcal{L}^n. \]
\end{definizione}
\begin{teorema}
	Sia $U$ connesso. Allora esiste $u \in H^1(U)$ soluzione debole del problema di Neumann
	\[ \begin{cases}
		-\Delta u = f & \textup{in } U,\\
		\frac{\partial u}{\partial \nu} = 0 & \textup{su } \partial U,
	\end{cases} \]
	se e solo se 
	\[ \int_U f d\mathcal{L}^n = 0. \]
\end{teorema}
\begin{proof}
	Sia $B: H^1(U) \times H^1(U) \to \R$ l'applicazione bilineare tale che
	\[ B(u,v) = \int_U Du \cdot Dv d\mathcal{L}^n. \]
	Siccome
	\[ \abs{B(u,v)} \le \norm{Du}_2\norm{Dv}_2 \le \norm{u}_{H^1(U)} \norm{v}_{H^1(U)},\]  
	$B$ è anche continua. Chiaramente poi l'applicazione $\left\lbrace v \to \int_U fvd\mathcal{L}^n = (f|v)_2\right\rbrace $ è lineare e 
	\[ \abs*{\int_U fvd\mathcal{L}^n} \le \norm{f}_2\norm{v}_{H^1(U)}, \]
	quindi è continua.
	
	Consideriamo il sottospazio chiuso di $H^1(U)$
	\[ X = \Set{ u \in H^1(U): \int_U u d\mathcal{L}^n = 0} \]
	e lo dotiamo della norma
	\[ \norm{u}_X = \left( \int_U \abs{Du}^2d\mathcal{L}^n\right) ^{\frac{1}{2}}. \]
	Notiamo subito che, per il Teorema di Poincaré a media nulla, esiste $C>0$ tale che per ogni $u \in X$
	\[ \norm{u}_2 \le C\norm{Du}_2 \]
	allora la norma $\norm{\;}_X$ è equivalente a quella dello spazio ambiente. Inoltre, se $v \in X$
	\[ \abs*{\int_U fvd\mathcal{L}^n} \le \norm{f}_2\norm{v}_{H^1(U)} \le C\norm{f}_2\norm{v}_{X}\]
	pertanto l'applicazione $\left\lbrace v \to \int_U fvd\mathcal{L}^n = (f|v)_2\right\rbrace $ è lineare e continua anche su $X$. Definiamo l'applicazione bilineare e continua $\tilde{B}: X\times X \to \R$ tale che
	\[ \tilde{B}(u,v) = B(u,v). \] 
	Osserviamo che
	\[ \tilde{B}(u,u) = \norm{u}_X^2, \]
	ossia $\tilde{B}$ è coercitiva. Allora per il Teorema di Lax--Milgram esiste ed è unica $u \in X$ tale che per ogni $v \in X$
	\[ B(u,v) = (f|v)_2. \]
	In particolare, $u \in H^1(U)$. Sia quindi $v \in H^1(U)$. Consideriamo $\tilde{v}= v- \fint_Uvd\mathcal{L}^n \in X$. Risulta
	\[ B(u,v) = B(u,\tilde{v}) = (f|\tilde{v})_2 = \int_U fv d\mathcal{L}^n - \fint_U v d\mathcal{L}^n \int_U f d\mathcal{L}^n.   \]
	Se 
	\[ \int_U f d\mathcal{L}^n = 0 \]
	allora $u$ è soluzione debole del problema dato.
	
	Viceversa, se esiste $u \in H^1(U)$ soluzione debole del problema dato, allora per ogni $v \in H^1(U)$
	\[ \fint_U v d\mathcal{L}^n \int_U f d\mathcal{L}^n = 0. \]
	La posizione $v = 1 \in H^1(U)$ restituisce
	\[ \int_U f d\mathcal{L}^n = 0. \qedhere \]
\end{proof}
\section{Regolarità}
In questa sezione studiamo la regolarità di soluzioni deboli di problemi di Dirichlet. In particolare, la fiducia che abbiamo riposto nella formulazione debole verrà ripagata.

Per dimostrare il seguente risultato faremo uso del \emph{metodo di Stampacchia}.
\begin{teorema}
	Supponiamo per ogni $i,j=1,\dots,n$, $a_{ij} \in C^1(U)$. Se $u \in H^1(U)$ soluzione debole di 
	\[ Lu= f \; \textup{in } U, \]
	allora $u \in H^2_{loc}(U)$. In particolare, per ogni aperto $V \subseteq U$ tale che $\overline{V} \subseteq U$, esiste $C>0$, dipendente solo da $U,V$ ed $L$, tale che 
	\[ \norm{u}_{H^2(V)} \le C\left( \norm{f}_{L^2(U)} + \norm{u}_{L^2(U)}\right).   \]
\end{teorema}
\begin{proof}
	Sia $W$ un aperto tale che $\overline{V} \subseteq W \subseteq \overline{W} \subseteq U$. Consideriamo la funzione di cut-off\footnote{Necessaria siccome non abbiamo informazioni sul comportamento al bordo.} $\zeta \in C^\infty_c(U)$ tale che $0\le \zeta \le 1, \zeta = 1$ su $V$, $\zeta = 0$ su $\R^n \setminus W$.
	Siccome $u \in H^1(U)$ è soluzione debole di 
	\[ Lu= f \; \textup{in } U, \]
	allora per ogni $v \in H^1_0(U)$ 
	\[ B(u,v) = (f|v)_2. \]
	In particolare,
	\[ \sum_{i,j=1}^n \int_U a_{ij}D_i u D_j v d\mathcal{L}^n = \int_U \tilde{f}v d\mathcal{L}^n \]
	dove $\tilde{f} = f - \sum_{i=1}^nb_i D_i u - cu. $
	
	Sia $h \ne 0$ sufficientemente piccolo, $ k =1,\dots, n$ e $v = -D_k^{-h}(\zeta^2D^h_ku) \in H^1_0(U)$. Allora
	\[ \sum_{i,j=1}^n \int_U a_{ij}D_i u D_j \left(-D_k^{-h}(\zeta^2D^h_ku) \right)  d\mathcal{L}^n = \int_U \tilde{f}\left(-D_k^{-h}(\zeta^2D^h_ku) \right)  d\mathcal{L}^n \]
	Denotiamo con 
	\begin{align*}
		A&=-\sum_{i,j=1}^n \int_U a_{ij}D_i u D_j \left(D_k^{-h}(\zeta^2D^h_ku) \right)  d\mathcal{L}^n, & B&=-\int_U \tilde{f}\left(D_k^{-h}(\zeta^2D^h_ku) \right)  d\mathcal{L}^n.
	\end{align*}
	Innanzitutto, potendo scambiare derivata debole e rapporto incrementale, applicando poi la formula di integrazione per parti con rapporti incrementali,
	\[ A=-\sum_{i,j=1}^n \int_U a_{ij}D_i u D_k^{-h} \left(D_j(\zeta^2D^h_ku) \right)  d\mathcal{L}^n =\sum_{i,j=1}^n \int_U D_k^{h}\left( a_{ij}D_i u\right)D_j(\zeta^2D^h_ku)d\mathcal{L}^n, \]
	ed utilizzando le formule di Leibniz per rapporti incrementali\footnote{Si tratta di 
		\[ D_k^h(vw) = v^hD_k^hw + w d^h_kv, \]
		con $v^h(x) = v(x+he_k)$.} e per derivate deboli
	\[ A =\sum_{i,j=1}^n \int_U \left[a^h_{ij}D^h_kD_iu + D^h_ka_{ij}D_iu \right]\left[2\zeta D_j \zeta D^h_ku+\zeta^2D^h_kD_ju \right]  d\mathcal{L}^n = A_1 + A_2 \]
	dove
	\[ A_1 = \sum_{i,j=1}^n \int_U \zeta^2a^h_{ij}D^h_kD_iuD^h_kD_ju d\mathcal{L}^n \]
	\[ A_2 = \sum_{i,j=1}^n \int_U \left(a^h_{ij}D^h_kD_iu2\zeta D_j \zeta D^h_ku + D^h_ka_{ij}D_iu\zeta^2D^h_kD_ju + D^h_ka_{ij}D_iu2\zeta D_j \zeta D^h_ku \right)  d\mathcal{L}^n. \]
	Innanzitutto, per uniforme ellitticità, esiste $\nu > 0$ tale che
	\[ A_1\ge \nu \int_U \zeta^2 \abs{D^h_kDu}^2 d\mathcal{L}^n.  \]
	Ora,
	\[ \abs*{\int_U a^h_{ij}D^h_kD_iu2\zeta D_j \zeta D^h_ku d\mathcal{L}^n} \le C \int_U \zeta \abs{D^h_kDu}\abs{D^h_ku}d\mathcal{L}^n = C \int_W \zeta\abs{D^h_kDu}\abs{D^h_ku}d\mathcal{L}^n,  \]
	per la disuguaglianza di Young pesata con $\varepsilon = \frac{\nu}{4C}$,
	\begin{align*}
		\abs*{\int_U a^h_{ij}D^h_kD_iu2\zeta D_j \zeta D^h_ku d\mathcal{L}^n} &\le \frac{\nu}{4}\int_W \zeta^2 \abs{D^h_kDu}^2d\mathcal{L}^n + C \int_W \abs{D^h_ku}^2d\mathcal{L}^n = \\
		&=\frac{\nu}{4}\int_U \zeta^2 \abs{D^h_kDu}^2d\mathcal{L}^n + C \int_W \abs{D^h_ku}^2d\mathcal{L}^n
	\end{align*}
	e, per il Teorema sui rapporti incrementali,
	\[ \abs*{\int_U a^h_{ij}D^h_kD_iu2\zeta D_j \zeta D^h_ku d\mathcal{L}^n} \le  \frac{\nu}{4}\int_U \zeta^2 \abs{D^h_kDu}^2d\mathcal{L}^n + C \int_U \abs{Du}^2d\mathcal{L}^n.\]
	Invece,
	\[ \abs*{\int_U D^h_ka_{ij}D_iu\zeta^2D^h_kD_ju d\mathcal{L}^n} \le C \int_U \zeta \abs{D^h_kDu}\abs{Du}d\mathcal{L}^n\]
	e, per la disuguaglianza di Young pesata con $\varepsilon = \frac{\nu}{4C}$,
	\[ \abs*{\int_U D^h_ka_{ij}D_iu\zeta^2D^h_kD_ju d\mathcal{L}^n} \le \frac{\nu}{4} \int_U \zeta^2 \abs{D^h_kDu}^2 d\mathcal{L}^n + C\int_U \abs{Du}^2d\mathcal{L}^n.\]
	Inoltre,
	\[ \abs*{\int_U D^h_ka_{ij}D_iu2\zeta D_j \zeta D^h_ku d\mathcal{L}^n} \le C \int_W \abs{D^h_k u}\abs{Du}d\mathcal{L}^n\]
	e, per la disuguaglianza di Young ed il Teorema sui rapporti incrementali,
	\begin{align*}
		\abs*{\int_U D^h_ka_{ij}D_iu2\zeta D_j \zeta D^h_ku d\mathcal{L}^n} &\le  C\int_W \abs{Du}^2 d\mathcal{L}^n + C\int_{W}\abs{D^h_ku}^2d\mathcal{L}^n \le \\
		&\le  C\int_U \abs{Du}^2 d\mathcal{L}^n.
	\end{align*}
	In definitiva,
	\[ A \ge  \frac{\nu}{2} \int_U \zeta^2 \abs{D^h_kDu}^2d\mathcal{L}^n - C \int_U \abs{Du}^2d\mathcal{L}^n. \]
	Ora, per il Teorema sui rapporti incrementali,
	\begin{align*}
		\int_U \abs{v}^2 d\mathcal{L}^n &= \int_U \abs*{D^{-h}_k\left( \zeta^2 D^h_k u\right)}^2d\mathcal{L}^n = \int_W \abs*{D^{-h}_k\left( \zeta^2 D^h_k u\right)}^2d\mathcal{L}^n \le \\
		&\le \int_U \abs*{D\left( \zeta^2 D^h_ku\right) }^2d\mathcal{L}^n \le C \int_U\left(  \zeta^2\abs*{D^h_k u}^2 + \zeta^2\abs*{D^h_kDu}^2\right) d\mathcal{L}^n \le \\
		&\le C \int_U \abs{Du}^2 d\mathcal{L}^n + C\int_U \zeta^2 \abs*{D^h_kDu}^2d\mathcal{L}^n.
	\end{align*}
	Allora, per la disuguaglianza di Young pesata con $\varepsilon = \frac{\nu}{4C}$,
	\begin{align*}
		\abs{B} &= \abs*{\int_U \left( f - \sum_{i=1}^{n}b_iD_iu - cu\right)\! vd\mathcal{L}^n} \le \frac{\nu}{4C} \int_U \abs{v}^2 d\mathcal{L}^n + C \!\! \int_U \!\! \left( \abs{f}^2 + \abs{u}^2 + \abs{Du}^2\right) d\mathcal{L}^n \le \\
		&\le \frac{\nu}{4} \int_U \zeta^2 \abs*{D^h_kDu}^2 d\mathcal{L}^n + C\int_U\left( \abs{f}^2 + \abs{u}^2 + \abs{Du}^2\right) d\mathcal{L}^n.
	\end{align*}
	In particolare,
	\begin{align*}
		\frac{\nu}{2} \int_U \zeta^2 \abs{D^h_kDu}^2d\mathcal{L}^n - C \int_U \abs{Du}^2d\mathcal{L}^n &\le A = B \le \\
		&\le  \frac{\nu}{4} \int_U \zeta^2 \abs*{D^h_kDu}^2 d\mathcal{L}^n + \\
		&+ C\int_U\left( \abs{f}^2 + \abs{u}^2 + \abs{Du}^2\right) d\mathcal{L}^n.
	\end{align*}
	Riassumendo, per ogni $k =1,\dots, n$ e per $h$ sufficientemente piccolo, 
	\[  \frac{\nu}{4} \int_U \zeta^2 \abs*{D^h_kDu}^2 d\mathcal{L}^n \le C\int_U\left( \abs{f}^2 + \abs{u}^2 + \abs{Du}^2\right) d\mathcal{L}^n. \]
	Per il Teorema sui rapporti incrementali otteniamo quindi che $Du \in H^1_{loc}(U)$, ossia $u \in H^2_{loc}(U)$. In particolare,
	\[ \norm{u}_{H^2(V)} \le C\left( \norm{f}_{L^2(U)} + \norm{u}_{H^1(U)}\right). \]
	\indent Ora, ripercorrendo i passaggi svolti, osserviamo che vale in particolare la stima
	\[ \norm{u}_{H^2(V)} \le C\left( \norm{f}_{L^2(W)} + \norm{u}_{H^1(W)}\right). \]
	Consideriamo ora la funzione di cut-off $\zeta \in C^\infty_c(U)$ tale che $0\le \zeta \le 1, \zeta = 1$ su $W$, $\zeta = 0$ su $\R^n \setminus U$. Posto $v = \zeta^2 u \in H^1_0(U)$ nella formulazione variazionale,
	\[ \sum_{i,j=1}^{n}\int_U a_{ij}D_iuDj\left( \zeta^2u\right) d\mathcal{L}^n = \int_U \left( \zeta^2fu - \zeta^2 u \sum_{i=1}^{n}b_iD_iu -c\zeta^2u^2\right) d\mathcal{L}^n \]
	\[ \sum_{i,j=1}^{n}\int_U a_{ij}D_iu\left( 2\zeta D_j \zeta u + \zeta^2 D_ju\right) d\mathcal{L}^n = \int_U \left( \zeta^2fu - \zeta^2 u \sum_{i=1}^{n}b_iD_iu -c\zeta^2u^2\right) d\mathcal{L}^n \]
	\[ \sum_{i,j=1}^{n}\int_U \left( 2a_{ij}\zeta D_j \zeta u D_iu +a_{ij}\zeta^2D_iuD_ju\right)  d\mathcal{L}^n = \int_U \left( \zeta^2fu - \zeta^2 u \sum_{i=1}^{n}b_iD_iu -c\zeta^2u^2\right) d\mathcal{L}^n \]
	da cui, utilizzando l'uniforme ellitticità e la disuguaglianza di Young, si ottiene
	\[ \int_U \zeta^2\abs{Du}^2d\mathcal{L}^n \le C\left(  \int_U \abs{f}^2d\mathcal{L}^n +\int_U \abs{u}^2d\mathcal{L}^n\right)  \]
	ma, per come è stato scelto il cut-off, ciò significa
	\[ \norm{u}_{H^1(W)} \le C\left( \norm{f}_{L^2(U)} + \norm{u}_{L^2(U)}\right)  \]
	da cui
	\[ \norm{u}_{H^2(V)} \le C\left( \norm{f}_{L^2(U)} + \norm{u}_{L^2(U)}\right)  \]
	che fornisce la tesi.
\end{proof}
\begin{osservazione}
	Nel Teorema precedente non abbiamo richiesto che $u \in H^1_0(U)$ intendendo che il problema di Dirichlet potrebbe avere anche dato al bordo diverso da $0$.
\end{osservazione}

Il processo può dunque essere iterato.
\begin{teorema}
	Siano $m \in \N$ ed $f \in H^m(U)$. Supponiamo per ogni $i,j=1,\dots,n$, $a_{ij},b_i,c \in C^{m+1}(U)$. Se $u \in H^1(U)$ soluzione debole di 
	\[ Lu= f \; \textup{in } U, \]
	allora $u \in H^{m+2}_{loc}(U)$. In particolare, per ogni aperto $V \subseteq U$ tale che $\overline{V} \subseteq U$, esiste $C>0$, dipendente solo da $m,U,V$ ed $L$, tale che 
	\[ \norm{u}_{H^{m+2}(V)} \le C\left( \norm{f}_{H^m(U)} + \norm{u}_{L^2(U)}\right).   \]
\end{teorema}
\begin{proof}
	Omettiamo la dimostrazione.
\end{proof}
\begin{teorema}
	Sia $f \in C^\infty(U)$. Supponiamo per ogni $i,j=1,\dots,n$, $a_{ij},b_i,c \in C^\infty(U)$. Se $u \in H^1(U)$ soluzione debole di 
	\[ Lu= f \; \textup{in } U, \]
	allora $u \in C^\infty(U)$.
\end{teorema}
\begin{proof}
	Omettiamo la dimostrazione.
\end{proof}
Aggiungendo altre ipotesi sui coefficienti, si può produrre anche regolarità sul bordo. Riportiamo per completezza i risultati, senza dimostrazione.

\begin{teorema}
	Supponiamo $\partial U$ di classe $C^2$ e per ogni $i,j=1,\dots,n$ che $a_{ij} \in C^1(\overline{U})$. Se $u \in H^1_0(U)$ soluzione debole di 
	\[ \begin{cases}
		Lu= f & \textup{in } U, \\
		u=0& \textup{su } \partial U,
	\end{cases} \]
	allora $u \in H^2(U)$. In particolare, esiste $C>0$, dipendente solo da $U$ ed $L$, tale che 
	\[ \norm{u}_{H^2(U)} \le C\left( \norm{f}_{L^2(U)} + \norm{u}_{L^2(U)}\right).   \]
	In più, se $u$ è l'unica soluzione, 
	\[ \norm{u}_{H^2(U)} \le C\norm{f}_{L^2(U)}.   \]
\end{teorema}
\begin{proof}
	Omettiamo la dimostrazione.
\end{proof}
\begin{teorema}
	Sia $m \in \N$. Supponiamo $\partial U$ di classe $C^{m+2}$, $f \in H^m(U)$ e per ogni $i,j=1,\dots,n$ che $a_{ij},b_i,c \in C^{m+1}(\overline{U})$. Se $u \in H^1_0(U)$ soluzione debole di 
	\[ \begin{cases}
		Lu= f & \textup{in } U, \\
		u=0& \textup{su } \partial U,
	\end{cases} \]
	allora $u \in H^{m+2}(U)$. In particolare, esiste $C>0$, dipendente solo da $m,U$ ed $L$, tale che 
	\[ \norm{u}_{H^{m+2}(U)} \le C\left( \norm{f}_{H^m(U)} + \norm{u}_{L^2(U)}\right).   \]
	In più, se $u$ è l'unica soluzione, 
	\[ \norm{u}_{H^{m+2}(U)} \le C\norm{f}_{H^m(U)}.   \]
\end{teorema}
\begin{proof}
	Omettiamo la dimostrazione.
\end{proof}
\begin{teorema}\label{smoothuptotheboundary}
	Supponiamo $\partial U$ di classe $C^\infty$, $f \in C^\infty(\overline{U})$  e per ogni $i,j=1,\dots,n$ che $a_{ij},b_i,c \in C^\infty(\overline{U})$. Se $u \in H^1_0(U)$ soluzione debole di 
	\[ \begin{cases}
		Lu= f & \textup{in } U, \\
		u=0& \textup{su } \partial U,
	\end{cases} \]
	allora $u \in C^\infty(\overline{U})$.
\end{teorema}
\begin{proof}
	Omettiamo la dimostrazione.
\end{proof}

Per giustizia morale, i risultati che non abbiamo dimostrato in questa sezione si basano tutti sull'idea di iterare il Metodo di Stampacchia adattandolo opportunamente. In questo modo il cerchio è chiuso: siamo partiti da problemi formulati classicamente, li abbiamo poi reinterpretati allargando il concetto di soluzione con la formulazione debole: in questo contesto siamo riusciti a dimostrare l'esistenza di soluzioni. Ora, avendo dimostrato che in realtà le soluzioni deboli che avevamo trovato possono essere reinterpretate classicamente, possiamo affermare effettivamente di aver risolto il problema da cui eravamo partiti.
\section{Principi del massimo}
In questa sezione siamo interessati a proprietà puntuali. Lavoreremo con $L$ non in forma di divergenza a coefficienti continui. Le altre ipotesi di struttura, salvo diversa specificazione, rimangono invariate.
\begin{teorema}\emph{\textbf{(principio del massimo debole)}}\index[analitico]{principio!del massimo!debole}
	Sia $u \in C^2(U)\cap C(\overline{U}) $. Supponiamo che $c = 0 $ in $U$. Valgono i seguenti fatti:
	\begin{enumerate}[$(a)$]
		\item se $Lu \le 0$ in $U$, allora
		\[ \underset{\overline{U}}{\max} \; u = \underset{\partial U}{\max} \; u. \]
		In altre parole, le sottosoluzioni raggiungono il loro massimo su $\partial U$.
		\item se $Lu \ge 0$ in $U$, allora
		\[ \underset{\overline{U}}{\min} \; u = \underset{\partial U}{\min} \; u. \]
		In altre parole, le supersoluzioni raggiungono il loro minimo su $\partial U$.
	\end{enumerate}
	
	\begin{proof}\hfill
		\par\noindent
		$(a)$ Supponiamo inizialmente che $Lu < 0$ in $U$. Se, per assurdo, esiste $\xi \in U$ tale che $u(\xi) = \underset{\overline{U}}{\textup{max}}\; u$, allora
		\begin{align*}
			Du(\xi) &= 0, & D^2u(\xi) &\le 0.
		\end{align*}
		
		Essendo $A = (a_{ij})$ simmetrica e definita positiva, esiste una matrice ortogonale $O=(o_{ij})$ tale che 
		\begin{align*}
			OAO^T &= \textup{diag}(d_1,\dots,d_n), & OO^T&=Id,
		\end{align*}
		con $d_k> 0$ per ogni $k=1,\dots,n$. Introduciamo il cambio di variabile $y = \xi + O(x-\xi)$. Allora $x = \xi + O^T(y-\xi)$. Sia $v$ tale che $u(x) = v(\xi + O(x-\xi)).$ Allora,
		\begin{align*}
			D_i u &= \sum_{k=1}^{n} D_kv o_{ik}, & D^2_{i,j}u &= \sum_{k,l=1}^{n}D^2_{k,l}v o_{ik}o_{jl}
		\end{align*}
		e quindi
		\begin{align*}
			\sum_{i,j=1}^{n}a_{ij}D^2_{i,j}u(\xi) &= \sum_{i,j=1}^{n}\sum_{l,k=1}^{n}a_{ij}D^2_{k,l}v(\xi) o_{ik}o_{jl} = \sum_{l,k=1}^{n}\sum_{i,j=1}^{n}a_{ij}D^2_{k,l}v(\xi) o_{ik}o_{jl} = \\
			&= \sum_{l,k=1}^{n} d_l\delta_k^l D^2_{k,l}v(\xi) = \sum_{k=1}^{n}d_kD^2_{k,k}v(\xi).
		\end{align*}
		A questo punto, dal fatto che $D^2u(\xi) \le 0$, segue che anche $D^2v(\xi) \le 0$, ossia per ogni $\eta \in \R^n$
		\[ \sum_{k,l=1}^{n} D^2_{k,l} v(\xi) \eta_k\eta_l \le 0. \]
		Scegliendo come $\eta$ gli elementi della base canonica si ottiene che per ogni $k =1,\dots,n$ 
		\[ D^2_{k,k}v(\xi) \le 0. \]
		Ne segue che $Lu(\xi) \ge 0$, assurdo.
		
		Nel caso generale, sia $w(x) = u(x) + \varepsilon e^{\lambda x_1}$ con $\varepsilon,\lambda>0$. Per uniforme ellitticità, $a_{ii}(x) \ge \nu$ per ogni $i=1,\dots, n$ e per ogni $x \in U$. In particolare, $a_{11}(x) \ge \nu$. Inoltre,
		\begin{align*}
			D_i(e^{\lambda x_1}) &= \lambda e^{\lambda x_1} \delta_{i}^1, & D^2_{i,j}(e^{\lambda x_1}) &= \lambda^2 e^{\lambda x_1}\delta_{i}^1\delta_{j}^1. 
		\end{align*}
		Allora,
		\begin{align*}
			Lw &= Lu + \varepsilon L(e^{\lambda x_1}) \le \varepsilon L(e^{\lambda x_1}) = -\sum_{i,j=1}^{n}a_{ij}\varepsilon\lambda^2e^{\lambda x_1}\delta_{i}^1\delta_{j}^1 + \sum_{i=1}^{n}\varepsilon b_i\lambda e^{\lambda x_1} \delta_{i}^1 = \\
			&=-a_{11}\varepsilon \lambda^2 e^{\lambda x_1} + b_1\varepsilon e^{\lambda x_1} \le -\nu\varepsilon \lambda^2 e^{\lambda x_1} + b_1\varepsilon e^{\lambda x_1} \le \varepsilon e^{\lambda x_1}\left( -\nu\lambda^2 + \lambda \norm{b}_\infty\right) <0
		\end{align*}
		per $\lambda$ sufficientemente grande. Allora, per il passo precedente,
		\[ \underset{\overline{U}}{\max} \; w = \underset{\partial U}{\max} \; w. \]
		Passando al limite per $\varepsilon \to 0^+$ si ottiene la tesi.\footnote{Siano $f,g$ due funzioni continue su $\overline{U}$ tali che 
			\[ \underset{\overline{U}}{\max} \; (f+\varepsilon g) = \underset{\partial U}{\max} \; (f+\varepsilon g). \]
			Allora, se $\varepsilon \to 0^+$,
			\[ \underset{\overline{U}}{\max} \; f = \underset{\partial U}{\max} \; f. \]
			Infatti, dato $x \in \overline{U}$,
			\[ f(x) = f(x) + \varepsilon g(x) - \varepsilon g(x) \le \underset{\overline{U}}{\max} \; (f+\varepsilon g) + \varepsilon \norm{g}_\infty = \underset{\partial U}{\max} \; (f+\varepsilon g) + \varepsilon \norm{g}_\infty \le \underset{\partial U}{\max} \; f + 2\varepsilon\norm{g}_\infty.  \]
			Quindi 
			\[ \underset{\partial U}{\max} \; f \le \underset{\overline{U}}{\max} \; f \le \underset{\partial U}{\max} \; f + 2\varepsilon\norm{g}_\infty \]
			e passando al limite per $\varepsilon \to 0^+$ si conclude.}
		\par\noindent
		$(b)$ Si tratta di scambiare $u$ con $-u$ nella $(a)$.
	\end{proof}
\end{teorema}
\begin{lemma}\textbf{\emph{(di Hopf)}}\index[analitico]{lemma!di Hopf}
	Sia $u \in C^2(U)\cap C^1(\overline{U})$. Supponiamo che $Lu \le 0$ in $U$, che esista $\xi \in \partial U$ tale che per ogni $x \in U$ si abbia
	\[ u(\xi ) > u(x) \]
	e che esista una palla aperta $\textup{B}$ tale che $\xi \in \partial \textup{B}$.\footnote{Tale condizione è automaticamente soddisfatta se $\partial U$ è di classe $C^2$.} Sia $\nu$ il versore normale esterno a $\textup{B}$ in $\xi$.  Valgono allora i seguenti fatti:
	\begin{enumerate}[$(a)$]
		\item se $c = 0 $ in $U$, allora
		\[ \frac{\partial u}{\partial \nu}(\xi) > 0, \]
		\item se $c \ge 0$ in $U$ e $u(\xi) \ge 0$, allora
		\[ \frac{\partial u}{\partial \nu}(\xi) > 0. \]
	\end{enumerate} 
\end{lemma}
\begin{proof}
	Osserviamo innanzitutto che 
	\[ \frac{\partial u}{\partial \nu}(\xi) \ge 0. \]
	Infatti, 
	\[ \frac{\partial u}{\partial \nu}(\xi) = \lim\limits_{h \to 0^-} \frac{u(\xi +h\nu) -u(\xi)}{h} \ge 0\]	
	in quanto $h < 0$ e $u(\xi +h\nu) -u(\xi)<0$.
	\par\noindent
	$(b)$ A meno di traslazione, possiamo supporre che $\textup{B} = \textup{B}(0,r)$ per qualche $r>0$. Poniamo, per ogni $x \in \textup{B}$, $v(x) = e^{-\lambda \abs{x}^2} -e^{-\lambda r^2}$ con $\lambda >0$. Allora $v(x) = 0$ per ogni $x \in \partial \textup{B}$ e 
	\begin{align*}
		D_i v(x) &= -2\lambda x_i e^{-\lambda \abs{x}^2}, & D^2_{i,j}v(x) =-2\lambda\delta_i^j e^{-\lambda \abs{x}^2} + 4\lambda^2 x_ix_je^{-\lambda \abs{x}^2}.
	\end{align*}
	Pertanto, utilizzando l'uniforme ellitticità con costante $\vartheta$,
	\begin{align*}
		Lv &= -\sum_{i,j=1}^{n}a_{ij}D^2_{i,j}v + \sum_{i=1}^{n} b_iD_iv +cv = \\
		&= -\sum_{i,j=1}^{n} a_{ij}\left( 4\lambda^2x_ix_je^{-\lambda \abs{x}^2}-2\lambda\delta_i^j e^{-\lambda \abs{x}^2}\right) -2\sum_{i=1}^{n}b_i\lambda x_ie^{-\lambda \abs{x}^2} +c\left( e^{-\lambda \abs{x}^2}-e^{-\lambda r^2}\right) \le \\
		&\le e^{-\lambda \abs{x}^2}\left( -4\vartheta \lambda^2\abs{x}^2 +2\lambda n\sup_i\norm{a_{_{ii}}}_\infty +2\lambda\norm{b}_\infty\abs{x} +\norm{c}_\infty\right).
	\end{align*}
	A questo punto sia $R = \textup{B} \setminus \textup{B}(0,\frac{r}{2})$. 
	
	
	\begin{center}
		\begin{tikzpicture}[x=0.8pt,y=0.8pt,yscale=-1,xscale=1]
			%uncomment if require: \path (0,300); %set diagram left start at 0, and has height of 300
			
			%Shape: Polygon Curved [id:ds20637273862932082] 
			\draw  [color={rgb, 255:red, 255; green, 255; blue, 255 }  ,draw opacity=1 ][fill={rgb, 255:red, 155; green, 155; blue, 155 }  ,fill opacity=0.6 ] (132,214) .. controls (131,227.4) and (214,254.6) .. (254,224.6) .. controls (294,194.6) and (379,214.6) .. (413,195.6) .. controls (447,176.6) and (492.48,125.93) .. (495.24,92.27) .. controls (498,58.6) and (444,95.6) .. (411,83.6) .. controls (378,71.6) and (260,55.6) .. (176,68.6) .. controls (92,81.6) and (109,175.6) .. (114,189.6) .. controls (119,203.6) and (133,200.6) .. (132,214) -- cycle ;
			%Shape: Ellipse [id:dp8510547001841648] 
			\begin{scope}
				\tikzset{
					render blur shadow/.prefix code={
						\colorlet{black}{rgb, 255:red, 255; green, 255; blue, 255 }  
					}
				}
				
				\draw  [fill={rgb, 255:red, 255; green, 255; blue, 255 }  ,fill opacity=0.5 ] (236,158.8) .. controls (236,130.52) and (259.06,107.6) .. (287.5,107.6) .. controls (315.94,107.6) and (339,130.52) .. (339,158.8) .. controls (339,187.08) and (315.94,210) .. (287.5,210) .. controls (259.06,210) and (236,187.08) .. (236,158.8) -- cycle ;
			\end{scope}
			%Shape: Ellipse [id:dp2967342374206603] 
			\draw  [fill={rgb, 255:red, 155; green, 155; blue, 155 }  ,fill opacity=1 ] (260.97,158.8) .. controls (260.97,144.23) and (272.85,132.42) .. (287.5,132.42) .. controls (302.15,132.42) and (314.03,144.23) .. (314.03,158.8) .. controls (314.03,173.37) and (302.15,185.18) .. (287.5,185.18) .. controls (272.85,185.18) and (260.97,173.37) .. (260.97,158.8) -- cycle ;
			
			%Shape: Circle [id:dp46699756212502197] 
			\draw  [fill={rgb, 255:red, 0; green, 0; blue, 0 }  ,fill opacity=1 ] (283.2,158.8) .. controls (283.2,156.43) and (285.13,154.5) .. (287.5,154.5) .. controls (289.87,154.5) and (291.8,156.43) .. (291.8,158.8) .. controls (291.8,161.17) and (289.87,163.1) .. (287.5,163.1) .. controls (285.13,163.1) and (283.2,161.17) .. (283.2,158.8) -- cycle ;
			%Shape: Circle [id:dp92244155475394] 
			\draw  [fill={rgb, 255:red, 0; green, 0; blue, 0 }  ,fill opacity=1 ] (294.2,208.8) .. controls (294.2,206.43) and (296.13,204.5) .. (298.5,204.5) .. controls (300.87,204.5) and (302.8,206.43) .. (302.8,208.8) .. controls (302.8,211.17) and (300.87,213.1) .. (298.5,213.1) .. controls (296.13,213.1) and (294.2,211.17) .. (294.2,208.8) -- cycle ;
			%Curve Lines [id:da6201100550341501] 
			\draw    (132,214) .. controls (131,227.4) and (214,254.6) .. (254,224.6) .. controls (294,194.6) and (379,214.6) .. (413,195.6) .. controls (447,176.6) and (492.48,125.93) .. (495.24,92.27) ;
			
			% Text Node
			\draw (291,141.4) node [anchor=north west][inner sep=0.75pt]    {$0$};
			% Text Node
			\draw (304.8,212.2) node [anchor=north west][inner sep=0.75pt]    {$\xi$};
			% Text Node
			\draw (135,191.4) node [anchor=north west][inner sep=0.75pt]    {$U$};
			% Text Node
			\draw (465,151.4) node [anchor=north west][inner sep=0.75pt]    {$\partial U$};
			% Text Node
			\draw (245,144.4) node [anchor=north west][inner sep=0.75pt]    {$R$};
		\end{tikzpicture}
	\end{center}
	Per $\lambda$ abbastanza grande abbiamo, in $R$,
	\[ Lv \le C\left( -\lambda^2 +\lambda + 1\right) \le 0.\]
	
	Mostriamo che esiste $\varepsilon > 0 $ tale che per ogni $x \in \partial \textup{B}(0,\frac{r}{2})$ si abbia
	\[ u(\xi) \ge u(x) +\varepsilon v(x). \]
	Se, per assurdo, per ogni $j \in \N\setminus \left\lbrace 0\right\rbrace $ esistesse $x_j \in \partial \textup{B}(0,\frac{r}{2})$ tale che
	\[ u(\xi) < u(x_j) + \frac{1}{j}v(x_j), \]
	dalla compattezza di $\partial \textup{B}(0,\frac{r}{2})$, esisterebbe $(x_{j_k})$ tale che $x_{j_k} \to \tilde{x}$ in $U$ per qualche $\tilde{x} \in U$. Allora $u(\xi) \le u(\tilde{x})$, assurdo.
	
	Osserviamo poi che per ogni $x \in \partial \textup{B}(0,r)$ si ha
	\[ u(\xi) \ge u(x) +\varepsilon v(x), \]
	quindi in definitiva per ogni $x \in \partial R = \partial \textup{B}(0,\frac{r}{2}) \cup \partial \textup{B}(0,r)$ si ha
	\[  u(\xi) \ge u(x) +\varepsilon v(x). \]
	
	Ora, su $R$
	\[ L\left(  u + \varepsilon v -u(\xi)\right)  = Lu +\varepsilon Lv - Lu(\xi) \le -cu(\xi) \le 0 \]
	e quindi, per il Principio del massimo debole,
	\[ \underset{\overline{R}}{\max} \; \left(  u + \varepsilon v -u(\xi)\right) = \underset{\partial R}{\max} \; \left(  u + \varepsilon v -u(\xi)\right) \le 0 \]
	da cui 
	\[   u + \varepsilon v -u(\xi) \le 0\]
	su $\overline{R}$, e quindi in particolare su $R$. In particolare, essendo $\varepsilon v(\xi) = 0$, 
	\[  u + \varepsilon v  \le u(\xi) + \varepsilon v(\xi), \]
	quindi, procedendo come nell'osservazione iniziale,
	\[ \frac{\partial u}{\partial \nu}(\xi) + \varepsilon \frac{\partial v}{\partial \nu}(\xi) \ge 0 \]
	allora
	\begin{align*}
		\frac{\partial u}{\partial \nu}(\xi) &\ge - \varepsilon \frac{\partial v}{\partial \nu}(\xi) = -\varepsilon Dv(\xi) \cdot \nu(\xi) =-\varepsilon Dv(\xi) \cdot \frac{\xi}{r} = \\
		&= -\varepsilon Dv(\xi) \cdot \frac{\xi}{r} = 2\varepsilon \lambda e^{-\lambda r^2}\frac{\xi \cdot \xi}{r} = 2\lambda \varepsilon re^{-\lambda r^2} > 0 
	\end{align*}
	da cui la tesi.
	\par\noindent
	$(a)$ Si tratta di una semplice variante di $(b)$.
\end{proof}

\begin{teorema}\textbf{\emph{(principio del massimo forte)}}\index[analitico]{principio!del massimo!forte}
	Siano $U$ connesso e $u \in C^2(U)\cap C(\overline{U}) $. Supponiamo $c= 0$ in $U$. Valgono i seguenti fatti:
	\begin{enumerate}[$(a)$]
		\item se $Lu \le 0$ in $U$ e $u$ raggiunge il suo massimo su $\overline{U}$ in un punto interno, allora $u$ è costante in $U$,
		\item se $Lu \ge 0$ in $U$ e $u$ raggiunge il suo minimo su $\overline{U}$ in un punto interno, allora $u$ è costante in $U$.
	\end{enumerate}
\end{teorema}
\begin{proof}\hfill
	\par\noindent
	$(a)$ Denotiamo con 
	\[ M = \underset{\overline{U}}{\textup{max}}\; u \]
	\[ C = \left\lbrace x \in U : u(x) = M\right\rbrace. \]
	Per ipotesi $C \ne \emptyset$. Supponiamo, per assurdo, che $u \ne M$. Allora sia 
	\[ V = \left\lbrace x \in U : u(x)< M\right\rbrace.\]
	Sia $y \in V$ tale che $\textup{dist}(y,C) < \textup{dist}(y,\partial U)$ e chiamiamo $\textup{B}$ la più grande palla aperta di centro $y$ che giace in $V$. 
	
	Esiste $\xi \in C$ tale che $\xi \in \partial \textup{B}$. Chiaramente $V$ soddisfa le ipotesi del Lemma di Hopf, pertanto
	\[ \frac{\partial u}{\partial \nu}(\xi) >0, \]
	ma essendo $\xi$ punto di massimo per $u$, $Du(\xi) = 0$, assurdo.
	\par\noindent
	$(b)$ Si tratta di scambiare $u$ con $-u$ nella $(a)$.
\end{proof}

Di particolare interesse è la disuguaglianza di Harnack. Vediamo dapprima il caso $L = \Delta$.
\begin{proposizione}
	Sia $u$ una funzione armonica. Allora per ogni aperto connesso $V$ tale che $\overline{V} \subseteq U$, esiste $C > 0$, dipendente solo da $V$, tale che 
	\[ \underset{V}{\sup} \; u \le C \underset{V}{\inf} \; u. \]
\end{proposizione}
\begin{proof}
	Sia $r = \frac{1}{4}\textup{dist}(V,\partial U)$. Siano $x,y \in V$ tali che $\abs{x-y}\le r$. Per il Teorema della media\footnote{Si rimanda alla 
		\href{https://www.dmf.unicatt.it/~musesti/EDFM/}{dispensa del corso di Equazioni differenziali della fisica matematica} del professor Alessandro Musesti per i dovuti accertamenti.}, 
	\begin{align*}
		u(x) = \fint_{\textup{B}(x,2r)} ud\mathcal{L}^n &= \frac{1}{\mathcal{L}^n(\textup{B}(0,1))(2r)^n} \int_{\textup{B}(x,2r)}ud\mathcal{L}^n \ge \frac{1}{\mathcal{L}^n(\textup{B}(0,1))(2r)^n} \int_{\textup{B}(y,r)}ud\mathcal{L}^n = \\
		&= \frac{1}{2^n} \frac{1}{\mathcal{L}^n(\textup{B}(0,1))(r)^n} \int_{\textup{B}(y,r)}ud\mathcal{L}^n = \frac{1}{2^n} \fint_{\textup{B}(y,r)} = \frac{1}{2^n} u(y). 
	\end{align*}
	Essendo $V$ connesso e $\overline{V}$ compatto, possiamo ricoprire $\overline{V}$ con una catena finita di palle aperte $\textup{B}_i$, con $i = 1,\dots, N$ ed $N$ dipendente solo da $V$, di raggio $r$ con intersezione non vuota e quindi per ogni $x,y \in V$
	\[ u(x) \ge \frac{1}{2^{nN}} u(y). \]
	In altre parole, esiste $C >0$, dipendente solo da $V$, tale che
	\[ \underset{V}{\textup{sup}}\; u \le Cu(x) \]
	e quindi 
	\[ \underset{V}{\textup{sup}}\; u \le C  \underset{V}{\textup{inf}}\; u.\qedhere\]
\end{proof}
Veniamo quindi al caso generale.
\begin{teorema}\textbf{\emph{(disuguaglianza di Harnack)}}\index[analitico]{disuguaglianza!di Harnack}
	Sia $u \in C^2(U)$ una soluzione di 
	\[ Lu = 0 \qquad \textup{in } U \]
	tale che $u\ge 0$ in $U$. Allora per ogni aperto connesso $V$ tale che $\overline{V} \subseteq U$, esiste $C > 0$, dipendente solo da $V,U$ ed $L$, tale che 
	\[ \underset{V}{\sup} \; u \le C \underset{V}{\inf} \; u. \]
\end{teorema}
\begin{proof}
	Limitiamo la dimostrazione al caso $a_{ij} \in C^\infty(U), b = 0, c=0$, ossia
	\[ \sum_{i,j=1}^{n}a_{ij}D^2_{i,j}u = 0 \qquad \textup{in } U. \]
	In particolare, allora $u \in C^\infty(U)$.
	
	Supponiamo inizialmente che $u > 0$ in $U$. Allora possiamo porre
	\[ v = \ln u. \]
	Pertanto, abbiamo che $u = e^v$ e per ogni $i,j = 1,\dots, n$
	\begin{align*}
		D_j u &= e^v D_j v & D^2_{ij}u &= e^v\left( D_i v D_j v + D^2_{i,j}v\right), 
	\end{align*}
	quindi, 
	\[ \sum_{i,j=1}^{n} a_{ij}D^2_{i,j}u = e^v \sum_{i,j=1}^n a_{ij}\left( D_ivD_jv + D^2_{ij}v\right)  = 0 \]
	ossia
	\[ \sum_{i,j=1}^{n}a_{ij}\left( D_ivD_jv + D^2_{i,j}v\right)  = 0 . \]
	In altre parole,
	\[ -\sum_{i,j=1}^{n}a_{ij}D^2_{i,j}v = \sum_{i,j=1}^{n}a_{ij}D_ivD_jv \]
	che, posto 
	\[ w =  \sum_{i,j=1}^{n}a_{ij}D_ivD_jv,\]
	diventa 
	\begin{equation}\label{eq1Harnack}
		-\sum_{i,j=1}^{n}a_{ij}D^2_{i,j}v = w.
	\end{equation}
	Ora, per uniforme ellitticità, esiste $\nu>0$ tale che 
	\[ w \ge \nu \abs{Dv}^2 \ge 0, \]
	allora $\abs{Dv}^2 \le \frac{1}{\nu}w$. Sia $V$ un aperto connesso tale che $\overline{V} \subseteq U$. Mostriamo che $w \in L^\infty(V)$. Innanzitutto, per ogni $k=1,\dots,n$,
	\[ D_k w = \sum_{i,j=1}^{n} D_k a_{ij} D_i v D_j v + \sum_{i,j=1}^{n} a_{ij} D^2_{k,i}v D_j v + \sum_{i,j=1}^{n} a_{ij} D_i v D^2_{k,j} v \]
	e dal fatto che $a_{ij} = a_{ji}$,
	\[ D_k w = \sum_{i,j=1}^{n} D_k a_{ij} D_i v D_j v + \sum_{i,j=1}^{n} a_{ij} D^2_{k,i}v D_j v + \sum_{i,j=1}^{n} a_{ji} D_i v D^2_{k,j} v, \]
	ma essendo nel terzo addendo $i,j$ indici muti, possiamo effettuare la permutazione che li scambia ottenendo
	\[ D_k w = \sum_{i,j=1}^{n} D_k a_{ij} D_i v D_j v + 2\sum_{i,j=1}^{n} a_{ij} D^2_{k,i}v D_j v. \]
	In maniera del tutto simile, per ogni $l = 1,\dots, n$,
	\[ D^2_{l,k} w = 2\sum_{i,j=1}^{n} a_{ij} D^3_{l,k,i}v D_j v + 2\sum_{i,j=1}^{n} a_{ij} D^2_{k,i}v D^2_{l,j}v + R_1 \]
	dove 
	\[ R_1 = \sum_{i,j=1}^{n} D^2_{l,k} a_{ij} D_i v D_j v + 2\sum_{i,j=1}^{n} D_k a_{ij} D^2_{l,i} v D_j v + 2\sum_{i,j=1}^{n} D_l a_{ij} D^2_{k,i} v D_j v. \]
	In particolare, esiste $C>0$, dipendente solo dagli $a_{ij}$ e da $V$, tale che
	\[ \abs{R_1} \le C\left( \abs{D^2 v}\abs{Dv} + \abs{Dv}^2\right)  \]
	e, per la disuguaglianza di Young pesata, per ogni $\varepsilon > 0$, 
	\[ \abs{R_1} \le \varepsilon \abs{D^2v}^2 + \frac{C}{\varepsilon}\abs{Dv}^2. \]
	Pertanto, 
	\[ -\sum_{k,l=1}^{n}a_{kl}D^2_{k,l}w = -2 \sum_{i,j=1}^na_{ij}D_j v \left( \sum_{k,l=1}^n a_{kl}D^3_{l,k,i} v\right)  - 2 \sum_{i,j,k,l=1}^{n} a_{ij}a_{kl}D^2_{k,i}v D^2_{l,k} v + R_2 \]
	dove 
	\[ R_2 = -\sum_{k,l=1}^n a_{kl} R_1 \]
	e chiaramente esiste $C>0$, dipendente solo dagli $a_{ij}$ e da $V$, tale che
	\[ \abs{R_2} \le \varepsilon \abs{D^2 v}^2 + \frac{C}{\varepsilon}\abs{Dv}^2. \]
	D'altra parte, per l'equazione \eqref{eq1Harnack},
	\[ D_k w = D_k \left(-\sum_{i,j=1}^{n}a_{ij}D^2_{i,j}v\right) = - \sum_{i,j=1}^na_{ij}D^3_{k,i,j}v - \sum_{i,j=1}^n D_ka_{ij} D^2_{i,j} v.  \]
	Chiamiamo 
	\[ R_3 = - \sum_{i,j=1}^n D_ka_{ij} D^2_{i,j} v, \]
	allora esiste $C>0$, dipendente solo dagli $a_{ij}$ e da $V$, tale che
	\[ \abs{R_3} \le C \abs{D^2v}. \]
	A questo punto, posto
	\[ \bar{b}_k = -2\sum_{l=1}^{n}a_{kl}D_l v, \]
	abbiamo subito che esiste $C>0$, dipendente solo dagli $a_{ij}$ e da $V$, tale che  $\abs{\bar{b}} \le C \abs{Dv}$. Inoltre, risulta che 
	\[ \sum_{k=1}^{n} \bar{b}_k D_k w = -\sum_{i,j,k=1}^{n} a_{ij}\bar{b}_k D^3_{k,i,j} v + R_4 \]
	dove 
	\[ R_4 = -\sum_{k=1}^{n}\bar{b}_k R_3 \]
	e quindi esiste $C>0$, dipendente solo dagli $a_{ij}$ e da $V$, tale che
	\[ \abs{R_4} \le \sum_{k=1}^{n}\abs{\bar{b}_k} \abs{R_3} \le 2C\sum_{k,l=1}^{n} \abs{a_{kl}} \abs{D_l v}\abs{D^2v} \le C\abs{D v} \abs{D^2 v}   \]
	e, per la disuguaglianza di Young pesata, per ogni $\varepsilon > 0$
	\[ \abs{R_4} \le \varepsilon \abs{D^2 v}^2 + \frac{C}{\varepsilon}\abs{Dv}^2. \]
	Osserviamo invece che 
	\[ \sum_{i,j,k=1}^{n} a_{ij}\bar{b}_k D^3_{k,i,j} v = 2 \sum_{i,j=1}^n a_{ij}\left( \sum_{k,l=1}^n a_{kl} D_l v D^3_{k,i,j} v\right)  \]
	ed effettuando la permutazione di indici che scambia $k,i$ e $l,j$ si ottiene
	\[  \sum_{i,j,k=1}^{n} a_{ij}\bar{b}_k D^3_{k,i,j} v = 2 \sum_{i,j=1}^na_{ij}D_j v \left( \sum_{k,l=1}^n a_{kl}D^3_{l,k,i} v\right). \]
	Alla luce di ciò, 
	\[ -\sum_{k,l=1}^{n}a_{kl}D^2_{k,l}w + \sum_{k=1}^{n} \bar{b}_k D_k w =  - 2 \sum_{i,j,k,l=1}^{n} a_{ij}a_{kl}D^2_{k,i}v D^2_{l,k} v + R_2 + R_4.  \]
	Ora, sempre per uniforme ellitticità, 
	\begin{align*}
		\sum_{i,j,k,l=1}^{n} a_{ij}a_{kl}D^2_{k,i}v D^2_{l,k} v &= \sum_{i,j=1}^{n}  a_{ij} \left( \sum_{k,l=1}^n a_{kl} D^2_{k,i}v D^2_{l,k} v \right) \ge \sum_{i,j=1}^{n}  a_{ij} \left( \nu \abs{D^2 v}^2\right) = \\
		&= \nu \abs{D^2 v}^2 \sum_{i,j=1}^{n}  a_{ij} \ge   \nu \abs{D^2 v}^2  \nu = \nu^2 \abs{D^2 v}^2,
	\end{align*}
	quindi scelto $\varepsilon = \frac{\nu^2}{4}$, esiste $C>0$, dipendente solo dagli $a_{ij}$ e da $V$, tale che
	\[ -\sum_{k,l=1}^{n}a_{kl}D^2_{k,l}w + \sum_{k=1}^{n} \bar{b}_k D_k w  \le -\frac{\nu^2}{2}\abs{D^2 v}^2 + C\abs{Dv}^2. \]
	Sia $W$ un aperto tale che $\overline{V} \subseteq W \subseteq \overline{W} \subseteq U$. Consideriamo la funzione di cut-off $\zeta \in C^\infty_c(U)$ tale che $0\le \zeta \le 1, \zeta = 1$ su $V$, $\zeta = 0$ su $\R^n \setminus W$. Consideriamo $z = \zeta ^4 w \in C_c(U)$.\footnote{In realtà è più regolare di così.} Allora,
	\[ D_k z = 4 \zeta^3 D_k \zeta w + \zeta^4 D_k w \]
	e 
	\[ D^2_{l,k} z = 12\zeta^2 D_l \zeta D_k \zeta w + 4 \zeta^3 D^2_{l,k} \zeta w + 4\zeta^3 D_k\zeta D_lw + 4\zeta^3D_l\zeta D_kw + \zeta^4 D^2_{l,k}w. \]
	Inoltre, esisterà un qualche $\xi \in U$ tale che $z$ ammette massimo su $U$ in $\xi$. In particolare, essendo $w \ge 0$, $\zeta(\xi)> 0$ e $Dz(\xi) = 0$, ossia
	\[ 4\zeta^3(\xi)D_k\zeta(\xi) w(\xi) + \zeta^4(\xi)D_kw(\xi)  =0 \]
	e dividendo per $\zeta^3(\xi)$
	\begin{equation}\label{eq2Harnack}
		4D_k\zeta(\xi) w(\xi) + \zeta(\xi)D_kw(\xi)  =0.
	\end{equation}
	Ora, 
	\begin{align*}
		-\sum_{k,l=1}^n a_{kl}D^2_{k,l}z = &-\sum_{k,l=1}^n a_{kl} \zeta^4 D^2_{l,k}w + \\
		& -\sum_{k,l=1}^n a_{kl}\left(12\zeta^2 D_l \zeta D_k \zeta w + 4 \zeta^3 D^2_{l,k} \zeta w + 4\zeta^3 D_k\zeta D_lw + 4\zeta^3D_l\zeta D_kw\right)
	\end{align*}
	mentre
	\[ \sum_{k=1}^n \bar{b}_k D_k z =  \sum_{k=1}^n \bar{b}_k \zeta^4 D_k w + \sum_{k=1}^n \bar{b}_k 4 \zeta^3 D_k \zeta w. \]
	Posto
	\begin{align*}
		R_5 =  &-\sum_{k,l=1}^n a_{kl}\left(12\zeta^2 D_l \zeta D_\zeta z w + 4 \zeta^3 D^2_{l,k} \zeta w + 4\zeta^3 D_k\zeta D_lw + 4\zeta^3D_l\zeta D_kw\right) + \\
		&+ \sum_{k=1}^n \bar{b}_k 4 \zeta^3 D_k \zeta w,
	\end{align*}
	esiste $C>0$, dipendente solo da $U,V$ ed $L$, tale che
	\[ \abs{R_5} \le C\left( \zeta^2 \abs{w} + \zeta^3 \abs{w} + \zeta^3\abs{Dw} + \abs{\bar{b}} \zeta^3 \abs{w} \right) \le C\left( \zeta^2 \abs{w} + \zeta^3\abs{Dw} + \zeta^3\abs{Dv}  \abs{w} \right).  \]
	Essendo quindi $\xi$ un punto di massimo, e siccome il prodotto scalare tra una matrice definita positiva ed una semidefinita negativa è negativo,
	\begin{align*}
		0\le  -\sum_{k,l=1}^n a_{kl}(\xi)D^2_{k,l}z(\xi) + \sum_{k=1}^n \bar{b}_k(\xi) D_k z(\xi) = &\zeta^4(\xi)\left(-\sum_{k,l=1}^n a_{kl}(\xi) D^2_{l,k}w(\xi) + \right. \\
		& +\left.\sum_{k=1}^n \bar{b}_k(\xi) D_k w(\xi)\right) +R_5(\xi),
	\end{align*}
	ma, per l'equazione \eqref{eq2Harnack}, $\abs{Dw(\xi)} \le C \abs{w(\xi)}$ e quindi 
	\[ \abs{R_5(\xi)} \le C \left( \zeta^2(\xi)\abs{w(\xi)} + \zeta^3(\xi)\abs{Dv(\xi)}\abs{w(\xi)}\right). \]
	Allora,
	\begin{align*}
		0 &\le \zeta^4(\xi)\left(-\sum_{k,l=1}^n a_{kl}(\xi) D^2_{l,k}w(\xi) +\sum_{k=1}^n \bar{b}_k(\xi) D_k w(\xi)\right) +R_5(\xi) \le \\
		&\le \zeta^4(\xi)\left( -\frac{\nu^2}{2}\abs{D^2 v(\xi)}^2 + C\abs{Dv(\xi)}^2\right) + C \left( \zeta^2(\xi)\abs{w(\xi)} + \zeta^3(\xi)\abs{Dv(\xi)}\abs{w(\xi)}\right)
	\end{align*}
	da cui, ricordando che $\abs{Dv}^2 \le \abs{w}$ e osservando che, per \eqref{eq1Harnack}, $\abs{w}\le \abs{D^2v}$ ,
	\begin{align*}
		\frac{\nu^2}{2}\zeta^4(\xi) \abs{D^2v(\xi)}^2 &\le C\zeta^4(\xi) \abs{Dv(\xi)}^2 + C\zeta^2(\xi)\abs{w(\xi)}+\zeta^3(\xi)\abs{Dv(\xi)}\abs{w(\xi)} \le \\
		&\le C\zeta^4(\xi) \abs{w(\xi)} + C\zeta^2(\xi)\abs{w(\xi)}+\zeta^3(\xi)\abs{Dv(\xi)}\abs{D^2v(\xi)}.
	\end{align*}
	Ora, per la disuguaglianza di Young pesata con $\varepsilon = \frac{\nu^2}{4}\zeta(\xi)$,
	\begin{align*}
		\zeta^3(\xi)\abs{Dv(\xi)}\abs{D^2v(\xi)} &\le \frac{\nu^2}{4}\zeta^4(\xi) \abs{D^2v(\xi)}^2 + C\zeta^3(\xi)\abs{Dv(\xi)}^2 \le \\
		&\le \frac{\nu^2}{4}\zeta^4(\xi) \abs{D^2v(\xi)}^2 + C\zeta^3(\xi)\abs{w(\xi)},
	\end{align*}
	quindi
	\[ \frac{\nu^2}{4}\zeta^4(\xi) \abs{D^2v(\xi)}^2 \le C\zeta^2(\xi)\abs{w(\xi)} \le C\abs{w(\xi)}. \]
	Allora, essendo $\abs{w(\xi)}^2\le C\abs{D^2v}^2$,
	\[ \zeta^4(\xi)\abs{w(\xi)}^2 \le C\abs{w(\xi)}. \]
	In definitiva, 
	\[ \underset{V}{\max} \; z \le C\]
	ed essendo $\zeta = 1$ su $V$ possiamo concludere che $\norm{w}_{L^\infty(V)} < +\infty$, ossia $w \in L^\infty(V)$. 
	
	Allora, anche $Dv \in L^\infty(V)$. A questo punto, dati $x,y \in V$, per la disuguaglianza di Lagrange\footnote{Qui si usa la connessione.}, esiste $\eta$ nel segmento congiungente $x,y$ tale che
	\[ \frac{u(x)}{u(y)} = e^{v(x)-v(y)} \le e^{\abs{x-y}\abs{Dv(\xi)}} \le e^{\textup{diam}(V)\norm{Dv}_{L^\infty(V)}} < +\infty.  \]
	Possiamo quindi affermare che per ogni $x,y \in V$
	\[ u(x) \le C u(y) \]
	ossia 
	\[ \underset{V}{\textup{sup}}\; u \le Cu(y) \]
	e quindi 
	\[ \underset{V}{\textup{sup}}\; u \le C  \underset{V}{\textup{inf}}\; u.\]
	
	Per quanto riguarda il caso generale, sia $\varepsilon >0 $ tale che 
	\[ \tilde{u} = u + \varepsilon >0. \]
	Applicando il caso precedente a $\tilde{u}$ e poi passando al limite per $\varepsilon \to 0$ si ha la tesi.
\end{proof}
\begin{osservazione}
	La disuguaglianza di Harnack mostra che i valori di soluzioni non negative di problemi ellittici, in una regione lontana dal bordo, sono comparabili: per ogni $x,y \in V$ tale che $\overline{V} \subseteq U$, 
	\[ \frac{1}{C} u(y) \le u(x) \le Cu(y). \]
	In particolare, se esiste $x \in V$ tale che $u(x) >0$, allora $u> 0$ in $V$.
\end{osservazione}
\section{Problemi agli autovalori}
I risultati di questa sezione sono gli analoghi nel caso di problemi ellittici di alcune affermazioni classiche di Algebra Lineare riguardo matrici simmetriche reali.

Nel corso di questa sezione consideriamo operatori in forma di divergenza del tipo 
\[ Lu = -\sum_{i,j=1}^{n}D_{j}(a_{ij}D_iu)\]
con $a_{ij} \in C^\infty(\overline{U})$. Oltre alle usuali ipotesi di struttura, considereremo anche $U$ connesso.
\begin{teorema}
	Valgono i seguenti fatti:
	\begin{enumerate}[$(a)$]
		\item ogni autovalore di $L$ è reale,
		\item contando ogni autovalore con la relativa molteplicità (finita), 
		\[ \Sigma = \left\lbrace \lambda_k : k \in \N\setminus\left\lbrace 0\right\rbrace \right\rbrace  \]
		dove 
		\[ 0 < \lambda_1 \le \lambda_2 \le \dots \]
		e $\lambda_k \to +\infty$ se $k \to +\infty$,
		\item esiste una base hilbertiana $(w_k)$ di $L^2(U)$ costituita da autovettori $w_{k} \in H^1_0(U)$ relativi agli autovalori $\lambda_k$, ossia per ogni $k \in \N\setminus\left\lbrace 0\right\rbrace$
		\[ \begin{cases}
			Lw_k = \lambda_k w_k & \textup{in } U,\\
			w_k = 0 & \textup{su } \partial U.
		\end{cases} \]
	\end{enumerate}
\end{teorema}
\begin{proof}
	In modo analogo quanto visto nei Teoremi di esistenza, si dimostra che
	\[ S = L^{-1} : L^2(U) \to L^2(U)\]
	è un operatore compatto.
	
	Mostriamo che $S$ è simmetrico. Siano $f,g \in L^2(U)$. Allora esistono e sono uniche $u,v \in H^1_0(U)$ soluzioni deboli, rispettivamente, di 
	\begin{align*}
		&\begin{cases}
			Lu = f & \textup{in } U,\\
			u = 0 & \textup{su } \partial U,
		\end{cases} & &\begin{cases}
			Lu = g & \textup{in } U,\\
			u = 0 & \textup{su } \partial U,
		\end{cases} 
	\end{align*}
	ossia $Sf = u$ e $Sg = v$. Per la formulazione debole, 
	\[ (Sf|g)_2 = (u|g)_2 = B(v,u) \]
	\[ (f|Sg)_2 = (f|v)_2 = B(u,v) \]
	e per simmetria di $B$, si ha $(Sf|g)_2 =(f|Sg)_2$, da cui $S$ è simmetrico. Inoltre, per ogni $f \in L^2(U)$
	\[ (Sf|f)_2 = (u|f)_2 = B(u,u) \ge 0. \]
	Dal Teorema Spettrale per operatori compatti, segue dunque che gli autovalori di $S$ sono tutti reali, strettamente positivi e le corrispondenti autofunzioni formano una base hilbertiana di $L^2(U)$. 
	
	Dal fatto che per $\eta \ne 0$, si ha che 
	\[ Sw = \eta w \Longleftrightarrow Lw = \lambda w \]
	con $\lambda = \frac{1}{\eta}$, segue la tesi.
\end{proof}
\begin{definizione}
	Chiamiamo $\lambda_1$ \emph{autovalore principale}\index[analitico]{autovalore!principale} o \emph{primo autovalore} di $L$.
\end{definizione}

\begin{lemma}
	Sia $H$ uno spazio di Hilbert ed $M \subseteq H$ non vuoto. Allora $\overline{\textup{span}(M)} = H$ se e solo se $M^\perp = \left\lbrace 0\right\rbrace $.
\end{lemma}
\begin{proof}
	Omettiamo la dimostrazione.
\end{proof}

\begin{teorema}\textbf{\emph{(principio variazionale per l'autovalore principale)}}\index[analitico]{principio! variazionale! per l'autovalore principale}
	Valgono i seguenti fatti:
	\begin{enumerate}[$(a)$]
		\item si ha che
		\[ \lambda_1 = \min \Set{B(u,u): u \in H^1_0(U), \norm{u}_{L^2(U)}=1}, \]
		\item il minimo è raggiunto in $w_1>0$ tale che 
		\[ \begin{cases}
			Lw_1 = \lambda_1 w_1 & \textup{in } U,\\
			w_1 = 0 & \textup{su } \partial U,
		\end{cases} \]
		\item se $u \in H^1_0(U)$ è una soluzione debole di 
		\[ \begin{cases}
			Lu = \lambda_1 u & \textup{in } U,\\
			u = 0 & \textup{su } \partial U,
		\end{cases} \]
		allora esiste $k \in \R$ tale che $u = k w_1$.
	\end{enumerate}
\end{teorema}
\begin{proof}
	Per la $(c)$ del Teorema precedente e la formulazione variazionale, per ogni $k,l \in \N\setminus\left\lbrace 0\right\rbrace $, 
	\[ B(w_k,w_l) = \lambda_k \delta_k^l \]
	e quindi 
	\[ B\left( \frac{w_k}{\sqrt{\lambda_k}},\frac{w_l}{\sqrt{\lambda_l}}\right)  =\delta_k^l. \]
	Inoltre, se $u \in H^1_0(U)$ e $\norm{u}_{L^2(U)}=1$, posto $d_k = (u|w_k)_{L^2(U)}$, si ha
	\begin{equation}\label{serie1primoautovalore}
		u = \sum_{k=1}^{\infty}d_k w_k\footnote{La serie converge in $L^2(U)$.}
	\end{equation}
	e, per la disuguaglianza di Bessel--Parseval,
	\[ \sum_{k=1}^{\infty} d_k^2 = \norm{u}_{L^2(U)}=1. \]
	
	Consideriamo $H^1_0(U)$ munito del prodotto scalare $B(\;,\;)$. Mostriamo che la norma indotta $\norm{\;}_B$ è equivalente alla norma $\norm{\;}_{H^1_0(U)}$. Innanzitutto, per ogni $u \in H^1_0(U)$, per uniforme ellitticità
	\[ \norm{u}_B^2 = \sum_{i,j=1}^n \int_{U} a_{ij}D_iuD_ju d\mathcal{L}^n \ge \nu \int_{U} \abs{Du}^2d\mathcal{L}^n = \nu \norm{u}^2_{H^1_0(U)}. \]
	Viceversa, 
	\[ \norm{u}^2_{B} \le C \int_{U} \abs{Du}^2d\mathcal{L}^n = C\norm{u}^2_{H^1_0(U)}. \]
	In definitiva, anche $\left( H^1_0(U),\norm{\;}_B\right)  $ è uno spazio di Hilbert su $\R$.
	
	Per quanto osservato inizialmente, $\left( \frac{w_k}{\sqrt{\lambda_k}}\right) $ è un sottoinsieme ortonormale di $H^1_0(U)$. Mostriamo che è una base hilbertiana di $H^1_0(U)$ con il prodotto scalare sopra introdotto. Questo fatto è equivalente ad affermare che 
	\[ \forall\; u \in H^1_0(U): \bigl( B(w_k,u) = 0 \textup{ per ogni } k \in \N\setminus \left\lbrace 0\right\rbrace \bigr) \Longrightarrow u = 0.\footnote{Sia $u \in H^1_0(U)$ tale che per ogni $k \in \N\setminus \left\lbrace 0\right\rbrace$, $B(w_k,u) = 0$. Se $\left( \frac{w_k}{\sqrt{\lambda_k}}\right) $ è una base hilbertiana,
		\begin{align*}
			u &= \sum_{k=1}^\infty l_k \frac{w_k}{\sqrt{\lambda_k}} & l_k &=  B\left( u,\frac{w_k}{\sqrt{\lambda_k}}\right) =\frac{1}{\sqrt{\lambda_k}} B(w_k,u)= 0,
		\end{align*}
		allora $u= 0$. Viceversa, sia $M = \left\lbrace \frac{w_k}{\sqrt{\lambda_k}} : k \in \N\setminus \left\lbrace 0\right\rbrace\right\rbrace  .$ Per ipotesi, $M^\perp = \left\lbrace 0\right\rbrace $, allora per il Lemma precedente si ha $\overline{\textup{span}(M)} = H^1_0(U)$. Se $u \in H^1_0(U)$, allora esiste $(u_h)$ in $\textup{span}(M)$ tale che $u_h \to u$ in $H^1_0(U)$. Deve essere 
		\[ u_h = \sum_{k=1}^{\infty} a_k^{(h)}\frac{w_k}{\sqrt{\lambda_k}} \]
		dove solo un numero finito di $\alpha_k^{(h)}$ è non nullo. Siccome $B\left( u_h,\frac{w_l}{\sqrt{\lambda_l}}\right)  = \alpha_l^{(h)}$, sfruttando la continuità del prodotto scalare, $B\left( u,\frac{w_l}{\sqrt{\lambda_l}}\right)  = B\left( \lim\limits_{h} u_h,\frac{w_l}{\sqrt{\lambda_l}}\right)  =  \lim\limits_{h} \alpha_l^{(h)}$ e per il Teorema della convergenza dominata (versione discreta), si ha $u = \lim\limits_{h} u_h = \sum\limits_{k=1}^{\infty} B\left( u,\frac{w_k}{\sqrt{\lambda_k}}\right) \frac{w_k}{\sqrt{\lambda_k}}$.} \]
	Ma quest'ultima affermazione è vera in quanto, essendo $(w_k)$ una base hilbertiana di $L^2(U)$, allora se per ogni $k \in \N\setminus \left\lbrace 0\right\rbrace$, 
	\[ B(w_k,u) = \lambda_k (w_k | u)_{L^2(U)} = 0 \]
	allora $u = 0$. 
	
	Di conseguenza, 
	\[ u = \sum_{k=1}^{\infty} \mu_k \frac{w_k}{\sqrt{\lambda_k}}\footnote{La serie converge in $H^1_0(U)$.} \]
	dove $\mu_k = B\left( u,\frac{w_k}{\sqrt{\lambda_k}}\right) $.
	
	Ora, siccome $\mu_k = d_k \sqrt{\lambda_k}$, la serie \eqref{serie1primoautovalore} converge anche in $H^1_0(U)$. Allora, sfruttando la continuità del prodotto scalare,
	\[ B(u,u) = B\left( \sum_{k=1}^{\infty} d_k w_k , \sum_{i=1}^{\infty} d_i w_i\right)  = \lim\limits_{m} \sum_{k,i=1}^{m} d_kd_iB(w_k,w_i) = \lim\limits_{m} \sum_{k=1}^{m} d_k^2\lambda_k \ge \lambda_1 \sum_{k=1}^{\infty} d_k^2 = \lambda_1. \]
	Siccome $B(w_1,w_1) = \lambda_1$, allora $w_1$ realizza il precedente minimo. 
	
	Mostriamo che se $u \in H^1_0(U)$ e $\norm{u}_{L^2(U)}=1$, $u$ è soluzione debole di 
	\[ \begin{cases}
		Lu = \lambda_1 u & \textup{in } U,\\
		u = 0 & \textup{su } \partial U,
	\end{cases} \]
	se e solo se $B(u,u) = \lambda_1$. Se $u$ è soluzione debole del problema precedente, per ogni $v \in H^1_0(U)$, $B(u,v) = \lambda_1 (u|v)_{L^2(U)}$. Scelto $v= u$, si ha $B(u,u) = \lambda_1 \norm{u}^2_{L^2(U)} = \lambda_1$. Viceversa,
	\[ \sum_{k=1}^\infty d^2_k \lambda_1 = \lambda_1 = B(u,u) = \sum_{k=1}^{\infty} d_k^2 \lambda_k \]
	allora
	\[ \sum_{k=1}^{\infty} (\lambda_k - \lambda_1) d_k^2 = 0 \]
	e di conseguenza deve essere $d_k = 0$ se $\lambda_k > \lambda_1$. Siccome $\lambda_1$ ha molteplicità finita, 
	\[ u = \sum_{k=1}^{m} (u|w_k)_{L^2(U)} w_k \]
	per qualche $m$, dove i $w_k$ sono autovettori relativi all'autovalore $\lambda_1$. Allora
	\[ Lu = \sum_{k=1}^{m} (u|w_k)_{L^2(U)} Lw_k = \lambda_1 u. \] 
	
	Mostriamo che se $u \in H^1_0(U)$ è una soluzione debole del problema precedentemente introdotto allora $u = 0$ oppure ha segno definito\footnote{Si intende che è strettamente positiva oppure strettamente negativa in $U$.}. Per fare ciò, supponiamo $u$ non nulla, allora senza perdita di generalità possiamo prenderla tale che $\norm{u}_{L^2(U)} = 1$. Posto,
	\begin{align*}
		\alpha &= \int_{U} (u^+)^2 d\mathcal{L}^n, & \beta &= \int_{U} (u^-)^2 d\mathcal{L}^n
	\end{align*}
	allora $\alpha + \beta = 1$. Osserviamo che $u^{\pm} \in H^1_0(U)$ con 
	\begin{align*}
		Du^+ &= \begin{cases}
			Du & \textup{q.o. su } \left\lbrace u \ge 0 \right\rbrace ,\\
			0 & \textup{q.o. su } \left\lbrace u \le 0 \right\rbrace ,
		\end{cases}
		&
		Du^- &= \begin{cases}
			0 & \textup{q.o. su } \left\lbrace u \ge 0 \right\rbrace ,\\
			-Du & \textup{q.o. su } \left\lbrace u \le 0 \right\rbrace.
		\end{cases}
	\end{align*}
	e 
	\begin{align*}
		B(u^+,u^-) &= \sum_{i,j=1}^n \int_U a_{ij} D_i u^+ D_j u^- d\mathcal{L}^n = \\
		& =\sum_{i,j=1}^n \int_{U \cap \left\lbrace u \ge 0 \right\rbrace} a_{ij} D_i u^+ D_j u^- d\mathcal{L}^n + \sum_{i,j=1}^n \int_{U \cap \left\lbrace u \le 0 \right\rbrace} a_{ij} D_i u^+ D_j u^- d\mathcal{L}^n = 0.
	\end{align*}
	Ora,
	\begin{align*}
		\lambda_1 &= B(u,u) = B(u^+,u^+) + B(u^-,u^-) + 2B(u^+,u^-) \ge \\
		&\ge \lambda_1 \norm{u^+}^2_{L^2(U)} + \lambda_1\norm{u^-}_{L^2(U)} = (\alpha + \beta) \lambda_1 = \lambda_1,
	\end{align*}
	allora deve essere
	\begin{align*}
		B(u^+,u^+) &= \lambda_1 \norm{u^+}^{2}_{L^2(U)}, & B(u^-,u^-) &= \lambda_1 \norm{u^-}^{2}_{L^2(U)}.
	\end{align*}
	Per il passaggio precedente quindi $u^+$ e $u^-$ sono soluzioni deboli dei problemi
	\begin{align*}
		&\begin{cases}
			Lu^+ = \lambda_1 u^+ & \textup{in } U,\\
			u^+ = 0 & \textup{su } \partial U,
		\end{cases} & &\begin{cases}
			Lu^- = \lambda_1 u^- & \textup{in } U,\\
			u^- = 0 & \textup{su } \partial U.
		\end{cases}
	\end{align*} 
	Ma essendo gli $a_{ij} \in C^\infty(U)$, $u^+ \in C^\infty(U)$ e 
	\[ Lu^+ = \lambda_1 u^+ \ge 0 \textup{ in } U. \]
	In definitiva, $u^+$ è una supersoluzione. Per il Principio del massimo forte quindi $u^+ >0$ in $U$ oppure $u^+ = 0$ in $U$. Un argomento analogo vale per $u^-$.
	
	Sia ora $u$ una soluzione debole non nulla del problema introdotto. Per quanto visto sopra, 
	\[ \int_U w_1 d\mathcal{L}^n \ne 0 \]
	e quindi esiste $k \in \R$ tale che 
	\[ \int_U \left( u-kw_1\right) d\mathcal{L}^n = 0. \]
	Ma 
	\[ L\left( u-kw_1\right)  = Lu -kLw_1 = \lambda_1 u - k Lw_1 = \lambda_1 \left( u-kw_1\right)  \]
	quindi anche $u-kw_1$ è soluzione del problema. Ma allora deve essere $u = kw_1$ in $U$.
\end{proof}
\begin{osservazione}
	L'affermazione $(c)$ del Principio variazionale per l'autovalore principale corrisponde a dire che $\lambda_1$ è \emph{semplice}. In particolare,
	\[ 0<\lambda_1 < \lambda_2 \le \lambda_3 \le \dots \]
\end{osservazione}
\begin{osservazione}
	L'affermazione $(a)$ del Principio variazionale per l'autovalore principale è una caratterizzazione di $\lambda_1$ attraverso un problema di minimo vincolato ed è nota come \emph{formula di Rayleigh}. Si può dimostrare, grazie all'omogeneità di $B(u,u)$, che è equivalente al problema di minimo libero
	\[ \lambda_1 = \underset{\underset{u \ne 0}{u \in H^1_0(U)}}{\min}\; \frac{B(u,u)}{\norm{u}_{L^2(U)}^2}. \] 
	Osserviamo infatti che 
	\[ \frac{B(u,u)}{\norm{u}_{L^2(U)}^2} = B\left( \frac{u}{\norm{u}_{L^2(U)}},\frac{u}{\norm{u}_{L^2(U)}}\right) \]
	e  
	\[ \norm*{\frac{u}{\norm{u}_{L^2(U)}}}_{L^2(U)}=1, \]
	allora
	\[ \lambda_1 \le \frac{B(u,u)}{\norm{u}_{L^2(U)}^2} \]
	e quindi 
	\[ \lambda_1 \le \underset{\underset{u \ne 0}{u \in H^1_0(U)}}{\min}\; \frac{B(u,u)}{\norm{u}_{L^2(U)}^2}. \]
	Viceversa, se $u \in H^1_0(U)$ e $\norm{u}_{L^2(U)} = 1$, allora
	\[ B(u,u) = \frac{B(u,u)}{\norm{u}^2_{L^2(U)}} \ge \underset{\underset{u \ne 0}{u \in H^1_0(U)}}{\min}\; \frac{B(u,u)}{\norm{u}_{L^2(U)}^2} \]
	da cui 
	\[ \lambda_1 \ge \underset{\underset{u \ne 0}{u \in H^1_0(U)}}{\min}\; \frac{B(u,u)}{\norm{u}_{L^2(U)}^2}.\]
\end{osservazione}

\begin{teorema}
	Siano $X$ uno spazio di Hilbert su $\K$ separabile e $S: X \to X$ un operatore lineare compatto e autoaggiunto. Allora esiste una base hilbertiana di $X$ composta da autovettori di $S$.
\end{teorema}
\begin{proof}
	Omettiamo la dimostrazione.
\end{proof}

Consideriamo l'operatore $S: L^2(U) \to L^2(U)$ tale che 
\[ Sf= u, \]
dove $u \in H^1_0(U)$ è l'unica soluzione debole del problema di Dirichlet
\[ \begin{cases}
	-\Delta u = f & \textup{in } U, \\
	u = 0 & \textup{su } \partial U.
\end{cases} \footnote{In altre parole $S$ è l'operatore inverso di $-\Delta$ (debole).}\]
In particolare, $\textup{Im}(S) \subseteq H^1_0(U)$. Chiaramente $S$ è lineare e per ogni $f,g \in L^2(U)$, dette $u = Sf$ e $v = Sg$,
\[ (Sf|g)_{L^2(U)} = (u|g)_{L^2(U)} = \int_U Du \cdot Dv d\mathcal{L}^n = (f|v)_{L^2(U)} = (f|Sg)_{L^2(U)},\]
quindi è pure autoaggiunto. Sia ora $(f_j)$ in $L^2(U)$ limitata. Dette $u_j = Sf_j \in H^1_0(U)$, dalla formulazione variazionale combinata con la disuguaglianza di Holder e la disuguaglianza di Poincaré,
\[ \norm{u_j}^2_{H^1_0(U)} = \int_U \abs{Du_j}^2 d\mathcal{L}^n = \int_U f_j u_j d\mathcal{L}^n \le \norm{f_j}_{L^2(U)} \norm{u_j}_{L^2(U)} \le C\norm{f_j}_{L^2(U)}\norm{u_j}_{H^1_0(U)},\]
quindi $(u_j)$ limitata in $H^1_0(U)$. Per il Teorema di Rellich--Kondrachov, esistono $(u_{j_k})$ e $u \in H^1_0(U)$ tale che $u_{j_k} \to u$ in $L^2(U)$. In particolare, $S$ è compatto. Allora, per il Teorema precedente, esiste $(w_k)$ base hilbertiana di $L^2(U)$ tale che 
\[ Sw_k = \lambda_k w_k. \]
In particolare, $(w_k)$ in $H^1_0(U)$ e
\[ -\Delta w_k = \frac{1}{\lambda_k} w_k \]
in senso debole, quindi
\[ \int_U Dw_h \cdot Dw_k d\mathcal{L}^n = \frac{1}{\lambda_k}  (w_h|w_k)_{L^2(U)},\]
segue che $(w_k)$ è una base hilbertiana di $H^1_0(U)$ costituita da autovettori di $-\Delta$. In particolare, da quanto visto sulla teoria della regolarità, le $w_k$ sono lisce.
		
		% Commenti speciali

% !TEX encoding = UTF-8 Unicode
% !TEX TS-program = pdflatex
% !TEX spellcheck = it-IT
% !TEX root = Dispensa/afed.tex

% Capitolo 4
\chapter{Alcuni complementi di calcolo delle variazioni}\label{Alcuni complementi di calcolo delle variazioni}
Benvenuti nel magico mondo del Calcolo delle Variazioni! Prima di proseguire si consiglia la lettura di
\begin{center}
	\href{https://www.dmf.unicatt.it/~musesti/FM/}{FISICA MATEMATICA\\
		appunti a cura di Alessandro Musesti}
\end{center}
non ve ne pentirete! :-)

Nel seguito, salvo diversa specificazione, $U$ indicherà un generico sottoinsieme aperto limitato di $\R^n$ con $\partial U$ di classe $C^1$. 
\section{Variazione prima e variazione seconda}
\begin{definizione}
	Chiamiamo \emph{lagrangiana}\index[analitico]{lagrangiana} ogni applicazione $L: \overline{U} \times \R \times \R^n \to \R$ di classe $C^\infty$ tale che
	\[ L(x,s,\xi) = L(x_1,\dots,x_n, s, \xi_1,\dots,\xi_n) \]
\end{definizione}
Nel seguito considereremo funzionali della forma
\[ I[u] = \int_U L(x,u(x),Du(x))d\mathcal{L}^n(x) \]
definiti inizialmente per $u \in C^\infty(\overline{U})$ tali che $u= g \textup{ su } \partial U,$ con $g$ un'applicazione nota. Il funzionale $I$ è anche detto \emph{energia} in quanto molto spesso può essere interpretato come energia meccanica di un particolare sistema fisico.
\begin{definizione}
	Sia $v \in C^\infty_c(U)$. Consideriamo l'applicazione $i: \R \to \R$ tale che
	\[  i(\tau) = I[u+\tau v].\]
	Chiamiamo \emph{variazione prima}\index[analitico]{variazione!prima} di $I$ la derivata prima di $i$, $i'(\tau)$, e \emph{variazione seconda}\index[analitico]{variazione!seconda} di $I$ la derivata seconda di $i$, $i''(\tau)$.
\end{definizione}
\begin{proposizione}\emph{\textbf{(condizioni necessarie per l'esistenza di minimi)}}
	Se $I$ ammette minimo in $u$, allora $i$ ammette minimo in $0$. In particolare, 
	\begin{align*}
		i'(0) &= 0, & i''(0)\ge 0.
	\end{align*}
	
	\begin{proof}
		Innanzitutto, per ogni $v \in C^\infty_c(U), \tau \in \R$, 
		\[ u + \tau v = g \textup{ su } \partial U \]
		in quanto $I$ ammette minimo in $u$. Allora, per ogni $v \in C^\infty_c(U), \tau \in \R$
		\[ I[u] \le I[u + \tau v] \]
		ossia $i(0) \le i(\tau)$ per ogni $\tau \in \R$, da cui la tesi.
	\end{proof}
\end{proposizione}
Calcoliamo esplicitamente in queste iniziali condizioni di regolarità la variazione prima di $I$ in un suo minimo $u$. 

Mostriamo innanzitutto che l'applicazione $f: \R \times \overline{U} \to \R $ tale che
\[ f(\tau, x) = L(x,u(x)+\tau v(x),Du(x)+\tau Dv(x)) \]
soddisfa le ipotesi del Teorema di derivazione sotto il segno di integrale. Innanzitutto, per ogni $\tau \in \R$, essendo $L$ di classe $C^\infty$, $u \in C^\infty(\overline{U})$ e $v \in C^\infty_c(U)$, allora sicuramente l'applicazione ${x \mapsto f(\tau,x)}$ è almeno continua in $\overline{U}$. In particolare è sommabile. Poi per ogni $x \in \overline{U}$, sempre per il fatto che $L$ è di classe $C^\infty$, allora l'applicazione ${\tau \mapsto f(\tau,x)}$ è derivabile. Ora,
\begin{align*}
	\frac{\partial f}{\partial \tau}(\tau,x) &= D_sL(x,u(x)+\tau v(x),Du(x)+ \tau Dv(x))v(x) + \\
	&+D_{\xi}L(x,u(x)+\tau v(x),Du(x)+\tau Dv(x))\cdot Dv(x)
\end{align*}
e non essendo restrittivo supporre $\tau \in [-1,1]$ si ha 
\[ \abs*{\frac{\partial f}{\partial \tau}(\tau,x)} \le A\norm{Dv}_\infty + B\norm{v}_\infty \]
dove 
\[ A = \max\Set{\abs{D_\xi L(x,s,\xi)} : \abs{\xi} \le \norm{Du}_\infty + \norm{Dv}_\infty, \abs{s} \le \norm{u}_\infty + \norm{v}_\infty, x \in \overline{U}} \]
e 
\[ B = \max\Set{\abs{D_s L(x,s,\xi)} : \abs{\xi} \le \norm{Du}_\infty + \norm{Dv}_\infty, \abs{s} \le \norm{u}_\infty + \norm{v}_\infty, x \in \overline{U}}. \]
Allora, detto $g = \left( A\norm{Dv}_\infty + B\norm{v}_\infty\right)\chi_{\overline{U}} : \overline{U} \to \R$, si ha per ogni $x \in \overline{U}$
\[ \abs*{\frac{\partial f}{\partial \tau}(\tau,x)} \le g(x) \]
con $g$ sommabile. 

Per il Teorema di derivazione sotto il segno di integrale quindi 
\begin{align*}
	i'(\tau) &=\int_U \Bigl( D_sL(x,u(x)+\tau v(x),Du(x)+\tau Dv(x))v(x) + \\
	& + D_{\xi}L(x,u(x)+\tau v(x),Du(x)+\tau Dv(x))\cdot Dv(x)\Bigr) d\mathcal{L}^n(x) 
\end{align*}
e posto $\tau = 0$ otteniamo 
\[ \int_U \Bigl( D_sL(x,u(x),Du(x))v(x) + D_{\xi}L(x,u(x),Du(x))\cdot Dv(x)\Bigr) d\mathcal{L}^n(x) = 0. \]
A questo punto, siccome $v$ ha supporto compatto, per la formula di Gauss--Green
\[ \int_U \Bigl( D_sL(x,u(x),Du(x)) -\sum_{i=1}^{n} D_{x_i}\bigl( D_{\xi_i}L(x,u(x),Du(x))\bigr) \Bigr)v(x) d\mathcal{L}^n(x) = 0 \]
ed il Lemma di Du Bois--Reymond si conclude che 
\[  D_sL(x,u(x),Du(x)) -\sum_{i=1}^{n} D_{x_i}\bigl( D_{\xi_i}L(x,u(x),Du(x))\bigr) = 0 \textup{ in } U. \]
\begin{definizione}
	Chiamiamo \emph{equazione di Eulero--Lagrange}\index[analitico]{equazione!di Eulero--Lagrange} associata al funzionale $I$, l'equazione differenziale alle derivate parziali quasilineare del secondo ordine in forma di divergenza
	\[ -\sum_{i=1}^{n} D_{x_i}\bigl( D_{\xi_i}L(x,u(x),Du(x))\bigr)  + D_sL(x,u(x),Du(x)) = 0. \]
\end{definizione}
Se $u$ è un minimo di $I$, allora risolve l'equazione di Eulero--Lagrange associata ad $I$. Molto spesso, si cerca proprio di ricondurre una PDE ad un funzionale in quanto è solitamente estremamente più semplice trovare i punti critici di un funzionale rispetto che risolvere direttamente l'equazione.
\begin{esempio}
	Consideriamo
	\[ L(x,s,\xi) = \frac{1}{2}\sum_{i,j=1}^{n}a_{ij}(x)\xi_i\xi_j - sf(x) \]
	con $a_{ij}(x) = a_{ji}(x)$ per ogni $i,j=1,\dots,n$. Allora, siccome per ogni $i,j=1,\dots,n$
	\begin{align*}
		D_{\xi_i} L &= \sum_{j=1}^n a_{ij}\xi_j, & D_s L = -f,
	\end{align*}
	l'equazione di Eulero--Lagrange associata al funzionale 
	\[ I[u] = \int_U  \left( \frac{1}{2}\sum_{i,j=1}^{n}a_{ij}(x)\xi_i\xi_j - sf(x)\right) d\mathcal{L}^n(x)\]
	è 
	\[ -\sum_{i,j=1}^n D_{x_i}\left( a_{ij}(x)D_{x_j}u(x)\right)  = f(x) \textup{ in } U, \]
	ossia un'equazione alle derivate parziali lineare del secondo ordine in forma di divergenza.
\end{esempio}

Non è detto, almeno a priori, che un minimo sia una soluzione in senso classico.
\begin{esempio}
	Consideriamo il problema
	\[ 	\begin{cases}
		-\Delta u = 0 & \textup{in } U,\\
		u= g & \textup{su } \partial U.
	\end{cases} \]
	Il funzionale associato è 
	\[ I[u] = \frac{1}{2}\int_U \abs{Du}^2 d\mathcal{L}^n. \]
	Se fossimo alla ricerca di soluzioni classiche, dovremmo cercare di dimostrare l'esistenza di minimi per $I$ su
	\[ C^2_g(U) = \left\lbrace u \in C^2(U) \cap C(\overline{U}): u = g \textup{ su } \partial U\right\rbrace. \]
	Provando ad applicare il Metodo diretto, prendiamo $(u_h)$ in $C^2_g(U)$ tale che 
	\[ I[u_h] \le \underset{u \in  C^2_g(U)}{\min}\;I[u] + \frac{1}{h+1}. \]
	Per andare avanti avremmo bisogno di un risultato di compattezza, tuttavia nel dominio dove stiamo lavorando non lo abbiamo. Possiamo esibire il seguente controesempio: siano $U = \left] -1,1\right[ $, $g = 1$ e 
	\[ u_h(x) = \sqrt{\frac{1}{h} +x^2} - \left( \sqrt{\frac{1}{h} +1}-1\right).  \]
	Chiaramente $u_h$ è di classe $C^\infty$ per ogni $h \in \N$. Inoltre,
	\[ I[u_h] = \frac{1}{2}\int_{-1}^{1} \frac{x^2}{\left( \frac{1}{h}+x^2\right) }dx \le \frac{1}{2} \int_{-1}^{1}dx = 1 < +\infty. \]
	Eppure, detto $\overline{u}(x) = \abs{x}$, da due razionalizzazioni inverse emerge che
	\[ \abs{u_h(x) - \overline{u}(x)} = \abs*{\left( \sqrt{ \frac{1}{h} + x^2} - \sqrt{x^2}\right) + \left( 1-  \sqrt{ \frac{1}{h} + 1}\right) }\le \frac{2}{\sqrt{h}}, \]
	quindi $u_h \to \overline{u}$ uniformemente, ma $\overline{u} \notin C^2_g(U)$.
\end{esempio}
\begin{esempio}\textbf{\emph{(superfici minime)}}
	Consideriamo il funzionale d'area 
	\[ I[u] = \int_U \sqrt{1 + \abs{Du}^2}d\mathcal{L}^n. \]
	L'equazione di Eulero--Lagrange associata è 
	\[ \sum_{i=1}^nD_{x_i}\left( \frac{D_{x_i} u}{\sqrt{1 + \abs{Du}^2}}\right)  = 0 \textup{ in } U \]
	ed è detta \emph{equazione delle superfici minime}. L'espressione a primo membro è $n$ volte la curvatura media del grafico di $u$, quindi le superfici minime sono caratterizzate da curvatura media nulla.
\end{esempio}

Calcoliamo esplicitamente in queste iniziali condizioni di regolarità anche la variazione seconda di $I$ in un suo minimo $u$. In modo analogo a quanto fatto per il calcolo della variazione prima, si verifica che è applicabile il Teorema di derivazione sotto il segno di integrale, quindi 
\begin{align*}
	i''(\tau) &= \int_U \Biggl( \sum_{i,j=1}^nD^2_{\xi_i,\xi_j}L(x,u(x)+\tau v(x), Du(x)+\tau Dv(x)) D_{x_i}v(x)D_{x_j}v(x) +  \\
	&+  2 \sum_{i=1}^nD^2_{\xi_i,s}L(x,u(x)+\tau v(x), Du(x)+\tau Dv(x)) D_{x_i} v(x) v(x) + \\
	&+ D^2_{s,s} L(x,u(x)+\tau v(x), Du(x)+\tau Dv(x)) v(x)^2\Biggr) d\mathcal{L}^n(x)
\end{align*}
e posto $\tau = 0$ si ha
\begin{align*}
	\int_U \Biggl( &\sum_{i,j=1}^nD^2_{\xi_i,\xi_j}L(x,u(x), Du(x)) D_{x_i}v(x)D_{x_j}v(x) + \\
	&+2 \sum_{i=1}^nD^2_{\xi_i,s}L(x,u(x), Du(x)) D_{x_i} v(x) v(x) + \\
	&+ D^2_{s,s} L(x,u(x), Du(x)) v(x)^2\Biggr) d\mathcal{L}^n(x) \ge 0
\end{align*}
per ogni $v \in C^\infty_c(U)$.

Cerchiamo di estrarre delle informazioni dalla disuguaglianza precedente. Innanzitutto, mediante approssimazione è possibile dimostrare la validità della disuguaglianza per $v$ lipschitziane che si annullano su $\partial U$. 

Consideriamo ora  $\varrho : \R \to \R$ l'applicazione tale che per ogni $x \in \R$ 
\[ \varrho(x+1) = \varrho(x) \]
e 
\[  \varrho(x) = \begin{cases}
	x & \textup{se } 0\le x \le \frac{1}{2}, \\
	1-x & \textup{se } \frac{1}{2}\le x \le 1. \\
\end{cases} \]
Evidentemente $\abs{\varrho'} = 1 $ q.o. in $\R$. Ora, dati $\varepsilon>0, \eta \in \R^n $, $\zeta \in C^\infty_c(U)$ definiamo per ogni $x \in U$
\[ v(x) = \varepsilon \varrho\left( \frac{x\cdot \eta}{\varepsilon}\right) \zeta(x) \]
e osserviamo che 
\[ D_{x_i} v(x) = \varrho'\left( \frac{x\cdot \eta}{\varepsilon}\right) \eta_i \zeta(x) + \varepsilon \varrho\left( \frac{x\cdot \eta}{\varepsilon}\right)  D_{x_i} \zeta(x). \]
Pertanto,
\[ 0 \le \int_U  \sum_{i,j=1}^n D^2_{\xi_i,\xi_j}L(x,u(x), Du(x)) \varrho'(x)^2\eta_i\eta_j \zeta(x)^2 d\mathcal{L}^n(x) + C_1 \varepsilon + C_2 \varepsilon^2  \]
e, se $\varepsilon \to 0^+$, ricordando che $\abs{\varrho'} = 1 $ q.o. in $\R$,
\[ 0 \le \int_U \sum_{i,j=1}^n D^2_{\xi_i,\xi_j}L(x,u(x), Du(x))\eta_i\eta_j \zeta(x)^2 d\mathcal{L}^n(x). \]
In definitiva\footnote{Si tratta di applicare una versione del Lemma di Du Bois--Reymond per le disuguaglianze.}, per ogni $\eta \in \R^n$ e per ogni $x \in U$
\[ \sum_{i,j=1}^n D^2_{\xi_i,\xi_j}L(x,u(x), Du(x))\eta_i\eta_j \ge 0,\]
da cui capiamo che sarà naturale supporre, nella prossima sezione, che $L$ sia convesso nei gradienti.

Prima di procedere, però, vediamo alcuni accenni al caso vettoriale.
\begin{definizione}
	Chiamiamo \emph{lagrangiana}\index[analitico]{lagrangiana} ogni applicazione $L: \overline{U} \times \R^m \times \textup{Mat}_{m,n}(\R) \to \R$ di classe $C^\infty$ tale che
	\[ L(x,s,\xi) = L(x_1,\dots, x_j,\dots,x_n, s_1,\dots,s_i,\dots,s_m, \xi_{11},\dots,\xi_{ij},\dots,\xi_{mn}) \]
\end{definizione}
Si considerano funzionali della forma
\[ I[u] = \int_U L(x,u(x),Du(x))d\mathcal{L}^n(x) \]
definiti per $u \in C^\infty(\overline{U};\R^m)$ tali che $u= g \textup{ su } \partial U,$ con $g: \partial U \to \R^m$ un'applicazione nota. Come nel caso scalare, il funzionale $I$ è anche detto \emph{energia} in quanto molto spesso può essere interpretato come energia meccanica di un particolare sistema fisico.
\begin{definizione}
	Sia $v \in C^\infty_c(U;\R^m)$. Consideriamo l'applicazione $i: \R \to \R$ tale che
	\[  i(\tau) = I[u+\tau v].\]
	Chiamiamo \emph{variazione prima}\index[analitico]{variazione!prima} di $I$ la derivata prima di $i$, $i'(\tau)$, e \emph{variazione seconda}\index[analitico]{variazione!seconda} di $I$ la derivata seconda di $i$, $i''(\tau)$.
\end{definizione}
\begin{proposizione}\emph{\textbf{(condizioni necessarie per l'esistenza di minimi)}}
	Se $I$ ammette minimo in $u$, allora $i$ ammette minimo in $0$. In particolare, 
	\begin{align*}
		i'(0) &= 0, & i''(0)\ge 0.
	\end{align*}
	
	\begin{proof}
		Si tratta di una semplice variante del caso scalare.
	\end{proof}
\end{proposizione}
In modo analogo a quanto fatto per il caso scalare, possiamo calcolare esplicitamente la variazione prima di $I$ in un suo minimo $u$. Da questo calcolo emerge, in tutta analogia, la seguente definizione.
\begin{definizione}
	Chiamiamo \emph{(sistema di) equazioni di Eulero--Lagrange}\index[analitico]{equazione!di Eulero--Lagrange} associate al funzionale $I$, il sistema di equazioni differenziali alle derivate parziali quasilineari del secondo ordine in forma di divergenza
	\[ \begin{cases}
		D_{s_1}L(x,u(x),Du(x))-\displaystyle \sum_{i=j}^{n} D_{x_j}\bigl( D_{\xi_{1j}}L(x,u(x),Du(x))\bigr) = 0 & \textup{in } U,\\
		& \vdots,\\
		D_{s_m}L(x,u(x),Du(x))-\displaystyle\sum_{j=1}^{n} D_{x_j}\bigl( D_{\xi_{mj}}L(x,u(x),Du(x))\bigr) = 0 & \textup{in } U.
	\end{cases}\]
\end{definizione}

\begin{definizione}
	Diciamo che $L$ è una \emph{lagrangiana nulla}\index[analitico]{lagrangiana!nulla} se ogni $u \in C^2(\overline{U};\R^m)$ è soluzione delle equazioni di Eulero--Lagrange.
\end{definizione}

\begin{teorema}
	Sia $L$ una lagrangiana nulla. Se $u,\overline{u}\in C^2(\overline{U};\R^m)$ tali che $u = \overline{u}$ su $\partial U$, allora $I[u] = I[\overline{u}]$.
\end{teorema}
\begin{proof}
	Consideriamo l'applicazione $j: [0,1] \to \R$ tale che	
	\[ j(\tau) = I[\tau u + (1-\tau)\overline{u}]. \]
	Allora, in modo analogo a quanto visto per il calcolo della variazione prima nel caso scalare, 
	\begin{align*}
		&j'(\tau) = \int_U \left( \sum_{i=1}^{m}D_{s_i}L(x,\tau u(x)+(1-\tau)\overline{u}(x), \tau Du(x) + (1-\tau)D\overline{u}(x))(u_i(x)-\overline{u}_i(x)) + \right. \\
		&\left. \hspace{-5pt}+ \hspace{-2pt}\sum_{j=1}^{n} \sum_{i=1}^{m} \hspace{-2pt} D_{\xi_{ij}}L(x,\tau u(x)\hspace{-2pt}+\hspace{-2pt}(1-\tau)\overline{u}(x), \tau Du(x)\hspace{-2pt} + \hspace{-2pt}(1-\tau)D\overline{u}(x)) D_{x_j}(u_i(x)-\overline{u}_i(x))\hspace{-2pt} \right) \hspace{-4pt}d\mathcal{L}^n(x)
	\end{align*}
	e, combinando il fatto che $u - \overline{u} = 0$ su $\partial U$ con le formule di Gauss--Green,
	\begin{align*}
		&j'(\tau) = \sum_{i=1}^{m}\int_U \biggl( D_{s_i}L(x,\tau u(x)+(1-\tau)\overline{u}(x), \tau Du(x) + (1-\tau)D\overline{u}(x)) + \\
		&\left.\hspace{-2pt} - \hspace{-2pt}\sum_{j=1}^{n}\hspace{-2pt}D_{x_j}\hspace{-2pt}\left( \hspace{-2pt} D_{\xi_{ij}}L(x,\tau u(x)+(1-\tau)\overline{u}(x), \tau Du(x) + (1-\tau)D\overline{u}(x))\hspace{-2pt}\right)  \right)\hspace{-3pt}(u_i(x)-\overline{u}_i(x)) \hspace{-1pt}d\mathcal{L}^n(x)
	\end{align*}
	Siccome $\tau u +(1-\tau)\overline{u} \in C^2(\overline{U};\R^m)$ e $L$ una lagrangiana nulla, segue che $j'(\tau) = 0$ per ogni $\tau \in [0,1]$. In particolare, $j$ è costante su $[0,1]$, da cui $I[u] = j(1) = j(0) = I[\overline{u}]$, ossia la tesi.
\end{proof}

Se $m=1$, ossia nel caso scalare, le uniche lagrangiane nulle sono le funzioni lineari in $\xi$. 

\begin{definizione}
	Sia $A \in \textup{Mat}_{n}(\R)$. Chiamiamo \emph{matrice dei cofattori} di $A$\index[analitico]{matrice!dei cofattori}\index[simboli]{$\textup{cof}(A)$} la matrice 
	\[ \textup{cof}A \]
	dove $\textup{cof}(A)_{ij} = (-1)^{i+j}\det(A^{(ij)})$ e $A^{(ij)} \in \textup{Mat}_{n-1}(\R)$ è ottenuta da $A$ eliminando la $i$-esima riga e la $j$-esima colonna.
\end{definizione}
\begin{lemma}
	Se $u:\R^n \to \R^n$ di classe $C^2$, allora per ogni $i=1,\dots, n$
	\[ \sum_{j=1}^{n} D_{x_j}(\textup{cof}(Du)_{ij}) = 0,\]
	ossia, in forma vettoriale, $\div (\textup{cof}(Du)) = 0$.
\end{lemma}
\begin{proof}
	Innanzitutto, per ogni $\xi \in \textup{Mat}_{n}(\R)$
	\[ (\det \xi)\textup{Id} = \xi^T\textup{cof}(\xi), \]
	ossia per ogni $i,j=1,\dots,n$
	\begin{equation}\label{serve}
		(\det \xi)\delta_{ij} = \sum_{k=1}^{n} \xi_{ki}\textup{cof}(\xi)_{kj}.
	\end{equation}
	Scelto $\xi = Du$ in \eqref{serve}, deriviamo ambo i membri per $x_j$ e sommando su $j$ otteniamo per ogni $i=1,\dots, n$
	\[ \sum_{j,p,q=1}^{n} \delta_{ij}\textup{cof}(Du)_{pq}D^2_{x_j,x_q}u_p = \sum_{k,j=1}^{n}\left(  D^2_{x_j,x_i}u_k \textup{cof}(Du)_{kj} + D_{x_i}u_k D_{x_j}(\textup{cof}(Du)_{kj}) \right) \]
	che può essere riscritta, scambiando $q$ con $j$ e $p$ con $k$ a primo membro, come 
	\[ \sum_{k=1}^n D_{x_i}u_k \left( \sum_{j=1}^{n}D_{x_j}\textup{cof}((Du)_{kj})\right) =0, \]
	ossia come un sistema lineare omogeneo che, scambiando $k$ con $i$, è 
	\[ \sum_{i=1}^n D_{x_k}u_i \left( \sum_{j=1}^{n}D_{x_j}(\textup{cof}(Du)_{ij})\right) =0. \]
	
	A questo punto, se $\det Du(x_0) \ne 0$, il sistema ha solo la soluzione banale, ossia per ogni  $k=1,\dots,n$
	\[ \sum_{j=1}^{n}D_{x_j}(\textup{cof}(Du(x_0))_{kj})=0. \]
	Se, invece, $\det Du(x_0) = 0$, sia $\varepsilon >0$ tale che $\det(Du(x_0) + \varepsilon\textup{Id}) \ne 0$. Applicando quanto fatto sopra a $u + \varepsilon x$ e mandando poi $\varepsilon \to 0^+$ segue la tesi.
\end{proof}
\begin{teorema}
	L'applicazione 
	\[ L(\xi) = \det(\xi) \]
	è una lagrangiana nulla.
\end{teorema}
\begin{proof}
	Per \eqref{serve}, per ogni $i,j = 1,\dots,n$
	\[ D_{ \xi_{ij}} \det (\xi) = (\textup{cof}\xi)_{ij}, \]
	da cui per ogni $u \in C^2(\overline{U};\R^n)$ e per ogni $i =1,\dots,n$
	\[ \sum_{j=1}^n D_{x_j}\left( D_{\xi_{ij}}L(Du)\right) =0, \]
	che è la tesi.
\end{proof}
\begin{osservazione}
	Grazie al fatto che il determinante è una lagrangiana nulla, possiamo dare una dimostrazione alternativa del Teorema di Brouwer rispetto a quella riportata nel Capitolo \ref{Teoremi di punto fisso}. 
	
	Vediamo innanzitutto il caso $K = \overline{\B(0,1)}$. Mostriamo che non esiste alcuna applicazione $f: \overline{\B(0,1)} \to \partial \B(0,1)$ di classe $C^2$ tale che per ogni $x \in \partial \B(0,1)$ si abbia $f(x) = x$. Supponiamo, per assurdo, che una tale applicazione esista. Allora, siccome $\det$ è una lagrangiana nulla,
	\[ \int_{\B(0,1)} \det(Df) d\mathcal{L}^n = \int_{\B(0,1)} \det(Dx)d\mathcal{L}^n(x) = \int_{\B(0,1)}d\mathcal{L}^n = \mathcal{L}^n(\B(0,1)) \ne 0.\]
	D'altra parte, $f(x) \cdot f(x) = \abs{f(x)}^2 = 1$ per ogni $x \in \overline{\B(0,1)}$ quindi, differenziando ambo i membri,
	\[ (Df)^T f = 0. \]
	Allora, siccome $\abs{f} = 1$ su $\overline{\B(0,1)}$, $0$ è un autovalore per $(Df)^T$, quindi $\det(Df) = \det((Df)^T) = 0$, assurdo.
	
	Mostriamo quindi che non esiste alcuna applicazione $f: \overline{\B(0,1)} \to \partial \B(0,1)$ continua tale che per ogni $x \in \partial \B(0,1)$ si abbia $f(x) = x$. Supponiamo, per assurdo, che una tale applicazione esista. Estendiamo $f$ con continuità fuori da $ \overline{\B(0,1)}$ in modo che sui segmenti radiali assuma il valore che ha su $\partial \B(0,1)$. Chiaramente, $f$ è non nulla su tutto $\R^n$. Sia $\varepsilon >0$ tale che l'$\varepsilon$-mollificatore\footnote{Perché abbiamo bisogno di un'applicazione almeno di classe $C^2$, così è pure $C^\infty$.} $\varrho_\varepsilon$ realizzi $ \varrho_\varepsilon \ast f \ne 0$ su $\R^n$ e $\varrho_\varepsilon \ast f(x) = x$ su $\R^n\setminus{\B(0,2)}$. Detto ciò, l'applicazione $\frac{2\varrho_\varepsilon \ast f}{\abs{\varrho_\varepsilon \ast f}}$, a meno di un riscalamento, è un'applicazione liscia che è l'identità su $\partial \B(0,1)$, da cui l'assurdo.
	
	A questo punto, la conclusione segue in modo simile a come riportato nel Capitolo \ref{Teoremi di punto fisso}.
\end{osservazione}

\section{Esistenza ed unicità di minimi}
In questa sezione ci occupiamo del caso scalare. Nel seguito considereremo $q \in \left] 1,\infty\right[ $.

Se $L$ è coercitiva, esistono $\alpha>0,\beta \ge 0$ tali che per ogni $x \in U, s \in \R,\xi \in \R^n$
\[ L(x,s,\xi) \ge \alpha\abs{\xi}^q-\beta. \]
Allora, 
\[ I[u] = \int_U L(x,u,Du)d\mathcal{L}^n(x) \ge \int_U \left(\alpha\abs{Du}^q-\beta\right)d\mathcal{L}^n = \alpha \norm{Du}_{L^q(U)}^q -\gamma  \]
dove $\gamma = \beta \mathcal{L}^n(U) >0$. Riassumendo,
\begin{equation}
	I[u] \ge \alpha \norm{Du}_{L^q(U)}^q -\gamma,
\end{equation}
detta \emph{condizione di coercitività}\index[analitico]{condizione!di coercitività} di $I$. Osserviamo subito che se $\norm{Du}_{L^q(U)} \to +\infty,$ allora $I[u] \to +\infty$. La condizione di coercitività ha significato anche per $u$ meno regolari rispetto alle ipotesi iniziali. 

Non è detto, a priori, che il funzionale ammetta minimo.
\begin{esempio}
	Dato $ q \in \left] 2,2^\ast\right[$, sia $I: H^1_0(U) \to \R$ tale che 
	\[ I[u] = \frac{1}{2}\int_U \abs{Du}^2 d\mathcal{L}^n - \frac{1}{q}\int_U \abs{u}^q d\mathcal{L}^n. \]
	Dato $u_0 \in H^1_0(U)$, consideriamo la successione $(hu_0)$. Allora,
	\[ I[hu_0] = \frac{h^2}{2}\int_U \abs{Du}^2 d\mathcal{L}^n - \frac{h^q}{q}\int_U \abs{u}^q d\mathcal{L}^n \le C(h^2 - h^q) \to -\infty, \]
	quindi $I$ non può ammettere minimo su $H^1_0(U)$.
\end{esempio}
Considereremo, da questo punto in poi, $I: W^{1,q}_0(U) \to \overline{\R}$. Osserviamo che, anche se è ben definito, potrebbe comunque fare $+\infty$.


\begin{teorema}
	Supponiamo che $L$ sia inferiormente limitata e convessa nella variabile $\xi$. Allora $I$ è debolmente inferiormente semicontinuo.
\end{teorema}
\begin{proof}
	Innanzitutto, senza perdita di generalità\footnote{Siccome $L \ge \beta$ e siamo in un dominio limitato, si tratta di applicare gli argomenti seguenti all'applicazione $\hat{L}= L - \beta \ge 0$.}, siccome $L$ è inferiormente limitata, possiamo supporre che 
	\[ L \ge 0. \]
	Siano $(u_h)$ in $W^{1,q}_0(U)$ ed $u \in W^{1,q}_0(U)$ tali che $u_h \rightharpoonup u$ in $W^{1,q}_0(U)$. Chiamiamo 
	\[ l = \underset{n}{\textup{lim inf}\;} I[u_h]. \]
	
	Mostriamo che $I[u] \le l.$ Innanzitutto, a meno di una sottosuccessione, possiamo supporre che $l = \lim\limits_{h} I[u_h]$. Inoltre, per il Teorema di Rellich, sicuramente $u_h \to u$ in $L^q(U)$ e quindi, a meno di una ulteriore sottosuccessione, $u_h \to u $ q.o. in $U$. Se $\varepsilon >0$, per il Teorema di Severini--Egorov esiste allora $E_{\varepsilon} \subseteq U$ tale che $u_h \to u$ uniformemente su $E_{\varepsilon}$ e 
	\[ \mathcal{L}^n(U\setminus E_\varepsilon) < \varepsilon. \]
	Senza perdita di generalità possiamo supporre $E_\varepsilon \subseteq E_{\varepsilon'}$ se $0 < \varepsilon' < \varepsilon$. Poniamo 
	\[ F_\varepsilon = \Set{x \in U : \abs{u(x)}+\abs{Du(x)}\le \frac{1}{\varepsilon}}. \]
	Osserviamo subito che 
	\[ N= \int_U \left( \abs{u}^q + \abs{Du}^q\right) d\mathcal{L}^n < +\infty,  \]
	ma 
	\begin{align*}
		N &= \int_{F_\varepsilon}\left( \abs{u}^q + \abs{Du}^q\right) d\mathcal{L}^n + \int_{U\setminus F_\varepsilon}\left( \abs{u}^q + \abs{Du}^q\right) d\mathcal{L}^n \ge \int_{U\setminus F_\varepsilon}\left( \abs{u}^q + \abs{Du}^q\right) d\mathcal{L}^n \ge \\
		&\ge \int_{U\setminus F_\varepsilon} \left( \frac{1}{2\varepsilon}\right) ^q d\mathcal{L}^n = \left( \frac{1}{2\varepsilon}\right) ^q \mathcal{L}^n(U\setminus F_\varepsilon)
	\end{align*}
	e quindi 
	\[ \mathcal{L}^n(U\setminus F_\varepsilon)  \le (2\varepsilon)^q N \]
	da cui
	\[ \lim\limits_{\varepsilon \to 0^+} \mathcal{L}^n(U\setminus F_\varepsilon)  = 0. \]
	Consideriamo anche 
	\[ G_\varepsilon = E_\varepsilon \cap F_\varepsilon. \]
	Siccome, utilizzando le Leggi di De Morgan,
	\[ \mathcal{L}^n(U\setminus G_\varepsilon) = \mathcal{L}^n((U\setminus E_\varepsilon) \cup (U\setminus F_\varepsilon)) \le \mathcal{L}^n(U\setminus E_\varepsilon) + \mathcal{L}^n(U\setminus F_\varepsilon) \]
	abbiamo subito che 
	\[ \lim\limits_{\varepsilon \to 0^+} \mathcal{L}^n(U\setminus G_\varepsilon)  = 0. \]
	Ora, essendo $L$ di classe $C^\infty$ e convessa nella variabile $\xi$, per ogni $x \in U$, per ogni $s \in \R$ e per ogni $\xi, \eta \in \R^n$
	\[ L(x,s,\xi) \ge L(x,s,\eta) + D_\xi L(x,s,\eta) \cdot (\xi - \eta)\footnote{Ossia $L$ sta sopra il suo piano tangente in $\xi$.} \]
	e quindi 
	\begin{align*}
		I[u_h] &= \int_U L(x,u_h(x),Du_h(x))d\mathcal{L}^n(x) \ge \int_{G_\varepsilon} L(x,u_h(x),Du_h(x))d\mathcal{L}^n(x) \ge \\
		&\ge \int_{G_\varepsilon} L(x,u_h(x),Du(x))d\mathcal{L}^n(x) + \\
		&+ \int_{G_\varepsilon}D_\xi L(x,u_h(x),Du(x))\cdot (Du_h(x) - Du(x))d\mathcal{L}^n(x) = \\
		&= \int_{G_\varepsilon} L(x,u_h(x),Du(x))d\mathcal{L}^n(x) + \\
		&+ \int_{G_\varepsilon}\left( D_\xi L(x,u_h(x),Du(x))-D_\xi L(x,u(x),Du(x))\right) \cdot (Du_h(x) - Du(x))d\mathcal{L}^n(x) + \\
		&+ \int_{G_\varepsilon}D_\xi L(x,u(x),Du(x))\cdot (Du_h(x) - Du(x))d\mathcal{L}^n(x).
	\end{align*}
	Ora, utilizzando la disuguaglianza di Holder,
	\begin{align*}
		&\abs*{\int_{G_\varepsilon}\left( D_\xi L(x,u_h(x),Du(x))-D_\xi L(x,u(x),Du(x))\right) \cdot (Du_h(x) - Du(x))d\mathcal{L}^n(x)} \le \\
		& \le \int_{G_\varepsilon}\abs*{ D_\xi L(x,u_h(x),Du(x))-D_\xi L(x,u(x),Du(x))}\abs*{Du_h(x) - Du(x)}d\mathcal{L}^n(x) \le \\
		&\le \norm{D_\xi L(\cdot,u_h,Du)-D_\xi L(\cdot,u,Du)}_{L^{q'}(G_\varepsilon)}\norm{Du_h - Du}_{L^q(G_\varepsilon)},
	\end{align*}
	ma $Du_h - Du \rightharpoonup 0$ in $L^q(U)$ e quindi in particolare $(Du_h - Du)$ è limitata in $L^q(U)$, quindi in $L^q(G_\varepsilon)$. Ne risulta che
	\begin{align*}
		&\abs*{\int_{G_\varepsilon}\left( D_\xi L(x,u_h(x),Du(x))-D_\xi L(x,u(x),Du(x))\right) \cdot (Du_h(x) - Du(x))d\mathcal{L}^n(x)} \le \\
		&\le C \norm{D_\xi L(\cdot,u_h,Du)-D_\xi L(\cdot,u,Du)}_{L^{q'}(G_\varepsilon)}.
	\end{align*}
	Dal Teorema di Lagrange ora, esiste $\sigma \in \left[ 0,1\right] $ tale che
	\begin{align*}
		&\abs*{\int_{G_\varepsilon}\left( D_\xi L(x,u_h(x),Du(x))-D_\xi L(x,u(x),Du(x))\right) \cdot (Du_h(x) - Du(x))d\mathcal{L}^n(x)} \le \\
		&\le C \left( \int_{G_\varepsilon}\abs*{D^2_{s,\xi}L(x,\sigma u_h(x) +(1-\sigma)u(x),Du(x))}^{q'}\abs*{u_h(x)-u(x)}^{q'}d\mathcal{L}^n(x)\right) ^{\frac{1}{q'}} \le \\
		&\le C \norm{u_h-u}_{L^\infty(G_\varepsilon)}\left( \int_{G_\varepsilon}\abs*{D^2_{s,\xi}L(x,\sigma u_h(x) +(1-\sigma)u(x),Du(x))}^{q'}d\mathcal{L}^n(x)\right) ^{\frac{1}{q'}}.
	\end{align*}
	Siccome
	\[ \abs{\sigma u_h - (1-\sigma)u} \le \sigma \abs{u_h} + (1-\sigma)\abs{u} \le \sigma \abs{u_h - u} + \abs{u} \le \sigma \norm{u_h-u}_{L^\infty(G_\varepsilon)} + \frac{1}{\varepsilon} \le \frac{2}{\varepsilon}\]
	si ha che $\norm{D^2_{s,\xi}L(\cdot,u,Du)}_{L^\infty(G_\varepsilon)} < +\infty$ e quindi 
	\begin{align*}
		&\abs*{\int_{G_\varepsilon}\left( D_\xi L(x,u_h(x),Du(x))-D_\xi L(x,u(x),Du(x))\right) \cdot (Du_h(x) - Du(x))d\mathcal{L}^n(x)} \le \\
		& \le C \norm{u_h-u}_{L^\infty(G_\varepsilon)}.
	\end{align*}
	In definitiva, se $h \to +\infty$
	\[ \int_{G_\varepsilon}\left( D_\xi L(x,u_h(x),Du(x))-D_\xi L(x,u(x),Du(x))\right) \cdot (Du_h(x) - Du(x))d\mathcal{L}^n(x) \to 0. \]
	Analogamente\footnote{Aggiungi e togli $L(x,u(x),Du(x))$ e poi ragiona come sopra.}, si vede che se $h \to +\infty$
	\[  \int_{G_\varepsilon} L(x,u_h(x),Du(x))d\mathcal{L}^n(x) \to  \int_{G_\varepsilon} L(x,u(x),Du(x))d\mathcal{L}^n(x). \]
	Invece, 
	\begin{align*}
		\int_{G_\varepsilon}D_\xi L(x,u(x),Du(x))\cdot (Du_h(x) - Du(x))d\mathcal{L}^n(x)  &= \int_{U} \chi_{G_\varepsilon}(x)D_\xi L(x,u(x),Du(x))\cdot \\
		&\cdot (Du_h(x) - Du(x))d\mathcal{L}^n(x), 
	\end{align*}
	ma, su $G_\varepsilon$, $D_\xi L(x,u(x),Du(x))$ è limitato e quindi in particolare sta in $L^{q'}(U)$. Per convergenza debole si conclude quindi che se $h \to +\infty$
	\[ \int_{G_\varepsilon}D_\xi L(x,u(x),Du(x))\cdot (Du_h(x) - Du(x))d\mathcal{L}^n(x)  \to 0. \]
	Riassumendo, per ogni $\varepsilon > 0$
	\[ l = \lim\limits_{h} I[u_h] \ge \int_{G_\varepsilon} L(x,u(x),Du(x))d\mathcal{L}^n(x). \]
	Ora, 
	\[ \int_{G_\varepsilon} L(x,u(x),Du(x))d\mathcal{L}^n(x) = \int_{U}g_\varepsilon(x) d\mathcal{L}^n(x), \]
	dove 
	\[ g_\varepsilon(x) = \chi_{G_\varepsilon}(x) L(x,u(x),Du(x)) \ge 0. \]
	Inoltre, se $\varepsilon' < \varepsilon$, 
	\[ \frac{1}{\varepsilon} < \frac{1}{\varepsilon'} \]
	e quindi $F_\varepsilon \subseteq F_{\varepsilon'}$. Ciò, combinato con il fatto che $E_\varepsilon \subseteq E_{\varepsilon'}$, permette di dedurre che 
	\[ g_\varepsilon \le  g_{\varepsilon'}.\]
	Siccome poi $\mathcal{L}^n(U\setminus G_\varepsilon) \to 0$ per $\varepsilon \to 0$, si ha anche che se $\varepsilon \to 0$, $g_\varepsilon \to L(x,u(x),Du(x))$ q.o. in U. Dal Teorema di Beppo--Levi segue dunque che se $\varepsilon \to 0$
	\[  l \ge \int_U L(x,u(x,Du(x)))d\mathcal{L}^n(x) = I[u].\qedhere\]
\end{proof}
\begin{osservazione}
	Vale la pena di notare che la difficoltà della dimostrazione sta nel fatto che $Du_h \rightharpoonup Du$ in $L^{q}(U)$ non implica, nemmeno passando a sottosuccessioni, $Du_h \to Du$ \emph{q.o.} in $U$. La chiave per superare questo ostacolo risiede proprio nella convessità nei gradienti.
\end{osservazione}
%\newpage
\begin{teorema}\textbf{\emph{(di esistenza)}}
	Supponiamo che $L$ sia coercitiva e convessa in $\xi$. 
	
	Allora esiste $u \in W^{1,q}_0(U)$ tale che 
	\[ I[u] = \underset{w \in W^{1,q}_0(U)}{\min \;}I[w]. \]
\end{teorema}
\begin{proof}
	Si tratta di applicare il Metodo Diretto del Calcolo delle Variazioni.
\end{proof}
Senza pretesa di completezza, la seguente è una possibile condizione sufficiente per l'unicità del minimo.
\begin{teorema}\textbf{\emph{(di unicità)}}
	Supponiamo che $L$ sia coercitiva e convessa nella variabile $\xi$. Supponiamo inoltre che $L$ non dipenda da $s$ e che esista $\nu > 0 $ tale che per ogni $x \in U$, per ogni $\xi,\eta \in \R^n$
	\[ \sum_{i,j=1}^n D^2_{\xi_i,\xi_j}L(x,\xi)\eta_i\eta_j \ge \nu \abs{\eta}^2. \qquad \textup{ (uniforme convessità)} \]
	Allora esiste ed è unico $u \in W^{1,q}_0(U)$ tale che 
	\[ I[u] = \underset{w \in W^{1,q}_0(U)}{\min \;}I[w]. \]
\end{teorema}
\begin{proof}
	Omettiamo la dimostrazione.
\end{proof}

\begin{definizione}
	Diciamo che $u \in W^{1,q}_0(U)$ è \emph{soluzione debole dell'equazione di Eulero--Lagrange}\index[analitico]{soluzione!debole!dell'equazione di Eulero--Lagrange} associata ad $L$ 
	\[ \begin{cases}
		-\sum_{i=1}^n D^2_{x_i,\xi_i}\left( L(x,u(x),Du(x))\right)  + D_{s}L(x,u(x),Du(x)) = 0 & \textup{in } U,\\
		u = 0 & \textup{su } \partial U,
	\end{cases} \]
	se per ogni $v \in W^{1,q}_0(U)$ si ha
	\[ \int_U \left( \sum_{i=1}^n D_{\xi_i} L(x,u(x),Du(x)) D_{x_i}v(x) + D_sL(x,u(x),Du(x)) v(x) \right) d\mathcal{L}^n(x) = 0. \]
\end{definizione}
\begin{teorema}
	Supponiamo esista $C>0$ tale che per ogni $x \in U$, per ogni $s \in \R$ e per ogni $\xi \in \R^n$ 
	\begin{align*}
		\abs{L(x,s,\xi)} &\le C (1+ \abs{s}^q+\abs{\xi}^q), \\
		\abs{D_\xi L(x,s,\xi)} &\le C (1+ \abs{s}^{q-1}+\abs{\xi}^{q-1}), \\
		\abs{D_s L(x,s,\xi)} &\le C (1+ \abs{s}^{q-1}+\abs{\xi}^{q-1}).
	\end{align*} 
	Se $u \in W^{1,q}_0(U)$ è tale che 
	\[ I[u] = \underset{w \in W^{1,q}_0(U)}{\min \;}I[w], \]
	allora è soluzione debole dell'equazione di Eulero--Lagrange associata ad $L$.
\end{teorema}
\begin{proof}
	Sia $v \in W^{1,q}_0(U)$. Dalla prima stima di crescita ricaviamo subito che $\abs{i(\tau)} < +\infty$ per ogni $\tau \in \R$. Ora, se $\tau \ne 0$, 
	\[ \dfrac{i(\tau)-i(0)}{\tau - 0} = \int_U \dfrac{L(x,u(x)+\tau v(x),Du(x)+\tau Dv(x)) - L(x,u(x),Du(x))}{\tau} d\mathcal{L}^n(x). \]
	Per $\tau \to 0$,
	\[  \dfrac{L(\cdot,u+\tau v,Du+\tau Dv) - L(\cdot,u,Du)}{\tau} \to \sum_{i=1}^n D_{\xi_i}L(\cdot,u,Du)D_{x_i}v + D_s L(\cdot,u,Du)v \]
	q.o. in $U$. Inoltre, per il Teorema fondamentale del calcolo integrale, per q.o. $x \in U$,
	\begin{align*}
		&\dfrac{L(x,u(x)+\tau v(x),Du(x)+\tau Dv(x)) - L(x,u(x),Du(x))}{\tau} = \\ &=\frac{1}{\tau}\int_{0}^{\tau} \frac{d}{dt} L(x,u(x)+tv(x,Du(x+tDv(x))dt = \\
		&= \frac{1}{\tau} \int_{0}^\tau\left( \sum_{i=1}^n D_{\xi_i}L(x,u(x),Du(x))D_{x_i}v(x) + D_s L(x,u(x),Du(x))v(x)\right) d t .
	\end{align*}
	e quindi combinando le stime di crescita con la disuguaglianza di Young si ottiene che 
	\begin{align*}
		&\abs*{\dfrac{L(x,u(x)+\tau v(x),Du(x)+\tau Dv(x)) - L(x,u(x),Du(x))}{\tau}} \le \\
		& \le C\left( \abs{Du}^q +\abs{u}^q + \abs{Dv}^q +\abs{v}^q + 1\right)  \in L^1(U).
	\end{align*}
	Il Teorema della convergenza dominata restituisce quindi
	\[ 0 = i'(0) = \int_{U}\left( \sum_{i=1}^n D_{\xi_i}L(x,u(x),Du(x))D_{x_i}v(x) + D_\xi L(x,u(x),Du(x))v(x)\right)d\mathcal{L}^n(x),  \]
	da cui la tesi.
\end{proof}
\begin{osservazione}
	Le stime di crescita possono essere indebolite. I dettagli esulano però dai nostri scopi.
\end{osservazione}
\begin{proposizione}
	Supponiamo che per ogni $x \in U$ l'applicazione $\left\lbrace (s,\xi) \mapsto L(x,s,\xi)\right\rbrace $ sia convessa. Allora ogni soluzione debole dell'equazione di Eulero--Lagrange associata ad $L$ è un minimo per $I$.
\end{proposizione}
\begin{proof}
	Sia $u$ una soluzione debole dell'equazione di Eulero--Lagrange associata ad $L$. Sia $w \in W^{1,q}_0(U)$. Dalla convessità dell'applicazione $\left\lbrace (s,\xi) \mapsto L(x,s,\xi)\right\rbrace $ si ha che per ogni $x \in U$, per ogni $s \in \R$ e per ogni $\xi, \eta \in \R^n$
	\[ L(x,s,\xi) + D_\xi L(x,s,\xi) \cdot (\eta - \xi) + D_s L(x,s,\xi) \cdot (w-s) \le L(x,w,\eta). \]
	Allora, scelto $\xi = Du(x)$, $\eta = Dw(x)$, $s = u(x)$ ed integrando su $U$
	\begin{align*}
		I[u] &+ \int_U \Bigl(D_\xi L(x,u(x),Du(x)) \cdot (Dw(x) - Du(x)) + \\
		&+D_s L(x,u(x),Du(x)) \cdot (w(x)-u(x)) \Bigr) d\mathcal{L}^n(x) \le I[w]
	\end{align*}
	ma dal fatto che $u$ è soluzione debole e che $w-u \in W^{1,q}_0(U)$ si ottiene
	\[ I[u] \le I[w] \]
	da cui la tesi.
\end{proof}

\section{Vincoli}
In questa sezione, $U$ indicherà un generico sottoinsieme aperto limitato di $\R^n$ e con $\partial U$ di classe $C^\infty$. 

Ci poniamo l'obiettivo di studiare alcune tipologie di problemi vincolati. Iniziamo dal caso dei \emph{vincoli di tipo integrale}\index[analitico]{vincolo!integrale}, di cui studiamo il seguente problema modello: consideriamo $U$ connesso ed il funzionale $I: H^1_0(U) \to \R$ tale che
\[ I[w] = \frac{1}{2}\int_U \abs{Dw}^2 d\mathcal{L}^n. \]
In aggiunta, data un'applicazione $G: \R\to \R$ liscia, tale che, detta $g= G'$, si abbia per ogni $z \in R$
\[ \abs{g(z)}\le C(1+\abs{z}), \]
e quindi, di conseguenza,
\[ \abs{G(z)}\le C(1+\abs{z}^2), \]
per qualche $C >0$, consideriamo anche il funzionale $J: H^1_0(U) \to \R$ tale che
\[ J[w] = \int_U G(w)d\mathcal{L}^n. \]
Detto 
\[ \mathcal{A} = \left\lbrace w \in H^1_0(U): J[w] = 0\right\rbrace, \]
siamo interessati a 
\[ \underset{w \in \mathcal{A}}{\inf}\; I[w]. \]
\begin{teorema}
	Se $\mathcal{A}$ è non vuoto, allora esiste $u \in \mathcal{A}$ tale che 
	\[ I[u] =  \underset{w \in \mathcal{A}}{\min}\; I[w]. \] 
\end{teorema}
\begin{proof}\footnote{Si tratta di usare il Metodo Diretto del Calcolo delle Variazioni.}
Sia $(u_h)$ in $\mathcal{A}$ una successione minimizzante, ossia 
\[ I[u_h] \to \underset{\mathcal{A}}{\inf}\; I. \]
In particolare, esiste $K>0$ tale che per ogni $h \in \N$
\[ \norm{u_h}_{H^1_0(U)} = \sqrt{2 I[u_h]} \le K <+\infty, \]
quindi, per il Teorema di Eberlein--Smulian, esistono $u \in H^1_0(U)$ e $(u_{h_k})$ tale che $u_{h_k} \rightharpoonup u$ in $H^1_0(U)$. In particolare, per la debole semicontinuità inferiore della norma di $H^1_0(U)$, 
\[ I[u] \le \underset{k}{\liminf} I[u_{h_k}] = \underset{\mathcal{A}}{\inf}\; I. \]
Per il Teorema di Rellich--Kondrachov e l'unicità del limite, a meno di sottosuccessioni, $u_{h_k} \to u$ in $L^2(U)$, quindi\footnote{Si usa il fatto che 
\[ \abs{G(t)-G(s)} \le \int_s^t \abs{g(\tau)} d\tau \le C \int_s^t (1+\abs{\tau}) d\tau \le C\int_s^t (1+\abs{s}+\abs{t})d\tau \le C\abs{t-s}(1+\abs{s}+\abs{t}).\]
} 
\[ \abs{J[u]} = \abs{J[u]-J[u_{h_k}]}\le \int_U \abs{G(u)-G(u_{h_k})}d\mathcal{L}^n \le C \int_U \abs{u-u_{h_k}}(1+\abs{u}+\abs{u_{h_k}})d\mathcal{L}^n,\]
da cui, per la Disuguaglianza di Holder,
\[ \abs{J[u]}  \le C \norm{u-u_{h_k}}_{L^2(U)} \to 0, \]
ossia $J[u] = 0$, quindi $u \in \mathcal{A}$ e la tesi segue.
\end{proof}
Vediamo come scrivere l'equazione di Eulero--Lagrange.
\begin{teorema}
Sia $u \in \mathcal{A}$ tale che 
	\[ I[u] =  \underset{w \in \mathcal{A}}{\min}\; I[w]. \]
Esiste $\lambda \in \R$ tale che per ogni $v \in H^1_0(U)$ 
\[ \int_U Du\cdot Dv d\mathcal{L}^n = \lambda \int_U g(u)v d\mathcal{L}^n. \]
\end{teorema}
\begin{proof}
Fissiamo, innanzitutto, $v \in H^1_0(U)$. Vediamo dapprima il caso in cui $g(u)\ne 0$ q.o. in $U$. In particolare, esiste $w \in H^1_0(U)$ tale che 
\[ \int_U g(u)wd\mathcal{L}^n \ne 0: \]
infatti se, per assurdo, per ogni $w \in H^1_0(U)$ si avesse 
\[ \int_U g(u)wd\mathcal{L}^n = 0, \]
in particolare varrebbe per ogni $w \in C^\infty_c(U)$ quindi, per il Lemma di Du Bois-Reymond, si otterrebbe $g(u)= 0$ q.o. in $U$, che è assurdo. Consideriamo $j: \R^2 \to \R$ tale che
\[ j(\tau,\sigma) = J[u+\tau v + \sigma w] = \int_U G(u+\tau v + \sigma w)d\mathcal{L}^n. \]
Chiaramente $j(0,0) = J[u] = 0$. Inoltre, per il Teorema di passaggio al limite sotto il segno di integrale ed il Teorema della convergenza dominata, si può osservare che $j$ è di classe $C^1$ e 
\begin{align*}
\frac{\partial j}{\partial \tau}(\tau,\sigma) &= \int_U g(u+\tau v + \sigma w)v d\mathcal{L}^n, & \frac{\partial j}{\partial \sigma}(\tau,\sigma) &= \int_U g(u+\tau v + \sigma w)w d\mathcal{L}^n.
\end{align*}
In particolare, 
\[ \frac{\partial j}{\partial \sigma}(0,0) \ne 0. \]
Per il Teorema delle funzioni implicite, esiste $\tau_0>0$ e $\varphi: [-\tau_0,\tau_0]\to \R$ di classe $C^1$ tale che $\varphi(0) = 0$ e $j(\tau,\varphi(\tau)) = 0$ per ogni $\tau \in [-\tau_0,\tau_0]$. In particolare,
\[ \frac{\partial j}{\partial \tau}(\tau,\varphi(\tau)) + \frac{\partial j}{\partial \sigma}(\tau,\varphi(\tau))\varphi'(\tau) = 0, \]
da cui 
\[ \varphi'(0) = - \frac{\int_U g(u)v d\mathcal{L}^n}{\int_U g(u)w d\mathcal{L}^n}. \]
Sia $i: [-\tau_0, \tau_0]\to \R$ di classe $C^1$ tale che
\[ i(\tau) = I[u+\tau v + \varphi(\tau) w]. \]
Siccome $j(\tau,\varphi(\tau)) = 0$, allora $u+\tau v + \varphi(\tau) w \in \mathcal{A}$. In particolare, siccome $I$ ha minimo in $u$, $i$ ha minimo in $0$ quindi, per il Teorema di derivazione sotto il segno di integrale,
	\[ 0 = i'(0) = \int_U Du \cdot (Dv + \varphi'(0)Dw)d\mathcal{L}^n\]
 da cui, posto 
 \[ \lambda = \frac{\int_U Du \cdot Dw d\mathcal{L}^n}{\int_U g(u)w d \mathcal{L}^n}, \]
 si ha 
 \[ \int_U Du\cdot Dv d\mathcal{L}^n = \lambda \int_U g(u)v d\mathcal{L}^n. \]
 
 Se, invece, $g=0$ q.o. in $U$, in particolare, $DG(u) = g(u)Du = 0$ q.o. in $U$.\footnote{Si tratta di usare il seguente fatto generale.
 	
 \par\noindent 
 \textbf{Teorema} \textit{Sia $u \in W^{1,1}_{loc}(U;\R^k)$ e sia $g: \R^k \to \R$ un'applicazione di classe $C^1$. Valgono i seguenti fatti:
 \begin{enumerate}[$(a)$]
 	\item $g(u)$ è $\mathcal{L}^n$-misurabile e $Dg(u)\cdot D_j(u)$ è $\mathcal{L}^n$-misurabile per ogni $j=1,\dots,n$,
 	\item se $g(u) \in L^1_{loc}(U)$ e $Dg(u) \cdot D_j u \in L^1_{loc}(U)$ per ogni $j=1,\dots,n$, allora $g(u) \in W^{1,1}_{loc}(U)$ e $D_j(g(u)) = Dg(u)\cdot D_j(u)$ per ogni $j=1,\dots,n$.
 \end{enumerate}}}
Dal fatto che $U$ è connesso, segue che $G(u)$ è costante su q.o. su $U$, ma $J[u] = 0$, quindi $G(u)=0$ q.o. in $U$. Siccome $u=0$ su $\partial U$ nel senso delle tracce, segue che $G(0)=0$.\footnote{Per definizione, esiste $(u_h)$ in $C^\infty_c(U)$ tale che $u_h \to u$ in $H^1_0(U)$. Siccome $G$ è liscia, in particolare, $G(u_h)$ è di classe $C^\infty$ per ogni $h \in \N$ e $G(u_h) \to G(u)$ in $W^{1,1}(U)$. Allora, 
\[ T(G(u)) = \lim\limits_{h} G(u_h)|_{\partial U} = \lim\limits_{h} G(u_{h}|_{\partial U}) = \lim\limits_{h} G(0) = G(Tu), \]
quindi $0 = T(0) = TG(u) = G(Tu) = G(0)$.} In particolare, $u= 0$ q.o. in $U$: se così non fosse, siccome $0$ è ammissibile, verrebbe violata la minimalità di $u$ essendo $I[u] > I[0] = 0$. La tesi, in questo caso, segue quindi per un qualsiasi $\lambda \in \R$.
\end{proof}
\begin{osservazione}
Il Teorema precedente afferma che i minimi del funzionale $I$ su $\mathcal{A}$ sono soluzioni deboli del problema di Dirichlet\footnote{Che è un problema agli autovalori non-lineare nelle incognite $(u,\lambda)$, con $u \ne 0$.}
\[ \begin{cases}
	-\Delta u = \lambda g(u) & \textup{ in } U,\\
	u = 0 & \textup{ su } \partial U,
\end{cases} \]
dove $\lambda$ è il \emph{moltiplicatore di Lagrange}\index[analitico]{moltiplicatore!di Lagrange} corrispondente al vincolo integrale $J[u] = 0$. 
\end{osservazione}

Veniamo ora al caso dei \emph{vincoli unilateri}\index[analitico]{vincolo!unilatero}, di cui studiamo l'esempio notevole del \emph{problema dell'ostacolo}: consideriamo $f: \overline{U} \to \R$ liscia ed il funzionale $I: H^1_0(U) \to \R$ tale che
\[ I[w] = \int_U \left( \frac{1}{2}\abs{Dw}^2-fw\right)  d\mathcal{L}^n. \]
Data un'applicazione $h: \overline{U}\to \R$ liscia, detta \emph{ostacolo}\index[analitico]{ostacolo}, e detto
\[ \mathcal{A} = \left\lbrace w \in H^1_0(U): w \ge h \textup{ q.o. in } U\right\rbrace, \]
siamo interessati a 
\[ \underset{w \in \mathcal{A}}{\inf}\; I[w]. \]
Chiaramente, $\mathcal{A}$ è un insieme convesso.
\begin{teorema}
Se $\mathcal{A}$ è non vuoto, allora esiste ed è unico $u \in \mathcal{A}$ tale che 
\[ I[u] =  \underset{w \in \mathcal{A}}{\min}\; I[w]. \] 
\end{teorema}
\begin{proof}
Sia $(u_h)$ in $\mathcal{A}$ una successione minimizzante, ossia 
\[ I[u_h] \to \underset{\mathcal{A}}{\inf}\; I. \]
In particolare, per la Disuguaglianza di Holder e la disuguaglianza di Poincaré,
\[ \norm{u_h}_{H^1_0(U)}^2 \le 2I[u_h] + C\norm{f}_{L^\infty(U)}\norm{u_h}_{H^1_0(U)} \]
e quindi, per la Disuguaglianza di Young pesata,
\[ \norm{u_h}_{H^1_0(U)}^2 \le 4I[u_h] + C. \]
Esiste, dunque, $K>0$ tale che per ogni $h \in \N$
\[ \norm{u_h}_{H^1_0(U)} \le K <+\infty, \]
quindi, per il Teorema di Eberlein--Smulian, esistono $u \in H^1_0(U)$ e $(u_{h_k})$ tale che $u_{h_k} \rightharpoonup u$ in $H^1_0(U)$.
Per il Teorema di Rellich--Kondrachov e l'unicità del limite, a meno di sottosuccessioni, $u_{h_k} \to u$ in $L^2(U)$, quindi $u_{h_k} \rightharpoonup u$ in $L^2(U)$ e, per le proprietà degli Spazi di Lebesgue, a meno di sottosuccessioni, $u_{h_k} \to u$ q.o. in $U$, quindi $u_{h_k} \ge h$ q.o. in $U$. Inoltre, usando anche la debole semicontinuità inferiore della norma di $H^1_0(U)$, 
\[ I[u] \le \underset{k}{\liminf} I[u_{h_k}] = \underset{\mathcal{A}}{\inf}\; I, \]
quindi 
\[ I[u] =  \underset{w \in \mathcal{A}}{\min}\; I[w].  \]

Supponiamo ora, per assurdo, che esistano $u,\overline{u} \in \mathcal{A}$ tali che $u \ne \overline{u}$ q.o. in $U$ e 
\[ I[u] =  \underset{w \in \mathcal{A}}{\min}\; I[w] = I[\overline{u}]. \]
Siccome $\mathcal{A}$ è convesso, $\frac{1}{2}(u + \overline{u}) \in \mathcal{A}$. Ma, per il fatto che $\norm{\;}_{H^1_0(U)}$ è una norma su $H^1_0(U)$, 
\begin{align*}
I\left[ \frac{u+\overline{u}}{2}\right] &= \int_U \frac{1}{2} \left( \abs*{\frac{Du + D\overline{u}}{2}}^2 -\frac{1}{2}fu -\frac{1}{2}f\overline{u}\right) d\mathcal{L}^n = \\
&= \int_U \frac{1}{8}\left( \abs{Du}^2 + \abs{D\overline{u}}^2 +2Du\cdot D\overline{u} -\frac{1}{2}fu -\frac{1}{2}f\overline{u}\right) d\mathcal{L}^n = \\
&= \int_U \frac{1}{8}\left( 2\abs{Du}^2 + 2\abs{D\overline{u}}^2 -\abs{Du-D\overline{u}}^2 -\frac{1}{2}fu -\frac{1}{2}f\overline{u}\right) d\mathcal{L}^n < \\
&< \frac{1}{2}I[u] +  \frac{1}{2}I[\overline{u}] = \underset{\mathcal{A}}{\min}\; I,
\end{align*}
assurdo e la tesi è dimostrata.
\end{proof}
\begin{teorema}
Sia $u \in \mathcal{A}$ l'unica soluzione di
\[ I[u] =  \underset{w \in \mathcal{A}}{\min}\; I[w]. \]
Per ogni $v \in \mathcal{A}$ 
\[ \int_U Du\cdot D(v-u) d\mathcal{L}^n \ge \int_U f(v-u) d\mathcal{L}^n. \]
\end{teorema}
\begin{proof}
Fissiamo, innanzitutto, $v \in \mathcal{A}$. Per ogni $\tau \in [0,1]$, siccome $\mathcal{A}$ è convesso, $u+\tau(v-u) \in \mathcal{A}$. Consideriamo $i: [0,1]\to \R$ di classe $C^1$ tale che 
\[ i(\tau) = I[u+\tau(v-u)]. \]
Allora $i(0)\le i(\tau)$ per ogni $\tau \in [0,1]$ quindi, per il Teorema della convergenza dominata,
\[ 0 \le \lim\limits_{\tau \to 0} \frac{i(\tau)- i(0)}{\tau} = \int_U \left( Du\cdot D(v-u) -f(v-u)\right) d\mathcal{L}^n, \]
da cui la tesi.
\end{proof}

Potendo dimostrare che $u \in W^{2,\infty}(U)$\footnote{Noi lo prendiamo per buono.}, allora l'insieme
\[ O= \left\lbrace x \in U: u(x) > h(x)\right\rbrace  \]
è aperto. Mostriamo che $u \in C^\infty(O)$ ed è soluzione dell'equazione 
\[ -\Delta u = f \textup{ in } O. \]
Sulla base della teoria della regolarità dei problemi lineari, è sufficiente dimostrare che è soluzione debole. Fissiamo $v \in C^\infty_c(O)$. Prima di tutto, esiste $\delta_0 >0$ tale che per ogni $\tau \in \left] -\delta_0,\delta_0\right[$ si abbia $u+\tau v \in \mathcal{A}$: infatti se, per assurdo, esistesse $(\tau_j)$ convergente a $0$ tale che $u+\tau_j v \notin \mathcal{A}$ allora, siccome $v$ ha supporto in $O$ e $u\ge h$ q.o. su $U$, dovrebbe essere $u + \tau_j v <h$ q.o. in $O$ e, passando al limite su $j$, $u \le h$ q.o. in $O$, da cui l'assurdo. In particolare, 
\[ \tau \int_O\left( Du\cdot Dv - fv\right) d\mathcal{L}^n \ge 0. \]
Siccome, in questo caso, è possibile scambiare $\tau$ con $-\tau$, si ottiene
\[ \int_O\left( Du\cdot Dv - fv\right) d\mathcal{L}^n = 0. \]

Se ora $v \in C^\infty_c(U)$ tale che $v \ge 0$ e $\tau \in \left] 0,1\right] $, $u +\tau v \ge u \ge h$, quindi $u+\tau v \in \mathcal{A}$. In particolare, 
\[ \int_U \left( Du\cdot Dv - fv\right) d\mathcal{L}^n \ge 0. \]
Dal fatto che $u \in W^{2,\infty}(U)$, si deduce che 
\[ \int_U \left( -\Delta u - f\right)v d\mathcal{L}^n \ge 0, \]
ossia $-\Delta u \ge f$ q.o. in $U$.

Riassumendo, 
\[ \begin{cases}
	u\ge h, -\Delta u \ge f & \textup{ q.o. in } U,\\
	-\Delta u = f & \textup{ su } O.
\end{cases} \]
In particolare, chiamiamo \emph{frontiera libera}\index[analitico]{frontiera!libera} l'insieme $F = \partial O \cap U$.

\begin{osservazione}
Lo studio dei problemi di frontiera libera è importante in molti ambiti: si trovano applicazioni nello studio di problemi di controllo ottimo legato al tempo di arresto di un moto browniano, in idrologia, nella teoria della plasticità,\dots
\end{osservazione}

Vediamo, infine, il caso dei \emph{vincoli puntuali}\index[analitico]{vincolo!puntuale}, di cui studiamo l'esempio notevole delle \emph{mappe armoniche sulla sfera}: consideriamo il funzionale $I: H^1(U;\R^m) \to \R$ tale che
\[ I[w] = \frac{1}{2}\int_U \abs{Dw}^2  d\mathcal{L}^n. \]
Data $g \in L^2(\partial U;\R^m)$ e detto
\[ \mathcal{A} = \left\lbrace w \in H^1(U;\R^m): w = g \textup{ su } \partial U \textup{ nel senso delle tracce, } \abs{w} = 1 \textup{ q.o. in } U\right\rbrace, \]
siamo interessati a 
\[ \underset{w \in \mathcal{A}}{\inf}\; I[w]. \]
L'idea è di minimizzare il funzionale energia su tutte le applicazioni da $U\subseteq \R^n$ nella sfera $\partial\B(0,1)\subseteq \R^m$.
\begin{osservazione}
Il problema delle mappe armoniche sulla sfera è utile per studiare alcuni modelli per i cristalli liquidi: in particolare, è usato nella Teoria di Frank.
\end{osservazione}
\begin{teorema}
	Se $\mathcal{A}$ è non vuoto, allora esiste $u \in \mathcal{A}$ tale che 
	\[ I[u] =  \underset{w \in \mathcal{A}}{\min}\; I[w]. \] 
\end{teorema}
\begin{proof}\footnote{Si tratta di usare il Metodo Diretto del Calcolo delle Variazioni.}
Sia $(u_h)$ in $\mathcal{A}$ una successione minimizzante, ossia 
\[ I[u_h] \to \underset{\mathcal{A}}{\inf}\; I. \]
In particolare, per il vincolo puntuale,
\[ \norm{u_h}_{H^1(U;\R^m)}^2 = \int_U \abs{u_h}^2d\mathcal{L}^n + 2I[u_h]\le \mathcal{L}^n(U) + 2I[u_h]. \]
Esiste, dunque, $K>0$ tale che per ogni $h \in \N$
\[ \norm{u_h}_{H^1(U;\R^m)} \le K <+\infty, \]
quindi, per il Teorema di Eberlein--Smulian, esistono $u \in H^1(U;\R^m)$ e $(u_{h_k})$ tale che $u_{h_k} \rightharpoonup u$ in $H^1(U;\R^m)$.
Per il Teorema di Rellich--Kondrachov e l'unicità del limite, a meno di sottosuccessioni, $u_{h_k} \to u$ in $L^2(U;\R^m)$, quindi $u_{h_k} \rightharpoonup u$ in $L^2(U;\R^m)$ e, per le proprietà degli Spazi di Lebesgue, a meno di sottosuccessioni, $u_{h_k} \to u$ q.o. in $U$, quindi $\abs{u} =1 $ q.o. in $U$ e, per convergenza debole,
\begin{align*}
	\int_{\partial U} \abs{Tu-g}^2d\mathcal{H}^{n-1}&\le C \int_{\partial U} \abs{Tu-Tu_{h_k}}^2d\mathcal{H}^{n-1} + C \int_{\partial U} \abs{Tu_h-g}^2d\mathcal{H}^{n-1} = \\
	&= C\int_{\partial U} \abs{Tu-Tu_{h_k}}^2d\mathcal{H}^{n-1} \to 0,
\end{align*}
quindi $u=g$ su $\partial U$ nel senso delle tracce. Inoltre, usando anche la debole semicontinuità inferiore della norma di $H^1(U;\R^m)$, 
\[ I[u] \le \underset{k}{\liminf} I[u_{h_k}] = \underset{\mathcal{A}}{\inf}\; I, \]
quindi 
\[ I[u] =  \underset{w \in \mathcal{A}}{\min}\; I[w]. \qedhere \]
\end{proof}
Vediamo come scrivere l'equazione di Eulero--Lagrange.
\begin{teorema}
Sia $u \in \mathcal{A}$ tale che 
\[ I[u] =  \underset{w \in \mathcal{A}}{\min}\; I[w]. \]
Per ogni $v \in H^1_0(U;\R^m)\cap L^\infty(U;\R^m)$\footnote{In modo che l'integrale a secondo membro sia finito, infatti, per il vincolo puntuale,
\[ \abs*{\int_U \abs{Du}^2u\cdot v d\mathcal{L}^n} \le \int_U\abs{Du}^2\abs{u}\abs{v} d\mathcal{L}^n = \int_U\abs{Du}^2\abs{v} d\mathcal{L}^n \le \norm{v}_{L^\infty(U;\R^m)}\norm{u}_{H^1(U;\R^m)}<+\infty. \]
}
\[ \int_U Du \cdot Dv d\mathcal{L}^n = \int_U \abs{Du}^2u\cdot v d\mathcal{L}^n. \]
\end{teorema}
\begin{proof}
Fissiamo, innanzitutto, $v \in H^1_0(U;\R^m)\cap L^\infty(U;\R^m)$. Usando il vincolo puntuale, segue l'esistenza di $\delta>0$ e $\sigma>0$ tale che per ogni $\tau \in \left] -\delta,\delta\right[ $ si abbia $\abs{u+\tau v}\ge \sigma$ q.o. in $U$: infatti se, per assurdo, esistessero $(\tau_j), (\sigma_j)$ convergenti a zero tali che $\abs{u+\tau_j v} < \sigma_j$ q.o. in $U$, si avrebbe $u = 0$ q.o. in $U$, assurdo.

Mostriamo che $\frac{u+\tau v}{\abs{u+\tau v}} \in \mathcal{A}$ per ogni $\tau \in \left] -\delta,\delta\right[ $. Innanzitutto, $\frac{u+\tau v}{\abs{u+\tau v}} \in H^1(U;\R^m)$\footnote{La stima della norma $L^2$ è evidente perché esce l'integrale di $1$. La parte con i gradienti segue dal seguente fatto generale
	\[ D_j\left( \frac{w_i}{\abs{w}}\right)  = \frac{D_j w_i \abs{w} - \frac{1}{\abs{w}}\displaystyle\sum_{h=1}^m w_iw_hD_j w_h}{\abs{w}^2},   \]
	combinato con il fatto che $\abs{u+\tau v}\ge \sigma$ q.o. in $U$ ed il vincolo puntuale.} e $\abs*{\frac{u+\tau v}{\abs{u+\tau v}}} = 1$ q.o. in $U$. Inoltre, siccome\footnote{Osserviamo che 
	\[ T\left( \sum_{j=1}^{m}u_j^2\right) = \sum_{j=1}^{m}(Tu_j)^2 = \sum_{j=1}^{m} g_j^2, \]
	dove abbiamo sottinteso un buon comportamento della traccia, che già è stato evidenziato nel caso dei vincoli integrali, ossia $\abs{g} = 1$.} 
	\begin{align*}
	T(u+\tau v) &= Tu +\tau Tv = Tu = g, & T\abs{u+\tau v} &= \abs{Tu+\tau Tv} = \abs{Tu} = \abs{g} = 1
	\end{align*}
	allora $T\left( \frac{u+\tau v}{\abs{u+\tau v}}\right)  = g$ q.o. in $\partial U$.\footnote{Abbiamo qui sottinteso un buon comportamento della traccia rispetto al quoziente.}
	
	Consideriamo $i: \left] -\delta,\delta\right[ \to \R$ di classe $C^1$ tale che 
	\[ i(\tau) = I\left[ \frac{u+\tau v}{\abs{u+\tau v}}\right] . \]
	In particolare, siccome $I$ ha minimo in $u$, $i$ ha minimo in $0$ quindi, per il Teorema di derivazione sotto il segno di integrale,
	\[ 0 = i'(0) = \int_U Du \cdot D\left( \frac{d}{d\tau}\left(\frac{u+\tau v}{\abs{u+\tau v}}\right) |_{\tau=0}\right) d\mathcal{L}^n.\]
	Ora, 
	\[ \frac{d}{d\tau}\frac{u+\tau v}{\abs{u+\tau v}} = \frac{v}{\abs{u+\tau v}} - \frac{((u+\tau v)\cdot v)(u+\tau v)}{\abs{u+\tau v}^3}, \]
	quindi 
	\[ \frac{d}{d\tau}\left( \frac{u+\tau v}{\abs{u+\tau v}}\right) |_{\tau=0} = v- (u\cdot v)u, \]
	da cui
	\[ \int_U Du \cdot Dv d\mathcal{L}^n = \int_U Du \cdot D((u\cdot v)u) d\mathcal{L}^n.\]
	Essendo $\abs{u}^2 = 1$, $(Du)^T u = 0$, quindi 
	\begin{align*}
	Du\cdot D((u\cdot v)v) &= \sum_{i=1}^{m}\sum_{j=1}^{n} D_j u_i \left( \sum_{h=1}^{m}D_j(v_hu_hu_i)\right)  = \\
	&= \sum_{i,h=1}^{m} \sum_{j=1}^{n} \left( \abs{D_ju_i}^2v_hu_h + D_ju_i u_h u_iD_jv_h + D_ju_i v_h u_iD_ju_h\right) =\\
	&= \sum_{i=1}^m\sum_{j=1}^{n}\abs{D_ju_i}^2(v\cdot u) + \sum_{h=1}^{m}\sum_{j=1}^{n}\left( \sum_{i=1}^{m}D_ju_i u_i\right) u_hD_jv_h + \\
	&+\sum_{h=1}^{m}\sum_{j=1}^{n}\left( \sum_{i=1}^{m}D_ju_i u_i\right) D_ju_hv_h = \abs{Du}^2(v\cdot u),
	\end{align*}
	da cui la tesi.
\end{proof}
\begin{osservazione}
Il Teorema precedente afferma che i minimi del funzionale $I$ su $\mathcal{A}$ sono soluzioni deboli del problema di Dirichlet
	\[ \begin{cases}
		-\Delta u = \abs{Du}^2u & \textup{ in } U,\\
		u = g & \textup{ su } \partial U,
	\end{cases} \]
	dove $\lambda = \abs{Du}^2$ è il \emph{moltiplicatore di Lagrange}\index[analitico]{moltiplicatore!di Lagrange} corrispondente al vincolo puntuale $\abs{u} = 1$.
\end{osservazione}
Vale la pena osservare una differenza: per un vincolo di tipo integrale il moltiplicatore di Lagrange è un numero reale; per un vincolo puntuale è un'applicazione.
\section{Punti critici}
Vediamo la teoria nell'ambientazione hilbertiana, trascurando possibili generalizzazioni in spazi di Banach o solamente metrici.
\begin{definizione}
	Consideriamo $(X, (\,|\,)_X)$ uno spazio di Hilbert su $\R$ ed $I: X \to \R$. Diciamo che $I$ è differenziabile in $u \in X$ se esiste $v \in X$\footnote{Qui si sta implicitamente usando il Teorema della rappresentazione di Riesz.} tale che per ogni $w \in X$
	\[ I[w] = I[u] + (v|w-u)_X + o(\norm{w-u}). \]
\end{definizione}
Se $v$ esiste, è unico e viene denotato con $I'[u]$.
\begin{definizione}
	Consideriamo $(X, (\,|\,)_X)$ uno spazio di Hilbert su $\R$, $u \in X$ ed $I: X \to \R$ differenziabile in $u$. Diciamo che $u$ è un \emph{punto critico} per $I$ \index[analitico]{punto!critico} se $I'[u]=0$.
\end{definizione}
\begin{definizione}
	Consideriamo $(X, (\,|\,)_X)$ uno spazio di Hilbert su $\R$ ed $I: X \to \R$. diciamo che $I \in C^1(X;\R)$ se $I'[u]$ esiste per ogni $u \in X$ e l'applicazione $I':X \to X$ è continua.
\end{definizione}
\begin{notazione}
	Consideriamo $(X, (\,|\,)_X)$ uno spazio di Hilbert su $\R$. Denotiamo con $\mathcal{C}$ l'insieme delle funzioni $I \in C^1(X;\R)$ tali che $I':X \to X$ sia lipschitziana sui sottoinsiemi limitati di $X$.
\end{notazione}
\begin{definizione}
	Consideriamo $(X, (\,|\,)_X)$ uno spazio di Hilbert su $\R$ ed $I \in C^1(X;\R)$. Diciamo che $I$ soddisfa la \emph{condizione di Palais--Smale}\index[analitico]{condizione!di Palais--Smale} se per ogni $(u_h)$ in $X$ tale che $(I[u_h])$ limitata e $I'[u_h] \to 0$ in $X$, allora esistono $(u_{h_k})$ ed $u \in X$ tali che $u_{h_k} \to u$ in $X$. 
\end{definizione}
Per evitare di appesantire tecnicamente la trattazione, considereremo in questa sezione $I \in \mathcal{C}$.\footnote{Tutto può essere fatto anche senza questa ipotesi però.} 

\begin{definizione}
	Consideriamo $(X, (\,|\,)_X)$ uno spazio di Hilbert su $\R$ ed $I \in \mathcal{C}$. Se $c \in \R$, chiamiamo \emph{$c$-sottolivello}\index[analitico]{$c$-sottolivello} l'insieme \index[simboli]{$A_c$}
	\[ A_c = \left\lbrace u \in X : I[u]\le c\right\rbrace  \]
	e \emph{insieme dei punti critici a livello $c$}\index[analitico]{punto!critico!al livello $c$} l'insieme \index[simboli]{$K_c$}
	\[ K_c =\left\lbrace u \in X : I[u]= c, I'[u]=0\right\rbrace. \]
	In particolare, diciamo che $c$ è un \emph{valore critico}\index[analitico]{valore!critico} per $I$ se $K_c \ne \emptyset$.
\end{definizione}

\begin{teorema}\emph{\textbf{(di deformazione)}}\index[analitico]{teorema!di deformazione}
	Consideriamo $(X, (\,|\,)_X)$ uno spazio di Hilbert su $\R$ ed $I \in \mathcal{C}$ soddisfacente la condizione di Palais--Smale. Se esiste $c \in \R$ tale che $K_c = \emptyset$, allora per ogni $\varepsilon>0$ sufficientemente piccolo esistono $\delta \in \left] 0,\varepsilon\right[$ ed $\eta \in C([0,1]\times X; X)$ tale che, detta $\eta_t(u) = \eta(t,u)$, valgono i seguenti fatti:
	\begin{enumerate}[$(a)$]
		\item per ogni $u \in X$, $\eta_0(u) = u$,
		\item per ogni $u \notin I^{-1}(\left[ c-\varepsilon,c+\varepsilon\right] )$ e per ogni $t \in [0,1]$, $\eta_t(u) = u$,
		\item per ogni $u \in X$ e per ogni $t \in [0,1]$, $I[\eta_t(u)] \le I[u]$,
		\item $\eta_1(A_{c+\delta})\subseteq A_{c-\delta}$.
	\end{enumerate}
\end{teorema}
\begin{proof}
	Mostriamo, innanzitutto, che esistono $\sigma, \varepsilon \in \left] 0,1\right[ $ tali che per ogni $u \in I^{-1}(\left] c-\varepsilon,c+\varepsilon\right] )$ si abbia $\norm{I'[u]}_X \ge \sigma$. Supponiamo, per assurdo, che esistano $(\sigma_j),(\varepsilon_j)$ in $\left] 0,1\right[ $ convergenti a $0$ e $u_j \in I^{-1}(\left] c-\varepsilon_j,c+\varepsilon_j\right] )$ tali che $\norm{I'[u_j]}_X \le \sigma_j$. In particolare, $I[u_j] \to c$, quindi è limitata, e $I'[u_j]\to 0$ quindi, per la condizione di Palais--Smale, a meno di sottosuccessioni, esiste $u \in X$ tale che $u_j \to u$ in $X$. Per continuità, $I[u] = c$ e $I'[u]=0$, in contraddizione con il fatto che $K_c = \emptyset$.
	
	Prendiamo $\delta \in \left] 0,\varepsilon\right[ \cap \left] 0,\frac{\sigma^2}{2}\right[$ e consideriamo 
	\begin{align*}
		A &= \left\lbrace u \in X : I[u]\le c-\varepsilon \textup{ o } I[u]\ge c+\varepsilon\right\rbrace, & B&=\left\lbrace u \in X: c-\delta \le I[u]\le c+\delta \right\rbrace.
	\end{align*}
	Siccome $I \in \mathcal{C}$, allora l'applicazione $\left\lbrace u \mapsto \textup{dist}(u,A) + \textup{dist}(u,B)\right\rbrace $ è limitata dal basso da una costante strettamente positiva su ogni sottoinsieme limitato di $X$. Di conseguenza, l'applicazione $g: X \to [0,1]$ tale che
	\[ g(u) = \frac{\textup{dist}(u,A)}{\textup{dist}(u,A) + \textup{dist}(u,B)} \]
	è lipschitziana sui sottoinsiemi limitati di $X$. In più, $g=0$ su $A$ e $g=1$ su $B$. Data $h: \left[ 0,+\infty\right[  \to \R$ tale che
	\[ h(t) = \begin{cases}
		1 & \textup{se }0\le t \le 1,\\
		\frac{1}{t} & \textup{se } t >1,
	\end{cases} \]
	definiamo $V: X \to X$ tale che
	\[ V(u) = -g(u)h(\norm{I'[u]}_X)I'[u]. \]
	Dal fatto che $V$ è limitata e lipschitziana sui sottoinsiemi limitati di $X$, per ogni $u \in X$ il problema 
	\[ \begin{cases}
		\frac{d\eta}{dt}(t) = V(\eta(t)) & \textup{se } t>0,\\
		\eta(0) = u,
	\end{cases} \]
	ammette una ed una sola soluzione $\eta$ definita su $\left[0,+\infty\right[ $\footnote{Non entriamo nei dettagli riguardo questo fatto: si tratta di estendere al caso hilbertiano quanto visto nel corso di Approfondimenti di Analisi Matematica.}. A meno di una restrizione e di un abuso di notazione, rimane definita un'applicazione $\eta \in C([0,1]\times X; X)$.
	\par\noindent
	$(a)$ Evidente.
	\par\noindent
	$(b)$ Sia $u \notin I^{-1}(\left[ c-\varepsilon,c+\varepsilon\right] ).$ In altre parole, $u \in A$. Allora $g(u) = 0$, quindi $V(u)=0$, ossia $u$ è soluzione del problema di Cauchy enunciato all'inizio. Per unicità della soluzione, $\eta_t = u$ per ogni $t \in [0,1]$.
	\par\noindent
	$(c)$ Per ogni $u \in X$, per ogni $t \in [0,1]$, siccome
	\begin{align*}
		\frac{d}{dt}I[\eta_t(u)] &= \left( I'[\eta_t(u)] | \frac{d}{dt}\eta_t(u)\right)_X = (I'[\eta_t(u)]| V(\eta_t(u)))= \\
		&= -g(\eta_t(u))h(\norm{I'[\eta_t(u)]}_X)\norm{I'[\eta_t(u)]}_X^2 \le 0,
	\end{align*}
	allora $I[\eta_t(u)] \le I[u]$.
	\par\noindent
	$(d)$ Sia $u \in A_{c+\delta}$. Se esiste $t \in [0,1]$ tale che $\eta_t(u) \notin B$, deve essere $I[\eta_t(u)]< c-\delta$ e, per la $(c)$,
	\[ I[\eta_1(u)] \le I[\eta_t(u)]< c-\delta, \]
	quindi $\eta_1(u) \in A_{c-\delta}$.
	
	D'altra parte, mettiamoci nel caso in cui per ogni $t \in [0,1]$ si ha $\eta_t(u) \in B$. Allora $g(\eta_t(u)) = 1$ per ogni $t \in [0,1]$ e 
	\[ \frac{d}{dt}I[\eta_t(u)]  =-h(\norm{I'[\eta_t(u)]}_X)\norm{I'[\eta_t(u)]}_X^2. \]
	Se $\norm{I'[\eta_t(u)]}_X \ge 1$, otteniamo 
	\[ \frac{d}{dt}I[\eta_t(u)] \le -1 < -\sigma^2. \]
	Se $\norm{I'[\eta_t(u)]}_X \le 1$, invece $h(\norm{I'[\eta_t(u)]}_X) = 1$ e quindi
	\[ \frac{d}{dt}I[\eta_t(u)] \le -\norm{I'[\eta_t(u)]}_X^2 \le -\sigma^2. \]
	In ogni caso, per la disuguaglianza di Lagrange,
	\[ I[\eta_1(u)] < I[u] - \sigma^2 \le c+\delta - \sigma^2 \le c-\delta, \]
	da cui $\eta_1(u) \in A_{c-\delta}$.
\end{proof}

Vediamo un controesempio al Teorema precedente.
\begin{esempio}
	Sia $f: \R^2 \to \R$ tale che 
	\[ f(x,y) = x^2+y^2. \]
	Siccome $f(0,0) = 0$ e $Df(0,0) =0$, allora $0 \in K_0$. In particolare, $K_0 \ne \emptyset$. Siccome $A_\delta = \B(0,\sqrt{\delta})$ e $A_{-\delta} = \emptyset$, non può esistere una deformazione che porta $A_\delta$ in $A_{-\delta}$.
\end{esempio}

\begin{teorema}\emph{\textbf{(del passo montano)}}
	Consideriamo $(X, (\,|\,)_X)$ uno spazio di Hilbert su $\R$ ed $I \in \mathcal{C}$ soddisfacente la condizione di Palais--Smale. Supponiamo valgano i seguenti fatti:
	\begin{enumerate}[$(i)$]
		\item $I[0] = 0$,
		\item esistono $r,a>0$ tali che per ogni $u \in \partial\B(0,r) \subseteq X$
		\[ I[u] \ge a, \]
		\item esiste $v \in X$ tale che $\norm{v}_X > r$ e $I[v] \le 0$.
	\end{enumerate}
	Allora, detto 
	\[ \Gamma = \left\lbrace g \in C([0,1];X): g(0) = 0, g(1) = v\right\rbrace,  \]
	si ha che 
	\[ c = \underset{g \in \Gamma}{\inf}\; \underset{t \in [0,1]}{\max}\; I[g(t)] \]
	è un valore critico per $I$.
\end{teorema}
\begin{proof}
	Mostriamo, innanzitutto, che $c\ge a>0$. Data $g \in \Gamma$, sia $\varphi \in C([0,1];\R)$ tale che
	\[ \varphi(t) = \norm{g(t)}_X. \]
	Osserviamo che $\varphi(0) = 0$ e $\varphi(1) = \norm{v}_X>r$, quindi esisterà $\overline{t} \in \left] 0,1\right[ $ tale che $\varphi(\overline{t})=r$, da cui $\norm{g(\overline{t})}_X = r$, quindi, per $(ii)$,
	\[  \underset{t \in [0,1]}{\max}\; I[g(t)] \ge I[g(\overline{t})] \ge a, \]
	da cui, passando all'inf su $\Gamma$, $c\ge a>0$.
	
	Supponiamo, per assurdo, che $K_c = \emptyset$. Allora, per il Teorema di deformazione, per ogni $\varepsilon \in \left] 0,\frac{a}{2}\right[ \cap \left] 0, c\right[ $ esiste $\delta \in \left] 0,\varepsilon\right[ $ ed $\eta = \eta_1 : X \to X$ tale che $\eta(A_{c+\delta})\subseteq A_{c-\delta}$ e, per ogni $u \notin I^{-1}(\left[ c-\varepsilon,c+\varepsilon\right] )$, $\eta(u) = u$.
	\begin{center}
		\begin{tikzpicture}[x=0.75pt,y=0.75pt,yscale=-1,xscale=1]
			%uncomment if require: \path (0,300); %set diagram left start at 0, and has height of 300
			
			%Straight Lines [id:da351590380041958] 
			\draw    (270,110.6) -- (440,110.4) ;
			\draw [shift={(442,110.4)}, rotate = 179.93] [color={rgb, 255:red, 0; green, 0; blue, 0 }  ][line width=0.75]    (10.93,-3.29) .. controls (6.95,-1.4) and (3.31,-0.3) .. (0,0) .. controls (3.31,0.3) and (6.95,1.4) .. (10.93,3.29)   ;
			%Straight Lines [id:da6689641897087604] 
			\draw    (310,104.6) -- (310,116.6) ;
			%Straight Lines [id:da03212271146730217] 
			\draw    (340,104.6) -- (340,116.6) ;
			%Straight Lines [id:da4012489765266216] 
			\draw    (400,104.6) -- (400,116.6) ;
			
			% Text Node
			\draw (305,83.4) node [anchor=north west][inner sep=0.75pt]    {$0$};
			% Text Node
			\draw (325,82.4) node [anchor=north west][inner sep=0.75pt]    {$c-\varepsilon $};
			% Text Node
			\draw (382,82.4) node [anchor=north west][inner sep=0.75pt]    {$c+\varepsilon $};
		\end{tikzpicture}
	\end{center}
	Per definizione di estremo inferiore, esiste $g \in \Gamma$ tale che
	\[ \underset{t \in [0,1]}{\max}\; I[g(t)] \le c +\delta, \]
	ossia, per ogni $t \in [0,1]$, $g(t) \in A_{c+\delta}$. Poniamo $\hat{g} = \eta \circ g$. Chiaramente $\hat{g} \in C([0,1];X)$. Inoltre, $\hat{g}(0) = \eta(0) = 0$ e $\hat{g}(1)= \eta(v) = v$, quindi $\hat{g} \in \Gamma$. Siccome $g(t) \in A_{c+\delta}$, allora $\hat{g}(t) \in A_{c-\delta}$, quindi 
	\[ \underset{t \in [0,1]}{\max}\; I[\hat{g}(t)] \le c -\delta, \]
	da cui $c \le c-\delta$, assurdo.
\end{proof}
\begin{osservazione}
	Pensiamo al caso $X = \R^2$. In questa situazione il grafico di $I$ è un "lenzuolo" sopra il piano $xy$. Sappiamo che l'origine è circondato da punti ad altezza $a$ e che esiste un punto che sta sotto il piano $xy$. Se partendo dall'origine volessi raggiungere tale punto dovrò necessariamente "salire" e poi "scendere": nel percorrere questo passo, attraverserò un punto critico per $I$.
\end{osservazione}


Il lettore interessato può approfondire altri aspetti legati al Calcolo delle Variazioni cominciando da
\textsc{B. Dacorogna}, \emph{Introduction to the Calculus of Variations}, Imperial College Press, London, 2014. Un altro testo, più recente, può essere \textsc{L. Brasco}, \emph{Handbook of Calculus of Variations for Absolute Beginners}, Springer, Cham, 2025. 
		
		% Commenti speciali

% !TEX encoding = UTF-8 Unicode
% !TEX TS-program = pdflatex
% !TEX spellcheck = it-IT
% !TEX root = Dispensa/afed.tex

% Capitolo 4
\chapter{Teoremi di punto fisso}\label{Teoremi di punto fisso}
\section{Introduzione}
Possedere condizioni sufficienti riguardo l'esistenza di punti fissi è fondamentale per un analista; in particolare, comunque, sarà indispensabile nell'ambito dell'Analisi non lineare, come vedremo nel capitolo \ref{Alcuni elementi di Analisi non lineare}. Un primo risultato, già noto dal corso di Analisi Matematica II, è il Teorema delle contrazioni.

\begin{teorema}\emph{\textbf{(di Banach--Caccioppoli)}}\index[analitico]{teorema!delle contrazioni}\index[analitico]{teorema!di Banach--Caccioppoli}
	Siano $(X,d)$ uno spazio metrico completo non vuoto ed $f: X \to X$ un'applicazione lipschitziana di costante $C \in \left[ 0,1\right[ $. Allora esiste uno ed un solo $\xi \in X$ tale che $f(\xi) = \xi$.
\end{teorema}
\begin{proof}
	Omettiamo la dimostrazione.
\end{proof}
Miriamo, nel corso di questo capitolo a studiare nel dettaglio altri risultati di questo tipo, sia in dimensione finita che infinita, e ad analizzare alcune prime notevoli applicazioni.
\section{Risultati in dimensione finita}
Il Teorema seguente è stato pubblicato per la prima volta nel $1912$ su \emph{Mathematische Annalen} da Luitzen Egbertus Jan "Bertus" Brouwer. Qui riportiamo una dimostrazione semplificata rispetto all'originale dovuta a John Milnor e, successivamente, rivista da Claude Ambrose Rogers.
\begin{teorema}\emph{\textbf{(di Brouwer)}}\index[analitico]{teorema!di Brouwer}
	Siano $K \subseteq \R^n$ convesso, compatto e non vuoto ed $f: K \to K$ un'applicazione continua. Esiste $x \in K$ tale che $f(x) = x$.
\end{teorema}
\begin{proof}
	Vediamo innanzitutto il caso $K = \overline{\B(0,1)}$. Mostriamo che non esiste alcuna applicazione $f: \overline{\B(0,1)} \to \partial \B(0,1)$ di classe $C^1$ tale che per ogni $x \in \partial \B(0,1)$ si abbia $f(x) = x$. Supponiamo, per assurdo, che una tale applicazione esista. Per ogni $x \in \overline{\B(0,1)}$ e per ogni $t \in [0,1]$ sia $f_t(x) = (1-t)x + tf(x)$. Osserviamo che
	\[ \abs{f_t(x)} \le (1-t)\abs{x} +t\abs{f(x)} \le (1-t) +t = 1, \]
	perciò $\textup{Im}(f_t)\subseteq \overline{\B(0,1)}$. Inoltre, se $x \in \partial \B(0,1)$,
	\[ f_t(x) = (1-t)x + tx = x. \]
	Detta $g(x) = f(x) - x$, si ha $f_t(x) = x + tg(x)$. Essendo $f$ di classe $C^1$ su $\overline{\B(0,1)}$,  tale è pure $g$, quindi, per la disuguaglianza di Lagrange ed il Teorema di Weierstrass, per ogni $x_1,x_2 \in \overline{\B(0,1)}$
	\[ \abs{g(x_2) - g(x_1)} \le \norm{dg}_\infty \abs{x_2 - x_1} \le C \abs{x_2 - x_1}, \]
	quindi $g$ è lipschitziana su $\overline{\B(0,1)}$. Se ora $x_1,x_2 \in \overline{\B(0,1)}$ tali che $f_t(x_1) = f_t(x_2)$, allora 
	\[ x_2 - x_1 = t(g(x_1 - g(x_2))), \]
	quindi 
	\[ \abs{x_2 - x_1} = t\abs{g(x_1) - g(x_2)} \le Ct \abs{x_2 - x_1}. \]
	In particolare, se $t \in \left[ 0,\frac{1}{C}\right[$, l'unica possibilità è che $x_1 = x_2$, altrimenti otterrei $t \ge \frac{1}{C}$, che sarebbe assurdo. In altre parole, per $t \in \left[ 0,\frac{1}{C}\right[$ $f_t$ è iniettiva. Osserviamo anche che per ogni $x \in \B(0,1)$
	\[ df_t(x) = \textup{Id} + tdg(x) \]
	quindi, essendo $g$ di classe $C^1$, esiste $t_0 \in  \left[ 0,\frac{1}{C}\right[$ tale che per ogni $t \in \left[ 0,t_0\right] $ e per ogni $x \in \B(0,1)$ si abbia $\det(df_t(x)) > 0$ che, in altre parole, significa che $df_t(x)$ è biettivo.\footnote{In generale, siano $A,B \in \textup{Mat}_n(\R)$ tali che $\det(A) > 0$ ed $a \in \R$. Consideriamo $\varphi: \left[ 0, a\right[ \to \R$ tale che $\varphi(t) = \det(A + tB)$. Essendo $\det$ un'applicazione continua, per composizione anche $\varphi$ è una funzione continua. In particolare è continua in $0$. Dal fatto che $\varphi(0) = \det(A) > 0$, per il Teorema della permanenza del segno, esiste $t_0 \in  \left[ 0,a\right[$ tale che per ogni $t \in \left[ 0,t_0\right]$ si abbia $\varphi(t) > 0$.} Per il Teorema di inversione locale, per ogni $t \in \left[ 0,t_0\right]$ e per ogni $x \in \B(0,1)$ esiste un intorno aperto $U_x$ di $x$ in $\B(0,1)$ tale che $f_t(U_x)$ sia aperto. In particolare, per ogni $t \in \left[ 0,t_0\right] $, siccome
	\[ G_t = f_t(\B(0,1)) = \bigcup_{x \in \B(0,1)} f_t(U_x),\]
	possiamo affermare che $G_t$ è aperto. Osserviamo che, in realtà, per ogni $t \in \left[ 0,t_0\right] $ si ha $G_t = \B(0,1)$. Infatti, se per assurdo così non fosse, allora esisterebbe $y_0 \in \partial G_t \cap \B(0,1)$. Sia allora $(x_h)$ in $\B(0,1)$ tale che $f_t(x_h) \to y_0$. Per compattezza, esiste $x_0 \in \overline{\B(0,1)}$ tale che, a meno di sottosuccessioni, $x_h \to x_0$. Allora, per continuità di $f_t$ ed unicità del limite, $f_t(x_0) = y_0$. Dal fatto che $G_t$ è aperto dovrebbe essere allora $x_0 \in \partial \B(0,1)$, ma quindi $y_ 0 = f_t(x_0) = x_0 \in \partial \B(0,1)$, che da la contraddizione. Possiamo quindi affermare che $f_t: \overline{\B(0,1)} \to \overline{\B(0,1)}$ è biettiva per ogni $t \in [0,t_0]$. Detta $F: [0,1] \to \R$ tale che
	\[ F(t) = \int_{B(0,1)} \det(df_t(x)) d\mathcal{L}^n(x), \]
	per la formula dell'area, per $t \in [0,t_0]$ si ha 
	\[ F(t) = \int_{B(0,1)} \abs{\det(df_t)} d\mathcal{L}^n = \int_{f_t(\B(0,1))} d\mathcal{L}^n = \mathcal{L}^n(f_t(\B(0,1))) = \mathcal{L}^n(\B(0,1)),  \]
	ossia $F$ è costante su $[0,t_0]$. Essendo, per definizione di determinante, $F$ un polinomio in $t$, allora è costante, quindi per ogni $t \in [0,1]$
	\[ F(t) = \mathcal{L}^n(\B(0,1)). \]
	In particolare, $F(1)> 0$. Ma, per ogni $x \in \overline{\B(0,1)}$ si ha $f_1(x) = f(x) \in \partial \B(0,1)$, quindi 
	\[ (f_1(x) | f_1(x)) = \abs{f_1(x)}^2=1. \]
	Preso $v \in \R^n$ tale che $x+tv \in \overline{\B(0,1)}$, 
	\[ 2(df_1(x)v | f_1(x)) = \frac{d}{dt} (f_1(x+tv) | f_1(x+tv))|_{t=0} = \frac{d}{dt} 1|_{t=0} = 0. \]
	Allora, $\textup{Im}(df_1(x)) \subseteq \textup{span}(f(x))^\perp$. Siccome $\dim (\textup{span}(f(x))) = 1$, allora per ogni $x \in \B(0,1)$ si ha $\textup{rg}(df_1(x)) \le n-1$. In particolare, per ogni $x \in \B(0,1)$ deve essere $\det(df_1(x)) = 0$. Allora $F(1) = 0$, assurdo.
	
	Ora, per il Teorema di Stone--Weierstrass, esiste una successione $(p_h)$ di applicazioni a componenti polinomiali tale che $p_h \to f$ uniformemente su $\overline{\B(0,1)}$. In particolare, esiste $M>0$ tale che per ogni $h \in \N$
	\[ \abs{p_h(x)} \le \abs{p_h(x) - f(x)} + \abs{f(x)} \le M, \]
	quindi $\norm{p_h} \le M$. Posto $\psi_h(x) = M^{-1} p_h(x)$ si ha $\norm{\psi_h}_\infty \le 1$, ossia $\textup{Im}(\psi_h) \subseteq \overline{\B(0,1)}$. Per assurdo, supponiamo che $\psi_h(x) \ne x$ per ogni $x \in \overline{\B(0,1)}$. Definiamo $f_h: \overline{\B(0,1)} \to \partial \B(0,1)$ di classe $C^1$ tale che $f_h(x)$ sia l'intersezione della semiretta uscente da $\psi_h(x)$, e passante per $x$, con $\partial \B(0,1)$.\footnote{Il fatto che sia di classe $C^1$ viene dal fatto che $\psi_h$ è di classe $C^1$ e la costruzione geometrica non provoca perdita di regolarità.}
	
	\begin{center}
		\begin{tikzpicture}[x=0.5pt,y=0.5pt,yscale=-1,xscale=1]
			%uncomment if require: \path (0,316); %set diagram left start at 0, and has height of 316
			
			%Shape: Circle [id:dp5526126125166915] 
			\draw   (203.5,154.25) .. controls (203.5,94.96) and (251.56,46.9) .. (310.85,46.9) .. controls (370.14,46.9) and (418.2,94.96) .. (418.2,154.25) .. controls (418.2,213.54) and (370.14,261.6) .. (310.85,261.6) .. controls (251.56,261.6) and (203.5,213.54) .. (203.5,154.25) -- cycle ;
			%Shape: Circle [id:dp3136828089910624] 
			\draw  [fill={rgb, 255:red, 0; green, 0; blue, 0 }  ,fill opacity=1 ] (272.4,172.53) .. controls (272.4,172.02) and (272.82,171.6) .. (273.33,171.6) .. controls (273.85,171.6) and (274.27,172.02) .. (274.27,172.53) .. controls (274.27,173.05) and (273.85,173.47) .. (273.33,173.47) .. controls (272.82,173.47) and (272.4,173.05) .. (272.4,172.53) -- cycle ;
			%Shape: Circle [id:dp8432142275934196] 
			\draw  [fill={rgb, 255:red, 0; green, 0; blue, 0 }  ,fill opacity=1 ] (332.27,90.13) .. controls (332.27,89.62) and (332.68,89.2) .. (333.2,89.2) .. controls (333.72,89.2) and (334.13,89.62) .. (334.13,90.13) .. controls (334.13,90.65) and (333.72,91.07) .. (333.2,91.07) .. controls (332.68,91.07) and (332.27,90.65) .. (332.27,90.13) -- cycle ;
			%Straight Lines [id:da2597269370190327] 
			\draw    (357.6,57.2) -- (273.33,172.53) ;
			%Shape: Circle [id:dp9542475070772747] 
			\draw  [fill={rgb, 255:red, 0; green, 0; blue, 0 }  ,fill opacity=1 ] (356.67,57.2) .. controls (356.67,56.68) and (357.08,56.27) .. (357.6,56.27) .. controls (358.12,56.27) and (358.53,56.68) .. (358.53,57.2) .. controls (358.53,57.72) and (358.12,58.13) .. (357.6,58.13) .. controls (357.08,58.13) and (356.67,57.72) .. (356.67,57.2) -- cycle ;
			
			% Text Node
			\draw (275.33,175) node [anchor=north west][inner sep=0.75pt]  {$\psi _{h}( x)$};
			% Text Node
			\draw (337.33,86.4) node [anchor=north west][inner sep=0.75pt]  {$x$};
			% Text Node
			\draw (367.6,35) node [anchor=north west][inner sep=0.75pt]  {$f_{h}( x)$};
			
			
		\end{tikzpicture}
	\end{center}
	Se $x \in \partial \B(0,1)$ risulta $f_h(x)= x$, assurdo.\footnote{La costruzione geometrica effettuata qui gioca un ruolo fondamentale: se avessimo preso l'altra intersezione che la retta passante per $\psi_h(x)$ e $x$ ha con $\partial \B(0,1)$, non avremmo questa proprietà. Inoltre questo modo di arrivare ad un assurdo giustifica l'utilizzo del Teorema di Stone--Weierstrass in quanto necessitiamo di funzioni di classe $C^1$.} Quindi, per ogni $h \in \N$ esiste $x_h \in \overline{\B(0,1)}$ tale che $\psi_h(x_h) = x_h$. Dal fatto che $\overline{\B(0,1)}$ è compatto, a meno di sottosuccessioni, $x_h \to x \in \overline{\B(0,1)}$. Ora, siccome $p_h \to f$ uniformemente, anche $\psi_h \to f$ uniformemente. In particolare, da 
	\[ \abs{\psi_h(x_h) - f(x)} \le \abs{\psi_h(x_h) - f(x_h)} + \abs{f(x_h)-f(x)} \le \norm{\psi_h -f}_\infty + \abs{f(x_h)-f(x)}\]
	e dalla continuità di $f$, 
	\[ \lim\limits_{h} \psi_h(x_h) = f(x), \]
	ma 
	\[ \lim\limits_{h} \psi_h(x_h) = \lim\limits_{h} x_h = x \]
	da cui $f(x) = x$.
	
	Quanto appena dimostrato si generalizza immediatamente al caso $K= \overline{\B(0,R)}$ per $R > 0$. Nel caso generale, sia $R > 0$ tale che $K \subseteq \overline{B(0,R)}$. Essendo $\R^n$ uno spazio di Hilbert e $K$ convesso, chiuso, in quanto compatto, e non vuoto, l'applicazione $P_K: \R^n \to K$, tale che per ogni $x \in \R^n$ $P_K(x)$ è la proiezione ortogonale di $x$ su $K$, è lipschitziana di costante $1$. In particolare è continua. Se consideriamo l'applicazione $f \circ P_K : \overline{\B(0,R)} \to \overline{\B(0,R)}$, anch'essa continua per composizione, per quanto visto precedentemente, esiste $x \in \overline{\B(0,R)}$ tale che $f(P_k(x)) = x$. Allora $x \in \textup{Im}(f)$, ossia $x \in K$, quindi $P_K(x) = x$ che conduce a $f(x) = x$, che è la tesi.
\end{proof}

\begin{osservazione}
	Siano $a<b$ ed $f: [a,b] \to [a,b]$ un'applicazione continua. In questa situazione la dimostrazione può essere molto semplificata, riconducendo il Teorema di Brouwer al Teorema di esistenza degli zeri. Se consideriamo infatti la funzione ausiliaria $\varphi: [a,b] \to [a,b]$ tale che 
	\[ \varphi(x) = f(x) - x. \]
	Chiaramente $\varphi$ è continua e, siccome $\textup{Im}(f) \subseteq [a,b]$, 
	\begin{align*}
		\varphi(b) &= f(b) - b \le 0 & &\textup{e} &  \varphi(a) &= f(a) - a \ge 0.
	\end{align*} 
	Se $\varphi(a) = 0$ (risp. $\varphi(b) = 0$), allora $f(a) = a$ (risp. $f(b) = b$). Se invece $\varphi(a),\varphi(b) \ne 0$, allora $\varphi(a) \varphi(b) < 0$ e, per il Teorema di esistenza degli zeri, esiste $x \in \left] a,b\right[ $ tale che $\varphi(x) = 0$, ossia $f(x) = x$.
\end{osservazione}

\begin{corollario}
	Sia $\overline{\B(0,1)} \subseteq \R^n$. Non esiste alcuna applicazione $f: \overline{\B(0,1)} \to \partial \B(0,1)$ continua tale che per ogni $x \in \partial \B(0,1)$ si abbia $f(x) = x$.
\end{corollario}
\begin{proof}
	Supponiamo, per assurdo, che tale applicazione $f: \overline{\B(0,1)} \to \partial \B(0,1)$ esista. Sia allora $g: \overline{\B(0,1)} \to \overline{\B(0,1)}$ tale che
	\[ g(x) = -f(x). \]
	Per il Teorema di Brouwer, esiste $x \in \overline{\B(0,1)}$ tale che $g(x) = x$. Essendo $\partial \B(0,1)$ simmetrico rispetto a $0$, risulta $\textup{Im}(g) \subseteq \partial \B(0,1)$ allora $x \in \partial \B(0,1)$ e \[ x = g(x) = -f(x) = -x, \] 
	da cui l'assurdo.
\end{proof}

Analizziamo alcune importanti applicazioni del Teorema di Brouwer. La prima è il Teorema di invarianza del dominio, per cui è necessaria una premessa di carattere topologico.
\begin{definizione}
	Consideriamo $(X,\tau)$ uno spazio topologico. Distinguiamo le seguenti proprietà di separazione:
	\begin{itemize}
		\item diciamo che $(X,\tau)$ è $T_0$\index[analitico]{spazio!topologico!$T_0$}\index[simboli]{$T_0$} se per ogni $x,y \in X$ tali che $x\ne y$, esiste $U \in \tau$ tale che $x \in U$ e $y\notin U$,
		\begin{center}
			\begin{tikzpicture}[x=0.75pt,y=0.75pt,yscale=-1,xscale=1]
				%uncomment if require: \path (0,300); %set diagram left start at 0, and has height of 300
				
				%Shape: Circle [id:dp1732517118606005] 
				\draw  [fill={rgb, 255:red, 0; green, 0; blue, 0 }  ,fill opacity=1 ] (205,150) .. controls (205,149.45) and (205.45,149) .. (206,149) .. controls (206.55,149) and (207,149.45) .. (207,150) .. controls (207,150.55) and (206.55,151) .. (206,151) .. controls (205.45,151) and (205,150.55) .. (205,150) -- cycle ;
				%Shape: Polygon Curved [id:ds020722407137701016] 
				\draw  [dash pattern={on 0.84pt off 2.51pt}] (160,119) .. controls (149,85.6) and (206,95.4) .. (250,119) .. controls (294,142.6) and (230,149) .. (250,179) .. controls (270,209) and (185,218.6) .. (165,188.6) .. controls (145,158.6) and (171,152.4) .. (160,119) -- cycle ;
				%Shape: Circle [id:dp19705533391128083] 
				\draw  [fill={rgb, 255:red, 0; green, 0; blue, 0 }  ,fill opacity=1 ] (315,156) .. controls (315,155.45) and (315.45,155) .. (316,155) .. controls (316.55,155) and (317,155.45) .. (317,156) .. controls (317,156.55) and (316.55,157) .. (316,157) .. controls (315.45,157) and (315,156.55) .. (315,156) -- cycle ;
				
				% Text Node
				\draw (191,132.4) node [anchor=north west][inner sep=0.75pt]    {$x$};
				% Text Node
				\draw (220,84.4) node [anchor=north west][inner sep=0.75pt]    {$U$};
				% Text Node
				\draw (301,138.4) node [anchor=north west][inner sep=0.75pt]    {$y$};
				
				
			\end{tikzpicture}
		\end{center}
		\item diciamo che $(X,\tau)$ è $T_1$\index[analitico]{spazio!topologico!$T_1$}\index[simboli]{$T_1$} se per ogni $x,y \in X$ tali che $x\ne y$, esistono $U,V \in \tau$ tali che $x \in U$, $y \notin U$, $x \notin V$ e $y \in V$,
		\begin{center}
			\begin{tikzpicture}[x=0.75pt,y=0.75pt,yscale=-1,xscale=1]
				%uncomment if require: \path (0,300); %set diagram left start at 0, and has height of 300
				
				%Shape: Circle [id:dp1732517118606005] 
				\draw  [fill={rgb, 255:red, 0; green, 0; blue, 0 }  ,fill opacity=1 ] (205,150) .. controls (205,149.45) and (205.45,149) .. (206,149) .. controls (206.55,149) and (207,149.45) .. (207,150) .. controls (207,150.55) and (206.55,151) .. (206,151) .. controls (205.45,151) and (205,150.55) .. (205,150) -- cycle ;
				%Shape: Polygon Curved [id:ds020722407137701016] 
				\draw  [dash pattern={on 0.84pt off 2.51pt}] (160,119) .. controls (149,85.6) and (206,95.4) .. (250,119) .. controls (294,142.6) and (230,149) .. (250,179) .. controls (270,209) and (185,218.6) .. (165,188.6) .. controls (145,158.6) and (171,152.4) .. (160,119) -- cycle ;
				%Shape: Circle [id:dp19705533391128083] 
				\draw  [fill={rgb, 255:red, 0; green, 0; blue, 0 }  ,fill opacity=1 ] (315,156) .. controls (315,155.45) and (315.45,155) .. (316,155) .. controls (316.55,155) and (317,155.45) .. (317,156) .. controls (317,156.55) and (316.55,157) .. (316,157) .. controls (315.45,157) and (315,156.55) .. (315,156) -- cycle ;
				%Shape: Polygon Curved [id:ds3453941623544552] 
				\draw  [dash pattern={on 0.84pt off 2.51pt}] (281,117) .. controls (357,83.2) and (358,79.8) .. (371,117) .. controls (384,154.2) and (413,162.2) .. (371,177) .. controls (329,191.8) and (306,216.6) .. (286,186.6) .. controls (266,156.6) and (205,150.8) .. (281,117) -- cycle ;
				
				% Text Node
				\draw (191,132.4) node [anchor=north west][inner sep=0.75pt]    {$x$};
				% Text Node
				\draw (220,84.4) node [anchor=north west][inner sep=0.75pt]    {$U$};
				% Text Node
				\draw (301,138.4) node [anchor=north west][inner sep=0.75pt]    {$y$};
				% Text Node
				\draw (373,91.4) node [anchor=north west][inner sep=0.75pt]    {$V$};
				
				
			\end{tikzpicture}
		\end{center}
		\item diciamo che $(X,\tau)$ è $T_2$, o \emph{di Hausdorff},\index[analitico]{spazio!topologico!$T_2$}\index[analitico]{spazio!topologico!di Hausdorff}\index[simboli]{$T_2$} se per ogni $x,y \in X$ tali che $x\ne y$, esistono $U,V \in \tau$ tali che $x \in U$, $y \in V$ e $U \cap V = \emptyset$,
		\begin{center}
			\begin{tikzpicture}[x=0.75pt,y=0.75pt,yscale=-1,xscale=1]
				%uncomment if require: \path (0,300); %set diagram left start at 0, and has height of 300
				
				%Shape: Circle [id:dp1732517118606005] 
				\draw  [fill={rgb, 255:red, 0; green, 0; blue, 0 }  ,fill opacity=1 ] (205,150) .. controls (205,149.45) and (205.45,149) .. (206,149) .. controls (206.55,149) and (207,149.45) .. (207,150) .. controls (207,150.55) and (206.55,151) .. (206,151) .. controls (205.45,151) and (205,150.55) .. (205,150) -- cycle ;
				%Shape: Polygon Curved [id:ds020722407137701016] 
				\draw  [dash pattern={on 0.84pt off 2.51pt}] (160,119) .. controls (149,85.6) and (206,95.4) .. (250,119) .. controls (294,142.6) and (230,149) .. (250,179) .. controls (270,209) and (185,218.6) .. (165,188.6) .. controls (145,158.6) and (171,152.4) .. (160,119) -- cycle ;
				%Shape: Circle [id:dp19705533391128083] 
				\draw  [fill={rgb, 255:red, 0; green, 0; blue, 0 }  ,fill opacity=1 ] (315,156) .. controls (315,155.45) and (315.45,155) .. (316,155) .. controls (316.55,155) and (317,155.45) .. (317,156) .. controls (317,156.55) and (316.55,157) .. (316,157) .. controls (315.45,157) and (315,156.55) .. (315,156) -- cycle ;
				%Shape: Polygon Curved [id:ds3453941623544552] 
				\draw  [dash pattern={on 0.84pt off 2.51pt}] (293,130.2) .. controls (319,73.4) and (358,79.8) .. (371,117) .. controls (384,154.2) and (413,162.2) .. (371,177) .. controls (329,191.8) and (265,214.2) .. (286,186.6) .. controls (307,159) and (267,187) .. (293,130.2) -- cycle ;
				
				% Text Node
				\draw (191,132.4) node [anchor=north west][inner sep=0.75pt]    {$x$};
				% Text Node
				\draw (220,84.4) node [anchor=north west][inner sep=0.75pt]    {$U$};
				% Text Node
				\draw (301,138.4) node [anchor=north west][inner sep=0.75pt]    {$y$};
				% Text Node
				\draw (373,91.4) node [anchor=north west][inner sep=0.75pt]    {$V$};
				
				
			\end{tikzpicture}
		\end{center}
		\item diciamo che $(X,\tau)$ è $T_3$, o \emph{regolare},\index[analitico]{spazio!topologico!$T_3$}\index[analitico]{spazio!topologico!regolare}\index[simboli]{$T_3$} se per ogni $x,y \in X$ tali che $x\ne y$, per ogni $C \subseteq X$ chiuso tale che $x \in C$ e $y \notin C$, esistono $U,V \in \tau$ tali che $C \subseteq U$, $y \in V$ e $U \cap V = \emptyset$,
		\begin{center}
			\begin{tikzpicture}[x=0.75pt,y=0.75pt,yscale=-1,xscale=1]
				%uncomment if require: \path (0,300); %set diagram left start at 0, and has height of 300
				
				%Shape: Circle [id:dp1732517118606005] 
				\draw  [fill={rgb, 255:red, 0; green, 0; blue, 0 }  ,fill opacity=1 ] (205,150) .. controls (205,149.45) and (205.45,149) .. (206,149) .. controls (206.55,149) and (207,149.45) .. (207,150) .. controls (207,150.55) and (206.55,151) .. (206,151) .. controls (205.45,151) and (205,150.55) .. (205,150) -- cycle ;
				%Shape: Polygon Curved [id:ds020722407137701016] 
				\draw  [dash pattern={on 0.84pt off 2.51pt}] (160,119) .. controls (149,85.6) and (206,95.4) .. (250,119) .. controls (294,142.6) and (230,149) .. (250,179) .. controls (270,209) and (185,218.6) .. (165,188.6) .. controls (145,158.6) and (171,152.4) .. (160,119) -- cycle ;
				%Shape: Circle [id:dp19705533391128083] 
				\draw  [fill={rgb, 255:red, 0; green, 0; blue, 0 }  ,fill opacity=1 ] (315,156) .. controls (315,155.45) and (315.45,155) .. (316,155) .. controls (316.55,155) and (317,155.45) .. (317,156) .. controls (317,156.55) and (316.55,157) .. (316,157) .. controls (315.45,157) and (315,156.55) .. (315,156) -- cycle ;
				%Shape: Polygon Curved [id:ds3453941623544552] 
				\draw  [dash pattern={on 0.84pt off 2.51pt}] (293,130.2) .. controls (319,73.4) and (358,79.8) .. (371,117) .. controls (384,154.2) and (413,162.2) .. (371,177) .. controls (329,191.8) and (265,214.2) .. (286,186.6) .. controls (307,159) and (267,187) .. (293,130.2) -- cycle ;
				%Shape: Polygon Curved [id:ds4009492428387247] 
				\draw   (180,133) .. controls (212,97.8) and (255,112.8) .. (247,147) .. controls (239,181.2) and (223,147.2) .. (224,177.2) .. controls (225,207.2) and (160,192) .. (172,168) .. controls (184,144) and (148,168.2) .. (180,133) -- cycle ;
				
				% Text Node
				\draw (191,132.4) node [anchor=north west][inner sep=0.75pt]    {$x$};
				% Text Node
				\draw (220,84.4) node [anchor=north west][inner sep=0.75pt]    {$U$};
				% Text Node
				\draw (301,138.4) node [anchor=north west][inner sep=0.75pt]    {$y$};
				% Text Node
				\draw (373,91.4) node [anchor=north west][inner sep=0.75pt]    {$V$};
				% Text Node
				\draw (226,180.6) node [anchor=north west][inner sep=0.75pt]    {$C$};
				
				
			\end{tikzpicture}
		\end{center}
		\item diciamo che $(X,\tau)$ è $T_4$, o \emph{normale},\index[analitico]{spazio!topologico!$T_4$}\index[analitico]{spazio!topologico!normale}\index[simboli]{$T_4$} se per ogni $x,y \in X$ tali che $x\ne y$, per ogni $C,K \subseteq X$ chiusi tale che $x \in C$ e $y \notin C$, $x \notin K$ e $y \in K$, esistono $U,V \in \tau$ tali che $C \subseteq U$, $K \subseteq V$ e $U \cap V = \emptyset$.
		\begin{center}
			\begin{tikzpicture}[x=0.75pt,y=0.75pt,yscale=-1,xscale=1]
				%uncomment if require: \path (0,300); %set diagram left start at 0, and has height of 300
				
				%Shape: Circle [id:dp1732517118606005] 
				\draw  [fill={rgb, 255:red, 0; green, 0; blue, 0 }  ,fill opacity=1 ] (205,150) .. controls (205,149.45) and (205.45,149) .. (206,149) .. controls (206.55,149) and (207,149.45) .. (207,150) .. controls (207,150.55) and (206.55,151) .. (206,151) .. controls (205.45,151) and (205,150.55) .. (205,150) -- cycle ;
				%Shape: Polygon Curved [id:ds020722407137701016] 
				\draw  [dash pattern={on 0.84pt off 2.51pt}] (160,119) .. controls (149,85.6) and (206,95.4) .. (250,119) .. controls (294,142.6) and (230,149) .. (250,179) .. controls (270,209) and (185,218.6) .. (165,188.6) .. controls (145,158.6) and (171,152.4) .. (160,119) -- cycle ;
				%Shape: Circle [id:dp19705533391128083] 
				\draw  [fill={rgb, 255:red, 0; green, 0; blue, 0 }  ,fill opacity=1 ] (315,156) .. controls (315,155.45) and (315.45,155) .. (316,155) .. controls (316.55,155) and (317,155.45) .. (317,156) .. controls (317,156.55) and (316.55,157) .. (316,157) .. controls (315.45,157) and (315,156.55) .. (315,156) -- cycle ;
				%Shape: Polygon Curved [id:ds3453941623544552] 
				\draw  [dash pattern={on 0.84pt off 2.51pt}] (293,130.2) .. controls (319,73.4) and (358,79.8) .. (371,117) .. controls (384,154.2) and (413,162.2) .. (371,177) .. controls (329,191.8) and (265,214.2) .. (286,186.6) .. controls (307,159) and (267,187) .. (293,130.2) -- cycle ;
				%Shape: Polygon Curved [id:ds4009492428387247] 
				\draw   (180,133) .. controls (212,97.8) and (255,112.8) .. (247,147) .. controls (239,181.2) and (223,147.2) .. (224,177.2) .. controls (225,207.2) and (160,192) .. (172,168) .. controls (184,144) and (148,168.2) .. (180,133) -- cycle ;
				%Shape: Polygon Curved [id:ds362325771433208] 
				\draw   (313,135.2) .. controls (333,134.2) and (342,114.2) .. (350,131.2) .. controls (358,148.2) and (353,124.2) .. (373,154.2) .. controls (393,184.2) and (304,186.8) .. (299,157) .. controls (294,127.2) and (293,136.2) .. (313,135.2) -- cycle ;
				
				% Text Node
				\draw (191,132.4) node [anchor=north west][inner sep=0.75pt]    {$x$};
				% Text Node
				\draw (220,84.4) node [anchor=north west][inner sep=0.75pt]    {$U$};
				% Text Node
				\draw (301,138.4) node [anchor=north west][inner sep=0.75pt]    {$y$};
				% Text Node
				\draw (373,91.4) node [anchor=north west][inner sep=0.75pt]    {$V$};
				% Text Node
				\draw (226,180.6) node [anchor=north west][inner sep=0.75pt]    {$C$};
				% Text Node
				\draw (323,108.6) node [anchor=north west][inner sep=0.75pt]    {$K$};
				
				
			\end{tikzpicture}
		\end{center}
	\end{itemize}
\end{definizione}
\begin{lemma}\emph{\textbf{(di Urysohm)}}\index[analitico]{lemma!di Urysohm} Consideriamo $(X,\tau)$ uno spazio topologico. I seguenti fatti sono equivalenti:
	\begin{enumerate}[$(a)$]
		\item $(X,\tau)$ è normale,
		\item per ogni $C,K \subseteq X$ chiusi e disgiunti esiste  un'applicazione $f: X \to [0,1]$ continua tale che $f= 0$ su $C$ ed $f=1$ su $K$.
	\end{enumerate}
\end{lemma}
\begin{proof}
	Omettiamo la dimostrazione.
\end{proof}
\begin{osservazione}
	Il Lemma di Urysohm è considerato storicamente il primo risultato di topologia generale con una dimostrazione non banale: la difficoltà sta nelle ipotesi di lavoro molto generali. Se ci mettiamo però in un contesto metrico, basta considerare
	\[ f(x) = \frac{\textup{dist}(x,C)}{\textup{dist}(x,C)+ \textup{dist}(x,K)}. \]
\end{osservazione}
\begin{teorema}\emph{\textbf{(di Tietze)}}\index[analitico]{teorema!di Tietze}
	Consideriamo $(X,\tau)$ uno spazio topologico normale, $A \subseteq X$ chiuso e $f: A \to \R$ un'applicazione continua. Esiste un'applicazione $\tilde{f}: X \to \R$ continua tale che $\tilde{f}|_A = f$.
\end{teorema}
\begin{proof}
	Mostriamo innanzitutto che se esiste $C>0$ tale che, per ogni $x \in A$, $\abs{f(x)}\le C$, allora esiste un'applicazione $g:X\to \R$ continua tale che, per ogni $x \in X$, $\abs{g(x)}\le \frac{1}{3}C$ e, per ogni $x \in A$, $\abs{f(x)-g(x)}\le \frac{2}{3}C$. Siano $Y = f^{-1}([-C,-\frac{1}{3}C])$ $Z=f^{-1}([\frac{1}{3}C,C])$.
	\begin{center}
		\begin{tikzpicture}[x=0.75pt,y=0.75pt,yscale=-1,xscale=1]
			%uncomment if require: \path (0,300); %set diagram left start at 0, and has height of 300
			
			%Straight Lines [id:da351590380041958] 
			\draw    (180,110.4) -- (440,110.4) ;
			\draw [shift={(442,110.4)}, rotate = 180] [color={rgb, 255:red, 0; green, 0; blue, 0 }  ][line width=0.75]    (10.93,-3.29) .. controls (6.95,-1.4) and (3.31,-0.3) .. (0,0) .. controls (3.31,0.3) and (6.95,1.4) .. (10.93,3.29)   ;
			%Straight Lines [id:da6689641897087604] 
			\draw    (310,104.6) -- (310,116.6) ;
			%Straight Lines [id:da03212271146730217] 
			\draw    (340,104.6) -- (340,116.6) ;
			%Straight Lines [id:da4012489765266216] 
			\draw    (400,104.6) -- (400,116.6) ;
			%Straight Lines [id:da5134778826508875] 
			\draw    (280,104.6) -- (280,116.6) ;
			%Straight Lines [id:da8281691966789628] 
			\draw    (220,104.6) -- (220,116.6) ;
			%Straight Lines [id:da30953233561890436] 
			\draw [color={rgb, 255:red, 0; green, 49; blue, 83 }  ,draw opacity=1 ][fill={rgb, 255:red, 0; green, 49; blue, 83 }  ,fill opacity=1 ]   (220,110.6) -- (280,110.6) ;
			%Straight Lines [id:da7267107874600898] 
			\draw [color={rgb, 255:red, 0; green, 49; blue, 83 }  ,draw opacity=1 ][fill={rgb, 255:red, 0; green, 49; blue, 83 }  ,fill opacity=1 ]   (340,110.6) -- (400,110.6) ;
			
			% Text Node
			\draw (305,86.4) node [anchor=north west][inner sep=0.75pt]    {$0$};
			% Text Node
			\draw (328,86.4) node [anchor=north west][inner sep=0.75pt]    {$\frac{1}{3} C$};
			% Text Node
			\draw (394,86.4) node [anchor=north west][inner sep=0.75pt]    {$C$};
			% Text Node
			\draw (251,86.4) node [anchor=north west][inner sep=0.75pt]    {$-\frac{1}{3} C$};
			% Text Node
			\draw (206,86.4) node [anchor=north west][inner sep=0.75pt]    {$-C$};
			% Text Node
			\draw (235,120.4) node [anchor=north west][inner sep=0.75pt]    {$f( Y)$};
			% Text Node
			\draw (352,121.4) node [anchor=north west][inner sep=0.75pt]    {$f( Z)$};
			
			
		\end{tikzpicture}
	\end{center}
	Siccome $f$ è continua, $Y$ e $Z$ sono chiusi in $A$, quindi lo sono in $X$, essendo $A$ chiuso. Sono inoltre disgiunti: infatti, se, per assurdo, esistesse $x \in Y \cap Z$, allora $f(x) \in [-C,-\frac{1}{3}C] \cap [\frac{1}{3}C,C] = \emptyset$, da cui la contraddizione. Per il Lemma di Urysohm, esiste $h: X \to [0,1]$ continua tale che $h=0$ su $Y$ e $h=1$ su $Z$. Mostriamo che $g: X \to \R$ tale che 
	\[ g(x) = \frac{2}{3}C\left( h(x)-\frac{1}{2}\right)  \]
	soddisfa i requisiti richiesti. In primo luogo, è chiaramente continua. Poi, per ogni $x \in X$
	\begin{align*}
		g(x) &\le  \frac{1}{3}C, & g(x) \ge - \frac{1}{3}C,
	\end{align*}
	quindi $\abs{g(x)}\le \frac{1}{3}C$. Ora, dal fatto che $A = Z \cup Y \cup f^{-1}([-\frac{C}{3},\frac{C}{3}])$, possiamo distinguere tre casi: se $x \in Y$, 
	\begin{align*}
		f(x) - g(x) &=  f(x) + \frac{1}{3}C \le 0 \le \frac{2}{3}C, & f(x) - g(x) \ge -C + \frac{1}{3}C = -\frac{2}{3}C;
	\end{align*}
	se $x \in Z$,
	\begin{align*}
		f(x) - g(x) &\le C - \frac{1}{3}C = \frac{2}{3}C, & f(x) - g(x) \ge \frac{1}{3}C - \frac{1}{3}C = 0 \ge -\frac{2}{3}C;
	\end{align*}
	se $x\in f^{-1}([-\frac{C}{3},\frac{C}{3}])$, ricordando che $h(x) \in [0,1]$,
	\begin{align*}
		f(x) - g(x) &\le \frac{C}{3} \le \frac{2}{3}C, & f(x) - g(x) \ge -\frac{2}{3}C.
	\end{align*}
	In ogni caso, $\abs{f(x)-g(x)}\le \frac{2}{3}C$. 
	
	Vediamo, dapprima, il caso $f:A \to [a,b]$, con $a<b \in \R$, continua. A meno di operare una traslazione ed un riscalamento, possiamo supporre $[a,b] = [0,1]$. Costruiamo ricorsivamente delle applicazioni $g_n: X \to \R$ continue tali che 
	\begin{align*}
		\forall\; x \in X: \;\abs{g_n(x)} &\le \frac{1}{3}\left( \frac{2}{3}\right) ^{n-1},\\
		\forall \; x \in A: \; \abs*{f(x)- \sum_{i=1}^{n}g_i(x)} &\le \left( \frac{2}{3}\right)^{n}.
	\end{align*}
	Siccome $\abs{f(x)}\le 1$ per ogni $x \in A$, l'esistenza di $g_1$ segue dal passaggio precedente. Supponiamo quindi di possedere $g_1,\dots,g_{n}$ e applichiamo il passaggio precedente a  
	\[ f_n =f- \sum_{i=1}^n g_i\]
	con $C = \left( \frac{2}{3}\right)^{n}$. Allora esiste $g_{n+1}: X \to \R$ continua tale che
	\begin{align*}
		\forall\; x \in X: \;\abs{g_{n+1}(x)} &\le \frac{1}{3}\left( \frac{2}{3}\right) ^{n},\\
		\forall \; x \in A: \; \abs*{f(x)- \sum_{i=1}^{n+1}g_i(x)} &\le \abs*{f(x)- \sum_{i=1}^{n}g_i(x) - g_{n+1}(x)} = \abs*{f_n(x)- g_{n+1}(x)}\le  \left( \frac{2}{3}\right)^{n+1}.
	\end{align*}
	Detta 
	\[ S_n(x) = \sum_{i=1}^{n}g_i(x), \]
	osserviamo che 
	\[ \sum_{i=1}^{\infty}\norm{g_i}_\infty \le \frac{1}{3} \sum_{i=1}^{\infty}\left( \frac{2}{3}\right)^{i-1} = 1, \]
	quindi la serie 
	\[ \sum_{i=1}^{\infty}g_i(x) \]
	è normalmente convergente in $C(X;\R)$, che è uno spazio di Banach. Allora esiste $\bar{f} \in  C(X;\R)$ tale che $S_n \to \bar{f}$ uniformemente\footnote{A voler essere precisi, dalla convergenza normale tiriamo fuori la convergenza uniforme e il limite uniforme di funzioni continue è continuo.}. Dal fatto che per ogni $x \in A$
	\[ \abs*{f(x)- \sum_{i=1}^{n}g_i(x)} \le \left( \frac{2}{3}\right)^{n}, \]
	segue che $\bar{f}|_A = f$. In particolare, per ogni $x \in X$,
	\[\bar{f}(x) \le \sum_{i=1}^{\infty}\norm{g_i}_\infty \le 1. \]
	A questo punto, basta considerare $\tilde{f} = (\bar{f})^+$. Chiaramente è una funzione continua. Poi, $0\le \tilde{f} \le 1$ e, su $ A$, $\bar{f}=f$, quindi $\tilde{f} = f^+ = f$, essendo $f\ge 0$, e possiamo concludere che $\tilde{f}$ soddisfa i requisiti richiesti.
	
	Vediamo, poi, il caso $f:A \to \left] -1,1\right[ $ continua. Siccome $\left] -1,1\right[ \subseteq [-1,1]$, dal passaggio precedente, esiste $\overline{f}: X \to [-1,1]$ tale che $\overline{f}|_A = f$. Sia $B= \overline{f}^{-1}(\left\lbrace -1,1\right\rbrace )$, chiuso in $X$ perché $\overline{f}$ continua e $A$ chiuso. Siccome $f$ è a valori in $\left] -1,1\right[$, allora $A \cap B = \emptyset$, quindi per il Lemma di Urysohm, esiste $k : X \to [0,1]$ tale che $k=1$ in $A$ e $k=0$ in $B$. Vediamo che l'applicazione $\tilde{f} = k\overline{f}: X \to \left] -1,1\right[$ soddisfa i requisiti richiesti. Innanzitutto, se $\overline{f}(x) \in \left] -1,1\right[ $, allora $\tilde{f} \in \left] -1,1\right[$. Invece, se $\overline{f}(x) \in \left\lbrace -1,1\right\rbrace$, $\tilde{f}(x) = 0 \in \left] -1,1\right[$, in quanto $k=0$ in $B$. In particolare, $\tilde{f}$ è ben definita. Chiaramente è poi continua e $\tilde{f}|_A = f$, in quanto $k = 1$ in $A$.
	
	Nel caso generale, sia $\Phi: \R \to \left] -1,1\right[$ un omeomorfismo.\footnote{Possiamo, ad esempio, prendere $\Phi^{-1}:\left] -1,1\right[ \to \R$ tale che
		\[ \Phi^{-1}(x)= \begin{cases}
			\frac{x}{x+1} & \textup{se } -1<x\le 0,\\
			\frac{x}{1-x} & \textup{se } 0\le x< 1,\\
		\end{cases} \]
		oppure un opportuno riscalamento della funzione $\arctan$.
	} Applicando il caso precedente a  $ \Phi\circ f: A \to \left] -1,1\right[$, segue l'esistenza di $\overline{f}: X \to \left] -1,1\right[$ tale che $\overline{f}|_A = \Phi\circ f$. A questo punto basta considerare $\tilde{f} = \Phi^{-1}\circ \overline{f}$.
\end{proof}

\begin{teorema}\textbf{\emph{(di invarianza del dominio)}}\index[analitico]{teorema!di invarianza!del dominio}
	Sia $U \subseteq \R^n$ aperto. Se $f: U \to \R^n$ è un'applicazione iniettiva e continua, allora $f(U)$ è aperto. In particolare, $f: U \to f(U)$ è un omeomorfismo.
\end{teorema}
\begin{proof}
	È sufficiente dimostrare che se  $f: \overline{\B(0,1)} \to \R^n$ è un'applicazione iniettiva e continua, allora $f(0) \in \textup{int}(f(\overline{\B(0,1)}))$. Infatti, la tesi corrisponde a richiedere che
	\[ y \in f(U) \Longrightarrow y \in \textup{int}(f(U)), \]
	ossia
	\[ x_0 \in U \Longrightarrow f(x_0) \in \textup{int}(f(U)), \]
	ma, a meno di una traslazione, possiamo sempre supporre che $x_0=0$, quindi la tesi sarebbe
	\[ f(0) \in \textup{int}(f(U)). \]
	Ora, siccome $U$ è aperto, a meno di un riscalamento, $0 \in \overline{\B(0,1)}\subsetneq U$. %Se $f(x) \in f(\overline{\B(0,1)})$, allora $x \in \overline{\B(0,1)}$, quindi $x \in U$, allora $f(x) \in f(U)$: in altre parole $f(\overline{\B(0,1)}) \subseteq f(U)$. In particolare, $f(\overline{\B(0,1)}) \subsetneq f(U)$: supponiamo, per assurdo, che  $f(\overline{\B(0,1)}) = f(U)$. Siccome $\overline{\B(0,1)}\subsetneq U$, esiste $\xi \in U \setminus \overline{\B(0,1)}$, ma, per ipotesi assurda, trovo anche $\eta \in \overline{\B(0,1)}$ tale che $f(\eta) = f(\xi)$. Dovendo essere, $\xi \ne \eta$, l'assurdo deriva dalla violazione dell'iniettività di $f$. 
	In particolare, siccome, per iniettività, $f(\overline{\B(0,1)}) \subseteq f(U)$, è sufficiente dimostrare che $f(0) \in \textup{int}(f(\overline{\B(0,1)}))$ e possiamo anche pensare direttamente che $f: \overline{\B(0,1)} \to \R^n$.
	
	Supponiamo, per assurdo che $f(0) \notin \textup{int}(f(\overline{\B(0,1)}))$. Allora, deve essere $f(0) \in \partial (f(\overline{\B(0,1)}))$, perché $f(0) \in f(\overline{\B(0,1)})$. In particolare, per ogni $\varepsilon > 0$, esiste $c \in \R^n\setminus f(\overline{\B(0,1)})$ tale che $\abs{c-f(0)}<\varepsilon$. Detti 
	\begin{align*}
		\Sigma_1 &= \left\lbrace y \in f(\overline{\B(0,1)}): \abs{y-c}\ge \varepsilon\right\rbrace,\\
		\Sigma_2 &= \left\lbrace y \in \R^n: \abs{y-c}= \varepsilon\right\rbrace,
	\end{align*}
	consideriamo l'insieme $\Sigma = \Sigma_1 \cup \Sigma_2$: in primo luogo $\Sigma_1$ è chiuso in $f(\overline{\B(0,1)})$\footnote{Perchè il suo complementare è l'intersezione di $f(\overline{\B(0,1)})$ con la palla aperta $\B(c,\varepsilon)$.}, che è compatto, siccome $f$ è continua, quindi $\Sigma_1$ è compatto. Poi $\Sigma_2 = \partial \B(c,\varepsilon)$, quindi anche lui è compatto. Essendo $\Sigma$ l'unione di due compatti, è anch'esso compatto. In più, $f(0) \notin \Sigma$.
	
	\begin{center}
		\begin{tikzpicture}[x=0.70pt,y=0.70pt,yscale=-1,xscale=1]
			%uncomment if require: \path (0,300); %set diagram left start at 0, and has height of 300
			
			%Curve Lines [id:da21018532107485988] 
			\draw [fill={rgb, 255:red, 155; green, 155; blue, 155 }  ,fill opacity=0.33 ]   (122,244.6) .. controls (101,211.6) and (115,79.6) .. (181,127.6) .. controls (247,175.6) and (377,99.6) .. (417.28,91.38) .. controls (457.56,83.16) and (598,151.6) .. (574,243.6) ;
			%Shape: Circle [id:dp004821536589640019] 
			\draw  [fill={rgb, 255:red, 0; green, 0; blue, 0 }  ,fill opacity=1 ] (294.4,135.3) .. controls (294.4,134.58) and (294.98,134) .. (295.7,134) .. controls (296.42,134) and (297,134.58) .. (297,135.3) .. controls (297,136.02) and (296.42,136.6) .. (295.7,136.6) .. controls (294.98,136.6) and (294.4,136.02) .. (294.4,135.3) -- cycle ;
			%Shape: Circle [id:dp8290836366492667] 
			\draw   (237.7,136.6) .. controls (237.7,104.57) and (263.67,78.6) .. (295.7,78.6) .. controls (327.73,78.6) and (353.7,104.57) .. (353.7,136.6) .. controls (353.7,168.63) and (327.73,194.6) .. (295.7,194.6) .. controls (263.67,194.6) and (237.7,168.63) .. (237.7,136.6) -- cycle ;
			%Shape: Circle [id:dp2809219934901368] 
			\draw  [fill={rgb, 255:red, 0; green, 0; blue, 0 }  ,fill opacity=1 ] (303.4,95.3) .. controls (303.4,94.58) and (303.98,94) .. (304.7,94) .. controls (305.42,94) and (306,94.58) .. (306,95.3) .. controls (306,96.02) and (305.42,96.6) .. (304.7,96.6) .. controls (303.98,96.6) and (303.4,96.02) .. (303.4,95.3) -- cycle ;
			%Shape: Circle [id:dp9678593963925135] 
			\draw   (246.7,96.6) .. controls (246.7,64.57) and (272.67,38.6) .. (304.7,38.6) .. controls (336.73,38.6) and (362.7,64.57) .. (362.7,96.6) .. controls (362.7,128.63) and (336.73,154.6) .. (304.7,154.6) .. controls (272.67,154.6) and (246.7,128.63) .. (246.7,96.6) -- cycle ;
			%Straight Lines [id:da42967896812356576] 
			\draw  [dash pattern={on 0.84pt off 2.51pt}]  (295.7,194.6) -- (295.7,136.6) ;
			%Straight Lines [id:da2746572242277334] 
			\draw  [dash pattern={on 0.84pt off 2.51pt}]  (304.7,94) -- (310,39.6) ;
			
			% Text Node
			\draw (297.7,137.4) node [anchor=north west][inner sep=0.75pt]  [font=\footnotesize]  {$f( 0)$};
			% Text Node
			\draw (194,175.4) node [anchor=north west][inner sep=0.75pt]    {$\Sigma _{1}$};
			% Text Node
			\draw (464,197.4) node [anchor=north west][inner sep=0.75pt]    {$f(\overline{\B(0,1)})$};
			% Text Node
			\draw (305.4,98.7) node [anchor=north west][inner sep=0.75pt]    {$c$};
			% Text Node
			\draw (297.7,169) node [anchor=north west][inner sep=0.75pt]    {$\varepsilon $};
			% Text Node
			\draw (291.7,52) node [anchor=north west][inner sep=0.75pt]    {$\varepsilon $};
			% Text Node
			\draw (213,48.4) node [anchor=north west][inner sep=0.75pt]    {$\Sigma _{2}$};
		\end{tikzpicture}
	\end{center}
	
	Mostriamo che $f: \overline{\B(0,1)} \to f(\overline{\B(0,1)})$ è un omeomorfismo. Chiaramente è biettiva, quindi basta mostrare che $f^{-1}:f(\overline{\B(0,1)}) \to \overline{\B(0,1)}$ è continua. Siano quindi $y \in f(\overline{\B(0,1)})$ e $(y_h)$ in $f(\overline{\B(0,1)})$ tali che $y_h \to y$. Data $(f^{-1}(y_h))$, consideriamo una qualsiasi sua sottosuccessione $(f^{-1}(y_{h_k}))$. Sia $(x_{h_k})$ in $\overline{\B(0,1)}$ tale che $y_{h_k} = f(x_{h_k})$. Per compattezza, esistono $x \in \overline{\B(0,1)}$ e $(x_{h_{k_j}})$ tali che $x_{h_{k_j}} \to x$. Essendo $f$ continua, allora $f(x_{h_{k_j}}) \to f(x)$, ossia $y_{h_{k_j}} \to f(x)$ e, per unicità del limite, $y = f(x)$ o, in altre parole, $x = f^{-1}(y)$. Quindi $f^{-1}(y_{h_{k_j}}) \to f^{-1}(y)$. Allora $f^{-1}(y_h)\to f^{-1}(y)$, quindi $f^{-1}$ è continua.
	
	A questo punto, essendo, per la continuità di $f$, $f(\overline{\B(0,1)})$ compatto, quindi completo, quindi chiuso, per il Teorema di Tietze, applicato componente per componente, esiste un'applicazione $g: \R^n \to \R^n$ continua tale che $g = f^{-1}$ su $f(\overline{\B(0,1)})$. 
	
	Osserviamo che $g \ne 0$ su $\Sigma_1$: infatti se, per assurdo, esiste $y \in \Sigma_1 \subseteq f(\overline{\B(0,1)})$ tale che $g(y)=0$, allora esiste $x \in \overline{\B(0,1)}$ tale che $y=f(x)$, ma $0= g(y) = f^{-1}(y) = x$, ossia $y = f(0) \in \Sigma_1$, assurdo. Quindi, per la compattezza di $\Sigma_1$ e la continuità di $g$, esiste $\delta >0$, che, senza perdita di generalità, possiamo supporre $\delta < \frac{1}{2}$, tale che per ogni $y \in \Sigma_1$
	\[ \abs{g(y)} \ge \delta. \]
	Per il Teorema di Stone--Weierstrass, essendo $\Sigma$ compatto, esiste un'applicazione a componenti polinomiali $p:\R^n \to \R^n$ tale che per ogni $y \in \Sigma$
	\[ \abs{p(y)-g(y)}< \frac{\delta}{2}. \]
	Osserviamo che $p \ne 0$ su $\Sigma_1$\footnote{Potrebbe però annullarsi su $\Sigma_2$.}: se così non fosse, esisterebbe $y \in \Sigma_1$ tale che $\abs{g(y)}<\frac{\delta}{2}$, assurdo.
	
	Mostriamo che esiste $a_0 \in \B(0,\frac{\delta}{2})\setminus p(\Sigma_2)$. Innanzitutto, $p$ ha componenti polinomiali, quindi è almeno di classe $C^1$. In particolare, $p$ è lipschitziana su $\B(c,2\varepsilon)$ con costante di Lipschitz $A$. Allora,
	\[ \mathcal{L}^n(p(\Sigma_2))\le A^n \mathcal{L}^n(\Sigma_2) =0. \]
	Se, per assurdo, per ogni $a \in \B(0,\frac{\delta}{2})$ si avesse $a \in p(\Sigma_2)$, allora $\B(0,\frac{\delta}{2}) \subseteq p(\Sigma_2)$, quindi, per la monotonia della misura di Lebesgue,
	\[ \mathcal{L}^n(\B(0,\frac{\delta}{2})) \le \mathcal{L}^n(p(\Sigma_2)) = 0, \]
	assurdo.
	
	Poniamo $\tilde{p} = p - a_0$. Osserviamo che $\tilde{p} \ne 0$ su $\Sigma$: innanzitutto, per ogni $y \in \Sigma$
	\[ \abs{\tilde{p}(y)-g(y)} \le \abs{p(y)-g(y)} + \abs{a_0} < \delta, \]
	quindi, come sopra, $\tilde{p} \ne 0$ su $\Sigma_1$. Inoltre, se, per assurdo, esistesse $y \in \Sigma_2$ tale che $\tilde{p}(y) = 0$, allora avremmo $a_0 = p(y) \in p(\Sigma_2)$, assurdo.
	
	Consideriamo $\Phi: f(\overline{\B(0,1)}) \to \Sigma$ tale che
	\[ f(y) = \begin{cases}
		y & \textup{se } \abs{y-c}\ge \varepsilon,\\
		c + \varepsilon \frac{y-c}{\abs{y-c}} & \textup{se } \abs{y-c}<\varepsilon.
	\end{cases} \]
	Anzitutto, $\Phi$ è ben definita: $c \notin f(\overline{\B(0,1)})$, quindi il denominatore non si annulla mai; inoltre, se $\abs{y-c} \ge \varepsilon$, $\Phi(y) = y \in \Sigma_1 \subseteq \Sigma$ mentre se $\abs{y-c} < \varepsilon$, allora 
	\[ \abs{\Phi(y)-c} = \varepsilon,\]
	quindi $\Phi(y) \in \Sigma_2 \subseteq \Sigma$. In più, $\Phi$ è continua: se $\abs{y-c} = \varepsilon$, allora
	\[ c + \varepsilon \frac{y-c}{\abs{y-c}} = c + y - c = y. \]
	
	Presa $h = \tilde{p} \circ \Phi: f(\overline{\B(0,1)}) \to \R^n$, continua per composizione, si ha che $h \ne 0$. Mostriamo che per ogni $x \in \overline{\B(0,1)}$ 
	\[ \abs{h(f(x))-x}\le 1. \]
	Per continuità di $g$ in $f(0)$, scegliamo $\varepsilon$ tale che
	\[ \abs{y - f(0)} < 2\varepsilon \Longrightarrow \abs{g(y)} = \abs{g(y)-g(f(0))} <\frac{1}{4}. \]
	Distinguiamo due casi: se $f(x) \in \Sigma_1$, 
	\[ \abs{h(f(x))-x} = \abs{\tilde{p}(\Phi(f(x))) -x} = \abs{\tilde{p}(f(x))-x} = \abs{\tilde{p}(f(x))- g(f(x))} < \delta < \frac{1}{2}<1. \]
	Se, invece, $y = f(x) \notin \Sigma_1$, allora $\abs{y-c}<\varepsilon$, quindi 
	\[ \abs{y - f(0)} \le \abs{y-c}+\abs{c-f(0)} < 2\varepsilon, \] 
	da cui $\abs{g(y)}<\frac{1}{4}$. Inoltre, 
	\[ \abs{\Phi(y) - f(0)} \le \abs{\Phi(y)-c}+\abs{c-f(0)} < 2\varepsilon, \] 
	da cui $\abs{g(\Phi(y))}<\frac{1}{4}$. Perciò,
	\begin{align*}
		\abs{h(f(x))-x} &= \abs{h(y) -g(y)} \le \abs{\tilde{p}(\Phi(y)))-g(\Phi(y))} + \abs{g(\Phi(y)) - g(y)} < \\
		&< \delta + \abs{g(\Phi(y)) } +\abs{g(y)}< 1.
	\end{align*}
	
	A questo punto, detta $\tilde{h}:\overline{\B(0,1)} \to \R^n$ tale che
	\[ \tilde{h}(x) = x-h(f(x)), \]
	chiaramente continua, si ha che, per ogni $x \in  \overline{\B(0,1)}$, $\abs{\tilde{h}(x)}\le 1$, ossia $\textup{Im}(\tilde{h})\subseteq \overline{\B(0,1)}$. Per il Teorema di Brouwer, esiste quindi $x_0 \in \overline{\B(0,1)}$ tale che $\tilde{h}(x_0) = x_0$, ossia $h(f(x_0)) = 0$, assurdo e la dimostrazione è conclusa.
\end{proof}
Per comprendere l'importanza delle ipotesi del Teorema di invarianza del dominio vediamo due controesempi evocativi.
\begin{esempio}
	Consideriamo l'applicazione continua e iniettiva $f: \left] 0,1\right[ \to \R^2$ tale che
	\[ f(t) = (t,0). \]
	Siccome $ \left] 0,1\right[ \subseteq \R$ e $1 \ne 2$, il Teorema di invarianza del dominio non si applica. Ciò accade a buon titolo, infatti $f( \left] 0,1\right[) =  \left] 0,1\right[ \times \left\lbrace 0\right\rbrace$, che non è aperto in $\R^2$.
\end{esempio}

\begin{esempio}
	Consideriamo l'applicazione continua e iniettiva $T_r: l^2(\R) \to l^2(\R) $ tale che
	\[ T_r((a_n)) = (0,a_0,a_1,a_2,\dots). \]
	Siccome $ l^2(\R) $ non ha dimensione finita, il Teorema di invarianza del dominio non si applica. Ciò accade a buon titolo, infatti $T_r(l^2(\R))$ non è aperto. Se per assurdo lo fosse, per ogni $(b_n) \in T_r(l^2(\R))$, ossia
	\[ (b_n) = (0,b_1,b_2,\dots), \]
	esisterebbe $r >0$ tale che $B((b_n),r) \subseteq T_r(l^2(\R))$. Ma ciò è assurdo perché 
	\[ (c_n) = \left( \frac{r}{2},b_1,b_2,\dots\right)  \in B((b_n),r),\]
	in quanto $\norm{(c_n)-(b_n)}_{l^2(\R)} = \frac{r}{2}<r$, ma $c_0 = \frac{r}{2} \ne 0$, quindi $(c_n) \in l^2(\R) \setminus T_r(l^2(\R))$.
\end{esempio}
Una conseguenza diretta del Teorema di invarianza del dominio è il seguente risultato.
\begin{teorema}\textbf{\emph{(di invarianza della dimensione)}}\index[analitico]{teorema!di invarianza!della dimensione}
	Siano $U \subseteq \R^n$ e $V \subseteq \R^m$ aperti. Se $n \ne m$, allora $U$ non è omeomorfo a $V$. In particolare, $\R^n$ non è omeomorfo a $\R^m$.
\end{teorema}
\begin{proof}
	Senza perdita di generalità, supponiamo $m > n$. Sia $i: \R^n \to \R^m$ tale che
	\[ i(x) = (x_1,\dots,x_n,0,\dots,0). \]
	Supponiamo, per assurdo, che esista un omeomorfismo $f: V \to U$ e consideriamo $g = i \circ f: V \to \R^m$. Chiaramente $g$ è continua e, dati $x,y \in V $ tali che $g(x) = g(y)$, risulta
	\[ (f_1(x),\dots,f_n(x),0,\dots,0) = (f_1(y),\dots,f_n(y),0,\dots,0). \]
	In particolare, $f(x) = f(y)$ che, essendo $f$ un omeomorfismo, quindi iniettiva, conduce a $x=y$. In altre parole $g$ è pure iniettiva. Per il Teorema di invarianza del dominio, $g(V)$ è aperto in $\R^m$, quindi per ogni $y \in g(V)$ esiste $r>0$ tale che $B(y,r) \subseteq g(V)$. Ma ciò è assurdo perché $z = \left( y_1,\dots,y_n,\frac{r}{2},0,\dots,0\right)  \in B(y,r)$, ma $z \in \R^m \setminus g(V)$.
\end{proof}

Grazie al Teorema di Brouwer è anche possibile dare una dimostrazione alternativa del Teorema fondamentale dell'Algebra.
\begin{teorema}\textbf{\emph{(fondamentale dell'Algebra)}}\index[analitico]{teorema!fondamentale dell'Algebra}
	Sia 
	\[ P(z) = \sum_{k=0}^{n}a_kz^k \]
	una funzione polinomiale complessa con $n\ge 1$ ed $a_n \ne 0$. Allora esiste $z \in \C$ tale che $P(z) = 0$.
\end{teorema}
\begin{proof}
	Innanzitutto, senza perdita di generalità, possiamo supporre $a_n = 1$ e identificare $\C$ con $\R^2$. Sia $r = 2 + \displaystyle\sum_{k=0}^{n-1}\abs{a_k}$ e $g: \overline{\B(0,r)} \to \C$ tale che
	\[ g(z) = \begin{cases}
		z- \frac{P(z)}{r}e^{i(1-n)\vartheta} & \textup{se } \abs{z} \le 1,\\
		z- \frac{P(z)}{r}z^{(1-n)} & \textup{se } \abs{z} > 1,
	\end{cases} \]
	ben definita e continua in quanto se $\abs{z} = 1$, $z = e^{i\vartheta}$ per qualche $\vartheta$ e $e^{i(1-n)\vartheta} = (e^{i\vartheta})^{1-n}$. Mostriamo che $\textup{Im}(g) \subseteq \overline{\B(0,r)}$. Se $\abs{z} \le 1$, 
	\[ \abs{g(z)} \le \abs{z} + \frac{\abs{P(z)}}{r} \le 1 + \frac{1+ \displaystyle\sum_{k=0}^{n-1}\abs{a_k}}{2+ \displaystyle\sum_{k=0}^{n-1}\abs{a_k}} \le 1+ 1 = 2 \le r.\]
	Se, invece, $1<\abs{z}\le r$, allora
	\[ \abs{g(z)} = \abs*{z - \frac{z}{r} - \frac{1}{rz^{n-1}}\sum_{k=0}^{n-1}a_kz^k} = \abs*{\frac{r-1}{r}z -\frac{1}{rz^{n-1}}\sum_{k=0}^{n-1}a_kz^k} \le \frac{r-1}{r}\abs{z} + \sum_{k=0}^{n-1}\frac{\abs{a_k}\abs{z^k}}{r\abs{z}^{n-1}},\]
	quindi 
	\[ \abs{g(z)} \le r-1 + \sum_{k=0}^{n-1}\frac{\abs{a_k}}{r\abs{z}^{n-1-k}} \le r-1 + \frac{\displaystyle\sum_{k=0}^{n-1}\abs{a_k}}{r} \le r-1 + 1 = r.\]
	
	Da ciò, per il Teorema di Brouwer, esiste $z \in \overline{\B(0,r)} \subseteq \C$ tale che $g(z) = z$, ossia $P(z) = 0$, che è la tesi.
\end{proof}

\section{Risultati in dimensione infinita}
Passiamo invece ai risultati che valgono senza l'ipotesi di dimensione finita.
\begin{teorema}\emph{\textbf{(di Schauder)}}\index[analitico]{teorema!di Schauder}
	Consideriamo $(X,\norm{\;})$ uno spazio normato su $\R$. Siano $C \subseteq X$ convesso, chiuso, limitato e non vuoto
	ed un'applicazione $f: C \to C$ continua e compatta. Esiste $x \in C$ tale che $f(x) = x$.
\end{teorema}
\begin{proof}
	Omettiamo la dimostrazione.
\end{proof}
Vediamo subito un'applicazione del Teorema di Schauder.
\begin{teorema}\textbf{\emph{(di Peano)}}\index[analitico]{teorema!di Peano}
	Siano $ (t_{0},u_{0}) \in \R\times\R^{n}$, $ a,b,r >0 $ e  
	\[ f: \left[t_{0} - a,t_{0}+b\right] \times \overline{\textup{B}(u_{0},r)} \to \R^{n} \]
	una applicazione continua tale che 
	\[ (b-a)\max\left\lbrace \abs*{f(\tau,\xi)} \; : \; \tau \in \left[t_{0} - a,t_{0} +b\right], \; \xi \in \overline{\textup{B}(u_{0},r)}\right\rbrace \le r. \]
	Allora esiste 
	\[ u: \left[t_{0} - a,t_{0}+b\right] \to \R^{n}, \]
	una soluzione del problema di Cauchy
	\[ \begin{cases}
		u' = f(t,u), \\
		u(t_{0}) = u_{0}.
	\end{cases} \]
\end{teorema}
\begin{proof}
	Essendo $f$ continua, $u: \left[t_{0} - a,t_{0}+b\right] \to \R^{n}$ è soluzione del problema di Cauchy \[ \begin{cases}
		u' = f(t,u), \\
		u(t_{0}) = u_{0},
	\end{cases} \]
	se e solo se $u$ è continua e per ogni $t \in \left[t_{0} - a,t_{0}+b\right]$
	\[ u(t) = u_0 + \int_{t_0}^{t} f(\tau,u(\tau)) d\tau.\]
	Per ogni $u \in \overline{\B(u_0,r)} \subseteq C(\left[t_{0} - a,t_{0}+b\right];\R^n)$ consideriamo il funzionale
	\[ \Phi(u)(t) = u_0 + \int_{t_0}^{t} f(\tau,u(\tau)) d\tau. \]
	Osserviamo che 
	\begin{align*}
		\abs{\Phi(u)(t) - u_0} \le\int_{t_0}^{t}\abs{f(\tau,u(\tau))d\tau}\le \frac{r}{b-a}(b-a) = r,
	\end{align*}
	per cui $\norm{\Phi(u) - u_0}_\infty \le r$. In altre parole possiamo pensare che $\Phi:\overline{\B(u_0,r)} \to \overline{\B(u_0,r)}$. La continuità di $\Phi$ discende dal Teorema della convergenza dominata, mostriamo quindi che $\Phi$ è compatto. Sia $(u_h)$ in $\overline{\B(u_0,r)}$, chiaramente limitata. Siccome $\textup{Im}(\Phi) \subseteq \overline{\B(u_0,r)}$, chiaramente $(\Phi(u_h))$ è limitata. Dato $\varepsilon > 0$, sia $\delta = \frac{\varepsilon(b-a)}{r}$. Allora per ogni $h \in \N$ e per ogni $s,t \in \left[t_{0} - a,t_{0}+b\right]$ tali che $\abs{t-s} < \delta$
	\[ \abs{\Phi(u_h)(t) - \Phi(u_h)(s)} \le \abs*{\int_{s}^{t}\abs{f(\tau,u_h(\tau))}d\tau} \le \frac{r}{b-a}\abs{t-s} < \varepsilon.\]
	Per il Teorema di Ascoli--Arzelà, esiste una sottosuccessione $(u_{h_k})$ convergente, ossia $\Phi$ è compatto. 
	
	Dal Teorema di Schauder discende quindi che esiste $u \in \overline{\B(u_0,r)}$ tale che $\Phi(u) = u$, da cui la tesi.
\end{proof}

Il risultato gemello del Teorema di Schauder è il seguente.
\begin{teorema}\emph{\textbf{(di Tychonoff)}}\index[analitico]{teorema!di Tychonoff}
	Consideriamo $(X,\norm{\;})$ uno spazio normato su $\R$. Siano $K \subseteq X$ convesso, compatto e non vuoto ed un'applicazione $f : K \to K$ continua. Esiste $x \in K$ tale che $f(x) = x$.
\end{teorema}
\begin{proof}
	Innanzitutto, fissato $\varepsilon>0$,
	\[ \bigcup_{x \in K}B(x,\varepsilon) \]
	è un ricoprimento aperto di $K$. Per compattezza, esistono quindi $x_1,\dots,x_{n_\varepsilon} \in K$ tali che
	\[ K \subseteq \bigcup_{i=1}^{n_\varepsilon}\B(x_i,\varepsilon). \]
	Consideriamo l'inviluppo convesso chiuso di $x_1,\dots, x_{n_\varepsilon}$, ossia
	\[ K_\varepsilon = \left\lbrace \sum_{i=1}^{n_\varepsilon} \lambda_i x_i : \lambda_i \in [0,1], \sum_{i=1}^{n_\varepsilon} \lambda_i =1\right\rbrace. \]
	
	Osserviamo che $K_\varepsilon$ è chiuso e convesso. Sia $\sum_{i=1}^{n_\varepsilon}\lambda_i^{(h)}x_i \to y \in X$. La successione $((\lambda_1^{(h)},\dots,\lambda_{n_\varepsilon}^{(h)}))$ è limitata in $[0,1]^{n_\varepsilon}$ quindi, per il Teorema di Bolzano--Weierstrass ed il fatto che  $[0,1]^{n_\varepsilon}$ è chiuso in $\R^{n_\varepsilon}$, esiste $(\lambda_1,\dots,\lambda_{n_\varepsilon}) \in [0,1]^{n_\varepsilon}$ tale che, a meno di sottosuccessioni, $(\lambda_1^{(h)},\dots,\lambda_{n_\varepsilon}^{(h)}) \to (\lambda_1,\dots,\lambda_{n_\varepsilon})$ in $\R^{n_\varepsilon}$. Passando al limite in 
	\[ \sum_{i=1}^{n_\varepsilon}\lambda_i^{(h_k)} = 1, \]
	otteniamo anche che 
	\[ \sum_{i=1}^{n_\varepsilon}\lambda_i = 1. \]
	In particolare poi, per continuità delle operazioni di somma e prodotto per scalare, 
	\[ \sum_{i=1}^{n_\varepsilon}\lambda_i^{(h_k)}x_i \to \sum_{i=1}^{n_\varepsilon}\lambda_ix_i,  \]
	quindi per unicità del limite deve essere $y = \displaystyle\sum_{i=1}^{n_\varepsilon}\lambda_ix_i \in K_\varepsilon$, da cui la chiusura di $K_\varepsilon$. Per quanto riguarda la convessità invece, se $\displaystyle\sum_{i=1}^{n_\varepsilon}\lambda_i x_i, \displaystyle\sum_{i=1}^{n_\varepsilon}\overline{\lambda}_i x_i \in K_\varepsilon$ e $\lambda \in [0,1]$
	\[ \lambda \sum_{i=1}^{n_\varepsilon}\lambda_i x_i + (1-\lambda)\sum_{i=1}^{n_\varepsilon}\overline{\lambda}_i x_i =  \sum_{i=1}^{n_\varepsilon}[\lambda \lambda_i +(1-\lambda)\overline{\lambda}_i]x_i \in K_\varepsilon  \]
	siccome, per convessità di $[0,1]$, per ogni $i=1,\dots,n_{\varepsilon}$ si ha che $\lambda \lambda_i +(1-\lambda)\overline{\lambda}_i \in [0,1]$ e 
	\[ \sum_{i=1}^{n_\varepsilon}\left( \lambda \lambda_i +(1-\lambda)\overline{\lambda}_i\right) = \lambda  \sum_{i=1}^{n_\varepsilon} \lambda_i +(1-\lambda) \sum_{i=1}^{n_\varepsilon}\overline{\lambda}_i = \lambda + 1 -\lambda = 1. \]
	
	Mostriamo che $K_\varepsilon \subseteq K$. Chiaramente, per ogni $i= 1,\dots, n_\varepsilon$ si ha $x_i \in K$. Poi, essendo $K$ convesso, per ogni $i,j = 1,\dots,n_\varepsilon$, dato $\lambda_i \in [0,1]$, si ha $\lambda_ix_i +(1-\lambda_i)x_j \in K$. Ora, per ogni $i,j,k=1,\dots,n_\varepsilon$, consideriamo $\lambda_i,\lambda_j,\lambda_k \in [0,1]$ tali che $\lambda_i + \lambda_j + \lambda_k = 1$. Preso
	\[ y = \frac{\lambda_j}{\lambda_j + \lambda_k}x_j + \frac{\lambda_k}{\lambda_j + \lambda_k}x_k, \]
	dal fatto che 
	\begin{align*}
		\frac{\lambda_k}{\lambda_j + \lambda_k} &\in [0,1], & \frac{\lambda_k}{\lambda_j + \lambda_k} &= 1 - \frac{\lambda_j}{\lambda_j + \lambda_k}
	\end{align*}
	e dal passo precedente, segue che $y \in K$. In particolare, 
	\[ \lambda_i x_i + \lambda_j x_j + \lambda_k x_k = \lambda_i x_i + (1-\lambda_i)y \in K.\]
	Procedendo ricorsivamente, si ottiene $K_\varepsilon \subseteq K$. In particolare, essendo $K_\varepsilon$ chiuso in un compatto, è anch'esso compatto.
	
	Consideriamo ora $P_\varepsilon : K \to K_\varepsilon$ tale che
	\[ P_\varepsilon(x) = \dfrac{\displaystyle \sum_{i=1}^{n_\varepsilon}\textup{dist}(x,K\setminus \B(x_i,\varepsilon))x_i}{\displaystyle\sum_{i=1}^{n_\varepsilon}\textup{dist}(u,K\setminus \B(x_i,\varepsilon))}. \]
	Vediamo che è ben definita. Innanzitutto, il denominatore è sempre diverso da zero: infatti, essendo una somma di termini positivi, per essere nullo dovrebbe avere tutti gli addendi nulli e, se ciò si verificasse, per ogni $i=1,\dots,n_\varepsilon$ si avrebbe $x \in K\setminus \B(x_i,\varepsilon)$, ossia $x \notin \B(x_i,\varepsilon)$, che contraddirebbe il fatto che K è ricoperto dall'unione delle $\B(x_i,\varepsilon)$. Inoltre, per ogni $i=1,\dots,n_\varepsilon$ si ha
	\begin{align*}
		\lambda_i = \dfrac{\textup{dist}(x,K\setminus \B(x_i,\varepsilon))}{\displaystyle\sum_{i=1}^{n_\varepsilon}\textup{dist}(x,K\setminus \B(x_i,\varepsilon))} &\in [0,1], & \sum_{i=1}^{n_{\varepsilon}}\lambda_i &= 1, 
	\end{align*}
	quindi $P_\varepsilon(x) \in K_\varepsilon$. In più, $P_\varepsilon$ è continua per costruzione e per ogni $x \in K$
	\begin{align*}
		\norm{P_\varepsilon(x)-x} &=  \norm*{\dfrac{\displaystyle \sum_{i=1}^{n_\varepsilon}\textup{dist}(x,K\setminus \B(x_i,\varepsilon))x_i}{\displaystyle\sum_{i=1}^{n_\varepsilon}\textup{dist}(x,K\setminus \B(x_i,\varepsilon))} - \dfrac{\displaystyle \sum_{i=1}^{n_\varepsilon}\textup{dist}(x,K\setminus \B(x_i,\varepsilon))}{\displaystyle\sum_{i=1}^{n_\varepsilon}\textup{dist}(x,K\setminus \B(x_i,\varepsilon))}x} = \\
		&= \norm*{\dfrac{\displaystyle \sum_{i=1}^{n_\varepsilon}\textup{dist}(x,K\setminus \B(x_i,\varepsilon))(x_i-x)}{\displaystyle\sum_{i=1}^{n_\varepsilon}\textup{dist}(x,K\setminus \B(x_i,\varepsilon))}} \le \dfrac{\displaystyle \sum_{i=1}^{n_\varepsilon}\textup{dist}(x,K\setminus \B(x_i,\varepsilon))\norm{x_i-x}}{\displaystyle\sum_{i=1}^{n_\varepsilon}\textup{dist}(x,K\setminus \B(x_i,\varepsilon))},
	\end{align*}
	dove, se $\textup{dist}(x,K\setminus \B(x_i,\varepsilon)) = 0$, chiaramente non ho contributo nelle somme; se invece si ha $\textup{dist}(x,K\setminus \B(x_i,\varepsilon)) \ne 0$, allora $\norm{x_i-x}<\varepsilon$. In definitiva,
	\begin{equation}\label{superimportanteperTychonoff}
		\norm{P_\varepsilon(x)-x} < \varepsilon.
	\end{equation}
	
	Detto $m_\varepsilon$ il massimo numero di $x_1,\dots,x_{n_\varepsilon}$ linearmente indipendenti, a meno di permutazioni degli indici, possiamo pensare che $x_1,\dots, x_{m_\varepsilon}$ siano linearmente indipendenti. A questo punto, a meno di una traslazione che manda $x_1$ in $0$, $K_\varepsilon$ è un sottoinsieme dello spazio vettoriale, di dimensione $m_\varepsilon-1$, generato da $x_2,\dots,x_{m_\varepsilon}$, quindi non è restrittivo pensare $K_\varepsilon \subseteq \R^{m_\varepsilon-1}$. Detta $f_\varepsilon: K_\varepsilon \to K_\varepsilon$ tale che
	\[ f_\varepsilon(x) = P_\varepsilon(f(x)), \]
	continua per composizione, dal Teorema di Brouwer segue l'esistenza di $x_\varepsilon \in K_\varepsilon$ tale che $f_\varepsilon(x_\varepsilon) = x_\varepsilon$.
	
	A questo punto, dal fatto che $K$ è compatto, esistono una sottosuccessione $(x_{\varepsilon_j})$ e $x \in K$ tali che $x_{\varepsilon_j} \to x$ in $X$. Usando la disuguaglianza \eqref{superimportanteperTychonoff},
	\[ \norm{x_{\varepsilon_j}- f(x_{\varepsilon_j})} =  \norm{f_{\varepsilon_j}(x_{\varepsilon_j})- f(x_{\varepsilon_j})} = \norm{P_{\varepsilon_j}(f(x_{\varepsilon_j})) -f(x_{\varepsilon_j})} < \varepsilon_j,\]
	e, per continuità della norma e di $f$, passando al limite ambo i membri si ottiene 
	\[ \norm{x - f(x)}\le 0, \]
	ossia $f(x) = x$, che da la tesi.
\end{proof} 
		
		% Commenti speciali

% !TEX encoding = UTF-8 Unicode
% !TEX TS-program = pdflatex
% !TEX spellcheck = it-IT
% !TEX root = Dispensa/afed.tex

% Capitolo 5
\chapter{Alcuni elementi di Analisi non lineare}\label{Alcuni elementi di Analisi non lineare}

\section{Introduzione}
La linearità è una delle proprietà più amate in tutta la Matematica. Anche in Analisi Matematica, lavorare con operatori lineari garantisce una trattazione pressoché organica e molto generale: basti pensare a quanto visto per gli operatori differenziali lineari del secondo ordine studiati nel capitolo \ref{Equazioni ellittiche lineari del secondo ordine}. Rinunciare alla linearità ha in generale un costo non indifferente, come mostrato dal seguente esempio motivazionale.
\begin{esempio}\label{equazionialgebriche}
	Consideriamo un'applicazione reale di variabile reale $f$, definita nel modo più generico assegnando l'espressione matematica $f(x)$, di classe $C^\infty$ nel suo dominio naturale. Siamo interessati a determinare gli zeri di $f$, ossia a trovare le soluzioni del problema
	\[ f(x) = 0. \]
	Se $f$ è lineare, il problema è banale. Se $f$ è un polinomio di grado almeno $1$, grazie al Teorema fondamentale dell'Algebra, sappiamo esattamente il numero di soluzioni del problema, ma se il grado del polinomio supera $4$, dal Teorema di Abel--Ruffini, a priori non ne esiste una formula di rappresentazione. In generale, presa $f$ qualsiasi, il problema può rivelarsi insidioso: potrei non avere nemmeno un risultato di esistenza di soluzioni.
\end{esempio}

In questo capitolo siamo interessati a presentare alcuni metodi per lo studio di determinate classi di operatori non lineari: l'obiettivo principale sarà la determinazione di risultati di esistenza di soluzioni per problemi differenziali non lineari.

\section{Metodi di punto fisso}
Alcuni problemi non lineari possono essere riscritti come problemi di punto fisso: un primo caso è l'Esempio \eqref{equazionialgebriche}.
\begin{esempio}
	Consideriamo un'applicazione reale di variabile reale $f$, definita nel modo più generico assegnando l'espressione matematica $f(x)$, di classe $C^\infty$ nel suo dominio naturale. Il problema 
	\[ f(x) = 0 \]
	è equivalente a 
	\[ f(x) + x = x \]
	ossia, posto $g(x) = f(x) + x$, al problema di punto fisso
	\[ g(x) = x. \]
\end{esempio}
La prima applicazione del risultati di punto fisso che vogliamo proporre è legata al Teorema delle contrazioni. 
\begin{definizione}
Siano $U$ un aperto limitato di $\R^n$ con $\partial U$ di classe $C^\infty$, $T>0$, $g \in H^1_0(U;\R^m)$ ed $f: \R^m \to \R^m$ un'applicazione lipschitziana e sub-lineare\footnote{A voler ben vedere la sub-linearità è conseguenza della lipschitzianità su $\R^m$, infatti $\abs{f(x)}- \abs{f(0)}\le \abs{f(x)-f(0)} \le C\abs{x}$, da cui $\abs{f(x)}\le C(1+\abs{x})$. Nel seguito la preciseremo comunque.}. Diciamo che $u \in L^2(0,T;H^1_0(U;\R^m))$ tale che $u' \in L^2(0,T;H^{-1}(U;\R^m))$ è una \emph{soluzione debole}\index[analitico]{soluzione!debole} del problema di Dirichlet
	\[ 	\begin{cases}
		\partial_t u-\Delta u = f(u) & \textup{in } U \times \left] 0,T\right],\\
		u = 0 & \textup{su } \partial U \times [0,T],\\
		u=g & \textup{su }  U \times \left\lbrace 0\right\rbrace ,
	\end{cases} \]
	se $u(0) = g$ e per ogni $v \in H^1_0(U;\R^m)$ e per \emph{q.o.} $t \in [0,T]$
	\[ \braket{u',v} + \int_U Du \cdot Dv d\mathcal{L}^n = \int_U f(u)\cdot v d\mathcal{L}^n. \]
\end{definizione}
\begin{osservazione}
La precedente definizione è ben posta grazie alla sub-linearità di $f$ ed al fatto che $u$ ammette un rappresentante in $C([0,T];L^2(U;\R^m))$.
\end{osservazione}
\begin{osservazione}
Se $f: U \times [0,T] \to \R^m$, ossia nel caso lineare e scalare, si può provare esistenza ed unicità della soluzione debole grazie al \emph{Metodo di Galerkin}.
\end{osservazione}

\begin{teorema}
Siano $U$ un aperto limitato di $\R^m$ con $\partial U$ di classe $C^\infty$, $g \in H^1_0(U;\R^m)$ ed $f: \R^m \to \R^m$ un'applicazione lipschitziana e sub-lineare. Esiste $T>0$ tale che il problema di Dirichlet
\[ 	\begin{cases}
	\partial_t u-\Delta u = f(u) & \textup{in } U \times \left] 0,T\right],\\
	u = 0 & \textup{su } \partial U \times [0,T],\\
	u=g & \textup{su }  U \times \left\lbrace 0\right\rbrace ,
\end{cases} \]
ammette una ed una sola soluzione debole.
\end{teorema}
\begin{proof}
Data $u \in C([0,T];L^2(U;\R^m))$, sia $h(t) = f(u(t))$. Dal fatto che $f$ è sub-lineare, $h \in L^2(0,T;L^2(U;\R^m))$. In particolare, il problema di Dirichlet
\[ 	\begin{cases}
	\partial_t w-\Delta w = h & \textup{in } U \times \left] 0,T\right],\\
	w = 0 & \textup{su } \partial U \times [0,T],\\
	w=g & \textup{su }  U \times \left\lbrace 0\right\rbrace ,
\end{cases} \]
ammette un'unica soluzione debole $w \in L^2(0,T;H^1_0(U;\R^m))$ tale che $w' \in L^2(0,T;H^{-1}(U;\R^m))$, ossia $w(0) = g$ e per ogni $v \in H^1_0(U;\R^m)$ e per q.o. $t \in [0,T]$
\[ \braket{w',v} + \int_U Dw \cdot Dv d\mathcal{L}^n = \int_U h\cdot v d\mathcal{L}^n. \]
Dal fatto che $w$ ha un rappresentante in $C([0,T];L^2(U;\R^m))$, che chiamiamo comunque $w$, rimane definita un'applicazione $A: C([0,T];L^2(U;\R^m)) \to C([0,T];L^2(U;\R^m))$ tale che 
\[ A(u) = w. \]
Dati $u_1,u_2 \in C([0,T];L^2(U;\R^m))$, siano $w_1 = A(u_1), w_2 = A(u_2)$. Innanzitutto, per linearità,
\begin{align*}
\braket{(w_1-w_2)',w_1-w_2} + \int_U \abs{D(w_1-w_2)}^2d\mathcal{L}^n = \int_U (h_1-h_2)\cdot (w_1-w_2) d\mathcal{L}^n,
\end{align*}
quindi 
\begin{align*}
	\frac{1}{2} \frac{d}{dt}\norm{w_1-w_2}_{L^2(U;\R^m)}^2+ \int_U \abs{D(w_1-w_2)}^2d\mathcal{L}^n = \int_U (h_1-h_2)\cdot (w_1-w_2) d\mathcal{L}^n,
\end{align*}
da cui, per la disuguaglianza di Cauchy--Schwarz combinata con la disuguaglianza di Poincaré e la disuguaglianza di Young pesata,
\begin{align*}
\frac{d}{dt}\norm{w_1-w_2}_{L^2(U;\R^m)}^2+ 2\norm{w_1 - w_2}_{H^1_0(U;\R^m)}^2 &= 2\int_U (h_1-h_2)\cdot (w_1-w_2) d\mathcal{L}^n \le \\
&\le 2 \norm{w_1-w_2}_{L^2(U;\R^m)}\norm{h_1-h_2}_{L^2(U;\R^m)}\le \\
&\le 2\norm{w_1 - w_2}_{H^1_0(U;\R^m)}^2 + C\norm{h_1-h_2}_{L^2(U;\R^m)}^2,
\end{align*}
che conduce, usando la lipschitzianità di $f$, a
\[ \frac{d}{dt}\norm{w_1-w_2}_{L^2(U;\R^m)}^2 \le C\norm{h_1-h_2}_{L^2(U;\R^m)}^2 = C \norm{f(u_1)-f(u_2)}_{L^2(U;\R^m)}^2 \le C \norm{u_1 - u_2}_{L^2(U;\R^m)}^2. \]
Integrando in $t$,
\[ \norm{w_1(s)-w_2(s)}_{L^2(U;\R^m)}^2 \le C \int_0^s \norm{u_1(t) - u_2(t)}_{L^2(U;\R^m)}^2d\mathcal{L}^1(t) \le CT \norm{u_1-u_2}_{C([0,T];L^2(U;\R^m))}^2, \]
quindi 
\[ \norm{A(u_1)-A(u_2)}_{C([0,T];L^2(U;\R^m))} \le(CT)^{\frac{1}{2}} \norm{u_1-u_2}_{C([0,T];L^2(U;\R^m))}, \]
con $C$ che non dipende da $T$. Se $T$ tale che $(CT)^{\frac{1}{2}} <1$, la tesi segue dal Teorema delle contrazioni.
\end{proof}

Il risultato di punto fisso a cui vorremmo fare riferimento per la prossima applicazione è il Teorema di Tychonoff, di cui il seguente è un Corollario che si adatta ai nostri scopi.
\begin{corollario}\emph{\textbf{(Teorema di Schaefer)}}\index[analitico]{teorema!di Schaefer}
	Consideriamo $(X,\norm{\;})$ uno spazio normato su $\R$ e $f: X \to X$ un'applicazione continua e compatta. Se l'insieme
	\[ B = \left\lbrace  x \in X : x = \lambda f(x), \textup{ per qualche } \lambda \in [0,1]\right\rbrace  \]
	è limitato, allora esiste $x \in X$ tale che $f(x) = x$.
\end{corollario}
\begin{proof}
	Sia $M > 0$ tale che 
	\[ \norm{x}<M \]
	per ogni $x \in B$. Consideriamo $\tilde{f}: \overline{\B(0,M)} \to \overline{\B(0,M)}$ tale che
	\[ \tilde{f}(x) = \begin{cases} 
		f(x) & \textup{se } \norm{f(x)}\le M,\\
		\frac{Mf(x)}{\norm{f(x)}} & \textup{se } \norm{f(x)}> M.
	\end{cases} \]
	
	Osserviamo che $\tilde{f}$ è ben definita: se $\norm{f(x)}\le M$, allora $\norm{\tilde{f}(x)} \le M$; se invece $ \norm{f(x)}> M$, 
	\[ \norm{\tilde{f}(x)} = \frac{M \norm{f(x)}}{\norm{f(x)}} =M.\]
	In particolare, $\tilde{f}(\overline{\B(0,M)}) \subseteq \overline{\B(0,M)}$.
	
	Mostriamo anche che $\tilde{f}$ è compatta. Sia $(x_h)$ in $\overline{\B(0,M)}$, chiaramente limitata. Se, definitivamente in $h$, $\norm{f(x_h)}\le M$ (risp. $\norm{f(x_h)} > M$), allora, definitivamente in $h$, $\tilde{f}(x_h) = f(x_h)$ (risp. $\tilde{f}(x_h) = \frac{Mf(x_h)}{\norm{f(x_h)}}$) e la compattezza segue da quella di $f$ (risp. da quella di $f$ e dalla continuità della norma). A meno di sottosuccessioni, posso sempre supporre di essere nella situazione sopra descritta, quindi $\tilde{f}$ è compatta.
	
	Consideriamo $C_H(\tilde{f}(\overline{\B(0,M)})) \subseteq X$, l'inviluppo convesso chiuso di $\tilde{f}(\overline{\B(0,M)})$. $C_H(\tilde{f}(\overline{\B(0,M)}))$ è convesso e chiuso, mostriamo che è compatto. Per fare ciò, consideriamo $(x_h)$ in $C_H(\tilde{f}(\overline{\B(0,M)}))$. In particolare, per ogni $h \in \N$
	\[ x_h = \sum_{i=1}^{n_h}\lambda_i^{(h)}\tilde{f}(y_i^{(h)}). \]
	Fissato $i$, a meno di sottosuccessioni, per compattezza di $\tilde{f}$ ed il Teorema di Bolzano--Weierstrass, $\lambda_i^{h} \to \lambda$ e $\tilde{f}(y_i^{h}) \to \tilde{f}(y_i)$. Se, per assurdo, $n_{h} \to +\infty$, allora il limite della successione $(x_h)$ non starebbe in  $C_H(\tilde{f}(\overline{\B(0,M)}))$, violandone la chiusura, quindi, a meno di sottosuccessioni, $(x_h)$ è convergente in $C_H(\tilde{f}(\overline{\B(0,M)}))$.
	
	Vediamo che è ben definita $\tilde{f}: C_H(\tilde{f}(\overline{\B(0,M)})) \to C_H(\tilde{f}(\overline{\B(0,M)}))$. In primo luogo,
	\[ C_H(\tilde{f}(\overline{\B(0,M)})) \subseteq C_h(\overline{\B(0,M)}) = \overline{\B(0,M)}, \]
	e
	\[ \tilde{f}(C_H(\tilde{f}(\overline{\B(0,M)}))) \subseteq \tilde{f}(\overline{\B(0,M)}) \subseteq C_H(\tilde{f}(\overline{\B(0,M)})). \]
	
	Dal Teorema di Tychonoff, esiste quindi $x \in C_H(\tilde{f}(\overline{\B(0,M)}))$ tale che $\tilde{f}(x) = x$. Ora, se per assurdo, $x$ non fosse un punto fisso per $f$, allora si avrebbe, per definizione di $\tilde{f}$, 
	\[ \norm{f(x)} > M \]
	e, detto $\lambda = \frac{M}{\norm{f(x)}} <1$, 
	\[ x = \tilde{f}(x) = \lambda f(x), \]
	quindi $x \in B$, ma $\norm{x} = \norm{\tilde{f}(x)} = M$, da cui l'assurdo e la tesi.
\end{proof}
Il Teorema di Schaefer rientra nella filosofia del \emph{metodo delle stime a priori}: se, assumendo per ipotesi che una soluzione di un certo problema esista, si riescono a dimostrare determinate stime per le soluzioni, allora queste stesse stime si possono usare per dimostrare l'esistenza di soluzioni (infatti se non valessero sicuramente la soluzione non esisterebbe). In ogni caso, il vantaggio del Teorema di Schaefer rispetto al Teorema di Tychonoff sta nel fatto che non è richiesto di specificare l'insieme convesso e compatto su cui lavorare. Vediamone un'applicazione.
\begin{definizione}
	Siano $U$ un aperto limitato di $\R^n$ con $\partial U$ di classe $C^\infty$, $b: \R^n \to \R$ un'applicazione lipschitziana, sub-lineare e di classe $C^\infty$ e $\mu >0$. Diciamo che $u$ è una \emph{soluzione classica}\index[analitico]{soluzione!classica} del problema di Dirichlet
	\[ 	\begin{cases}
		-\Delta u + b(Du) + \mu u = 0 & \textup{in } U,\\
		u = 0 & \textup{su } \partial U,
	\end{cases} \]
	se $u \in C^2(U)\cap C(\overline{U})$, $-\Delta u(x) + b(Du(x)) + \mu u(x) = 0$ per ogni $x \in U$ e $u(x) = 0$ per ogni $x \in \partial U$.
\end{definizione}
Vediamo ora come indebolire la definizione di soluzione del problema 
\[ 	\begin{cases}
	-\Delta u + b(Du) + \mu u = 0 & \textup{in } U,\\
	u = 0 & \textup{su } \partial U,
\end{cases} \]
Se $u$ è una soluzione classica, allora per ogni $x \in U$
\[ -\Delta u(x) + b(Du(x)) + \mu u(x) = 0. \]
Moltiplicando ambo i membri per una certa $\varphi \in C^\infty_c(U)$ ed integrando otteniamo 
\[ \int_U \left( -\Delta u \varphi + \mu u\varphi\right) d\mathcal{L}^n = -\int_U b(Du) \varphi d\mathcal{L}^n. \]
Per le formule di Gauss--Green abbiamo quindi, per ogni $\varphi \in C^\infty_c(U)$
\[ \int_U \left( Du \cdot D\varphi +\mu u \varphi\right)  d\mathcal{L}^n = -\int_U b(Du) \varphi d\mathcal{L}^n. \]
Osserviamo che l'equazione precedente ha significato con $u,\varphi \in H^1_0(U)$.
\begin{definizione}
	Siano $U$ un aperto limitato di $\R^n$ con $\partial U$ di classe $C^\infty$, $b: \R^n \to \R$ un'applicazione lipschitziana, sub-lineare e di classe $C^\infty$ e $\mu >0$.	Diciamo che $u \in H^1_0(U)$ è una \emph{soluzione debole}\index[analitico]{soluzione!debole} del problema di Dirichlet
	\[ 	\begin{cases}
		-\Delta u + b(Du) + \mu u = 0 & \textup{in } U,\\
		u = 0 & \textup{su } \partial U,
	\end{cases} \]
	se per ogni $v \in H^1_0(U)$, 
	\[ \int_U \left( Du \cdot Dv +\mu u v\right)  d\mathcal{L}^n = -\int_U b(Du) v d\mathcal{L}^n. \]
\end{definizione}
\begin{osservazione}
	La precedente definizione è ben posta, infatti, dalla sub-linearità di $b$ combinata con le disuguaglianze di Holder e Poincaré,
	\[ \int_U b(Du) v d\mathcal{L}^n \le C\norm{v}_{H^1_0(U)}(1+\norm{u}_{H^!_0(U)})<+\infty.  \]
	e, chiaramente,
	\[  \int_U \left( Du \cdot Dv +\mu u v\right)  d\mathcal{L}^n < +\infty.\]
\end{osservazione}

Vediamo come, applicando il Teorema di Schaefer, riusciamo a dimostrare l'esistenza di soluzioni deboli al problema in esame.
\begin{teorema}
	Siano $U$ un aperto limitato di $\R^n$ con $\partial U$ di classe $C^\infty$, $b: \R^n \to \R$ un'applicazione lipschitziana, sub-lineare e di classe $C^\infty$ e $\mu >0$. Se $\mu$ è sufficientemente grande, allora esiste $u \in H^2(U)\cap H^1_0(U)$ soluzione debole
	del problema di Dirichlet
	\[ 	\begin{cases}
		-\Delta u + b(Du) + \mu u = 0 & \textup{in } U,\\
		u = 0 & \textup{su } \partial U.
	\end{cases} \] 
\end{teorema}
\begin{proof}
	Innanzitutto, per ogni $u \in H^1_0(U)$, posto $f = -b(Du)$, dalla sub-linearità di $b$, si ha
	\[ \norm{f}_{L^2(U)}^2 = \int_U f^2 d\mathcal{L}^n\le C\int_U(1+\abs{Du})^2 d\mathcal{L}^n \le C(1+ \norm{u}_{H^1_0(U)}^2)<+\infty, \]
	quindi $f \in L^2(U)$. Allora, per il primo Teorema di esistenza di soluzioni deboli, per ogni $\mu >0$ esiste ed è unica $w \in H^1_0(U)$ soluzione debole di 
	\[ \begin{cases}
		-\Delta w + \mu w = f & \textup{in } U,\\
		w=0 & \textup{su } \partial U.
	\end{cases} \]
	Per quanto visto sulla regolarità delle soluzioni deboli nel caso lineare, $w \in H^2(U)$ ed esiste $C>0$, dipendente solo da $U$ e $\mu$, tale che
	\[ \norm{w}_{H^2(U)} \le C \norm{f}_{L^2(U)}. \]
	Rimane definita $A: H^1_0(U) \to H^1_0(U)$ tale che
	\[ A(u) = w, \]
	dove $w$ è la soluzione sopra descritta. In particolare, per ogni $u,v \in H^1_0(U)$
	\begin{equation}\label{perilgranfinale}
		\int_U \left( D(A(u))\cdot Dv + \mu A(u)v\right) d\mathcal{L}^n = -\int_U b(Du)vd\mathcal{L}^n
	\end{equation}
	e, per la sub-linearità di $b$,
	\[ \norm{A(u)}_{H^2(U)} \le C (1+ \norm{u}_{H^1_0(U)}). \]
	
	Mostriamo che $A$ è continua.  Siano $(u_h)$ in $H^1_0(U)$ e $u \in H^1_0(U)$ tali che $u_h \to u$ in $H^1_0(U)$. Siccome $(u_h)$ è convergente, allora è limitata, quindi 
	\[ \underset{h}{\sup}\; \norm{A(u_h)}_{H^2(U)} \le C \underset{h}{\sup}\; \left( 1 + \norm{u_h}_{H^1_0(U)}\right)  < +\infty, \]
	ossia $(A(u_h))$ è limitata in $H^2(U)$. Allora esiste $w \in H^1(U)$ tale che, a meno di una sottosuccessione, $A(u_h) \to w$ in $H^1(U)$. Siccome $H^1_0(U)$ è completo, in particolare è chiuso perciò $w \in H^1_0(U)$ e $A(u_h) \to w$ in $H^1_0(U)$. Ora, per definizione di $A$, per ogni $v \in H^1_0(U)$
	\begin{equation}\label{formulazionedeboleA}
		\int_U \left( D(A(u_h))\cdot Dv + \mu A(u_h)v\right) d\mathcal{L}^n = -\int_U b(Du_h)vd\mathcal{L}^n.
	\end{equation}
	Siccome $A(u_h) \to w$ in $H^1_0(U)$, in particolare $D(A(u_h)) \to Dw$ in $L^2(U)$, quindi $D(A(u_h)) \rightharpoonup Dw$ in $L^2(U)$. Similmente, $A(u_h) \to w$ in $L^2(U)$, quindi $A(u_h) \rightharpoonup w$ in $L^2(U)$. In più, dal fatto che $u_h \to u$ in $H^1_0(U)$, si ha $Du_h \to Du$ in $L^2(U)$ quindi dal fatto che $b$ è sub-lineare,
	\begin{align*}
	\int_U(b(Du_h)-b(Du))vd\mathcal{L}^n &\le C \int_U (\abs{Du_h}-\abs{Du})vd\mathcal{L}^n \le C\int_U \abs{Du_h - Du}\abs{v}d\mathcal{L}^n \le \\
	&\le C\int_U \abs{Du_h - Du}^2d\mathcal{L}^n \int_U \abs{v}^2 d\mathcal{L}^n \to 0,
	\end{align*}
	quindi $b(Du_h) \rightharpoonup b(Du)$ in $L^2(U)$. Passando al limite in \eqref{formulazionedeboleA}, per ogni $u,v \in H^1_0(U)$
	\[ \int_U \left( D(w)\cdot Dv + \mu wv\right) d\mathcal{L}^n = -\int_U b(Du)vd\mathcal{L}^n, \]
	ossia $w$ è soluzione debole del problema di Dirichlet
	\[ \begin{cases}
		-\Delta w + \mu w = -b(Du) & \textup{in } U,\\
		w=0 & \textup{su } \partial U,
	\end{cases} \]
	che, in altre parole, significa $A(u) = w$, quindi $A(u_h) \to A(u)$ in $H^1_0(U)$ e la continuità è dimostrata.
	
	Mostriamo che $A$ è compatta. Sia $(u_h)$ in $H^1_0(U)$ limitata. Allora 
	\[ \underset{h}{\sup}\; \norm{A(u_h)}_{H^2(U)} \le C \underset{h}{\sup}\; \left( 1 + \norm{u_h}_{H^1_0(U)}\right)  < +\infty, \]
	ossia $(A(u_h))$ è limitata in $H^2(U)$. Allora esiste $w \in H^1(U)$ tale che, a meno di una sottosuccessione, $A(u_h) \to w$ in $H^1(U)$. Siccome $H^1_0(U)$ è completo, in particolare è chiuso perciò $w \in H^1_0(U)$ e $A(u_h) \to w$ in $H^1_0(U)$, da cui la compattezza.
	
	Mostriamo che, per $\mu$ sufficientemente grande, l'insieme
	\[ B= \left\lbrace  u \in H^1_0(U) : u = \lambda A(u), \textup{ per qualche } \lambda \in [0,1]\right\rbrace \]
	è limitato in $H^1_0(U)$. Se $u \in B$, allora $A(u) = \frac{u}{\lambda}$, quindi $u \in H^2(U)\cap H^1_0(U)$ e, preso $v= u$ nella formulazione variazionale che definisce $A$, 
	\[ \int_U \left( \abs{Du}^2 + \mu \abs{u}^2\right) d\mathcal{L}^n = -\int_U\lambda b(Du)u d\mathcal{L}^n.\]
	Essendo $b$ sub-lineare, usando la disuguaglianza di Young pesata con $\varepsilon=\frac{1}{2C}$ ed dal fatto che $\lambda \le 1$,
	\begin{align*}
		-\int_U\lambda b(Du)u d\mathcal{L}^n &\le \lambda C \int_U (1+\abs{Du})\abs{u}d\mathcal{L}^n\le C \int_U (1+\abs{Du})\abs{u}d\mathcal{L}^n\le \\
		&\le \frac{1}{2} \int_U \abs{Du}^2d\mathcal{L}^n + C\int_U \left( 1+\abs{u}^2\right) d\mathcal{L}^n,
	\end{align*}
	quindi 
	\[ \frac{1}{2} \int_U \abs{Du}^2d\mathcal{L}^n + (\mu-C) \int_U \abs{u}^2 d\mathcal{L}^n \le C  \]
	dove $C$ non dipende da $\lambda$ e nemmeno da $u$. Se quindi $\mu > C$, $\norm{u}_{H^1_0(U)}\le C$, ossia $B$ è limitato in $H^1_0(U)$.
	
	Possiamo concludere che, per il Teorema di Schaefer, esiste $u \in H^1_0(U)\cap H^2(U)$ tale che $A(u)=u$ che, combinato con \eqref{perilgranfinale}, permette di affermare che per ogni $v \in H^1_0(U)$
	\[ \int_U \left( Du\cdot Dv + \mu uv)\right) d\mathcal{L}^n = -\int_U b(Du)vd\mathcal{L}^n, \]
	da cui la tesi.
\end{proof}
\begin{osservazione}
	Dal Teorema precedente abbiamo direttamente l'esistenza di una soluzione forte. In più, la teoria della regolarità per questo problema differenziale può essere ricondotta a quella del caso lineare, e quindi al Metodo di Stampacchia: siccome $u \in H^2(U)$, allora $Du \in H^1(U)$, quindi $f = -b(Du) \in H^1(U)$. Aumentando le ipotesi di regolarità su $b$ si può proseguire. Non entriamo nel dettaglio: la morale è che la non linearità è più un problema per l'esistenza che per la regolarità.
\end{osservazione}

\section{Metodo di monotonia}
\begin{definizione}
	Siano $U$ un aperto limitato di $\R^n$ con $\partial U$ di classe $C^\infty$, $a: \R^n \to \R^n$ un'applicazione di classe $C^\infty$ ed $f: U \to \R$ un'applicazione continua. Diciamo che $u$ è una \emph{soluzione classica}\index[analitico]{soluzione!classica} del problema di Dirichlet
	\[ 	\begin{cases}
		-\div(a(Du)) = f & \textup{in } U,\\
		u = 0 & \textup{su } \partial U,
	\end{cases} \]
	se $u \in C^2(U)\cap C(\overline{U})$, $-\div(a(Du(x))) = f(x)$ per ogni $x \in U$ e $u(x) = 0$ per ogni $x \in \partial U$.
\end{definizione}
\begin{osservazione}
	Se $a = \textup{Id}$, allora ritroviamo il problema di Poisson, visto nel capitolo \ref{Equazioni ellittiche lineari del secondo ordine}.
\end{osservazione}
\begin{osservazione}
	Se esiste $F: \R^n \to \R$ tale che $a = DF$, ossia $a$ ammette potenziale scalare, allora l'equazione 
	\[ -\div(a(Du)) = f  \textup{ in } U \]
	è l'equazione di Eulero--Lagrange associata al funzionale
	\[ I[u] = \int_U (F(Du) - fu) d\mathcal{L}^n, \]
	quindi possiamo utilizzare le tecniche del Calcolo delle Variazioni, come visto nel capitolo \ref{Alcuni complementi di calcolo delle variazioni}, per studiare il problema. Infatti, con le convenzioni sulle notazioni introdotte nel capitolo \ref{Alcuni complementi di calcolo delle variazioni},
	\begin{align*}
		D_{\xi_i}L(x,s,\xi) &= D_\xi F(\xi), & D_s L(x,s,\xi) = -f
	\end{align*}
	per cui l'equazione di Eulero--Lagrange è 
	\[ -\sum_{i=1}^{n} D_{x_i}(D_{\xi_i}F(Du)) = f  \textup{ in } U \]
	ossia, in forma vettoriale,
	\[ -\div(DF(Du)) = f  \textup{ in } U. \]
\end{osservazione}

Consideriamo per un attimo la situazione in cui $a$ ammette potenziale scalare $F$ convesso. Allora, per ogni $\xi,\eta \in \R^n$
\begin{align*}
	(a(\xi) - a(\eta))\cdot (\xi - \eta) &= (DF(\xi)- DF(\eta))\cdot(\xi-\eta) = \\
	&= \sum_{i=1}^n (D_iF(\xi)- D_iF(\eta))\cdot(\xi_i-\eta_i)
\end{align*}
e, per il Teorema di Lagrange, esiste $t \in [0,1]$ tale che
\[ 	(a(\xi) - a(\eta))\cdot (\xi - \eta) = \sum_{i,j=1}^n D^2_{i,j}F(\xi + t(\eta - \xi))(\xi_j - \eta_j)(\xi_i - \eta_i).  \]
A questo punto, dalla convessità di $F$, si deduce che
\[ (a(\xi) - a(\eta))\cdot (\xi - \eta) \ge 0. \]
Motivati da questo calcolo, diamo la seguente definizione.
\begin{definizione}
	Diciamo che un'applicazione $a: \R^n \to \R^n$ è \emph{monotona}\index[index]{monotonia}, se per ogni $\xi,\eta \in \R^n$
	\[ (a(\xi) - a(\eta))\cdot (\xi - \eta) \ge 0. \]
\end{definizione}
\begin{osservazione}
	In $\R$, una funzione monotona, nel senso dato sopra, altro non è che una funzione crescente. Infatti, se $\xi \le \eta$, l'unica possibilità per rispettare la disuguaglianza di monotonia è $a(\xi) \le a(\eta)$.
\end{osservazione}

Nel seguito, salvo diversa specificazione, supporremo sempre che $U$ sia un aperto limitato di $\R^n$ non vuoto con $\partial U$ di classe $C^\infty$ e che siano assegnati $f \in L^2(U)$ ed $a: \R^n \to \R^n$ di classe $C^\infty$ tale che
\begin{itemize}
	\item per ogni $\xi,\eta \in \R^n$
	\[ (a(\xi) - a(\eta))\cdot (\xi - \eta) \ge 0. \]
	In altre parole, $a$ è monotona,
	\item esiste $C >0$ tale che per ogni $\xi \in \R^n$ 
	\[ \abs{a(\xi)} \le C(1+\abs{\xi}).\]
	In altre parole, $a$ è sub-lineare,
	\item esistono $\alpha > 0$ e $\beta \ge 0$ tali che per ogni $\xi \in \R^n$
	\[ a(\xi)\cdot \xi \ge \alpha \abs{\xi}^2 - \beta.\]
	In altre parole, $a$ è coercitiva.
\end{itemize}

Vediamo, a titolo di esempio, che esistono funzioni $a$ con le ipotesi richieste.
\begin{esempio}
	Consideriamo $a = \textup{Id}$. Allora per ogni $\xi,\eta \in \R^n$
	\[ (\xi - \eta) \cdot (\xi-\eta) = \abs{\xi - \eta}^2 \ge 0. \]
	Inoltre,
	\[ \abs{\xi}\le 1 + \abs{\xi} \]
	e
	\[ \xi \cdot \xi = \abs{\xi}^2. \]
	Quindi l'identità, che è ovviamente di classe $C^\infty$, è una scelta ammissibile.
\end{esempio}
Con lievi variazioni rispetto all'esempio precedente si può dimostrare che se $a$ è lineare, allora è una scelta ammissibile. Se però ci fermassimo qui, potremmo tranquillamente tornare al capitolo \ref{Equazioni ellittiche lineari del secondo ordine} perché avremmo dei riscalamenti dell'equazione di Poisson. Vediamo che in realtà esistono anche applicazioni non lineari ammissibili.
\begin{esempio}
	Consideriamo l'applicazione $a(\xi) = \xi + \arctan \left( \xi_1\right) e_1$, chiaramente di classe $C^\infty$. Allora per ogni $\xi,\eta \in \R^n$
	\begin{align*}
		&\left( \left( \xi+\arctan \left( \xi_1\right) e_1 \right)  - \left( \eta+\arctan \left( \eta_1\right)e_1 \right) \right)  \cdot (\xi-\eta) = \\
		&=\abs{\xi - \eta}^2 + (\arctan\xi_1-\arctan\eta_1)\cdot(\xi_1 - \eta_1) \ge 0.
	\end{align*}
	Inoltre,
	\[ \abs{\xi + \arctan \left( \xi_1\right)e_1 }\le \abs{\xi} + \frac{\pi}{2} \le (1+\frac{\pi}{2})(1+\abs{\xi})\]
	e, per la disuguaglianza di Young pesata con $\varepsilon = \frac{1}{2}$,
	\[ (\xi+\arctan \left( \xi_1\right)e_1) \cdot \xi = \abs{\xi}^2 +\arctan \left( \xi_1\right)\xi_1 \ge \abs{\xi}^2 - \frac{\pi}{2}\abs{\xi} \ge  \frac{1}{2} \abs{\xi}^2 - \beta\]
	con $\beta \ge 0$.
\end{esempio}

Vediamo ora come indebolire la definizione di soluzione del problema 
\[ 	\begin{cases}
	-\div(a(Du)) = f & \textup{in } U,\\
	u = 0 & \textup{su } \partial U.
\end{cases} \]
Se $f: U \to \R$ continua ed $u$ è una soluzione classica, allora per ogni $x \in U$
\[ -\div(a(Du(x))) = f(x). \]
Moltiplicando ambo i membri per una certa $\varphi \in C^\infty_c(U)$ ed integrando otteniamo 
\[ -\int_U \div(a(Du)) \varphi d\mathcal{L}^n = \int_U f \varphi d\mathcal{L}^n. \]
Per le formule di Gauss--Green abbiamo quindi, per ogni $\varphi \in C^\infty_c(U)$
\[ \int_U a(Du) \cdot D\varphi d\mathcal{L}^n = \int_U f \varphi d\mathcal{L}^n. \]
Osserviamo che l'equazione precedente ha significato con $f \in L^2(U)$ e $u,\varphi \in H^1_0(U)$.
\begin{definizione}
	Diciamo che $u \in H^1_0(U)$ è una \emph{soluzione debole}\index[analitico]{soluzione!debole} del problema di Dirichlet
	\[ 	\begin{cases}
		-\div(a(Du)) = f & \textup{in } U,\\
		u = 0 & \textup{su } \partial U,
	\end{cases} \]
	se per ogni $v \in H^1_0(U)$, 
	\[ \int_U a(Du) \cdot Dv d\mathcal{L}^n = \int_U f v d\mathcal{L}^n. \]
\end{definizione}
\begin{osservazione}
	La precedente definizione è ben posta, infatti, dalla sub-linearità di $a$,
	\[ \int_U a(Du) \cdot Dv d\mathcal{L}^n \le \int_U \abs{a(Du)} \abs{Dv} d\mathcal{L}^n \le C \norm{Dv}_{L^2(U)}\left( \mathcal{L}^n(U)+\norm{Du}_{L^2(U)}\right)<+\infty  \]
	e, chiaramente,
	\[ \int_U f v d\mathcal{L}^n \le \norm{f}_{L^2(U)}\norm{v}_{L^2(U)} < +\infty.\]
\end{osservazione}

Volendo affrontare ora il problema dell'esistenza di soluzioni deboli, iniziamo con un risultato di carattere tecnico che ci sarà utile nel seguito.
\begin{lemma}\label{tecnicopermetodidimonotonia}
	Sia $v: \R^n \to \R^n$ un'applicazione continua. Se esiste $r> 0$ tale che per ogni $x \in \partial B(0,r)$
	\[ v(x) \cdot x \ge 0, \]
	allora esiste $x \in \overline{\B(0,r)}$ tale che $v(x) = 0$.
\end{lemma}
\begin{proof}
	Per assurdo, supponiamo che per ogni $x \in \overline{\B(0,r)}$ si abbia $v(x) \ne 0$. Consideriamo l'applicazione $w: \overline{\B(0,r)} \to 	\partial \B(0,r)$ continua tale che 
	\[ w(x) = -\frac{r}{\abs{v(x)}}v(x). \]
	Per il Teorema di Brouwer, esiste $z \in \overline{\B(0,r)}$ tale che $w(z) = z$. In particolare deve essere $z \in \partial\B(0,r)$, quindi 
	\[ r^2 = z \cdot z = w(z) \cdot z =  -\frac{r}{\abs{v(z)}}v(z) \cdot z \le 0,\]
	che è assurdo.
\end{proof}
Siamo dunque pronti per dimostrare l'esistenza di soluzioni deboli per il problema che stiamo studiando: affronteremo la dimostrazione seguendo il cosiddetto  \emph{metodo di Browder e Minty}.
\begin{teorema}
	Esiste $u \in H^1_0(U)$ soluzione debole del problema di Dirichlet
	\[ 	\begin{cases}
		-\div(a(Du)) = f & \textup{in } U,\\
		u = 0 & \textup{su } \partial U.
	\end{cases} \]
\end{teorema}
\begin{proof}
	Fissiamo innanzitutto $(w_j)$, una base hilbertiana liscia di $H^1_0(U)$.\footnote{Ottenibile, ad esempio, dall'operatore $-\Delta$ (debole), come visto alla fine del capitolo \ref{Equazioni ellittiche lineari del secondo ordine}.}
	Fissato $m \in \N$, mostriamo che esiste $u_m \in H^1_0(U)$ tale che
	\[ u_m = \sum_{j=0}^{m}d_jw_j \]
	e per ogni $k = 0,\dots, m$
	\[ \int_U a(Du_m)\cdot Dw_k d\mathcal{L}^n = \int_U fw_k d\mathcal{L}^n.\footnote{In altre parole, vorremmo scegliere i coefficienti $d_j$ in modo che $u_m$ sia soluzione debole della proiezione del problema 
		\[ 	\begin{cases}
			-\div(a(Du)) = f & \textup{in } U,\\
			u = 0 & \textup{su } \partial U,
		\end{cases} \]
		sul sottospazio vettoriale di $H^1_0(U)$ di dimensione finita generato da $w_0,\dots,w_m$.}\]
	Per fare ciò, consideriamo l'applicazione $v: \R^{m+1} \to \R^{m+1}$ tale che per ogni $k =0,\dots, m$
	\[ v_k (d) = \int_U \left( a\left( \sum_{j=0}^{m}d_jDw_j\right) \cdot Dw_k -fw_k\right) d\mathcal{L}^n.\]
	Dal fatto che $a$ è sub-lineare, l'applicazione $v$ è ben definita. Inoltre è continua per il Teorema della convergenza dominata. Ora, per ogni $d \in \R^m$ 
	\begin{align*}
		v(d) \cdot d &= \sum_{j=0}^{m}v_j(d)d_j = \sum_{j=0}^{m} \int_U \left(  a\left( \sum_{k=0}^{m}d_kDw_k\right) \cdot d_jDw_j -d_jfw_j\right) d\mathcal{L}^n = \\
		&= \int_U \left( a\left( \sum_{k=0}^{m}d_kDw_k\right) \cdot \sum_{j=0}^{m}d_jDw_j -\sum_{j=0}^{m}d_jfw_j\right) d\mathcal{L}^n,
	\end{align*}
	quindi, per la coercitività di $a$,
	\[ v(d) \cdot d \ge \alpha \int_U \abs*{\sum_{j=0}^{m}d_jDw_j}^2 d\mathcal{L}^n -\beta\mathcal{L}^n(U) - \sum_{j=0}^{m} d_j \int_U fw_jd\mathcal{L}^n. \]
	Ora, nel primo integrale i doppi prodotti sono tutti nulli in quanto $w_0,\dots,w_m$ sono tra loro ortogonali in $H^1_0(U)$, quindi 
	\[ \alpha \int_U \abs*{\sum_{j=0}^{m}d_jDw_j}^2 d\mathcal{L}^n = \alpha \sum_{j=0}^{m} d_j^2 \int_U \abs{Dw_j}^2d\mathcal{L}^n = \alpha \abs{d}^2 \]
	Per quanto riguarda l'ultimo addendo invece, detta $u \in H^1_0(U)$ la soluzione debole del problema di Dirichlet
	\[ \begin{cases}
		-\Delta u = f & \textup{in } U,\\
		u= 0 & \textup{su } \partial U,
	\end{cases} \]
	per la disuguaglianza di Young pesata con $\varepsilon = \frac{\alpha}{2}$, la formulazione variazionale e la disuguaglianza di Holder,
	\begin{align*}
		d_j  \int_U fw_jd\mathcal{L}^n &\le \abs{d_j} \abs*{ \int_U fw_jd\mathcal{L}^n}\le \frac{\alpha}{2} \abs{d_j}^2 + C \left( \int_U fw_j d\mathcal{L}^n\right) ^2 \le \\
		&\le \frac{\alpha}{2} \abs{d_j}^2 + C\left( \int_U Du \cdot Dw_j d\mathcal{L}^n\right) ^2 \le \frac{\alpha}{2} \abs{d_j}^2 + C\norm{u}_{H^1_0(U)}^2\norm{w_j}_{H^1_0(U)}^2 \le \\
		&\le \frac{\alpha}{2} \abs{d_j}^2 + C, 
	\end{align*}
	quindi
	\[ - \sum_{j=0}^{m} d_j \int_U fw_jd\mathcal{L}^n \ge -\frac{\alpha}{2} \abs{d}^2 - C. \]
	In definitiva, 
	\[ v(d) \cdot d \ge  \frac{\alpha}{2} \abs{d}^2 - C.\]
	Scelto $r\ge \sqrt{\frac{2C}{\alpha}}$, per il Lemma \eqref{tecnicopermetodidimonotonia}, esiste $d \in \overline{\B(0,r)}$ tale che $v(d) = 0$, ossia per ogni $k = 0,\dots, m$
	\begin{equation}\label{projectionproblemmonotonia}
		\int_U a(Du_m)\cdot Dw_k d\mathcal{L}^n = \int_U fw_k d\mathcal{L}^n.
	\end{equation}
	
	Mostriamo ora che esiste $C$, dipendente solo da $U$ ed $a$, tale che per ogni $m \in \N$
	\[ \norm{u_m}_{H^1_0(U)} \le C(1+\norm{f}_{L^2(U)}). \]
	Per fare ciò, moltiplichiamo per $d_k$ ambo i membri di \eqref{projectionproblemmonotonia} e sommiamo per $k=0,\dots m$, ottenendo
	\begin{equation}\label{miserve}
		\int_U a(Du_m) \cdot Du_m d\mathcal{L}^n = \int_U fu_m d\mathcal{L}^n.
	\end{equation}
	Dalla coercitività di $a$, 
	\[ \int_U a(Du_m) \cdot Du_m d\mathcal{L}^n \ge \alpha \int_U \abs{Du_m}^2d\mathcal{L}^n -\beta \mathcal{L}^n(U), \]
	quindi combinando la disuguaglianza di Holder e la disuguaglianza di Poincaré,
	\[ \alpha \int_U \abs{Du_m}^2d\mathcal{L}^n \le C +\int_U fu_m d\mathcal{L}^n \le C +\norm{f}_{L^2(U)}\norm{u_m}_{L^2(U)} \le C+C\norm{f}_{L^2(U)}\norm{u_m}_{H^1_0(U)}, \]
	perciò, per la disuguaglianza di Young pesata con $\varepsilon = \frac{\alpha}{2C}$,
	\[ \alpha \int_U \abs{Du_m}^2d\mathcal{L}^n \le C + \frac{\alpha}{2} \int_U \abs{Du_m}^2d\mathcal{L}^n +C\norm{f}_{L^2(U)}^2, \]
	allora
	\[ \frac{\alpha}{2} \int_U \abs{Du_m}^2d\mathcal{L}^n \le C(1 + \norm{f}_{L^2(U)}^2) \]
	e, per la sub-additività della radice quadrata,
	\[ \norm{u_m}_{H^1_0(U)} \le C(1+\norm{f}_{L^2(U)}). \]
	
	Da quanto appena dimostrato, si deduce che $(u_m)$ è limitata in $H^1_0(U)$, quindi, per il Teorema di Eberlein--Smulian, esistono $(u_{m_j})$ e $u\in H^1_0(U)$ tale che $u_{m_j} \rightharpoonup u$ in $H^1_0(U)$. %Per il Teorema di Rellich--Kondrachov combinato con l'unicità del limite debole, $u_{m_j} \to u$ in $L^2(U)$. 
	In particolare, $u_{m_j} \rightharpoonup u$ in $L^2(U)$. Essendo $a$ sub-lineare, 
	\[ \int_U \abs{a(Du_{m_j})}^2 d\mathcal{L}^n \le C(\mathcal{L}^n(U) + \norm{u_{m_j}}_{H^1_0(U)}^2)<+\infty, \]
	quindi $(a(Du_{m_j}))$ limitata in $L^2(U;\R^n)$ e, per il Teorema di Eberlein--Smulian, esiste $\xi \in L^2(U;\R^n)$ tale che, a meno di un'ulteriore sottosuccessione, $a(Du_{m_j}) \rightharpoonup \xi$ in $L^2(U;\R^n)$. 
	
	Per ogni $v \in H^1_0(U)$, scrivendo l'equazione $\eqref{projectionproblemmonotonia}$ per $u_{m_j}$ e passando al limite per $j \to +\infty$ otteniamo per ogni $k \in \N$
	\[ \int_U \xi \cdot Dw_k d\mathcal{L}^n = \int_U fw_k d\mathcal{L}^n \]
	che moltiplicata ambo i membri per $(v|w_k)_{H^1_0(U)}$ e sommata su $k$ fornisce
	\begin{equation}\label{esistenzamonotonia2}
		\int_U \xi \cdot Dv d\mathcal{L}^n = \int_U fvd\mathcal{L}^n.
	\end{equation}
	
	Dalla monotonia di $a$, per ogni $m \in \N$ e per ogni $w \in H^1_0(U)$,
	\[ \int_U \left( a(Du_{m_j})- a(Dw)\right) \cdot (Du_{m_j}-Dw) d\mathcal{L}^n \ge 0,\]
	ossia, utilizzando l'equazione \eqref{miserve},
	\[ \int_U fu_{m_j} d\mathcal{L}^n  -  \int_U a(Dw)\cdot Du_{m_j} d\mathcal{L}^n - \int_U a(Du_{m_j})\cdot Dw d\mathcal{L}^n +\int_U a(Dw)\cdot Dw d\mathcal{L}^n\ge 0\]
	che per $j \to +\infty$ diventa
	\[ \int_U fu d\mathcal{L}^n  -  \int_U a(Dw)\cdot Du d\mathcal{L}^n - \int_U \xi\cdot Dw d\mathcal{L}^n +\int_U a(Dw)\cdot Dw d\mathcal{L}^n \ge 0,\]
	che, utilizzando l'equazione \eqref{esistenzamonotonia2} con $v = u$, può essere scritta come
	\[ \int_U (\xi -a(Dw)) \cdot (Du - Dw)d\mathcal{L}^n\ge 0.\]
	A questo punto, fissato $v \in H^1_0(U)$, preso $w = u-\lambda v$, con $\lambda \in \left]0,1\right] $, dalla precedente equazione abbiamo 
	\[ \int_U (\xi -a(Du-\lambda Dv)) \cdot Dvd\mathcal{L}^n\ge 0.\]
	Dal fatto che $a$ è continua, 
	\[ \lim\limits_{\lambda \to 0}a(Du-\lambda Dv)\cdot Dv = a(Du) \textup{ q.o. in } U, \]
	inoltre, per la sub-linearità di $a$ e per la disuguaglianza di Young,
	\begin{align*}
		\abs{a(Du-\lambda Dv)\cdot Dv}&\le \abs{a(Du-\lambda Dv)}\abs{Dv}\le C(1+\abs{Du - \lambda Dv})\abs{Dv} \le \\
		&\le  C(1 + \abs{Dv}^2 + \abs{Du}^2) \in L^1(U)
	\end{align*}
	quindi, per il Teorema della convergenza dominata, passando al limite per $\lambda \to 0$ si ottiene che per ogni $v \in H^1_0(U)$
	\[ \int_U (\xi - a(Du))\cdot Dv d\mathcal{L}^n \ge 0. \]
	A meno di scambiare $v$ con $-v$, vale anche
	\[ \int_U (\xi - a(Du))\cdot Dv d\mathcal{L}^n \le 0, \]
	quindi 
	\[ \int_U (\xi - a(Du))\cdot Dv d\mathcal{L}^n = 0. \]
	In particolare, dalla \eqref{esistenzamonotonia2}, per ogni $v \in H^1_0(U)$
	\[ \int_U a(Du)\cdot Dv d\mathcal{L}^n = \int_U \xi \cdot Dv d\mathcal{L}^n = \int_U fv d\mathcal{L}^n, \]
	da cui la tesi.
\end{proof}
\begin{osservazione}
	Vale la pena osservare la strategia dimostrativa utilizzata per il Teorema precedente nasce per risolvere il seguente problema: a causa della non linearità di $a$, anche sapendo che $u_{m_j} \rightharpoonup u$ in $H^1_0(U)$, non è possibile dedurre che $a(Du_{m_j}) \rightharpoonup a(Du)$ in un qualche senso. La chiave per superare questo stallo è proprio la monotonia.
\end{osservazione}
\begin{definizione}
	Diciamo che un'applicazione $a: \R^n \to \R^n$ è \emph{strettamente monotona}\index[index]{monotonia!stretta}, se esiste $\vartheta >0$ tale che per ogni $\xi,\eta \in \R^n$
	\[ (a(\xi) - a(\eta))\cdot (\xi - \eta) \ge \vartheta \abs{\xi -\eta}^2. \]
\end{definizione}
\begin{teorema}
	Se $a$ è strettamente monotona, allora la soluzione debole del problema di Dirichlet
	\[ 	\begin{cases}
		-\div(a(Du)) = f & \textup{in } U,\\
		u = 0 & \textup{su } \partial U,
	\end{cases} \]
	è unica.
\end{teorema}
\begin{proof}
	Supponiamo esistano $u,\tilde{u}$ soluzioni deboli del problema di Dirichlet
	\[ 	\begin{cases}
		-\div(a(Du)) = f & \textup{in } U,\\
		u = 0 & \textup{su } \partial U.
	\end{cases} \]
	Allora, per ogni $v \in H^1_0(U)$
	\[ \int_U a(Du)\cdot Dv d\mathcal{L}^n = \int_U fv d\mathcal{L}^n = \int_U a(D\tilde{u})\cdot Dv d\mathcal{L}^n, \]
	ossia 
	\[ \int_U(a(Du)-a(D\tilde{u}))\cdot Dv d\mathcal{L}^n = 0. \]
	Scelto $v = u - \tilde{u}$,
	\[ 0 = \int_U(a(Du)-a(D\tilde{u}))\cdot (Du - D\tilde{u})d\mathcal{L}^n \ge \vartheta \int_U \abs{Du- D\tilde{u}}^2 d\mathcal{L}^n = \norm{u-\tilde{u}}_{H^1_0(U)}^2.\]
	L'unica possibilità è $\norm{u-\tilde{u}}_{H^1_0(U)}^2 = 0$, ossia $u = \tilde{u}$ q.o. in $U$.
\end{proof}

Un'ampia classe di scelte per $a$ grazie alle quali abbiamo esistenza ed unicità della soluzione debole è la seguente.
\begin{esempio}
	Consideriamo $a(\xi) = \xi +b(\xi)$ dove 
	\[ b(\xi) = (b_1(\xi_1),\dots,b_n(\xi_n)) \]
	con $b_i : \R\to \R$ monotone crescenti, limitate e di classe $C^\infty$ per ogni $i=1,\dots,n$. Allora per ogni $\xi,\eta \in \R^n$
	\begin{align*}
		(a(\xi)-a(\eta))\cdot (\xi-\eta) &= \abs{\xi-\eta}^2 +(b(\xi)-b(\eta))\cdot(\xi-\eta) = \\
		&=\abs{\xi-\eta}^2 + \sum_{i=1}^{n}(b_i(\xi_1)-b_i(\eta_i))(\xi_i-\eta_i) \ge \abs{\xi-\eta}^2.
	\end{align*}
	Inoltre,
	\[ \abs{a(\xi)} \le \abs{\xi} + \abs{b(\xi)} \le \abs{\xi} + C \le (1+C)(1+\abs{\xi})\]
	e, per la disuguaglianza di Cauchy--Schwarz e la disuguaglianza di Young pesata con $\varepsilon = \frac{1}{2}$
	\[ a(\xi) \cdot \xi = \xi^2 + b(\xi) \cdot \xi \ge \xi^2 -\abs{b(x)} \abs{\xi} \ge  \frac{1}{2} \abs{\xi}^2 - \beta\]
	con $\beta \ge 0$.
\end{esempio}
\begin{osservazione}
	In ipotesi di stretta monotonia, è possibile ricondursi al \emph{Metodo di Stampacchia}, già usato nel capitolo \ref{Equazioni ellittiche lineari del secondo ordine}, per lo studio della regolarità delle soluzioni.
\end{osservazione}

\section{Metodo delle sopra/sotto-soluzioni}
\begin{proposizione}\label{monotonia}
Consideriamo $U\subseteq \R^n$ aperto limitato con bordo di classe $C^1$, $f_1,f_2 \in L^2(U)$ e $u_1,u_2 \in H^1_0(U)$ le soluzioni deboli dei problemi di Dirichlet
\begin{align*}
	&\begin{cases}
		-\Delta u = f_1 & \textup{in } U,\\
		u = 0 & \textup{su } \partial U,
	\end{cases} & &\begin{cases}
	-\Delta u = f_2 & \textup{in } U,\\
	u = 0 & \textup{su } \partial U.
	\end{cases}
\end{align*}
Se $f_2 \le f_1$ \emph{q.o.} in $U$, allora $u_2 \le u_1$ \emph{q.o.} in $U$.
\end{proposizione}
\begin{proof}
Osserviamo, anzitutto, che $(u_2 - u_1)^+ \in H^1_0(U)$ con 
\[ D(u_2-u_1)^+ = \begin{cases}
	D(u_2-u_1) & \textup{se } u_2 \ge u_1, \\
	0 & \textup{altrimenti}.
\end{cases} \]

Ora, per linearità, 
\[ \int_U D(u_2-u_1) \cdot D(u_2-u_1)^+ d\mathcal{L}^n = \int_U (f_2 - f_1)(u_2-u_1)^+ d\mathcal{L}^n \le 0, \]
ma
\[ \int_U D(u_2-u_1) \cdot D(u_2-u_1)^+ d\mathcal{L}^n = \int_U \abs{D(u_2-u_1)^+}^2 d\mathcal{L}^n, \]
allora $\norm{(u_2-u_1)^+}_{H^1_0(U)}\le 0$, da cui $(u_2-u_1)^+ = 0$ q.o. in $U$, cioè $u_2 \le u_1$ q.o. in $U$.
\end{proof}

\begin{definizione}
	Consideriamo $U\subseteq \R^n$ aperto limitato con bordo di classe $C^\infty$, $C >0$ e $f: \R \to \R$ di classe $C^\infty$ tale che, per ogni $ x\in \R$, $\abs{f'(x)}\le C$.\footnote{In particolare, grazie alla disuguaglianza di Lagrange, $f$ è lipschitziana e sub-lineare.}
	Diciamo che $u \in H^1_0(U)$ è una \emph{soluzione debole}\index[analitico]{soluzione!debole} del problema di Dirichlet
	\[ 	\begin{cases}
		-\Delta u = f(u) & \textup{in } U,\\
		u = 0 & \textup{su } \partial U,
	\end{cases} \]
	se per ogni $v \in H^1_0(U)$
	\[ \int_U Du \cdot Dv d\mathcal{L}^n = \int_U f(u)vd\mathcal{L}^n. \]
\end{definizione}
\begin{definizione}
Consideriamo $U\subseteq \R^n$ aperto limitato con bordo di classe $C^\infty$, $C >0$ e $f: \R \to \R$ di classe $C^\infty$ tale che, per ogni $ x\in \R$, $\abs{f'(x)}\le C$. Diciamo che $\overline{u} \in H^1(U)$ è una \emph{sopra-soluzione debole}\index[analitico]{sopra-soluzione debole} del problema di Dirichlet
\[ 	\begin{cases}
	-\Delta u = f(u) & \textup{in } U,\\
	u = 0 & \textup{su } \partial U,
\end{cases} \]
se per ogni $v \in H^1_0(U)$ tale che $v\ge 0$ \emph{q.o.} in $U$
\[ \int_U D\overline{u} \cdot Dv d\mathcal{L}^n \ge \int_U f(\overline{u})vd\mathcal{L}^n. \]
\end{definizione}
\begin{definizione}
	Consideriamo $U\subseteq \R^n$ aperto limitato con bordo di classe $C^\infty$, $C >0$ e $f: \R \to \R$ di classe $C^\infty$ tale che, per ogni $ x\in \R$, $\abs{f'(x)}\le C$.
	Diciamo che $\underline{u} \in H^1(U)$ è una \emph{sotto-soluzione debole}\index[analitico]{sotto-soluzione debole} del problema di Dirichlet
	\[ 	\begin{cases}
		-\Delta u = f(u) & \textup{in } U,\\
		u = 0 & \textup{su } \partial U,
	\end{cases} \]
	se per ogni $v \in H^1_0(U)$ tale che $v\ge 0$ \emph{q.o.} in $U$
	\[ \int_U D\underline{u} \cdot Dv d\mathcal{L}^n \le \int_U f(\underline{u})vd\mathcal{L}^n. \]
\end{definizione}
In particolare, una soluzione debole è sia sopra-soluzione che sotto-soluzione.

Il principale risultato della sezione è il seguente: l'idea di base è che è più facile trovare sopra/sotto-soluzioni deboli rispetto che soluzioni deboli.
\begin{teorema}\label{metodosubsuper}
Consideriamo $U\subseteq \R^n$ aperto limitato con bordo di classe $C^\infty$, $C >0$ e $f: \R \to \R$ di classe $C^\infty$ tale che, per ogni $ x\in \R$, $\abs{f'(x)}\le C$. Se esistono una sopra-soluzione debole $\overline{u} \in H^1(U)$ ed una sotto-soluzione debole $\underline{u} \in H^1(U)$ del problema di Dirichlet
\[ 	\begin{cases}
	-\Delta u = f(u) & \textup{in } U,\\
	u = 0 & \textup{su } \partial U,
\end{cases} \]
 tali che $\underline{u} \le 0$ e $\overline{u} \ge 0$ su $\partial U$ nel senso delle tracce e $\underline{u} \le \overline{u}$ \emph{q.o.} in $U$, allora esiste una soluzione debole $u \in H^1_0(U)$ tale che $\underline{u} \le u \le \overline{u}$ \emph{q.o.} in $U$.
\end{teorema}
\begin{proof}
Innanzitutto, essendo $f' \ge -C$, esiste $\lambda >0$ sufficientemente grande tale che 
\[ f' + \lambda \ge 0, \]
quindi l'applicazione $\left\lbrace x \mapsto f(x) + \lambda x\right\rbrace$ è crescente.

Detta $u_0 = \underline{u}$, costruiamo ricorsivamente una successione $(u_k)$ in $H^1_0(U)$ tale che $u_{k+1}$ è la soluzione debole del problema di Dirichlet
\[ 	\begin{cases}
	-\Delta u_{k+1} + \lambda u_{k+1} = f(u_k) + \lambda u_k & \textup{in } U,\\
	u_{k+1} = 0 & \textup{su } \partial U.
\end{cases} \]
Mostriamo, per induzione, che $(u_k)$ è crescente q.o. in $U$. Per prima cosa, detto $T$ l'operatore di traccia su $U$, siccome $T(u_0-u_1) = Tu_0 - Tu_1 \le 0$ q.o. in $U$, allora $T(u_0-u_1)^+ = 0$ q.o. in $U$, ossia $(u_0-u_1)^+ \in H^1_0(U)$ con 
\[ D(u_0-u_1)^+ = \begin{cases}
	D(u_0-u_1) & \textup{se } u_0 \ge u_1, \\
	0 & \textup{altrimenti}.
\end{cases} \]
In particolare, sottraendo alla definizione di sotto-soluzione debole $u_0$ la formulazione variazionale del problema di Dirichlet che ha soluzione debole $u_1$, entrambe valutate per $v= (u_0-u_1)^+$, si ottiene
\[ \int_U \left( D(u_0 - u_1) \cdot D(u_0-u_1)^+ + \lambda(u_0-u_1)(u_0-u_1)^+\right)  d\mathcal{L}^n \le 0,\]
ma
\begin{align*}
	\int_U D(u_0-u_1) \cdot D(u_0-u_1)^+ d\mathcal{L}^n &= \int_U \abs{D(u_0-u_1)^+}^2 d\mathcal{L}^n, \\
	 \int_U (u_0-u_1) \cdot (u_0-u_1)^+ d\mathcal{L}^n &= \int_U \abs{(u_0-u_1)^+}^2 d\mathcal{L}^n,
\end{align*}
quindi 
\[ \int_{U} \left( \abs{D(u_0-u_1)^+}^2 + \lambda \abs{(u_0-u_1)^+}^2\right)  d\mathcal{L}^n \le 0, \]
da cui $u_0 \le u_1$ q.o. in $U$. Supponiamo ora che $u_{k-1} \le u_k$ q.o. in $U$. Siccome $\left\lbrace x \mapsto f(x) + \lambda x\right\rbrace$ è crescente, $f(u_{k-1}) + \lambda u_{k-1} \le f(u_k) + \lambda u_k$ che, ragionando in modo analogo a quanto fatto sopra, da $u_k \le u_{k+1}$ q.o. in $U$.

Mostriamo ora, per induzione, che per ogni $k \in \N$ 
\[ u_k \le \overline{u} \textup{ q.o. in } U. \]
Il caso $k = 0$ segue direttamente dalle ipotesi. Supponiamo quindi che $u_k \le \overline{u}$ q.o. in $U$. Sempre ragionando in modo analogo a quanto fatto sopra segue che $u_{k+1}\le \overline{u}$ q.o. in $U$.

Riassumendo,
\[ \underline{u}\le \dots \le u_k \le u_{k+1} \le \overline{u} \textup{ q.o. in } U,\]
quindi per q.o. $x \in U$ esiste $u(x) = \lim\limits_{k} u_k(x)$. In particolare, $\underline{u} \le u \le \overline{u}$ q.o. in $U$, da cui 
\[ \int_U \abs{u}^2 d\mathcal{L}^n \le \int_U \max\left\lbrace \abs{\underline{u}}^2,\abs{\overline{u}}^2\right\rbrace d\mathcal{L}^n < +\infty, \]
che vuol dire $u \in L^2(U)$, e 
\[ \forall\; k \in N: \abs{u_k-u}^2 \le 2 \max\left\lbrace \abs{\underline{u}}^2,\abs{\overline{u}}^2\right\rbrace \in L^1(U)\]
che, combinato con il Teorema della convergenza dominata, da $u_k \to u$ in $L^2(U)$. Ora, per la sub-linearità di $f$,
\[ \norm{f(u_k)}_{L^2(U)}\le C(1+ \norm{u_k}_{L^2(U)}), \]
quindi, per la formulazione variazionale del problema di Dirichlet che ha soluzione debole $u_k$, valutata per $v=u_k$, e per la disuguaglianza di Young,
\begin{align*}
	\norm{u_k}_{H^1_0(U)}^2 &\le \int_{U} \left( \abs{Du_k}^2 +\lambda \abs{u_k}^2 \right) d\mathcal{L}^n = \int_U \left( f(u_{k-1})+\lambda u_{k-1}\right) u_k d\mathcal{L}^n \le \\
	&\le \int_U \left( f(u_{k})+\lambda u_{k}\right) u_k d\mathcal{L}^n \le C(\norm{f(u_k)}^2_{L^2(U)} + \norm{u_k}_{L^2(U)}^2) \le \\
	&\le C(1 + \norm{u_k}_{L^2(U)}^2) \le C<+\infty.
\end{align*}
In altre parole, $(u_k)$ è limitata in $H^1_0(U)$, quindi per il Teorema di Eberlein--Smulian, a meno di sottosuccessioni, per l'unicità del limite si ha $u_k \rightharpoonup u$ in $H^1_0(U)$. In modo analogo, siccome $(f(u_k))$ è limitata in $L^2(U)$, $f(u_k) \rightharpoonup f(u)$ in $L^2(U)$. Passando quindi al limite per $k \to +\infty$ nella formulazione variazionale del problema di Dirichlet che ha soluzione debole $u_k$ otteniamo per ogni $v \in H^1_0(U)$
\[ \int_{U} \left( Du \cdot Dv +\lambda uv \right) d\mathcal{L}^n = \int_U \left( f(u)v + \lambda uv\right)  d\mathcal{L}^n, \]
ossia
\[ \int_{U} Du \cdot Dv  d\mathcal{L}^n = \int_Uf(u)v d\mathcal{L}^n, \]
che è la tesi.
\end{proof}

Vediamo di seguito un esempio di applicazione del Teorema precedente.
\begin{esempio}
Consideriamo $U\subseteq \R^n$ aperto limitato con bordo di classe $C^\infty$, $C >0$ e $f: \R \to [1,2]$ di classe $C^\infty$ tale che, per ogni $x \in \R$, $\abs{f'(x)}\le C$. Sia $\underline{u} \in H^1_0(U)$ la soluzione debole del problema di Dirichlet
\[ 	\begin{cases}
	-\Delta \underline{u} = 1 & \textup{in } U,\\
	\underline{u} = 0 & \textup{su } \partial U,
\end{cases} \]
allora per ogni $v \in H^1_0(U)$ con $v \ge 0$ \emph{q.o.} in $U$
\[ \int_U D\underline{u}\cdot Dv d\mathcal{L}^n \le \int_U vd\mathcal{L}^n \le \int_U fv d\mathcal{L}^n, \]
quindi $\underline{u}$ è sotto-soluzione debole per il problema di Dirichlet
\[ 	\begin{cases}
	-\Delta u = f(u) & \textup{in } U,\\
	u = 0 & \textup{su } \partial U.
\end{cases} \]
Analogamente, $\overline{u} \in H^1_0(U)$ tale che 
\[ 	\begin{cases}
	-\Delta \overline{u} = 2 & \textup{in } U,\\
	\overline{u} = 0 & \textup{su } \partial U,
\end{cases} \]
è sopra-soluzione debole per il lo stesso problema di Dirichlet. Siccome $1\le 2$, dalla Proposizione \eqref{monotonia}, discende che $\underline{u} \le \overline{u}$ \emph{q.o.} in $U$. Inoltre, $\underline{u} = \overline{u} = 0$ su $\partial U$ nel senso delle tracce. Dal Teorema \eqref{metodosubsuper}, il problema di Dirichlet 
\[ 	\begin{cases}
	-\Delta u = f(u) & \textup{in } U,\\
	u = 0 & \textup{su } \partial U.
\end{cases} \]
ammette soluzione debole.
\end{esempio}

\section{Metodo dei punti critici}
Nel seguito, salvo diversa specificazione, supporremo sempre che $n\ge3$, $U$ sia un aperto limitato di $\R^n$ non vuoto con $\partial U$ di classe $C^\infty$ e che siano assegnate due applicazioni $g: U \times \R \to \R$ di Carathéodory e 
\[ G(x,s) = \int_0^sg(x,t)dt, \]
tali che 
\begin{enumerate}[$(i)$]
	\item  per ogni $\varepsilon>0$ esista $a_\varepsilon \in L^{\frac{2n}{n+2}}(U)$ tale che per q.o. $x \in U$ e per ogni $s \in \R$
	\[ \abs{g(x,s)} \le a_\varepsilon(x) + \varepsilon\abs{s}^{\frac{n+2}{n-2}},\footnote{Attenzione: $\frac{n+2}{n-2} = 2^\ast -1$.} \]
	\item esistono $R>0$ e $q>2$ tali che per q.o. $x \in U$ e per ogni $s \in \R$ tale che $\abs{s} \ge R$
	\[ 0 <qG(x,s) \le g(x,s)s, \]
	\item per ogni $\varepsilon>0$ esista $C_\varepsilon \in \R$ tale che per q.o. $x \in U$ e per ogni $s \in \R$
	\[ G(x,s) \le \varepsilon\abs{s}^2+C_\varepsilon\abs{s}^{2^\ast},\]
\end{enumerate}
ed un operatore differenziale $L$ uniformemente ellittico in forma di divergenza tale che
\begin{itemize}
	\item per ogni $i,j = 1,\dots, n$, 
	\[ a_{ij}=a_{ji}. \]
	In altre parole, $A$ è simmetrica,
	\item per ogni $i,j = 1,\dots, n$, 
	\[ a_{ij} \in L^\infty(U), b_i,c = 0. \]
\end{itemize}

\begin{definizione}
Diciamo che $u \in H^{1}_0(U)$ è una \emph{soluzione debole}\index[analitico]{soluzione!debole} del problema di Dirichlet
	\[ 	\begin{cases}
		-\displaystyle\sum_{i,j=1}^n D_j(a_{ij}D_iu) = g(x,u) & \textup{in } U,\\
		u = 0 & \textup{su } \partial U,
	\end{cases} \]
	se per ogni $v \in H^1_0(U)$ 
	\[ \int_U \sum_{i,j=1}^{n}a_{ij}D_iuD_jv d\mathcal{L}^n = \int_U g(x,u)vd\mathcal{L}^n. \]
\end{definizione}

Consideriamo il funzionale continuo $J: H^1_0(U) \to \R$ tale che
\[ J[u] = \frac{1}{2}\int_U\sum_{i,j=1}^{n}a_{ij}D_iuD_ju d\mathcal{L}^n - \int_U G(x,u)d\mathcal{L}^n. \]
Innanzitutto, per ogni $u,v \in H^1_0(U)$
\begin{align*}
\lim\limits_{\tau \to 0} \frac{J[u+\tau v]- J[u]}{\tau} &= \lim\limits_{\tau \to 0}\left(  \dfrac{\displaystyle\int_U \sum_{i,j=1}^{n}a_{ij}(D_i(u+\tau v)D_j(u+\tau v)-D_iuD_ju) d\mathcal{L}^n}{2\tau} \right. +\\
&-\left. \dfrac{\displaystyle\int_U (G(x,u+\tau v)-G(x,u)) d\mathcal{L}^n}{\tau}\right),
\end{align*}
ma 
\[ \lim\limits_{\tau \to 0}  \dfrac{\displaystyle\int_U \sum_{i,j=1}^{n}a_{ij}(D_i(u+\tau v)D_j(u+\tau v)-D_iuD_ju) d\mathcal{L}^n}{2\tau} = \int_U \sum_{i,j=1}^{n}a_{ij}D_iuD_jv d\mathcal{L}^n \]
e, detta $\varphi(t) = G(u+t\tau v)$, per il Teorema di Lagrange, esiste $\sigma \in \left] 0,1\right[ $ tale che
\begin{align*}
\lim\limits_{\tau \to 0}\dfrac{\displaystyle\int_U (G(x,u+\tau v)-G(x,u)) d\mathcal{L}^n}{\tau} &= \lim\limits_{\tau \to 0}\dfrac{\displaystyle\int_U (\varphi(1)-\varphi(0)) d\mathcal{L}^n}{\tau} = \lim\limits_{\tau \to 0}\dfrac{\displaystyle\int_U \varphi'(\sigma) d\mathcal{L}^n}{\tau} = \\
&= \lim\limits_{\tau \to 0}\int_U g(u+\sigma \tau v) d\mathcal{L}^n = \int_U g(x,u)vd\mathcal{L}^n,
\end{align*}
dove nell'ultimo passaggio abbiamo usato la continuità dell'operatore di Nemytskij, quindi
\[ \braket{dJ[u],v} = J'[u](v) = \lim\limits_{\tau \to 0} \frac{J[u+\tau v]- J[u]}{\tau} = \int_U \sum_{i,j=1}^{n}a_{ij}D_iuD_jv d\mathcal{L}^n - \int_U g(x,u)vd\mathcal{L}^n.\footnote{Attenzione al cambio di notazione rispetto ai capitoli precedenti.}\]
In particolare, $J$ è derivabile ed è il funzionale associato al problema di Dirichlet in esame, nel senso delle equazioni di Eulero--Lagrange. Inoltre, $J$ è differenziabile, infatti, detta $\varphi(t) = G(x,u+tv)-tg(x,u)v$,
\begin{align*}
	&\abs{J[u+v]-J[u]-\braket{dJ[u],v}} = \\
	&\left| \frac{1}{2}\left(\sum_{i,j=1}^{n} \int_U a_{ij}\left( D_i(u+v)D_j(u+v)-D_iuD_ju-2D_iuD_jv\right) d\mathcal{L}^n\right) + \right.\\
	&\left. -\int_U \left( G(x,u+v)-G(x,u)-g(x,u)v\right) d\mathcal{L}^n\right|  \le \\
	&\le \frac{1}{2} \abs*{\sum_{i,j=1}^{n}\int_U a_{ij}D_iuD_jvd\mathcal{L}^n} + \abs*{\int_U (\varphi(1)-\varphi(0))d\mathcal{L}^n},
\end{align*}
quindi, per il Teorema di Lagrange, esiste $\sigma \in \left] 0,1\right[ $ tale che
\begin{align*}
\abs{J[u+v]-J[u]-\braket{dJ[u],v}} &\le \frac{1}{2} \abs*{\sum_{i,j=1}^{n}\int_U a_{ij}D_iuD_jvd\mathcal{L}^n} + \abs*{\int_U \varphi'(\sigma)d\mathcal{L}^n} \le \\
&\le C\norm{v}_{H^1_0(U)}^2 + \int_U \abs{g(u+\sigma v)-g(u)}\abs{v}d\mathcal{L}^n
\end{align*}
e, per la Disuguaglianza di Holder ed il Teorema di Sobolev,
\begin{align*}
	\abs{J[u+v]-J[u]-\braket{dJ[u],v}} &\le C\norm{v}_{H^1_0(U)}^2 +\norm{g(u+tv)-g(u)}_{L^{\frac{2n}{n+2}}(U)} \norm{v}_{L^{\frac{2n}{n-2}}(U)}\le  \\
	&\le C\norm{v}_{H^1_0(U)} \left( \norm{v}_{H^1_0(U)} + \norm{g(u+tv)-g(u)}_{L^{\frac{2n}{n+2}}(U)}\right),
\end{align*}
quindi, per la continuità dell'operatore di Nemytskij,
\[ \lim\limits_{v \to 0} \frac{J[u+v]-J[u]-\braket{dJ[u],v}}{\norm{v}_{H^1_0(U)}}  =0. \]
In modo simile, si vede che $J$ è di classe $C^1$.

\begin{lemma}\label{lemma_tech_minorazioneG}
Per \emph{q.o.} $x \in U$ e per ogni $s \in \R$ tale che $\abs{s}\ge R$, esiste $k(x)$ tale che
\[ G(x,s) \ge k \abs{s}^q. \]
\end{lemma}
\begin{proof}
Vediamo dapprima il caso $s\ge R$. Detta $\varphi(s) = \frac{G(x,s)}{s^q}$, si ha, per la $(ii)$,
\[ \varphi'(s) = \frac{g(x,s)s-qG(x,s)}{s^{q+1}} \ge0, \]
quindi $\varphi(s) \ge \varphi(R)$, ossia
\[ \frac{G(x,s)}{s^q} \ge \frac{G(x,R)}{R^q} \]
o, in altre parole,
\[ G(x,s) \ge k s^q. \]
Il caso $s\le -R$ può essere trattato in modo simile.
\end{proof}
\begin{lemma}\label{Lemma_JPS}
$J$ soddisfa la condizione di Palais--Smale.
\end{lemma}
\begin{proof}
Sia $(u_h)$ in $H^1_0(U)$ tale che $(J[u_h])$ sia limitata e $dJ[u_h] \to 0$ in $H^{-1}(U)$.\footnote{Attenzione: non stiamo applicando il Teorema della rappresentazione di Riesz; vedremo che è molto più comodo ragionare così.}
Per il Teorema di Bolzano--Weierstrass, esistono $(u_{h_k})$ e $m \in \R$ tali che 
\begin{align*}
	J[u_{h_k}] &\to m, & dJ[u_{h_k}] &\to 0.
\end{align*}
In particolare, esiste una successione $(\alpha_k)$ in $\R$ tale che $\alpha_k \to 0$ e
\[ \frac{1}{2} \int_U \sum_{i,j=1}^{n} a_{ij}D_iu_{h_k}D_ju_{h_k}d\mathcal{L}^n-\int_U G(x,u_{h_k}) d\mathcal{L}^n = m + \alpha_k.\]
Sottraendo, a quest'ultima equazione moltiplicata per $q$, 
\[ \int_U \sum_{i,j=1}^{n}a_{ij}D_iu_{h_k}D_ju_{h_k} d\mathcal{L}^n - \int_U g(x,u_{h_k})u_{h_k}d\mathcal{L}^n = \braket{dJ[u_{h_k}],u_{h_k}} \]
otteniamo, a meno di rinominare la successione $\alpha_k$,
\begin{align*}
	&\frac{q-2}{2} \int_U \sum_{i,j=1}^{n} a_{ij}D_i u_{h_k} D_j u_{h_k} d\mathcal{L}^n + \\
	&+ \int_U \left( g(x,u_{h_k})u_{h_k}-qG(x,u_{h_k})\right) d\mathcal{L}^n = qm + \alpha_k - \braket{dJ[u_{h_k}],u_{h_k}}, 
\end{align*}
che può anche essere scritta nella forma
\begin{align*}
	&\frac{q-2}{2} \int_U \sum_{i,j=1}^{n} a_{ij}D_i u_{h_k} D_j u_{h_k} d\mathcal{L}^n + \int_{U \cap \left\lbrace \abs{u_{h_k}}\ge R\right\rbrace}  \left( g(x,u_{h_k})u_{h_k}-qG(x,u_{h_k})\right) d\mathcal{L}^n\\
	&+ \int_{U \cap \left\lbrace \abs{u_{h_k}}\le R\right\rbrace } \left( g(x,u_{h_k})u_{h_k}-qG(x,u_{h_k})\right) d\mathcal{L}^n = qm + \alpha_k - \braket{dJ[u_{h_k}],u_{h_k}}. 
\end{align*}
Ora, per uniforme ellitticità,
\[ \frac{q-2}{2} \int_U \sum_{i,j=1}^{n} a_{ij}D_i u_{h_k} D_j u_{h_k} d\mathcal{L}^n  \ge  \frac{q-2}{2} \nu \norm{u_{h_k}}^2_{H^1_0(U)}, \]
per la $(ii)$,
\[ \int_{U \cap \left\lbrace \abs{u_{h_k}}\ge R\right\rbrace}  \left( g(x,u_{h_k})u_{h_k}-qG(x,u_{h_k})\right) d\mathcal{L}^n \ge 0, \]
per la $(i)$,
\[ \int_{U \cap \left\lbrace \abs{u_{h_k}}\le R\right\rbrace } \left( g(x,u_{h_k})u_{h_k}-qG(x,u_{h_k})\right) d\mathcal{L}^n \ge C_R, \]
quindi possiamo scrivere, utilizzando la continuità del differenziale,
\[  \frac{q-2}{2} \nu \norm*{u_{h_k}}^2_{H^1_0(U)} \le qm +\alpha_k + \norm{dJ[u_{h_k}]}_{H^{-1}}\norm{u_{h_k}}_{H^1_0(U)} + C_R \]
che, a seguito di un'applicazione della disuguaglianza di Young pesata sul terzo addendo a secondo membro, diventa
\[ \norm{u_{h_k}}^2_{H^1_0(U)} \le C <+\infty, \]
ossia $(u_{h_k})$ è limitata in $H^1_0(U)$. A meno di sottosuccessioni, per il Teorema di Eberlein--Smulian, il Teorema di Rellich--Kondrachov, le proprietà degli spazi $L^p$ e l'unicità del limite debole, esiste $u \in H^1_0(U)$ tale che per ogni $p < 2^\ast$ si abbia
\begin{align*}
	u_{h_k} \rightharpoonup u & \textup{ in } H^1_0(U), & u_{h_k} \to u & \textup{ in } L^p(U).
\end{align*}
Tuttavia, per uniforme ellitticità,
\begin{align*}
\nu \norm{Du_{h_k}-Du}_{L^2(U)}^2 &\le \int_U \sum_{i,j=1}^{n} a_{ij} D_i(u_{h_k}-u)D_j(u_{h_k}-u) d\mathcal{L}^n = \\
&= \int_U \sum_{i,j=1}^{n} a_{ij}D_iu_{h_k} D_j u_{h_k} d\mathcal{L}^n - 2 \int_U \sum_{i,j=1}^{n}D_iu_{h_k}D_ju d\mathcal{L}^n + \\
&+\int_U \sum_{i,j=1}^{n}D_iuD_ju d\mathcal{L}^n = \\
&= \int_U g(x,u_{h_k})u_{h_k}d\mathcal{L}^n + \braket{dJ[u_{h_k}],u_{h_k}} - 2 \int_U \sum_{i,j=1}^{n}D_iu_{h_k}D_ju d\mathcal{L}^n + \\
&+ \int_U \sum_{i,j=1}^{n}D_iuD_ju d\mathcal{L}^n.
\end{align*}
A questo punto, osservando che $\braket{dJ[u_{h_k}],u_{h_k}} \to 0$\footnote{Si tratta di un fatto generale. Consideriamo $(X,\norm{\;}_X)$ uno spazio normato su $\K$, $(x;h)$ in $X$ e $(\varphi_h)$ in $X'$. Se esistono $x \in X$ e $\varphi \in X'$ tali che $x_h \rightharpoonup x$ in $X$ e $\varphi_h \to \varphi$ in $X'$, allora $\braket{\varphi_h,x_h} \to \braket{\varphi,x}$.

Infatti, siccome la convergenza debole implica la limitatezza,
\begin{align*}
\abs{\braket{\varphi_h,x_h}-\braket{\varphi,x}}&\le \abs{\braket{\varphi_h,x_h} - \braket{\varphi,x_h}} + \abs{\braket{\varphi,x_h} - \braket{\varphi,x_h}} \le \\
&\le \norm{\varphi_h - \varphi}_{X'}\norm{x_h}_X + \abs{\braket{\varphi,x_h} - \braket{\varphi,x_h}} \le \\
&\le C\norm{\varphi_h - \varphi}_{X'} + \abs{\braket{\varphi,x_h} - \braket{\varphi,x_h}} \to 0.
\end{align*}
}
per ogni $\varepsilon >0$ esiste $N_1 \ge 0 $ tale che per ogni $k\ge N_1$
\[ \braket{dJ[u_{h_k}],u_{h_k}} < \frac{\nu \varepsilon}{3}.  \]
Siccome $u_{h_k} \rightharpoonup u$ in $H^1_0(U)$, essendo $a_{ij}D_ju \in L^2(U)$ per ogni $i,j=1,\dots,n$, esiste $N_2\ge0$ tale che per ogni $k\ge N_2$
\[ \int_U \sum_{i,j=1}^{n}D_iu_{h_k}D_ju d\mathcal{L}^n > \int_U \sum_{i,j=1}^{n}D_iuD_ju d\mathcal{L}^n - \frac{\nu \varepsilon}{6}, \] 
quindi, per $k \ge \max\left\lbrace N_1,N_2\right\rbrace$,
\[ \nu \norm{Du_{h_k}-Du}_{L^2(U)}^2 < \int_U g(x,u_{h_k})u_{h_k}d\mathcal{L}^n - \int_U \sum_{i,j=1}^{n}D_iuD_ju d\mathcal{L}^n + \frac{2}{3} \nu \varepsilon. \]
Siccome per ogni $v \in H^1_0(U)$ si ha $\braket{J[u_{h_k}],v} \to 0$, $u_{h_k} \rightharpoonup u$ in $H^1_0(U)$ e $a_{ij}D_ju \in L^2(U)$ per ogni $i,j=1,\dots,n$, per la continuità dell'operatore di Nemytskij (in $H^{-1}(U)$)\footnote{Prestiamo attenzione al fatto che $g$ non dipende dai gradienti, quindi le convergenze che abbiamo sono sufficienti.}, passando al limite nell'espressione di $\braket{J[u_{h_k}],v}$ si trova
\[ \int_U \sum_{i,j=1}^{n}D_iuD_jv d\mathcal{L}^n = \int_U g(x,u)vd\mathcal{L}^n.\footnote{Potrebbe sembrare che già qui abbiamo trovato una soluzione debole del nostro problema. Attenzione però: noi, già a priori, sappiamo che il problema ammette la soluzione nulla. Stiamo cercando soluzioni non banali. A questo livello nulla ci garantisce che $u$ non sia la soluzione banale.} \]
In particolare, 
\[ \nu \norm{Du_{h_k}-Du}_{L^2(U)}^2 < \int_U g(x,u_{h_k})u_{h_k}d\mathcal{L}^n - \int_U g(x,u)ud\mathcal{L}^n + \frac{2}{3} \nu \varepsilon. \]
Sempre per la continuità dell'operatore di Nemytskij (su $H^{-1}(U)$) e siccome $u_{h_k} \rightharpoonup u$ in $H^1_0(U)$, esiste $N_3 \ge 0$ tale che se $k \ge N_3$
\[ \int_U g(x,u_{h_k})u_{h_k}d\mathcal{L}^n -  \int_U g(x,u)ud\mathcal{L}^n < \frac{\nu \varepsilon}{3}, \]
quindi per ogni $k \ge \max\left\lbrace N_1,N_2,N_3\right\rbrace $
\[ \norm{Du_{h_k}-Du}_{L^2(U)}^2 < \varepsilon,  \]
ossia $u_{h_k} \to u$ in $H^1_0(U)$, da cui la tesi.
\end{proof}

\begin{teorema}
Esiste $u \in H^{1}_0(U)\setminus \left\lbrace 0\right\rbrace $ soluzione debole del problema di Dirichlet
\[ 	\begin{cases}
	-\displaystyle\sum_{i,j=1}^n D_j(a_{ij}(x)D_iu(x)) = g(x,u) & \textup{in } U,\\
	u = 0 & \textup{su } \partial U.
\end{cases} \]
\end{teorema}
\begin{proof}
Innanzitutto, per il Lemma \eqref{Lemma_JPS}, il funzionale $J$ soddisfa la condizione di Palais--Smale. Chiaramente, $J[0] = 0$.
	
Mostriamo che esistono $a,r >0$ tali che per ogni $u \in \partial \B(0,r) \subseteq H^1_0(U)$
\[ J[u] \ge a. \] 
In primo luogo, per uniforme ellitticità e per la $(iii)$,
\[ J[u] \ge \frac{1}{2} \nu \int_U \abs{Du}^2 d\mathcal{L}^n - \varepsilon \int_U \abs{u}^2 d\mathcal{L}^n - C_\varepsilon \int_U \abs{u}^{2^\ast}d\mathcal{L}^n. \]
A questo punto, per la disuguaglianza di Poincaré ed il Teorema di Sobolev,
\[ J[u] \ge \frac{\nu}{2} \norm{u}_{H^1_0(U)}^2 -\varepsilon C_p\norm{u}_{H^1_0(U)}^2 - C_\varepsilon C_s \norm{u}_{H^1_0(U)}^{2^\ast} = \frac{\nu}{2} r^2 -\varepsilon C_pr^2 - C_\varepsilon C_s r^{2^\ast}.  \]
Se scegliamo $\varepsilon$ tale che $\frac{\nu}{2}> \varepsilon C_p$, è sufficiente scegliere $a,r$ tali che 
\[ \left( \frac{\nu}{2} -\varepsilon C_p\right) r^2 - C_\varepsilon C_s r^{2^\ast} \ge a>0. \]

Mostriamo che esiste $v \in H^1_0(U)$ tale che $\norm{v}_{H^1_0(U)}>r$ e $J[v]\le 0$. Per il Lemma \eqref{lemma_tech_minorazioneG} e la $(i)$, fissato $\overline{u} \in H^1_0(U)\setminus \left\lbrace 0\right\rbrace$,
\begin{align*}
	J[t\overline{u}] &\le Ct^2 \norm{\overline{u}}_{H^1_0(U)}^2 - \int_{U \cap \left\lbrace \abs{t\overline{u}} \ge R\right\rbrace} G(x,t\overline{u})d\mathcal{L}^n -\int_{U \cap \left\lbrace \abs{t\overline{u}} \le R\right\rbrace} G(x,t\overline{u})d\mathcal{L}^n \le \\
	&\le C t^2\norm{\overline{u}}_{H^1_0(U)}^2 -\int_{U \cap \left\lbrace \abs{t\overline{u}} \ge R\right\rbrace} k\abs{t\overline{u}}^qd\mathcal{L}^n -\int_{U \cap \left\lbrace \abs{t\overline{u}} \le R\right\rbrace} G(x,t\overline{u})d\mathcal{L}^n = \\
	&=  C t^q\norm{\overline{u}}_{H^1_0(U)}^2 - t^q\int_U k\abs{\overline{u}}^q\mathcal{L}^n +\int_{U \cap \left\lbrace \abs{t\overline{u}} \le R\right\rbrace}(k\abs{t\overline{u}}^q-G(x,t\overline{u}))d\mathcal{L}^n \le \\
	&\le Ct^q\norm{\overline{u}}_{H^1_0(U)}^2 - t^q\int_U k\abs{\overline{u}}^q\mathcal{L}^n + C_R.
\end{align*}
Si tratta quindi di scegliere $v = t\overline{u}$ per $t$ sufficientemente grande da garantire 
\[ \begin{cases}
	Ct^q\norm{\overline{u}}_{H^1_0(U)}^2 - t^q\int_U k\abs{\overline{u}}^q\mathcal{L}^n + C_R \le 0, \\
	t \norm{\overline{u}}_{H^1_0(U)} > r.
\end{cases} \]

Per il Teorema del passo montano, detto 
\[ \Gamma = \left\lbrace g \in C([0,1];H^1_0(U)): g(0) = 0, g(1) = v\right\rbrace,  \]
si ha che 
\[ c = \underset{g \in \Gamma}{\inf}\; \underset{t \in [0,1]}{\max}\; J[g(t)] \]
è un valore critico per $J$, ossia esiste $u \in H^1_0(U)$ tale che $J[u]=c$ e $J'[u] = 0$. In altre parole, $u$ è soluzione debole non banale dell'equazione di Eulero--Lagrange di $J$, ossia del problema di Dirichlet
\[ 	\begin{cases}
	-\displaystyle\sum_{i,j=1}^n D_j(a_{ij}(x)D_iu(x)) = g(x,u) & \textup{in } U,\\
	u = 0 & \textup{su } \partial U,
\end{cases} \]
da cui la tesi.
\end{proof}

Un esempio modello è il seguente.
\begin{esempio}
Consideriamo il problema di Dirichlet
\[ 	\begin{cases}
	-\Delta u = \abs{u}^{p-1}u & \textup{in } U,\\
	u = 0 & \textup{su } \partial U,
\end{cases} \]
dove $p \in \left] 1,\frac{n+2}{n-2}\right[$, corrispondente alla scelta $a_{ij}= \delta_{ij}$, chiaramente ammissibile, e $g(x,s) = \abs{s}^{p-1}s$. 

Vediamo che la scelta di $g$ è ammissibile. Innanzitutto, chiaramente è di Carathéodory e possiamo scegliere $G(x,s) = \frac{\abs{s}^{p+1}}{p+1}$.
Poi, la $(i)$ segue dalla disuguaglianza di Young pesata\footnote{Con uno degli esponenti $\frac{2^\ast-1}{p}$.} e dal fatto che, essendo $U$ limitato, le costanti sono elementi di $L^{\frac{2n}{n+2}}(U)$, infatti  
\[ \abs{s}^p \le C + \varepsilon \abs{s}^{\frac{n+2}{n-2}}.\]
La $(ii)$ segue dalla scelta $q= p+1$. Per quanto riguarda la $(iii)$, segue dalla disuguaglianza di Young pesata, infatti presto $t \in \left] 0,1\right[ $ tale che $p+1 = 2t + 2^\ast(1-t)$, si ha
\[ \frac{\abs{s}^{p+1}}{p+1} = \frac{1}{p+1}\abs{s}^{2t}\abs{s}^{2^\ast(1-t)} \le \varepsilon \abs{s}^2 + C_\varepsilon \abs{s}^{2^\ast}.  \]

Il funzionale associato, in questo caso, è $J: H^1_0(U) \to \R$ tale che
\[ J[u] = \frac{1}{2}\int_U \abs{Du}^2 d\mathcal{L}^n - \frac{1}{p+1} \int_U \abs{u}^{p+1}d\mathcal{L}^n, \]
infatti, per ogni $u,v \in H^1_0(U)$
\[ J'[u](v) = \int_U Du \cdot Dvd\mathcal{L}^n - \int_U \abs{u}^{p-1}uv d\mathcal{L}^n, \]
da cui la corrispondenza tra soluzioni deboli dell'equazione di Eulero--Lagrange e del problema posto sopra.
\end{esempio}
\section{Non esistenza di soluzioni}
Nel seguito, salvo diversa specificazione, supporremo sempre che $n\ge3$, $U$ sia un aperto limitato di $\R^n$ non vuoto con $\partial U$ di classe $C^\infty$. Dato $p \in \left] 1,+\infty\right[$, consideriamo il problema di Dirichlet
\begin{equation}\label{modelloapotenza}
\begin{cases}
	-\Delta u = \abs{u}^{p-1}u & \textup{in } U,\\
	u = 0 & \textup{su } \partial U.
\end{cases}
\end{equation}

Nella sezione precedente abbiamo potuto osservare che nel caso $p \in \left] 1,\frac{n+2}{n-2}\right[$, mediante il Metodo dei punti critici, il problema \eqref{modelloapotenza} ammette una soluzione non banale. In realtà, siccome il funzionale $J$ associato è pari, è possibile dimostrare che esistono infinite soluzioni $u_j \in H^1_0(U)$ non banali e a segni alterni tali che $J[u_j]\to +\infty$.

Studiamo ora il caso  $p \in \left[ \frac{n+2}{n-2}, +\infty\right[$.

\begin{lemma}\label{normaleperPohozaev}
Supponiamo che $U$ sia stellato rispetto a $0$. Allora per ogni $x \in \partial U$ si ha
\[ x \cdot \nu(x) \ge 0. \]
\end{lemma}
\begin{proof}
Fissiamo $x \in \partial U$. Mostriamo, innanzitutto, che, per ogni $\varepsilon>0$ esiste $\delta >0$ tale che per ogni $y \in \overline{U}$ tale che $\abs{x-y}<\delta$ si ha 
\[ \nu(x)\cdot \frac{y-x}{\abs{y-x}} \le \varepsilon. \]
Se, per assurdo, esiste $\varepsilon >0$ tale che per ogni $j \in \N$ esiste $y_j \in \B(x, \frac{1}{j}) \cap U$ tale che
\[ \nu(x) \cdot \frac{y_j-x}{\abs{y_j -x}} >\varepsilon>0, \]
dal fatto che $y_j \to x$, si avrebbe $\varepsilon \le 0$, assurdo. In particolare, 
\[ \underset{\underset{y \in \overline{U}}{y \to x}}{\limsup}\; \nu(x)\cdot \frac{y-x}{\abs{y-x}} \le 0.  \]
Dal fatto che $U$ è stellato rispetto a $0$, per ogni $\lambda \in [0,1]$ $y = \lambda x \in U \subseteq \overline{U}$, allora
\[ \nu(x) \cdot \frac{x}{\abs{x}} = - \lim\limits_{\lambda \to 1^-} \nu(x) \cdot \frac{\lambda x-x}{\abs{\lambda x -x}} \ge 0, \]
da cui la tesi.  
\end{proof}
Vediamo che, sotto l'ipotesi che $U$ sia stellato rispetto a $0$, non esistono soluzioni non banali per il problema \eqref{modelloapotenza}.
\begin{teorema}
Supponiamo che $U$ sia stellato rispetto a $0$ e $p \in \left] \frac{n+2}{n-2}, +\infty\right[$. Se $u \in C^2(\overline{U})$ è una soluzione del problema \eqref{modelloapotenza}, allora $u = 0$ in $U$.
\end{teorema}
\begin{proof}
Innanzitutto, dal fatto che $u$ è una soluzione classica di \eqref{modelloapotenza}, per ogni $x \in U$ si ha
\[ 	-\Delta u(x) = \abs{u(x)}^{p-1}u(x). \]
Moltiplicando ambo i membri per $x \cdot Du$ ed integrando su $U$ otteniamo 
\[ \int_U (-\Delta u(x))(x \cdot Du(x))d\mathcal{L}^n(x) = \int_U \abs{u(x)}^{p-1}u(x) (x\cdot Du)d\mathcal{L}^n(x). \]
Per quanto riguarda il primo membro, per le formule di Gauss--Green, 
\begin{align*}
	&\int_U (-\Delta u(x))(x \cdot Du(x))d\mathcal{L}^n(x) = -\sum_{i,j=1}^{n}\int_U D^2_{i,i}u x_j D_ju d\mathcal{L}^n(x) = \\
	&= \sum_{i,j=1}^{n}\int_U D_iu D_i(x_j D_ju) d\mathcal{L}^n(x) - \sum_{i,j=1}^{n}\int_{\partial U} D_iu x_j D_ju \nu_i d\mathcal{H}^{n-1}(x)= \\
	&= \sum_{i,j=1}^{n} \int_U D_i u \delta_{ij} D_ju d\mathcal{L}^n + \sum_{i,j=1}^{n} \int_U D_iu x_j D^2_{i,j}u d\mathcal{L}^n(x) - \sum_{i,j=1}^{n}\int_{\partial U} D_iu x_j D_ju \nu_i d\mathcal{H}^{n-1}(x)= \\
	&= \sum_{i=1}^{n}\int_{U} \abs{D_iu}^2 d\mathcal{L}^{n} + \frac{1}{2}\sum_{j=1}^{n}\int_U D_j\left( \abs{Du}^2\right)  x_j d\mathcal{L}^n(x) - \sum_{i,j=1}^{n}\int_{\partial U} D_iu x_j D_ju \nu_i d\mathcal{H}^{n-1}(x)=\\
	&= \int_{U} \abs{Du}^2 d\mathcal{L}^{n} + \frac{1}{2}\left( -\sum_{j=1}^{n} \int_U \abs{Du}^2 d\mathcal{L}^n + \sum_{j=1}^{n}\int_{\partial U} \abs{Du}^2 x_j\nu_j d\mathcal{H}^{n-1}(x)\right) + \\
	&- \sum_{i,j=1}^{n}\int_{\partial U} D_iu x_j D_ju \nu_i d\mathcal{H}^{n-1}(x) = \\
	&= -\frac{n-2}{2} \int_U \abs{Du}^2 d\mathcal{L}^n + \frac{1}{2}\int_{\partial U} \abs{Du}^2 (x \cdot \nu)d\mathcal{H}^{n-1}(x) -  \sum_{i,j=1}^{n}\int_{\partial U} D_iu x_j D_ju \nu_i d\mathcal{H}^{n-1}(x),
\end{align*}
ma siccome $u = 0$ su $\partial U$, $Du$ è parallelo a $\nu$ su $\partial U$, ossia $Du = \pm \abs{Du}\nu$. In particolare, per ogni $i=1,\dots,n$ si ha $D_iu = \pm \abs{Du}\nu_i$, quindi 
\begin{align*}
	&\int_U (-\Delta u(x))(x \cdot Du(x))d\mathcal{L}^n(x) = \\
	&=\frac{2-n}{2} \int_U \abs{Du}^2 d\mathcal{L}^n + \frac{1}{2}\int_{\partial U} \abs{Du}^2 (x \cdot \nu)d\mathcal{H}^{n-1}(x) - \sum_{i,j=1}^{n}\int_{\partial U} \abs{Du}^2\nu_i x_j\nu_j \nu_i d\mathcal{H}^{n-1}(x) = \\
	&= \frac{2-n}{2} \int_U \abs{Du}^2 d\mathcal{L}^n - \frac{1}{2}\int_{\partial U} \abs{Du}^2 (x \cdot \nu)d\mathcal{H}^{n-1}(x). 
\end{align*}
Per quanto riguarda il secondo membro, invece, sempre per le formule di Gauss--Green ed il fatto che $u = 0$ su $\partial U$, 
\begin{align*}
\int_U \abs{u(x)}^{p-1}u(x) (x\cdot Du)d\mathcal{L}^n(x) &= \sum_{j=1}^{n} \int_U \abs{u}^{p-1} u x_j D_ju d\mathcal{L}^n(x) = \\
	&= \frac{1}{p+1} \sum_{j=1}^{n}\int_U D_j\left( \abs{u}^{p+1}\right)  x_j d\mathcal{L}^n(x) = \\
	&= -\frac{n}{p+1}\int_U \abs{u}^{p+1}d\mathcal{L}^n.
\end{align*}
Riassumendo, 
\begin{equation}\label{Derrick--Pohozaevidentity}
 \frac{n-2}{2} \int_U \abs{Du}^2 d\mathcal{L}^n + \frac{1}{2}\int_{\partial U} \abs{Du}^2 (x \cdot \nu)d\mathcal{H}^{n-1}(x) = \frac{n}{p+1}\int_U \abs{u}^{p+1}d\mathcal{L}^n.\footnote{Nota come \emph{identità di Derrick--Pohozaev}\index[analitico]{identità!di Derrick--Pohozaev}.}
\end{equation}

Combinando la \eqref{Derrick--Pohozaevidentity} con il Lemma \eqref{normaleperPohozaev}, possiamo scrivere
\[  \frac{n-2}{2} \int_U \abs{Du}^2 d\mathcal{L}^n \le \frac{n}{p+1}\int_U \abs{u}^{p+1}d\mathcal{L}^n.  \]
 
Tuttavia, la formulazione variazionale del problema \eqref{modelloapotenza}, per $v =u$, conduce a 
\[ \int_U \abs{Du}^2 d\mathcal{L}^n = \int_U \abs{u}^{p+1}d\mathcal{L}^n, \]
quindi 
\[ \left(  \frac{n-2}{2} - \frac{n}{p+1}\right) \int_U \abs{Du}^2 d\mathcal{L}^n \le 0.   \]
Siccome 
\[ \frac{n-2}{2} - \frac{n}{p+1} >0,\footnote{Si tratta di un calcolo diretto,
		\[ \frac{n-2}{2} - \frac{n}{p+1} >0 \Longleftrightarrow  (p+1)(n-2) > 2n \Longleftrightarrow  (n-2)p + n-2 > 2n \Longleftrightarrow  (n-2)p > n+2, \]
	quindi,
	\[ 	\frac{n-2}{2} - \frac{n}{p+1} >0 \Longleftrightarrow p > \frac{n+2}{n-2}. \]
} \]
l'unica possibilità è che $u = 0$, da cui la tesi.
\end{proof}

\begin{osservazione}
Nel caso $p = \frac{n+2}{n-2}$, la formulazione variazionale del problema \eqref{modelloapotenza}, per $v =u$, combinata con \eqref{Derrick--Pohozaevidentity}, conduce a 
\[\int_{\partial U} \abs{Du}^2 (x \cdot \nu)d\mathcal{H}^{n-1}(x) = 0. \]
Se $x \cdot \nu >0$ su una porzione di $\partial U$ di misura positiva,  dal fatto che $u = 0$ su $\partial U$, tramite una \emph{proprietà di continuazione unica}\index[analitico]{continuazione unica}, si può affermare anche in questo caso che $u=0$ su $U$. I dettagli della dimostrazione esulano dai nostri scopi.
\end{osservazione}

\begin{osservazione}
Nel caso $p= \frac{n+2}{n-2}$, se $U$ non è stellato rispetto a $0$, è possibile fornire, sotto opportune ipotesi tecniche, dei risultati di esistenza di soluzioni non banali: si tratta però, principalmente, di scenari molto specialistici che esulano dai nostri scopi.
\end{osservazione}

Vediamo, infine, un caso che esula dalla nostra trattazione, ma comunque degno di nota.
\begin{esempio}
Consideriamo il problema di Dirichlet
\[ 	\begin{cases}
	-\Delta u = u^{\frac{n+2}{n-2}} & \textup{in } \R^n,\\
	u > 0 & \textup{in } \R^n.\\
\end{cases} \]
Intorno al 1976, Giorgio Talenti ha lavorato su questo problema esibendo come soluzione
\[ \mathcal{T}_0(x) = \left( n(n-2)\right) ^{\frac{n-2}{4}} \frac{1}{(1+\abs{x}^2)^{\frac{n-2}{2}}}.\]
Addirittura, ogni elemento della \emph{famiglia delle talentiane}\index[analitico]{applicazione!talentiana}, ossia  
\[ \mathcal{T}_\varepsilon(x) = \varepsilon^{\frac{n-2}{2}}\mathcal{T}_0\left( \varepsilon(x-x_0)\right), \qquad \varepsilon>0,x_0\in \R^n, \]
è una soluzione. Infatti, 
\begin{align*}
 -\Delta \mathcal{T}_\varepsilon &= -\varepsilon^{\frac{n-2}{2}} \Delta \left(\mathcal{T}_0\left( \varepsilon(x-x_0)\right)  \right) = -\varepsilon^{\frac{n-2}{2}+ 2} (\Delta \mathcal{T}_0)(\varepsilon(x-x_0)) = \\
 &= \varepsilon^{\frac{n-2}{2}+2} \mathcal{T}_0\left( \varepsilon(x-x_0)\right)^{\frac{n+2}{n-2}} = \varepsilon^{\frac{n-2}{2}+2} \mathcal{T}_0\left( \varepsilon(x-x_0)\right)^{\frac{n+2}{n-2}} \left( \frac{\varepsilon^{\frac{n-2}{2}}}{\varepsilon^{\frac{n-2}{2}}}\right)^{\frac{n+2}{n-2}} = \\
 &=\varepsilon^{\frac{n-2}{2}+2 -\frac{n-2}{2}}\mathcal{T}_\varepsilon(x)^{\frac{n+2}{n-2}} = \mathcal{T}_\varepsilon(x)^{\frac{n+2}{n-2}}. 
\end{align*}
\end{esempio} 		
		\printindex[simboli]
		\printindex[analitico]
	\end{document}